%--------------------------------------------------------------
% thesis.tex 
%--------------------------------------------------------------
% Corso di Laurea in Informatica 
% - template for the main file of Informatica@Unifi Thesis 
% - based on Classic Thesis Style Copyright (C) 2008 
%   Andr\'e Miede http://www.miede.de   
%--------------------------------------------------------------
\documentclass[oneside,openany,titlepage,fleqn,
	headinclude,12pt,a4paper,footinclude]{scrbook}

\usepackage[nottoc]{tocbibind}
%--------------------------------------------------------------
% Inputs for the title page 

\newcommand{\myName}{Giosuè Aiello\xspace}
\newcommand{\myTime}{2026\xspace}
%--------------------------------------------------------------

\usepackage[utf8]{inputenc} 
\usepackage[T1]{fontenc} 

% switch italian and english, depending on the language of the thesis
% this will change the chapter titles
\usepackage[italian]{babel}
%\usepackage[english]{babel}

\usepackage[fleqn]{amsmath}
\usepackage{amssymb}
\usepackage{array}       % \extrarowheight e colonne personalizzate  
\usepackage{ellipsis}
\usepackage{float}
\usepackage{caption}
\usepackage{cancel}

\usepackage[pdftex]{graphicx} 
\graphicspath{{./Figure/}}

% Font dei titoli KOMA-Script di default reimpostato a Serif
\setkomafont{disposition}{\normalfont\rmfamily\bfseries}

\usepackage{tikz}
\usepackage{tikz-3dplot}
\usetikzlibrary{shapes,arrows,positioning,calc,fit,backgrounds,decorations.pathreplacing,decorations.pathmorphing,patterns,trees}
\definecolor{mygreen}{rgb}{0.65, 0.75, 0.1}
\definecolor{mypen}{rgb}{0.1, 0.1, 0.4}
\definecolor{notered}{RGB}{180,0,0}
\definecolor{notedarkred}{RGB}{120,0,0}
\usepackage{tcolorbox}
\tcbuselibrary{skins,breakable}
\usepackage{enumitem}
\usepackage{ccicons}

% --- Custom Styling Definitions ---
% Palette semantica dei riquadri didattici
\definecolor{boxdef}{RGB}{168,38,38}    % Definizioni      (rosso mattone)
\definecolor{boxteo}{RGB}{186,112,12}   % Teoremi          (ocra)
\definecolor{boxoss}{RGB}{28,82,158}    % Osservazioni     (blu)
\definecolor{boxese}{RGB}{24,112,84}    % Esempi/Esercizi  (verde bosco)
\definecolor{boxatt}{RGB}{158,48,20}    % Avvertenze       (terra di Siena)
\definecolor{boxnota}{RGB}{102,72,142}  % Note storiche    (viola)
\definecolor{boxform}{RGB}{16,102,112}  % Formulari        (teal)

% Stile base condiviso: un'unica geometria per tutti i riquadri,
% differenziati soltanto dal colore semantico passato come argomento.
\tcbset{
    fisicabox/.style={
        enhanced,
        breakable,
        sharp corners,
        frame hidden,
        boxrule=0pt,
        colback=#1!4!white,
        borderline west={2.2pt}{0pt}{#1},
        left=9pt, right=9pt, top=6pt, bottom=6pt,
        before skip=9pt, after skip=9pt,
        fonttitle=\bfseries,
        coltitle=#1!82!black,
        colbacktitle=#1!11!white,
        toptitle=3.5pt, bottomtitle=3.5pt,
        titlerule=0pt
    },
    % Variante senza titolo (riquadro "muto")
    fisicaboxplain/.style={
        enhanced,
        breakable,
        sharp corners,
        frame hidden,
        boxrule=0pt,
        colback=#1!4!white,
        borderline west={2.2pt}{0pt}{#1},
        left=9pt, right=9pt, top=7pt, bottom=7pt,
        before skip=9pt, after skip=9pt
    }
}

% Etichetta del titolo: maiuscoletto per la categoria, tondo per il nome
\newcommand{\boxlabel}[2]{{\scshape #1}\ \textemdash\ #2}

\newtcolorbox{definizione}[1]{fisicabox=boxdef,  title={\boxlabel{Definizione}{#1}}}
\newtcolorbox{teorema}[1]{fisicabox=boxteo,      title={\boxlabel{Teorema}{#1}}}
\newtcolorbox{osservazione}[1]{fisicabox=boxoss, title={\boxlabel{Osservazione}{#1}}}
\newtcolorbox{esercizio}[1]{fisicabox=boxese,    title={\boxlabel{Esercizio}{#1}}}
\newtcolorbox{esempio}[1]{fisicabox=boxese,      title={\boxlabel{Esempio}{#1}}}
\newtcolorbox{attenzione}[1]{fisicabox=boxatt,   title={\boxlabel{Attenzione}{#1}}}
\newtcolorbox{notastorica}[1]{fisicabox=boxnota, title={\boxlabel{Nota storica}{#1}}}
\newtcolorbox{formulario}[1]{fisicabox=boxform,  title={\boxlabel{Formulario}{#1}}}

% Riquadro di sintesi senza intestazione, per i risultati chiave
\newtcolorbox{risultato}{fisicaboxplain=boxform}

\usepackage{scrlayer-scrpage}

% Clear default headers/footers
\clearpairofpagestyles

% Define a new layer for our TikZ page number
\DeclareNewLayer[
    background,
    contents={%
        \begin{tikzpicture}[remember picture, overlay]
            \ifodd\value{page}
                % Odd pages: number at bottom-right corner
                \node[anchor=center, font=\small] at ($(current page.south east)+(-1cm,1cm)$) {\thepage};
            \else
                % Even pages: number at bottom-left corner
                \node[anchor=center, font=\small] at ($(current page.south west)+(1cm,1cm)$) {\thepage};
            \fi
        \end{tikzpicture}%
    }
]{pagenumberlayer}

% Add this layer to both plain (chapter) and headings (normal) styles
\AddLayersToPageStyle{scrheadings}{pagenumberlayer}
\AddLayersToPageStyle{plain.scrheadings}{pagenumberlayer}

% Define a custom command to toggle page numbering visibility
\newcommand{\suppresspagenumber}{\clearpairofpagestyles}
\newcommand{\activatepagenumber}{%
    \clearpairofpagestyles%
    \AddLayersToPageStyle{scrheadings}{pagenumberlayer}%
    \AddLayersToPageStyle{plain.scrheadings}{pagenumberlayer}%
}

%--------------------------------------------------------------
% Layout setting
%--------------------------------------------------------------
\setlength{\extrarowheight}{3pt} 
\captionsetup{format=hang,font=small}

\usepackage{geometry}
\geometry{
	a4paper,
	ignoremp,
	left=1.8cm,
	right=1.8cm,
	top=2cm,
	bottom=2cm,
	headheight=15pt
}


%--------------------------------------------------------------

\usepackage{booktabs}  % for \toprule, \midrule, \bottomrule
\usepackage{pifont}    % for \ding{55}
\usepackage{hyperref}  % load LAST - for \url, \texorpdfstring
\usepackage{bookmark}
\begin{document}
\frenchspacing
\raggedbottom
\setcounter{page}{1} % Title page is page 1
\suppresspagenumber % No page numbers in frontmatter
\pagestyle{scrheadings} % active style (but currently empty)
%--------------------------------------------------------------
% Frontmatter
%--------------------------------------------------------------
\begin{titlepage}
  \newgeometry{margin=2.5cm}
  \centering
  \vspace*{3.4cm}
  {\Large\scshape\color{notedarkred}Appunti di\par}
  \vspace{-0.6cm}
  {\fontsize{78}{85}\selectfont\bfseries\color{notered}FISICA\ {\color{notedarkred}1}\par}
  \vspace{1cm}
  {\color{notered}\rule{8cm}{1.2pt}\par}
  \vspace{1.2cm}
  {\Large\scshape Giosu\`e Aiello\par}
  \vfill
  \noindent
  \begin{minipage}[b]{0.5\linewidth}
    \textit{versione 1.0}
  \end{minipage}%
  \begin{minipage}[b]{0.5\linewidth}
    \raggedleft
    Universit\`a di Pisa\\
    anno \myTime
  \end{minipage}
  \thispagestyle{empty}
  \restoregeometry
\end{titlepage}


% Dedication at top
\begin{flushright}
\vspace*{1cm}
\emph{Alla mia famiglia:\\
Francesco Aiello, Rosaria Fassari,\\
Emanuele Antonino Aiello,\\
Alfio Fassari, Domenica Calvagna,\\
ai migliori amici che potessi avere,\\
i Cicci di cho cho:\\
Martina Marzione, Davide Lombardo,\\
Caterina Schinzi, Francesca Guido\\
e Agnese Micheletti.}
\end{flushright}

\vfill
\noindent
\textbf{Teoria di Fisica 1} \\
\copyright\ \myTime \ by \myName \\

\noindent
{\Large \ccbync} Licensed under Creative Commons Attribution-NonCommercial 4.0 International. \\
To view a copy of this license, visit \url{https://creativecommons.org/licenses/by-nc/4.0/}.

\clearpage
\thispagestyle{empty}
\vspace*{\fill}
\begin{center}
    {\Large\scshape\color{notedarkred} Nota di merito \par}
    \vspace{0.7cm}
    {\color{notered}\rule{6cm}{1pt}\par}
    \vspace{1.4cm}

    \begin{minipage}{0.78\textwidth}
        \setlength{\parskip}{0.7em}
        \noindent
        Questo testo non nasce dal nulla. Alla sua origine ci sono gli
        \textbf{appunti manoscritti di Chiara Rustici}, studentessa
        dell'Università di Pisa.

        \noindent
        A partire da quelle pagine mi sono divertito a costruire una versione
        più curata: ho riformulato diversi passaggi di teoremi e definizioni,
        ampliato le spiegazioni e arricchito il testo di figure oltre quanto
        presente nelle note originarie, il tutto per offrire al lettore
        un'esperienza di studio più comoda.

        \noindent
        Grazie, Chiara, per la base.
    \end{minipage}

    \vspace{1.8cm}
    {\color{notered}\rule{3cm}{0.6pt}\par}
\end{center}
\vspace*{\fill}

\clearpage
\addtocontents{toc}{\protect\thispagestyle{empty}} % Double check TOC is empty
{\hypersetup{linkcolor=black}
\pdfbookmark[0]{Indice}{indice}
\tableofcontents
}

\clearpage
%--------------------------------------------------------------
% Mainmatter
%--------------------------------------------------------------
\activatepagenumber % Start showing page numbers with our custom layer
% We don't reset counters, continuous from cover=1.
% Explicitly ensure chapter pages use the plain style which now has our layer
\renewcommand*{\chapterpagestyle}{plain.scrheadings}

%--------------------------------------------------------------
% Chapters
%--------------------------------------------------------------
% User notes will go here once provided
\chapter{Introduzione}

Alla fine del 1700, quella che veniva chiamata "Filosofia naturale" iniziò a scindersi in due grandi ambiti di studio:
\begin{itemize}
    \item \textbf{Scienze Biologiche}: dedicate allo studio degli esseri viventi.
    \item \textbf{Scienze Fisiche}: dedicate allo studio della materia inanimata.
\end{itemize}
All'interno delle scienze fisiche, si consolidò ulteriormente la distinzione tra:
\begin{itemize}
    \item \textbf{Chimica}: che studia la formazione e la trasformazione delle molecole.
\end{itemize}

\begin{definizione}{Fisica}
La Fisica (dal greco \emph{physis}, "natura") è una \textbf{scienza sperimentale} in continua evoluzione che indaga i costituenti fondamentali della materia e le loro interazioni attraverso il \textbf{Metodo Scientifico}. Si occupa dei fenomeni naturali, sia quelli "reali" (osservati in natura) che quelli "riprodotti" (in laboratorio).
\end{definizione}

\vspace{0.5cm}

\begin{figure}[ht]
\centering
\begin{tikzpicture}[scale=0.65, transform shape,
    node distance=1.5cm,
    auto,
    block/.style={rectangle, draw, fill=blue!8, text width=10em, text centered, rounded corners=3pt, minimum height=4em, thick},
    decision/.style={diamond, aspect=2, draw, fill=yellow!15, text width=5em, text centered, inner sep=2pt, thick},
    outcome/.style={rectangle, draw, thick, text width=9em, text centered, rounded corners=3pt, minimum height=3.5em},
    arrow/.style={->, >=stealth, thick, draw=gray!80}
]
    % Nodes aligned vertically
    \node [block] (osserv) {Osservazione};
    \node [block, below=of osserv] (riprod) {Riproduzione Semplificata\\(Esperimento)};
    \node [block, below=of riprod] (dati) {Raccolta Dati};
    \node [block, below=of dati] (interp) {Interpretazione\\(Teoria)};
    
    \node [decision, below=1.8cm of interp] (verifica) {Conferma?};

    % Outcomes (closer to diamond as requested)
    \node [outcome, right=1.5cm of verifica, fill=green!15] (tesi) {Tesi Confermata};
    \node [outcome, left=1.5cm of verifica, fill=red!15] (crisi) {Crisi\\(Correzioni)};

    % Edges
    \draw [arrow] (osserv) -- (riprod);
    \draw [arrow] (riprod) -- (dati);
    \draw [arrow] (dati) -- (interp);
    \draw [arrow] (interp) -- (verifica);
    
    \draw [arrow] (verifica) -- node [above, near start] {\small SÌ} (tesi);
    \draw [arrow] (verifica) -- node [above, near start] {\small NO} (crisi);
    
    % Feedback loop: detached from lilac blocks, reconnecting from the left to the first light block
    \draw [arrow, rounded corners=15pt] (crisi.west) -- ++(-1.5,0) |- node [pos=0.25, left, rotate=90, anchor=south] {\small Nuova Ipotesi} (osserv.west);

\end{tikzpicture}
\caption{Schema del Metodo Scientifico: un processo ciclico di verifica.}
\label{fig:metodo}
\end{figure}

\newpage

\section{Il Metodo Scientifico Sperimentale}

Il metodo scientifico è il pilastro dell'indagine fisica. È un ciclo continuo che permette di validare o confutare le teorie.
Partendo dall'\textbf{osservazione} di un fenomeno, si tenta una \textbf{riproduzione semplificata} (esperimento) per raccogliere \textbf{dati}. L'\textbf{interpretazione} di questi dati porta alla formulazione di una \textbf{teoria}. Se la teoria viene verificata sperimentalmente, diventa una tesi confermata; altrimenti, si entra in una fase di crisi che richiede correzioni o nuove ipotesi (vedi Fig. \ref{fig:metodo}).



\section{Evoluzione delle Branche della Fisica}

La Fisica si è evoluta storicamente passando dalla visione classica a quella moderna, spinte dalla necessità di spiegare fenomeni sempre più complessi o estremi.

\begin{figure}[ht]
\centering
\begin{tikzpicture}[scale=0.85, transform shape,
    box_classica/.style={rectangle, draw=black, thick, fill=orange!15, text width=3.8cm, text centered, rounded corners=3pt, minimum height=1.2cm},
    box_moderna/.style={rectangle, draw=black, thick, fill=cyan!15, text width=3.8cm, text centered, rounded corners=3pt, minimum height=1.2cm},
    box_ponte/.style={rectangle, draw=black, thick, fill=yellow!25, text width=4cm, text centered, rounded corners=3pt, minimum height=1.8cm},
    header/.style={font=\bfseries\large},
    label_style/.style={font=\itshape\small},
    arrow/.style={->, >=stealth, line width=2pt, draw=gray!60}
]

    % Column Headers
    \node[header] at (0, 1) {FISICA CLASSICA};
    \node[header] at (10, 1) {FISICA MODERNA};

    % Classica Nodes
    \node[box_classica] (mecc) at (0, 0) {Meccanica};
    \node[box_classica] (termo) at (0, -1.5) {Termodinamica};
    \node[box_classica] (acus) at (0, -3) {Acustica};
    \node[box_classica] (ott) at (0, -4.5) {Ottica};

    % Moderna Nodes
    \node[box_moderna] (rel) at (10, 0) {Relatività};
    \node[box_moderna] (atom) at (10, -1.5) {Fisica Atomica};
    \node[box_moderna] (nucl) at (10, -3) {Fisica Nucleare};
    \node[box_moderna] (sub) at (10, -4.5) {Fisica Subnucleare};

    % Ponte Elettromagnetismo (Vertically aligned with column centers at -2.25)
    \node[box_ponte] (em) at (5, -2.25) {Elettromagnetismo\\($\vec{E}, \vec{B}$)};
    
    % Passaggio Text (outside box, stylized)
    \node[below=0.3cm of em, label_style] {\textit{Passaggio}};

    % Clean, straight, centered arrows in the gaps
    % Gap Left: mecc.east (x=1.9) to em.west (x=3.0). Gap = 1.1cm.
    % We draw arrow from x=2.1 to x=2.8.
    \draw [arrow] (2.1, -2.25) -- (2.8, -2.25);

    % Gap Right: em.east (x=7.0) to rel.west (x=8.1). Gap = 1.1cm.
    % We draw arrow from x=7.2 to x=7.9.
    \draw [arrow] (7.2, -2.25) -- (7.9, -2.25);

\end{tikzpicture}
\caption{Lo sviluppo storico dalla Fisica Classica alla Moderna.}
\end{figure}

Dalla \textbf{Meccanica Newtoniana} (o Classica), valida per velocità basse rispetto alla luce ($v \ll c$) e dimensioni umane, si è passati a nuove teorie per spiegare gli estremi:
\begin{itemize}
    \item \textbf{Relatività Ristretta e Generale}: necessarie per velocità prossime a ($c$) o forti campi gravitazionali.
    \item \textbf{Meccanica Quantistica}: necessaria per descrivere il comportamento della materia su scala atomica ($L \sim$ atomo).
\end{itemize}

Inoltre, esiste un legame circolare tra \textbf{Fisica e Tecnologia}: la fisica ispira nuove tecnologie, e le nuove tecnologie forniscono strumenti più precisi per fare nuova fisica (es. acceleratori di particelle, laser, semiconduttori).

\section{La Misurazione}
\begin{definizione}{Misurazione}
La misurazione è il processo di confronto di una grandezza fisica con un'unità di misura predefinita. Poiché è impossibile ottenere un valore "esatto", il risultato si esprime come:
\[ x \pm \delta x \]
\end{definizione}

Le cause di incertezza possono essere:
\begin{itemize}
    \item \textbf{Variabilità del campione}.
    \item \textbf{Errore Sistematico}: precisione dello strumento.
    \item \textbf{Errore Accidentale}: fluttuazione delle misure ripetute sullo stesso campione.
\end{itemize}

\begin{esercizio}{Propagazione dell'Errore: Esempio Geometrico}
Un errore può propagarsi attraverso i calcoli, rendendo la determinazione dell'incertezza $\delta x$ complessa.
Consideriamo l'esempio semplice dell'area di un rettangolo misurato.

    \begin{center}
    \begin{tikzpicture}[scale=1.5]
        % Rettangolo Principale A = ab
        \draw[thick, fill=blue!5] (0,0) rectangle (3,2);
        \node at (1.5,1) {\Large $Area = a \cdot b$};
        
        % Striscia errore altezza (a * delta b)
        \draw[thick, dashed, fill=red!10] (0,2) rectangle (3,2.5);
        \node at (1.5, 2.25) {\small $a \cdot \delta b$};
        
        % Striscia errore base (b * delta a)
        \draw[thick, dashed, fill=red!10] (3,0) rectangle (3.5,2);
        \node[rotate=90] at (3.25, 1) {\small $b \cdot \delta a$};
        
        % Angolo trascurabile (delta a * delta b)
        \draw[thick, dashed, fill=gray!20] (3,2) rectangle (3.5,2.5);
        \node[scale=0.7, align=center] at (4.2, 2.25) {$\delta a \cdot \delta b$\\trascurabile};
        
        % Quote
        \draw[<->] (0,-0.2) -- (3,-0.2) node[midway, below] {$a$};
        \draw[<->] (3,-0.2) -- (3.5,-0.2) node[midway, below] {$\delta a$};
        
        \draw[<->] (-0.2,0) -- (-0.2,2) node[midway, left] {$b$};
        \draw[<->] (-0.2,2) -- (-0.2,2.5) node[midway, left] {$\delta b$};
    \end{tikzpicture}
    \captionof{figure}{Propagazione dell'errore nel calcolo dell'Area ($A \pm \delta A$).}
    \label{fig:errore_area}
    \end{center}

Se misuriamo i lati come $a \pm \delta a$ e $b \pm \delta b$, l'area calcolata sarà:
\[
A_{mis} = (a \pm \delta a)(b \pm \delta b) = \underbrace{ab}_{A} \pm \underbrace{a\delta b \pm b\delta a}_{\delta A} + \underbrace{\delta a \delta b}_{\approx 0}
\]
Poiché $\delta a \ll a$ e $\delta b \ll b$, il termine quadratico $\delta a \delta b$ è trascurabile. L'errore assoluto sull'area è quindi:
\[ \delta A \approx a\delta b + b\delta a \]
\end{esercizio}

\subsection{Rappresentazione dei Dati ed Errori}
A questo proposito introduciamo i concetti di:
\begin{itemize}
    \item \textbf{Cifra Significativa (CF)}: Indica la precisione della misura (es. misurare $3,0$ è diverso da misurare $3$, in quanto indica una precisione decimale).
    \item \textbf{Notazione Scientifica (NS)}: Si usa la forma $x \cdot 10^n$ per scrivere numeri molto grandi o piccoli mantenendo chiare le cifre significative.
\end{itemize}

L'errore può essere classificato come:
\begin{itemize}
    \item \textbf{Assoluto} ($x \pm \delta x$): Ha la stessa unità di misura della grandezza.
    \item \textbf{Relativo} ($x \pm \frac{\delta x}{|x|}$): È un numero puro (spesso espresso in \%).
\end{itemize}

\begin{definizione}{Grandezza Fisica}
La definizione di una grandezza fisica si dice \textbf{operativa}, poiché elenca le procedure da compiere per misurarne la quantità attraverso il confronto con un'\textbf{Unità di Misura} (UdM). Se le grandezze a confronto non sono della stessa specie (\textbf{non omogenee}), si opera una equivalenza tramite un \textbf{fattore di ragguaglio}.
\end{definizione}

Discende quindi la distinzione tra:
\begin{itemize}
    \item \textbf{Grandezze Fondamentali} (o Primitive): Misurate per via \textbf{diretta}.
    \item \textbf{Grandezze Derivate}: Misurate per via \textbf{indiretta} (tramite relazioni matematiche tra grandezze primitive).
\end{itemize}

Nel Sistema Internazionale (SI), le \textbf{Grandezze Fondamentali} sono 7:
\begin{enumerate}
    \item \textbf{Lunghezza} ($m$, metro).
    \item \textbf{Massa} ($kg$, chilogrammo) - misurata con bilancia o spettrometro di massa.
    \item \textbf{Tempo} ($s$, secondo) - misurato con orologio.
    \item \textbf{Temperatura} ($K$, Kelvin) - misurata con termometro.
    \item \textbf{Intensità di corrente} ($A$, Ampere) - misurata con amperometro.
    \item \textbf{Intensità luminosa} ($cd$, candela) - misurata con fotometro.
    \item \textbf{Quantità di sostanza} ($mol$, mole).
\end{enumerate}

\section{Sistemi di Unità di Misura}
Esistono diversi sistemi di unità di misura costituiti da grandezze primitive indipendenti tra loro. I principali sono:
\begin{itemize}
    \item \textbf{MKS}: Metro ($m$), Chilogrammo ($kg$), Secondo ($s$). È alla base del Sistema Internazionale (SI).
    \item \textbf{CGS}: Centimetro ($cm$), Grammo ($g$), Secondo ($s$). Usato spesso in chimica o elettromagnetismo teorico.
    \item \textbf{Pratico}: Metro ($m$), Chilogrammo-peso ($kg_p$), Secondo ($s$). In questo sistema la forza è fondamentale e la massa derivata.
\end{itemize}

Tutti questi sistemi sono caratterizzati da tre dimensioni fondamentali: Lunghezza $[L]$, Massa $[M]$ e Tempo $[T]$.

\begin{osservazione}{Rivoluzione della Metrologia (2019)}
Dal 2019, i campioni fisici materiali (come il cilindro di platino-iridio per il chilogrammo) sono stati definitivamente abbandonati. Le unità di misura sono ora definite esclusivamente in base a \textbf{costanti fondamentali della fisica} (es. costante di Planck $h$, velocità della luce $c$, carica elementare $e$), garantendo invarianza e riproducibilità universale indipendentemente dal luogo o dallo strumento.
\end{osservazione}

\section{Analisi Dimensionale}
L'analisi dimensionale permette di studiare le dimensioni fisiche di una grandezza e verificare la correttezza delle equazioni. Ogni grandezza $X$ può essere espressa come prodotto di potenze delle grandezze fondamentali:
\[ [X] = [L]^\alpha \cdot [M]^\beta \cdot [T]^\gamma \quad \text{con } \alpha, \beta, \gamma \in \mathbb{Q} \]

\paragraph{Principio di Omogeneità:} In un'equazione fisica $A = B + C$, tutti i termini devono avere le stesse dimensioni ($[A] = [B] = [C]$).
Se $\alpha=\beta=\gamma=0$, la grandezza si dice \textbf{adimensionale} (numero puro).

\begin{esercizio}{Il Pendolo Semplice}
Vogliamo determinare da quali grandezze fisiche dipende il **Periodo di oscillazione ($T$)** di un pendolo semplice.

    \begin{center}
    \begin{tikzpicture}
        % Ceiling
        \draw[thick] (-2,0) -- (2,0);
        \fill[pattern=north east lines] (-2,0) rectangle (2,0.2);
        
        % Vertical dashed line
        \draw[dashed] (0,0) -- (0,-3);
        
        % Pendulum string
        \draw[thick] (0,0) -- (1.5,-2.5) node[midway, right] {$l$};
        
        % Mass
        \draw[thick, fill=gray!30] (1.5,-2.5) circle (0.3cm) node[right=0.4cm] {$m$};
        
        % Arc angle
        \draw[->] (0,-1) arc (270:300:1);
        \node at (0.3,-1.2) {$\theta$};
        
        % Gravity vector
        \draw[->, thick] (-1.5, -1) -- (-1.5, -2) node[right] {$\vec{g}$};
    \end{tikzpicture}
    \captionof{figure}{Pendolo semplice di lunghezza $l$ e massa $m$.}
    \end{center}

Ipotizziamo che il periodo $T$ dipenda dalla massa $m$, dalla lunghezza del filo $l$ e dall'accelerazione di gravità $g$. Scriviamo una relazione di proporzionalità generica:
\[ T = k \cdot m^\alpha \cdot l^\beta \cdot g^\gamma \]
dove $k$ è una costante adimensionale. Analizziamo le dimensioni di ogni termine:
\begin{itemize}
    \item Periodo: $[T] = [T]^1$
    \item Massa: $[m] = [M]$
    \item Lunghezza: $[l] = [L]$
    \item Gravità ($m/s^2$): $[g] = [L][T]^{-2}$
\end{itemize}

Sostituiamo nell'equazione dimensionale:
\[ [M]^0 [L]^0 [T]^1 = [M]^\alpha \cdot [L]^{\beta+\gamma} \cdot [T]^{-2\gamma} \]

Per il principio di omogeneità, impostiamo il sistema:
\[
\begin{cases}
    \alpha = 0 \\
    \beta + \gamma = 0 \\
    -2\gamma = 1
\end{cases}
\implies \alpha = 0, \beta = 1/2, \gamma = -1/2
\]

Otteniamo così la legge del periodo:
\[ \boxed{T = k \sqrt{\frac{l}{g}}} \]
\textbf{Conclusione}: Il periodo è indipendente dalla massa e dipende solo dalla lunghezza del filo e dalla gravità locale.
\end{esercizio}

\section{Multipli e Sottomultipli}
Ecco una tabella completa dei prefissi utilizzati nel Sistema Internazionale per esprimere multipli (lettere maiuscole, tranne $k$, $h$, $da$) e sottomultipli (lettere minuscole).

\begin{table}[h]
    \centering
    \renewcommand{\arraystretch}{1.2}
    \begin{tabular}{llc | llc}
        \toprule
        \multicolumn{3}{c|}{\textbf{Multipli (Grandi)}} & \multicolumn{3}{c}{\textbf{Sottomultipli (Piccoli)}} \\
        \midrule
        \textbf{Prefisso} & \textbf{Simb.} & \textbf{Fattore} & \textbf{Prefisso} & \textbf{Simb.} & \textbf{Fattore} \\
        \midrule
        deca  & $da$ & $10^1$  & deci  & $d$ & $10^{-1}$ \\
        etto  & $h$  & $10^2$  & centi & $c$ & $10^{-2}$ \\
        chilo & $k$  & $10^3$  & milli & $m$ & $10^{-3}$ \\
        Mega  & $M$  & $10^6$  & micro & $\mu$ & $10^{-6}$ \\
        Giga  & $G$  & $10^9$  & nano  & $n$ & $10^{-9}$ \\
        Tera  & $T$  & $10^{12}$ & pico  & $p$ & $10^{-12}$ \\
        Peta  & $P$  & $10^{15}$ & femto & $f$ & $10^{-15}$ \\
        Exa   & $E$  & $10^{18}$ & atto  & $a$ & $10^{-18}$ \\
        \bottomrule
    \end{tabular}
    \caption{Tabella completa dei multipli e sottomultipli SI.}
    \label{tab:multipli}
\end{table}

\section{Recap: Legami cinematici tra posizione, velocità e accelerazione}

Abbiamo definito le relazioni di derivazione, adesso introduciamo quelle di integrazione.
Possiamo schematizzare i legami tra posizione $\vec{x}(t)$, velocità $\vec{v}(t)$ e accelerazione $\vec{a}(t)$ come segue:

\begin{center}
\begin{tikzpicture}[node distance=2.5cm, auto, >=stealth]
    \node (x) {$\vec{x}$};
    \node (v) [right=of x] {$\vec{v}$};
    \node (a) [right=of v] {$\vec{a}$};

    \draw[->] (x) to [bend left=30] node {$v = \frac{d\vec{x}}{dt}$} (v);
    \draw[->] (v) to [bend left=30] node {$a = \frac{d\vec{v}}{dt}$} (a);
    
    \draw[->] (a) to [bend left=30] node {$v = \int a \, dt$} (v);
    \draw[->] (v) to [bend left=30] node {$x = \int v \, dt$} (x);
\end{tikzpicture}
\end{center}

Possiamo suddividere un intervallo di tempo $\Delta t$ in molti intervalli $\delta t = t_{i+1} - t_i$. Inoltre, la velocità media nell'intervallo può essere approssimata come la media aritmetica delle velocità agli estremi:
\[ v_m(t_i, t_{i+1}) = \frac{x(t_{i+1}) - x(t_i)}{t_{i+1} - t_i} \approx \frac{v(t_{i+1}) + v(t_i)}{2} \]

Lo spazio percorso tra un istante $t_0$ e un istante $t_n$ è dato dalla somma degli spostamenti nei singoli intervallini:
\begin{align*}
x_n - x_0 &= x(t_n) - x(t_0) = [x(t_n) - x(t_{n-1})] + \dots + [x(t_1) - x(t_0)] \\
&= v_m(t_{n-1}, t_n)\delta t + \dots + v_m(t_0, t_1)\delta t \\
&\approx \sum_{i=0}^{n-1} \left[ \frac{v(t_{i+1}) + v(t_i)}{2} \right] \delta t
\end{align*}
Geometricamente, notiamo che il termine $\frac{v(t_i) + v(t_{i+1})}{2} \delta t$ corrisponde all'\textbf{Area del trapezio} sotteso dalla curva $v(t)$ nell'intervallo $[t_i, t_{i+1}]$.
Pertanto:
\[ x_n - x_0 \approx \sum \text{Aree trapezi} \]

Passando al limite per $\delta t \to 0$:
\[ x_n - x_0 = \lim_{\delta t \to 0} \sum_{i=0}^{n-1} \left( \frac{v(t_i) + v(t_{i+1})}{2} \right) \delta t = \int_{t_0}^{t_n} v(t) \, dt \]

\begin{figure}[h]
    \centering
    \begin{tikzpicture}[scale=0.8]
        % Axes
        \draw[->] (0,0) -- (6,0) node[right] {$t$};
        \draw[->] (0,0) -- (0,4) node[above] {$v(t)$};
        
        % Curve
        \draw[thick, blue] plot [smooth] coordinates {(0.5,1) (1.5,2.5) (3,2.8) (4.5,3.2) (5.5,3.5)};
        
        % Trapezoids
        \foreach \x/\y in {0.5/1, 1.5/2.5, 3/2.8, 4.5/3.2} {
            \draw[dashed] (\x,0) -- (\x, \y);
            \node[below] at (\x,0) {\tiny $t_{i}$}; % Placeholder indices
        }
        \draw[dashed] (5.5,0) -- (5.5, 3.5);
        
        % Pattern area
        \fill[pattern=north east lines, pattern color=gray!40] (1.5,0) -- (1.5,2.5) -- (3,2.8) -- (3,0) -- cycle;
        \draw (1.5,2.5) -- (3,2.8); % Top segment of trapezoid
        
        \node at (2.25, 1) {Area};
        \node[below] at (0.5,0) {$t_0$};
        \node[below] at (5.5,0) {$t_n$};
    \end{tikzpicture}
    \caption{Interpretazione geometrica dell'integrale come area sottesa al grafico.}
\end{figure}

In sintesi, le relazioni integrali fondamentali sono:

\begin{itemize}
    \item \textbf{Spostamento vettoriale}:
    \[ \Delta \vec{x} = \vec{x}(t_n) - \vec{x}(t_0) = \int_{t_0}^{t_n} \vec{v}(t) \, dt \]
    
    \item \textbf{Variazione di velocità vettoriale}:
    \[ \Delta \vec{v} = \vec{v}(t_n) - \vec{v}(t_0) = \int_{t_0}^{t_n} \vec{a}(t) \, dt \]
    
    \item \textbf{Spazio percorso (Ascissa Curvilinea)}:
    \[ \Delta S = S(t_n) - S(t_0) = \int_{t_0}^{t_n} |\vec{v}(t)| \, dt \]
    \textit{Nota: qui si integra il modulo della velocità (la velocità scalare istantanea).}
    
    \item \textbf{Variazione di velocità scalare}:
    \[ \Delta v = v(t_n) - v(t_0) = \int_{t_0}^{t_n} |\vec{a}(t)| \, dt \]
    \textit{Nota: analogamente, si considera il contributo del modulo dell'accelerazione.}
\end{itemize}

\section{Moti Rettilinei Fondamentali}

Prima di analizzare moti più complessi, definiamo i due pilastri della cinematica unidimensionale.

\subsection{Moto Rettilineo Uniforme (MRU)}
È il moto di un punto che percorre una traiettoria rettilinea con velocità $\vec{v}$ costante in modulo, direzione e verso.
\[ a = 0 \implies v(t) = \text{costante} \]
La sua legge oraria è rappresentata da una retta:
\[ x(t) = x_0 + v(t - t_0) \]

\subsection{Moto Rettilineo Uniformemente Accelerato (MRUA)}
È il moto di un punto che si muove lungo una retta con accelerazione $\vec{a}$ costante.
Le leggi del moto sono:
\begin{align}
    v(t) &= v_0 + a(t - t_0) \\
    x(t) &= x_0 + v_0(t - t_0) + \frac{1}{2}a(t - t_0)^2
\end{align}

\begin{osservazione}{Derivazione Formula di Torricelli}
La formula di Torricelli mette in relazione velocità e spazio percorso eliminando la dipendenza esplicita dal tempo. Si ottiene isolando l'intervallo di tempo $\Delta t$ dalla legge della velocità e sostituendolo nella legge della posizione.
Ponendo $\Delta t = (t - t_0)$ e $\Delta x = (x - x_0)$:
\begin{enumerate}
    \item Dalla (1): $\Delta t = \frac{v - v_0}{a}$
    \item Sostituiamo in (2): $\Delta x = v_0 \left( \frac{v - v_0}{a} \right) + \frac{1}{2} a \left( \frac{v - v_0}{a} \right)^2$
\end{enumerate}
Sviluppando i calcoli:
\[ \Delta x = \frac{v_0 v - v_0^2}{a} + \frac{1}{2} a \frac{(v - v_0)^2}{a^2} = \frac{v_0 v - v_0^2}{a} + \frac{v^2 - 2v v_0 + v_0^2}{2a} \]
Moltiplicando tutto per $2a$:
\[ 2a \Delta x = 2v_0 v - 2v_0^2 + v^2 - 2v v_0 + v_0^2 \]
\[ 2a \Delta x = v^2 - v_0^2 \implies \boxed{v^2 = v_0^2 + 2a \Delta x} \]
\end{osservazione}

\section{Altri Moti della Cinematica}
Oltre ai moti rettilinei uniformi e accelerati, la cinematica studia:
\begin{itemize}
    \item Moto vario accelerato
    \item Moto armonico (semplice, smorzato, forzato)
    \item Moto dei gravi (caduta libera e parabolico)
    \item Moti curvilinei (circolare, ellittico)
\end{itemize}

\section{Moto Armonico Semplice (MAS)}

Il moto armonico descrive i sistemi che oscillano attorno a una posizione di equilibrio stabile. A differenza del MRUA, nel moto armonico l'accelerazione \textbf{NON è costante}.

\subsection{Funzioni Periodiche vs Armoniche}
È fondamentale distinguere tra una funzione periodica generica e una funzione armonica.

\begin{figure}[ht]
    \centering
    \begin{minipage}{0.48\textwidth}
        \centering
        \begin{tikzpicture}[scale=0.6]
            \draw[->] (0,0) -- (4.5,0) node[right] {$t$};
            \draw[->] (0,-1.5) -- (0,1.5) node[above] {$f(t)$};
            \draw[thick, blue] (0,1) -- (1,1) -- (1,-1) -- (2,-1) -- (2,1) -- (3,1) -- (3,-1) -- (4,-1) -- (4,1);
            \node[below, font=\small] at (2.25,-1.6) {Funzione Periodica};
        \end{tikzpicture}
    \end{minipage}
    \hfill
    \begin{minipage}{0.48\textwidth}
        \centering
        \begin{tikzpicture}[scale=0.6]
            \draw[->] (0,0) -- (4.5,0) node[right] {$t$};
            \draw[->] (0,-1.5) -- (0,1.5) node[above] {$x(t)$};
            \draw[thick, red, smooth, domain=0:4.2, samples=100] plot (\x, {sin(\x*180/1)});
            \node[below, font=\small] at (2.25,-1.6) {Funzione Armonica};
        \end{tikzpicture}
    \end{minipage}
\end{figure}

\begin{itemize}
    \item \textbf{Funzione Periodica}: Una funzione che si ripete identica dopo un periodo $T$ (es. l'onda quadra a sinistra). Qualunque forma è ammessa purché $f(t) = f(t+T)$.
    \item \textbf{Funzione Armonica}: Una specifica funzione periodica con andamento sinusoidale ($A\sin(\omega t + \phi)$). Matematicamente, è legata a una forza di richiamo proporzionale allo spostamento.
\end{itemize}

\begin{definizione}{Moto Armonico Semplice (MAS)}
Il moto armonico è un moto la cui legge oraria è una \textbf{funzione armonica}. Esso deve soddisfare l'equazione differenziale specifica:
\begin{equation}
    \ddot{x}(t) = -\omega^2 x(t)
\end{equation}
Tale equazione indica che in ogni istante l'accelerazione è proporzionale allo spostamento e rivolta verso la posizione di equilibrio.
\end{definizione}

\begin{osservazione}{Determinismo e Teorema di Cauchy}
Per determinare completamente l'evoluzione di un moto armonico, non basta l'equazione differenziale; è necessario conoscere le \textbf{condizioni iniziali} ($x_0, v_0$). Secondo il \textbf{Teorema di Cauchy} (esistenza e unicità), date tali condizioni, la soluzione è univocamente determinata.
\end{osservazione}

\subsection{Legge Oraria}
La soluzione dell'equazione differenziale è del tipo:
\begin{equation}
    x(t) = A \sin(\omega t + \phi)
\end{equation}
Dove:
\begin{itemize}
    \item $[A]$: \textbf{Ampiezza} dell'oscillazione (massimo spostamento dalla posizione di equilibrio).
    \item $[\omega]$: \textbf{Pulsazione} del moto (rad/s).
    \item $[\phi]$: \textbf{Fase iniziale} (o sfasamento rispetto a $\omega t$).
    \item $x(t)$: \textbf{Elongazione} al tempo $t$.
\end{itemize}

Poiché il moto è periodico, definiamo:
\begin{itemize}
    \item \textbf{Periodo} $[T]$: Il tempo impiegato per compiere un'oscillazione completa $[s]$.
    \[ T = \frac{2\pi}{\omega} \]
    \item \textbf{Frequenza} $[f]$: Numero di oscillazioni compiute nell'unità di tempo $[Hz]$ ($s^{-1}$).
    \[ f = \frac{1}{T} = \frac{\omega}{2\pi} \]
\end{itemize}

\subsection{Velocità e Accelerazione}
Derivando la legge oraria otteniamo:
\begin{align}
    v(t) &= \frac{dx}{dt} = A\omega \cos(\omega t + \phi) \\
    a(t) &= \frac{dv}{dt} = -A\omega^2 \sin(\omega t + \phi)
\end{align}
Osserviamo che:
\[ a(t) = -\omega^2 [A \sin(\omega t + \phi)] = -\omega^2 x(t) \]
Ritroviamo così la caratteristica fondamentale: l'accelerazione è proporzionale a $-x(t)$.

\subsection{Forme Alternative}
Le leggi orarie si possono trovare scritte anche in modi equivalenti, sfruttando le proprietà trigonometriche (come $\sin(\theta) = \cos(90^\circ - \theta)$):
\begin{itemize}
    \item $x(t) = A \cos(\omega t + \psi)$ (dove $\psi \neq \phi$)
    \item Combinazione lineare: $x(t) = \alpha \sin(\omega t) + \beta \cos(\omega t)$
\end{itemize}
Esistono correlazioni tra i coefficienti per passare da una forma all'altra:
\[ A = \sqrt{\alpha^2 + \beta^2}, \quad \tan(\phi) = \frac{\beta}{\alpha} \]

\paragraph{Nota Geometrica:}
Si può vedere il moto armonico semplice come la \textbf{proiezione di un moto circolare uniforme} su un diametro (una retta).

\section{Moto Smorzato Esponenzialmente}
È un moto la cui velocità decresce esponenzialmente nel tempo.
Le leggi orarie sono:
\[ x(t) = A \left(1 - e^{-\frac{t}{\tau}}\right), \quad v(t) = \frac{A}{\tau} e^{-\frac{t}{\tau}}, \quad a(t) = -\frac{A}{\tau^2} e^{-\frac{t}{\tau}} \]
Dove $\tau$ è la costante di tempo e $A$ è il valore asintotico della posizione (la quota a cui tende).

\begin{figure}[h]
    \centering
    \begin{tikzpicture}[scale=0.9]
        \draw[->] (0,0) -- (5,0) node[right] {$t$};
        \draw[->] (0,0) -- (0,3) node[above] {$x(t)$};
        \draw[dashed] (0,2.5) -- (5,2.5) node[right] {$A$};
        \draw[thick, blue] plot[domain=0:4.8, samples=100] (\x, {2.5*(1 - exp(-\x))});
    \end{tikzpicture}
    \caption{Andamento della posizione nel moto smorzato: tende asintoticamente ad $A$.}
\end{figure}

\section{Moto dei Gravi}
È un moto rettilineo uniformemente accelerato con accelerazione costante $\vec{a} = -\vec{g}$ (diretta verso il basso), dove $g \approx 9.81 \, m/s^2$.

\subsection{Caduta da $y_0 \neq 0$}
Consideriamo un corpo lasciato cadere da un'altezza $y_0$ con velocità iniziale nulla ($v_0 = 0$).

Le equazioni sono:
\begin{align*}
    a(t) &= -g \\
    v(t) &= -g(t - t_0) \\
    y(t) &= y(t_0) - \frac{1}{2}g(t - t_0)^2
\end{align*}

Possiamo calcolare il \textbf{tempo di caduta} $t_f$ imponendo $y(t_f) = 0$:
\[ t_f = t_0 + \sqrt{\frac{2y_0}{g}} \]

La \textbf{velocità finale} (di impatto) sarà:
\[ v_f = - \sqrt{2y_0 g} \]
(negativa perché il vettore velocità è diretto verso il basso, concorde a $\vec{g}$).

\begin{center}
    \begin{tikzpicture}[scale=0.8]
        \draw[->] (0,0) -- (3,0) node[right] {$t$};
        \draw[->] (0,0) -- (0,3) node[above] {$y(t)$};
        \draw[thick, red] plot[domain=0:2.2, samples=50] (\x, {2.5 - 0.5*\x*\x});
        \node[left] at (0,2.5) {$y_0$};
        \draw[dashed] (2.23,0) -- (2.23, 2.5); % impact time
        \node[below] at (2.23,0) {$t_f$};
        \draw[->, gray] (2.7, 2.5) -- (2.7, 1.5) node[midway, right] {$\vec{g}$};
    \end{tikzpicture}
\end{center}

\subsection{Lancio verso l'alto da $y_0 = 0$}
Consideriamo un corpo lanciato verso l'alto con velocità iniziale $v(t_0) > 0$.

Le equazioni sono:
\begin{align*}
    v(t) &= v(t_0) - g(t - t_0) \\
    y(t) &= v(t_0)(t - t_0) - \frac{1}{2}g(t - t_0)^2
\end{align*}

\textbf{Analisi del moto:}
\begin{itemize}
    \item $v(t) > 0$ nella fase di ascesa.
    \item $v(t) < 0$ nella fase di discesa.
    \item Quando $v(t) = 0$, il corpo raggiunge la \textbf{quota massima} $y_m$. Questo avviene all'istante:
    \[ t_{\text{apice}} = \frac{v(t_0)}{g} \]
\end{itemize}
Sostituendo il tempo nell'equazione della posizione, otteniamo la quota massima:
\[ y_m = \frac{v(t_0)^2}{2g} \]

Il tempo totale di volo $t_f$ (per tornare a terra) è il doppio del tempo di salita:
\[ t_f = t_0 + \frac{2v(t_0)}{g} \]

\begin{center}
    \begin{tikzpicture}[scale=0.8]
        \draw[->] (0,0) -- (4,0) node[right] {$t$};
        \draw[->] (0,0) -- (0,3) node[above] {$y(t)$};
        % Parabola starting at 0, peak at x=2, y=2, landing at x=4
        \draw[thick, green!60!black] plot[domain=0:4, samples=50] (\x, {2*\x - 0.5*\x*\x});
        
        \draw[dashed] (0,2) -- (2,2);
        \node[left] at (0,2) {$y_m$};
        
        \draw[dashed] (2,0) -- (2,2);
        \node[below] at (2,0) {$\frac{v_0}{g}$};
        
        \node[below] at (4,0) {$t_f$};
    \end{tikzpicture}
\end{center}

\section{Moto Parabolico}
Il moto parabolico è un moto in due dimensioni, composto dalla sovrapposizione di due moti indipendenti:
\begin{itemize}
    \item \textbf{Asse x}: Moto rettilineo uniforme ($a_x = 0$) $\Rightarrow v_x = \text{costante}$.
    \item \textbf{Asse y}: Moto uniformemente accelerato ($a_y = -g = -9.81 \, m/s^2$).
\end{itemize}

\subsection{Lancio con $\theta = 0^\circ$ (Lancio Orizzontale)}
Consideriamo un corpo lanciato orizzontalmente da un'altezza $y_0$ con velocità $v_0$.
In questo caso, l'angolo iniziale è nullo, quindi:
\[ v_{0x} = v_0 \cos(0^\circ) = v_0, \quad v_{0y} = v_0 \sin(0^\circ) = 0 \]

\textbf{Equazioni della velocità:}
\[
\begin{cases}
    v_x(t) = v_0 \\
    v_y(t) = -gt
\end{cases}
\]

\textbf{Leggi orarie:}
\[
\begin{cases}
    x(t) = v_0 t \\
    y(t) = y_0 - \frac{1}{2} g t^2
\end{cases}
\]

Il corpo atterra quando $y(t_f) = 0$.
\begin{itemize}
    \item \textbf{Tempo di caduta}: $t_f = \sqrt{\frac{2y_0}{g}}$
    \item \textbf{Gittata} ($x_f$): Sostituendo $t_f$ in $x(t)$:
    \[ x_f = v_0 \sqrt{\frac{2y_0}{g}} \]
    \item \textbf{Angolo finale d'impatto} $\theta_f$:
    \[ \tan \theta_f = \frac{v_y(t_f)}{v_x(t_f)} = \frac{-g t_f}{v_0} = \frac{-\sqrt{2y_0 g}}{v_0} \]
    \item \textbf{Velocità d'impatto verticale}: $v_{yf} = - \sqrt{2y_0 g}$
\end{itemize}

\begin{center}
    \begin{tikzpicture}[scale=0.8]
        \draw[->] (0,0) -- (4,0) node[right] {$x(t)$};
        \draw[->] (0,0) -- (0,3) node[above] {$y(t)$};
        
        \node[left] at (0,2.5) {$y_0$};
        \draw[dashed] (0,2.5) -- (4,2.5);
        \draw[->, thick] (0,2.5) -- (1,2.5) node[above] {$\vec{v}_0$};
        
        % Parabola starting at (0, 2.5) going down
        \draw[thick, blue] plot[domain=0:2.25, samples=50] (\x, {2.5 - 0.49*\x*\x});
        
        \node[below] at (2.25,0) {$x_f$};
    \end{tikzpicture}
\end{center}

\subsection{Lancio con \texorpdfstring{$\theta \neq 0^\circ$}{theta != 0} (Lancio Obliquo)}
Il corpo viene lanciato con una velocità iniziale $v_0$ che forma un angolo $\theta$ con l'orizzontale.
\[ v_{0x} = v_0 \cos\theta, \quad v_{0y} = v_0 \sin\theta \]

\textbf{Equazioni della velocità:}
\[
\begin{cases}
    v_x(t) = v_{0x} \quad (\text{costante}) \\
    v_y(t) = v_{0y} - gt
\end{cases}
\]

\textbf{Leggi orarie:}
\[
\begin{cases}
    x(t) = x_0 + v_{0x} t \\
    y(t) = y_0 + v_{0y} t - \frac{1}{2} g t^2
\end{cases}
\]

\textbf{Proprietà del moto:}
\begin{itemize}
    \item \textbf{Direzione istantanea}: $\tan \theta(t) = \frac{v_y(t)}{v_x(t)}$.
    \item \textbf{Quota Massima} ($y_m$): Si raggiunge quando $v_y = 0$, ovvero a $t = \frac{v_{0y}}{g}$.
    \[ y_m = y_0 + \frac{v_{0y}^2}{2g} = y_0 + \frac{(v_0 \sin\theta)^2}{2g} \]
    \item \textbf{Tempo di volo} (atterraggio alla stessa quota $y_0$): $t_f = \frac{2v_{0y}}{g}$.
    \item \textbf{Gittata} ($x_m$): Distanza orizzontale percorsa.
    \[ x_m = v_{0x} t_f = \frac{v_0^2 \sin(2\theta)}{g} \]
    Nota: La gittata è massima per $\theta = 45^\circ$ ($\sin 90^\circ = 1$) ed è uguale per angoli complementari ($\theta$ e $90^\circ - \theta$).
\end{itemize}

\begin{center}
    \begin{tikzpicture}[scale=0.7]
        \draw[->] (0,0) -- (4.5,0) node[right] {$x$};
        \draw[->] (0,0) -- (0,3) node[above] {$y$};
        
        % Parabola
        \draw[thick, green!60!black] plot[domain=0:4, samples=50] (\x, {-\x*(\x-4)/1.5});
        
        % Vector v0
        \draw[->, thick] (0,0) -- (1, 1.5) node[midway, left] {$\vec{v}_0$};
        \draw (0.5,0) arc (0:56:0.5);
        \node at (0.7, 0.3) {$\theta$};
        
        \draw[dashed] (2, 2.66) -- (0, 2.66) node[left] {$y_m$};
        \draw[dashed] (2, 2.66) -- (2, 0) node[below] {$x_m/2$};
        \node[below] at (4,0) {$x_m$};
    \end{tikzpicture}
\end{center}

\subsubsection{Equazione della Traiettoria}
Eliminando il tempo $t = x/v_{0x}$ dalle leggi orarie (ponendo $x_0=y_0=0$), otteniamo la relazione tra $y$ e $x$:
\[ y(x) = x \tan\theta - \frac{g}{2v_0^2 \cos^2\theta} x^2 \]
Questa è l'equazione di una parabola con concavità verso il basso.

\subsubsection{Geometria della Sicurezza}
\begin{itemize}
    \item \textbf{Luogo dei vertici}: Le coordinate dei vertici delle parabole al variare di $\theta$ descrivono un'ellisse:
    \[ x_v^2 + 4y_v^2 = \frac{2v_0^2}{g} y_v \]
    \item \textbf{Parabola di Sicurezza}: L'inviluppo di tutte le possibili traiettorie (al variare di $\theta$ con $v_0$ fissato) definisce la regione di spazio raggiungibile dal proiettile. Questa regione è delimitata dalla \textbf{parabola di sicurezza}.
\end{itemize}

\section{Moto Circolare Uniforme (MCU)}
Il moto circolare è caratterizzato da alcune nuove grandezze vettoriali:
\begin{itemize}
    \item $\vec{\omega}$: \textbf{Velocità Angolare}.
    \item $\vec{a}_n$: \textbf{Accelerazione Centripeta} (o Normale).
    \item $\vec{a}_t$: \textbf{Accelerazione Tangenziale} (presente solo se il moto non è uniforme).
    \item $\vec{\alpha}$: \textbf{Accelerazione Angolare} (presente solo se il moto non è uniforme).
\end{itemize}

\begin{definizione}{Componenti dell'Accelerazione}
\begin{itemize}
    \item \textbf{Accelerazione Centripeta / Normale} ($\vec{a}_n$): Esprime la \textbf{variazione di direzione} della velocità. È sempre presente in un moto curvilineo e punta verso il centro di curvatura.
    \item \textbf{Accelerazione Tangenziale} ($\vec{a}_t$): Esprime la \textbf{variazione di modulo} della velocità. È parallela alla velocità istantanea.
    \item \textbf{Accelerazione Angolare} ($\vec{\alpha}$): Derivata della velocità angolare rispetto al tempo ($\vec{\alpha} = d\vec{\omega}/dt$).
\end{itemize}
\end{definizione}

\begin{figure}[h]
    \centering
    \begin{tikzpicture}[scale=1]
        % Circle
        \draw[thick] (0,0) circle (2cm);
        \draw[dashed] (0,0) -- (1.414, 1.414) node[midway, above left] {$r$};
        \filldraw (1.414, 1.414) circle (2pt) node[above right] {$P$};
        
        % Vectors
        \draw[->, thick, blue] (1.414, 1.414) -- (-0.5, 3.3) node[above] {$\vec{v}$ (tangente)};
        \draw[->, thick, red] (1.414, 1.414) -- (0.5, 0.5) node[midway, right] {$\vec{a}_n \perp \vec{v}$};
        
        % Tangential acceleration (dashed imaginary for MCU)
        \draw[->, thick, orange, dashed] (1.414, 1.414) -- (0.4, 2.4) node[left] {$\vec{a}_t$};
        
        \node at (0,-2.5) {Scomposizione dell'accelerazione};
        
        % Angular velocity 3D illusion
        \begin{scope}[xshift=5cm]
            \draw (0,0) ellipse (1.5cm and 0.5cm);
            \draw[->, thick, purple] (0,0) -- (0, 2) node[right] {$\vec{\omega}$};
            \draw[->] (0.5, 0.2) arc (0:330:0.3 and 0.1);
            \node at (0,-1) {Direzione di $\vec{\omega}$};
        \end{scope}
    \end{tikzpicture}
    \caption{Cinematica del moto circolare.}
\end{figure}

\subsection{Relazioni Fondamentali (MCU)}
Nel moto circolare uniforme ($|v|=$ cost, $\alpha=0$), le grandezze sono così relazionate:
\begin{itemize}
    \item Velocità tangenziale: $|v| = \omega r$
    \item Periodo: $T = \frac{2\pi}{\omega} = \frac{2\pi r}{v}$
    \item Frequenza: $f = \frac{1}{T}$
    \item Velocità angolare: $|\omega| = \frac{2\pi}{T} = 2\pi f$
    \item Modulo accelerazione centripeta: $a_c = \omega^2 r = \frac{v^2}{r}$
\end{itemize}

\clearpage
\section{Approfondimento: Sistemi di Riferimento e Algebra dei Versori}
Prima di analizzare nel dettaglio la dinamica del moto circolare, è fondamentale definire formalmente i sistemi di coordinate e le proprietà dei versori (vettori unitari) associati.

\subsection{Versori e Derivazione Temporale}
Un sistema di riferimento è definito da un'origine $O$ e da una base di versori. La caratteristica cruciale per la cinematica è se questi versori sono \textbf{fissi} o \textbf{mobili} nel tempo.

\subsubsection{Coordinate Cartesiane $(x, y)$}
La base è costituita dai versori $\{\hat{x}, \hat{y}\}$:
\begin{itemize}
    \item Hanno modulo unitario ($|\hat{x}| = |\hat{y}| = 1$) e sono ortogonali.
    \item Hanno \textbf{direzione e verso fissi} nello spazio.
\end{itemize}
Pertanto, la loro derivata rispetto al tempo è nulla:
\[ \frac{d\hat{x}}{dt} = \vec{0}, \quad \frac{d\hat{y}}{dt} = \vec{0} \]
Il vettore posizione si scrive come: $\vec{r}(t) = x(t)\hat{x} + y(t)\hat{y}$.

\begin{center}
\begin{tikzpicture}[scale=1.5]
    % Assi
    \draw[->] (-2.5,0) -- (2.5,0) node[right] {$x$};
    \draw[->] (0,-2.5) -- (0,2.5) node[above] {$y$};
    
    % Traiettoria Circolare (MCU)
    \draw[thick, gray, dashed] (0,0) circle (2cm);
    
    % Versori fissi origin
    \draw[->, thick, red] (0,0) -- (1,0) node[below] {$\hat{x}$};
    \draw[->, thick, red] (0,0) -- (0,1) node[left] {$\hat{y}$};
    
    % Punto P a 45 gradi su r=2 -> P(1.414, 1.414)
    \filldraw (1.414, 1.414) circle (1.5pt) node[above right] {$P(x,y)$};
    \draw[dashed] (1.414,0) node[below] {$x$} -- (1.414,1.414) -- (0,1.414) node[left] {$y$};
    
    % Vettore posizione
    \draw[->, thick, blue] (0,0) -- (1.414, 1.414) node[midway, above left] {$\vec{r}(t)$};
\end{tikzpicture}
\end{center}

\subsubsection{Coordinate Polari Piane $(\rho, \theta)$}
La base è costituita dai versori $\{\hat{\rho}, \hat{\theta}\}$:
\begin{itemize}
    \item $\hat{\rho}$: Versore \textbf{radiale}, diretto dall'origine verso il punto $P$.
    \item $\hat{\theta}$: Versore \textbf{trasverso} (o azimutale), perpendicolare a $\hat{\rho}$ nel verso di crescita di $\theta$.
\end{itemize}

\begin{center}
\begin{tikzpicture}[scale=1.5]
    % Assi riferimento
    \draw[->, gray] (-2.5,0) -- (2.5,0) node[right] {$x$};
    \draw[->, gray] (0,-2.5) -- (0,2.5) node[above] {$y$};
    
    % Traiettoria Circolare (MCU)
    \draw[thick, gray, dashed] (0,0) circle (2cm);
    
    % Punto P
    \filldraw (1.414, 1.414) circle (1.5pt) node[above left] {$P(\rho,\theta)$};
    
    % Vettore rho (coincide con il raggio)
    \draw[->, thick, blue] (0,0) -- (1.414, 1.414) node[midway, below right] {$\vec{\rho}$};
    \draw[dashed, gray] (1.414,1.414) -- (2.121, 2.121); % Prolungamento
    
    % Angolo theta
    \draw (0.8,0) arc (0:45:0.8);
    \node at (1, 0.4) {$\theta$};
    
    % Versori mobili in P
    \draw[->, thick, red] (1.414, 1.414) -- (2.121, 2.121) node[right] {$\hat{\rho}$};
    \draw[->, thick, red] (1.414, 1.414) -- (0.707, 2.121) node[above] {$\hat{\theta}$};
    
    % Angolo retto tra i versori
    \draw (1.556, 1.556) -- (1.414, 1.697) -- (1.273, 1.556);
\end{tikzpicture}
\end{center}

Possiamo esprimerli in funzione della base cartesiana (fissa):
\[ \hat{\rho} = \cos\theta \, \hat{x} + \sin\theta \, \hat{y} \]
\[ \hat{\theta} = -\sin\theta \, \hat{x} + \cos\theta \, \hat{y} \]

\begin{teorema}{Formule di Poisson nel Piano}
Le derivate temporali dei versori polari mobili $\{\hat{\rho}, \hat{\theta}\}$ si esprimono in funzione della velocità angolare $\dot{\theta}$:
\begin{equation}
    \frac{d\hat{\rho}}{dt} = \dot{\theta} \hat{\theta}, \quad \frac{d\hat{\theta}}{dt} = -\dot{\theta} \hat{\rho}
\end{equation}
\end{teorema}
Queste relazioni (Formule di Poisson nel piano) sono fondamentali per derivare velocità e accelerazione in coordinate polari.

\subsubsection{Coordinate Intrinseche $(s)$}
La base è locale alla traiettoria, definita dai versori $\{\hat{\tau}, \hat{n}\}$:
\begin{itemize}
    \item $\hat{\tau}$: Versore \textbf{tangente}, parallelo alla velocità.
    \item $\hat{n}$: Versore \textbf{normale}, diretto verso il centro di curvatura locale.
\end{itemize}

\begin{center}
\begin{tikzpicture}[scale=1.5]
    % Traiettoria (Circolare per MCU)
    \draw[thick, gray, dashed] (0,0) circle (2cm);
    
    % Centro di curvatura (origine)
    \filldraw (0, 0) circle (1pt) node[below] {$C$};
    
    % Punto P a 45 gradi su r=2
    \filldraw (1.414, 1.414) circle (1.5pt) node[above right] {$P$};
    
    % Raggio di curvatura
    \draw[dashed] (0, 0) -- (1.414, 1.414) node[midway, left] {$R$};
    
    % Versori mobili in P
    % n punta verso il centro C (origine), theta/tau è radente
    \draw[->, thick, red] (1.414, 1.414) -- (0.707, 2.121) node[above left] {$\hat{\tau}$};
    \draw[->, thick, red] (1.414, 1.414) -- (0.707, 0.707) node[below right] {$\hat{n}$};
    
    % Vettore velocità (parallelo a tau)
    \draw[->, thick, blue] (1.414, 1.414) -- (0, 2.828) node[left] {$\vec{v}$};
\end{tikzpicture}
\end{center}

Nel caso del moto circolare, valgono le seguenti identificazioni geometriche con i versori polari:
\[ \hat{\tau} \equiv \hat{\theta}, \quad \hat{n} \equiv -\hat{\rho} \]
La relazione fondamentale è $\frac{d\hat{\tau}}{dt} = \frac{v}{R} \hat{n}$ (dove $R$ è il raggio di curvatura).

\section{Analisi del Moto Circolare nei vari Sistemi}
Applichiamo ora l'algebra dei versori appena definita al caso specifico del moto circolare di raggio $R$ costante.

\subsection{Caso 1: Coordinate Cartesiane}
Descriviamo il punto $P$ tramite la proiezione sugli assi $x$ e $y$.
Sia $\theta(t) = \omega t + \phi_0$ la legge oraria angolare.
\[
\begin{cases}
    x(t) = R \cos(\omega t + \phi_0) \\
    y(t) = R \sin(\omega t + \phi_0)
\end{cases}
\]

\textbf{Velocità}: Deriviamo le componenti rispetto al tempo ($\dot{x}, \dot{y}$):
\[
\begin{cases}
    v_x = \dot{x} = -R\omega \sin(\omega t + \phi_0) = -\omega y \\
    v_y = \dot{y} = R\omega \cos(\omega t + \phi_0) = \omega x
\end{cases}
\]

\textbf{Accelerazione}: Deriviamo nuovamente ($\ddot{x}, \ddot{y}$):
\[
\begin{cases}
    a_x = \ddot{x} = -R\omega^2 \cos(\omega t + \phi_0) = -\omega^2 x \\
    a_y = \ddot{y} = -R\omega^2 \sin(\omega t + \phi_0) = -\omega^2 y
\end{cases}
\]
Notiamo che $\vec{a} = a_x \hat{x} + a_y \hat{y} = -\omega^2 (x\hat{x} + y\hat{y}) = -\omega^2 \vec{r}$.
Questo conferma che l'accelerazione è \textbf{centripeta} (opposta al raggio vettore) e proporzionale al quadrato della velocità angolare. Inoltre, le proiezioni sugli assi sono \textbf{moti armonici} di pulsazione $\omega$.

\subsection{Caso 2: Coordinate Polari}
Il vettore posizione è definito semplicemente da:
\[ \vec{r}(t) = \rho(t) \hat{\rho} \]
Nel moto circolare, il raggio è costante $\rho(t) = R$, quindi $\dot{\rho} = 0, \ddot{\rho} = 0$.
L'angolo varia come $\theta(t)$, con velocità angolare $\dot{\theta} = \omega$ e accelerazione angolare $\ddot{\theta} = \alpha$.

\textbf{Velocità}: Deriviamo $\vec{r}$ usando la regola del prodotto:
\[ \vec{v} = \frac{d}{dt}(R \hat{\rho}) = \dot{R}\hat{\rho} + R\frac{d\hat{\rho}}{dt} \]
Poiché $\dot{R}=0$ e $\frac{d\hat{\rho}}{dt} = \dot{\theta}\hat{\theta} = \omega \hat{\theta}$:
\[ \vec{v} = R\omega \hat{\theta} \]
(La velocità è puramente tangenziale, diretta lungo $\hat{\theta}$).

\textbf{Accelerazione}: Deriviamo $\vec{v}$:
\[ \vec{a} = \frac{d}{dt}(R\omega \hat{\theta}) = R\dot{\omega}\hat{\theta} + R\omega\frac{d\hat{\theta}}{dt} \]
Ricordando che $\dot{\omega} = \alpha$ e $\frac{d\hat{\theta}}{dt} = -\omega \hat{\rho}$:
\[ \vec{a} = R\alpha \hat{\theta} - R\omega^2 \hat{\rho} \]
Abbiamo ottenuto le due componenti vettoriali in modo formale:
\begin{itemize}
    \item \textbf{Componente Tangenziale} ($\vec{a}_t = R\alpha \hat{\theta}$): Dipende da $\alpha$. Se il moto è uniforme ($\alpha=0$), questa componente è nulla.
    \item \textbf{Componente Centripeta} ($\vec{a}_n = -R\omega^2 \hat{\rho}$): Sempre presente, diretta lungo $-\hat{\rho}$ (verso il centro).
\end{itemize}

\subsection{Caso 3: Coordinate Intrinseche}
In questo sistema, i versori sono definiti localmente sulla curva. Per il moto circolare valgono le identità:
\[ \hat{\tau} \equiv \hat{\theta}, \quad \hat{n} \equiv -\hat{\rho} \]

Sappiamo che:
\[ \vec{v} = v \hat{\tau} \]
Deriviamo rispetto al tempo:
\[ \vec{a} = \frac{dv}{dt}\hat{\tau} + v\frac{d\hat{\tau}}{dt} \]
Usando la relazione geometrica $\frac{d\hat{\tau}}{dt} = \frac{v}{R}\hat{n}$:
\[ \vec{a} = \frac{dv}{dt}\hat{\tau} + \frac{v^2}{R}\hat{n} \]

Sostituendo i versori intrinseci con quelli polari:
\[ \vec{a} = (R\alpha)\hat{\theta} + \frac{(R\omega)^2}{R}(-\hat{\rho}) = R\alpha \hat{\theta} - R\omega^2 \hat{\rho} \]

\clearpage
\section{Moti Curvilinei Generali in Coordinate Polari}
Estendiamo l'analisi al caso più generale di un \textbf{moto curvilineo}, dove il raggio di curvatura non è necessariamente costante (quindi $\rho = \rho(t)$ varia nel tempo). Vogliamo derivare l'espressione completa di velocità e accelerazione in coordinate polari piane $(\rho, \theta)$.

\subsection{Posizione e Velocità}
Il vettore posizione è:
\[ \vec{r}(t) = \rho(t) \hat{\rho}(t) \]
Derivandolo rispetto al tempo per ottenere la velocità, dobbiamo applicare la regola del prodotto (poiché sia il modulo $\rho$ che il versore $\hat{\rho}$ variano):
\[ \vec{v} = \frac{d\vec{r}}{dt} = \frac{d\rho}{dt}\hat{\rho} + \rho\frac{d\hat{\rho}}{dt} \]
Ricordando la formula di Poisson $\frac{d\hat{\rho}}{dt} = \dot{\theta}\hat{\theta}$:
\[ \vec{v} = \dot{\rho}\hat{\rho} + \rho\dot{\theta}\hat{\theta} \]
Possiamo identificare i due termini (anche quelli nascosti dal foglio verde nella figura originale):
\begin{itemize}
    \item $v_\rho = \dot{\rho}$: \textbf{Velocità Radiale} (dovuta alla variazione della distanza dall'origine).
    \item $v_\theta = \rho\dot{\theta}$: \textbf{Velocità Trasversa/Angolare} (dovuta alla rotazione del raggio vettore).
\end{itemize}

\subsection{Accelerazione: Derivazione Completa}
Per ottenere l'accelerazione, deriviamo nuovamente la velocità:
\[ \vec{a} = \frac{d}{dt}(\dot{\rho}\hat{\rho} + \rho\dot{\theta}\hat{\theta}) \]
Applichiamo la derivata ai due termini separatamente:
\begin{enumerate}
    \item Primo termine ($\dot{\rho}\hat{\rho}$):
    \[ \frac{d}{dt}(\dot{\rho}\hat{\rho}) = \ddot{\rho}\hat{\rho} + \dot{\rho}\frac{d\hat{\rho}}{dt} = \ddot{\rho}\hat{\rho} + \dot{\rho}\dot{\theta}\hat{\theta} \]
    \item Secondo termine ($\rho\dot{\theta}\hat{\theta}$), regola del prodotto a tre fattori ($fgh \to f'gh + fg'h + fgh'$):
    \[ \frac{d}{dt}(\rho\dot{\theta}\hat{\theta}) = \dot{\rho}\dot{\theta}\hat{\theta} + \rho\ddot{\theta}\hat{\theta} + \rho\dot{\theta}\frac{d\hat{\theta}}{dt} \]
    Ricordando che $\frac{d\hat{\theta}}{dt} = -\dot{\theta}\hat{\rho}$, diventa:
    \[ = \dot{\rho}\dot{\theta}\hat{\theta} + \rho\ddot{\theta}\hat{\theta} - \rho\dot{\theta}^2\hat{\rho} \]
\end{enumerate}

Sommando tutti i contributi e raggruppando per versori:
\[ \vec{a} = (\ddot{\rho}\hat{\rho} + \dot{\rho}\dot{\theta}\hat{\theta}) + (\dot{\rho}\dot{\theta}\hat{\theta} + \rho\ddot{\theta}\hat{\theta} - \rho\dot{\theta}^2\hat{\rho}) \]

\begin{teorema}{Accelerazione in Coordinate Polari}
L'espressione completa dell'accelerazione per un moto curvilineo generale $\vec{r} = \rho \hat{\rho}$ è:
\begin{equation}
    \vec{a} = (\ddot{\rho} - \rho \dot{\theta}^2)\hat{\rho} + (2\dot{\rho}\dot{\theta} + \rho \ddot{\theta}) \hat{\theta}
\end{equation}
Identifichiamo i 4 contributi fisici:
\begin{itemize}
    \item $\ddot{\rho}$: \textbf{Accelerazione Radiale}.
    \item $-\rho \dot{\theta}^2$: \textbf{Accelerazione Centripeta}.
    \item $\rho \ddot{\theta}$: \textbf{Accelerazione Tangenziale}.
    \item $2\dot{\rho}\dot{\theta}$: \textbf{Accelerazione di Coriolis}.
\end{itemize}
\end{teorema}

\chapter{Dinamica del Punto Materiale}
La dinamica studia il movimento dei corpi in relazione alle cause che lo modificano, ovvero alle \textbf{FORZE}.
Nella dinamica le grandezze fisiche fondamentali che entrano in gioco sono:
\begin{enumerate}
    \item Lunghezza $[L]$
    \item Tempo $[T]$
    \item Massa $[M]$
\end{enumerate}

\section{La Forza (Definizione Generale)}
In generale, la \textbf{Forza} è definita come un'\textbf{INTERAZIONE} fra corpi.
Tutti i tipi di forze condividono alcune proprietà fondamentali:
\begin{enumerate}
    \item \textbf{Grandezze Vettoriali}: sono descritte da un modulo, una direzione e un verso.
    \item \textbf{Principio di Sovrapposizione}: vige l'indipendenza delle azioni simultanee. Se più forze agiscono su un corpo, la forza totale è la somma vettoriale delle singole forze: $\vec{F}_{tot} = \sum \vec{F}_i$.
    \item \textbf{Dipendenza dalla Distanza}: tendenzialmente il modulo della forza $|\vec{F}|$ diminuisce all'aumentare della distanza $d$ tra i corpi (spesso $|\vec{F}| \propto d^{-2}$).
    \item \textbf{Indipendenza dal Sistema di Riferimento}: la forza $\vec{F}$ è invariante per trasformazioni di Galileo (non cambia al variare del riferimento inerziale), così come non cambiano la distanza, il modulo della velocità relativa $|\vec{v}_{rel}|$, la massa e la carica.
\end{enumerate}

\section{Tipologie di Forze}
Esistono diversi tipi di forze, classificabili in due grandi categorie:

\subsection{Forze A Distanza (Fondamentali)}
Agiscono senza contatto diretto tra i corpi:
\begin{itemize}
    \item Forza Gravitazionale.
    \item Forza Elettromagnetica.
    \item Forze Nucleari (Forte e Debole) e Atomiche.
\end{itemize}

\subsection{Forze Di Contatto}
Sono manifestazioni macroscopiche delle forze fondamentali (principalmente elettromagnetiche) che si generano dal contatto tra superfici.

\begin{figure}[ht]
    \centering
    \begin{tikzpicture}[
        grow=right,
        level 1/.style={sibling distance=3.2cm, level distance=4.5cm},
        level 2/.style={sibling distance=1.8cm, level distance=5.5cm},
        edge from parent/.style={draw, thick, -stealth},
        every node/.style={font=\small},
        box/.style={draw, rounded corners, minimum height=0.9cm, minimum width=2.2cm, align=center, font=\small},
        leaf/.style={draw, rounded corners=3pt, fill=gray!8, minimum height=0.8cm, align=center, font=\small}
    ]
    \node[box, fill=mypen!15, font=\bfseries] {Forze di Contatto}
        child { 
            node[box, fill=green!15] {Rigido--Fluido}
            child { node[leaf] {\textbf{Attrito Viscoso}} }
            child { node[leaf] {\textbf{Pressione}} }
        }
        child { 
            node[box, fill=orange!15] {Deformabile--Def.}
            child { node[leaf] {\textbf{Anelastiche}} }
            child { node[leaf] {\textbf{Elastiche}} }
        }
        child { 
            node[box, fill=red!15] {Rigido--Rigido}
            child { node[leaf] {\textbf{Attrito} \\ \scriptsize{Statico, Dinamico}} }
            child { node[leaf] {\textbf{Vincolari} \\ \scriptsize{Tensioni, Reazioni}} }
        };
    \end{tikzpicture}
    \caption{Classificazione delle Forze di Contatto.}
\end{figure}

\section{Misurazione delle Forze}
La misurazione delle forze può avvenire in due modi principali:

\begin{enumerate}[label=\alph*)]
    \item \textbf{Misura Dinamica}:
    È una misura indiretta. Conoscendo la massa $m$ del corpo e potendo misurare la sua accelerazione $\vec{a}$, si sfrutta la seconda legge della dinamica:
    \[ \vec{F} = m \cdot \vec{a} \]
    
    \item \textbf{Misura Statica}:
    È una misura diretta basata sull'equilibrio. Si bilancia una forza nota $F_n$ con la forza incognita $F_i$, oppure si rapportano le deformazioni (lunghezze) causate dalle forze.
    
    Strumenti tipici:
    \begin{itemize}
        \item \textbf{Dinamometro}: Sfrutta la deformazione elastica di una molla ($F = k \Delta x$).
        \item \textbf{Bilancia}: Sfrutta l'equilibrio dei momenti o delle forze peso.
    \end{itemize}
    
    \begin{center}
    \begin{tikzpicture}[scale=0.8]
        % Dinamometro
        \begin{scope}[yshift=0cm]
             \draw[thick] (-0.5, 3) -- (0.5, 3);
             \draw[decoration={aspect=0.3, segment length=2mm, amplitude=2mm, coil}, decorate] (0, 3) -- (0, 1);
             \draw[->, thick] (0, 1) -- (0, 0) node[right] {$\vec{F}$};
             \draw (-0.2, 1) -- (0.2, 1); 
             \node[left] at (-0.5, 2) {\small Dinamometro};
        \end{scope}

        % Bilancia
        \begin{scope}[xshift=6cm]
            \draw[thick] (0, 0) -- (0, 2.5); % Piatto centrale
            \draw[thick] (-2, 2.5) -- (2, 2.5); % Braccio
            \draw (-2, 2.5) -- (-2, 1.5); % Filo sx
            \draw (-2.4, 1.5) -- (-1.6, 1.5); % Piatto sx
            
            \draw (2, 2.5) -- (2, 1.5); % Filo dx
            \draw (1.6, 1.5) -- (2.4, 1.5); % Piatto dx
            
            \filldraw (0, 2.5) circle (2pt); % Fulcro
            \node at (0, -0.5) {\small Bilancia};
        \end{scope}
    \end{tikzpicture}
    \end{center}
\end{enumerate}

\section{Principi della Dinamica}
I principi della dinamica furono enunciati formalmente da \textbf{Isaac Newton} nella sua opera monumentale \textit{"Philosophiae Naturalis Principia Mathematica"} (1687), riprendendo e sistematizzando gli studi precedenti di \textbf{Galileo Galilei}.

\subsection{Primo Principio: Principio di Inerzia}
\begin{definizione}{Primo Principio: Principio di Inerzia}
Un corpo permane nel suo stato di quiete o di moto rettilineo uniforme finché non interviene una forza esterna a variarne lo stato.
\[ \sum \vec{F} = 0 \implies \vec{v} = \text{cost} \]
Questa proprietà intrinseca della materia è detta \textbf{Inerzia}.
\end{definizione}

Pertanto, se la risultante delle forze è nulla ($\vec{F}_{tot} = 0$), il corpo si trova in uno degli \textbf{Stati di Equilibrio}:
\begin{itemize}
    \item \textbf{Quiete} ($v=0$).
    \item \textbf{Moto Rettilineo Uniforme} ($v = \text{cost} \neq 0$).
\end{itemize}

La condizione $\vec{F}_{tot} = 0$ si verifica in due casi esclusivi:
\begin{enumerate}[label=\alph*)]
    \item Il corpo non è soggetto ad alcuna forza ($F=0$).
    \item Il corpo è soggetto a più forze la cui somma vettoriale è nulla ($\sum \vec{F} = 0$).
\end{enumerate}

\begin{esercizio}{Il Paradosso Sperimentale e l'Esperimento Ideale di Galileo}
A livello concreto, è impossibile osservare il principio di inerzia in forma pura sulla Terra, poiché non otteniamo mai una reale prova sperimentale dell'assenza completa di forze. Infatti, le forze dissipative (come l'attrito dell'aria o del piano) sono sempre presenti.

Per risolvere questo paradosso, Galileo ideò un celebre \textbf{Esperimento Intellettuale} (\textit{Gedankenexperiment}):

    \begin{center}
    \begin{tikzpicture}[scale=0.85, transform shape] 
        % --- Caso A: Simmetrico Ripido ---
        \begin{scope}[xshift=-4cm, yshift=2.5cm]
            \coordinate (TopLeft) at (-2, 2);
            \coordinate (Bottom) at (0, 0);
            \coordinate (TopRight) at (2, 2);
            \draw[thick, rounded corners=15pt] (TopLeft) -- (Bottom) -- (TopRight);
            \draw[dashed, gray] (TopLeft) -- (TopRight);
            \draw[thick] (-2,0) -- (2,0);
            \draw[thick] (-2,2) -- (-2,0);
            \draw[thick] (2,2) -- (2,0);
            \node[left] at (-2.2, 1) {\textbf{A}};
            \node[right, gray] at (TopRight) {$h$};
            \coordinate (BallDown) at (-1.2, 1.2);
            \filldraw[fill=gray!30, draw=black] (BallDown) circle (0.18);
            \draw[->, thick] (BallDown) -- ++(-45:0.8);
            \coordinate (BallUp) at (1.2, 1.2);
            \filldraw[fill=gray!30, draw=black] (BallUp) circle (0.18);
            \draw[->, thick] (BallUp) -- ++(45:0.8);
        \end{scope}
        
        % --- Caso B: Simmetrico Largo ---
        \begin{scope}[xshift=4cm, yshift=2.5cm]
            \coordinate (TopLeftB) at (-3.5, 2);
            \coordinate (BottomB) at (0, 0);
            \coordinate (TopRightB) at (3.5, 2);
            \draw[thick, rounded corners=20pt] (TopLeftB) -- (BottomB) -- (TopRightB);
            \draw[dashed, gray] (TopLeftB) -- (TopRightB);
            \draw[thick] (-3.5,0) -- (3.5,0);
            \draw[thick] (-3.5,2) -- (-3.5,0);
            \draw[thick] (3.5,2) -- (3.5,0);
            \node[left] at (-3.7, 1) {\textbf{B}};
            \node[right, gray] at (TopRightB) {$h$};
            \coordinate (BallDownB) at (-2, 1.14);
            \filldraw[fill=gray!30, draw=black] (BallDownB) circle (0.18);
            \draw[->, thick] (BallDownB) -- ++(-29.7:0.8);
            \coordinate (BallUpB) at (2, 1.14);
            \filldraw[fill=gray!30, draw=black] (BallUpB) circle (0.18);
            \draw[->, thick] (BallUpB) -- ++(29.7:0.8);
        \end{scope}
        
        % --- Caso C: Piano Orizzontale ---
        \begin{scope}[xshift=0cm, yshift=-1.5cm]
            \draw[thick] (-4, 0) -- (4, 0);
            \draw[->, thick, dashed] (4, 0) -- (5.5, 0); 
            \node[left] at (-4.2, 0.5) {\textbf{C}};
            \filldraw[fill=gray!30, draw=black] (0, 0.18) circle (0.18);
            \draw[->, thick, mygreen] (0.3, 0.18) -- (1.8, 0.18) node[above] {$\vec{v} = \text{cost}$};
            \node[below, mygreen] at (1.5, -0.2) {MRU};
        \end{scope}
    \end{tikzpicture}
    \captionof{figure}{Esperimento dei piani inclinati di Galileo.}
    \end{center}

Immaginando di avere due piani inclinati contrapposti:
\begin{enumerate}
    \item Il corpo scivola sul primo piano \textbf{accelerando} e risale sul secondo \textbf{decelerando} fino a tornare alla quota iniziale $h$.
    \item Riducendo l'inclinazione del secondo piano, il corpo percorre un tratto più lungo.
    \item Se l'inclinazione del secondo piano diventa \textbf{nulla} (piano orizzontale), il corpo continuerà a muoversi indefinitamente a velocità costante (\textbf{MRU}).
\end{enumerate}
\end{esercizio}

\subsection{Sistemi di Riferimento Inerziali}
Un altro fatto fondamentale è legato alla relatività del moto.
Infatti:
\begin{itemize}
    \item \textbf{Caso A (Sistemi Fermi/Inerziali)}: Se l'osservatore è fermo, vede i corpi liberi (su cui non agiscono forze) muoversi di MRU o stare fermi.
    \item \textbf{Caso B (Sistemi Accelerati)}: Se l'osservatore si trova in un sistema accelerato (es. su una giostra in rotazione, $\omega \neq 0$), vedrà i corpi muoversi di moto accelerato (forze apparenti) anche se su di essi non agisce alcuna forza reale. In questo caso il Primo Principio \textbf{NON} è valido.
\end{itemize}

\begin{definizione}{Sistemi di Riferimento Inerziali}
Si definisce \textbf{Inerziale} un sistema di riferimento in cui vale il Primo Principio della Dinamica. Tutti i sistemi in moto rettilineo uniforme rispetto a un sistema inerziale sono anch'essi inerziali (\textbf{Classe di Inerzialità}).
\end{definizione}

\subsubsection{Proprietà dei Sistemi Inerziali}
\begin{enumerate}
    \item \textbf{Classe di Equivalenza}: Sono inerziali tutti e soli i sistemi che si muovono di \textbf{MRU} rispetto ad un sistema inerziale $S$.
    \item \textbf{Transitivitá}: Se $S$ è inerziale e $S'$ è in MRU rispetto a $S \Rightarrow S'$ è inerziale.
    \item \textbf{Le Stelle Fisse}: Un sistema di riferimento solidale con le "stelle fisse" (stelle enormemente lontane, percepite come "ferme") è un'ottima approssimazione di sistema inerziale.
    \item \textbf{Relatività di Galileo}: Le leggi della meccanica sono le stesse in \textbf{tutti} i sistemi di riferimento inerziali. Non esiste un sistema inerziale privilegiato.
\end{enumerate}

\begin{osservazione}{La Terra è un sistema inerziale?}
Rigorosamente \textbf{NO}, a causa dei moti di rotazione e rivoluzione. Tuttavia, per la maggior parte delle applicazioni locali, può essere approssimata come tale poiché queste accelerazioni sono trascurabili rispetto a $g$.
\end{osservazione}
Tuttavia, può essere considerata \textbf{approssimativamente inerziale} per la maggior parte degli esperimenti di laboratorio, poiché tali accelerazioni sono molto piccole rispetto alla gravità:
\[ a_{rot} \approx 2.5 \cdot 10^{-3} g, \quad a_{riv} \approx 10^{-4} g \]
Questi effetti diventano rilevanti solo per moti su grande scala (es. correnti atmosferiche) o misure di altissima precisione.

\section{Secondo Principio: Legge Fondamentale (Legge di Newton)}

% ---------------------------------------------------------------
\subsection{Prima Formulazione (Massa Costante)}
% ---------------------------------------------------------------

\noindent
\begin{minipage}[t]{0.52\textwidth}
\vspace{0pt}
dove $\sum\vec{F}$ è la somma delle forze applicate
ad un corpo, avente la stessa direzione dell'accelerazione
$\vec{a}$; la costante di proporzionalità è la massa $m$:
\begin{equation}\label{eq:2nd_law}
    \boxed{\sum \vec{F} = m \cdot \vec{a}}
\end{equation}
\end{minipage}%
\hfill
\begin{minipage}[t]{0.44\textwidth}
\vspace{4pt}
\begin{osservazione}{Unità di misura della Forza}
Nel Sistema Internazionale, la forza si misura in \textbf{Newton} [N]:
\[ 1 \, N = 1 \, kg \cdot m/s^2 \]
\end{osservazione}
\end{minipage}

\medskip
\begin{osservazione}{Relazione tra i Principi}
Affinché valga il 2° principio è necessario definire un sistema di riferimento inerziale tramite il 1° principio. In sistemi non inerziali compaiono forze apparenti che invalidano la forma $\sum \vec{F} = m\vec{a}$ se non corrette.
\end{osservazione}

\begin{definizione}{Secondo Principio della Dinamica}
In un sistema di riferimento \textbf{inerziale}, l'accelerazione di
un punto materiale è \emph{direttamente proporzionale} alla forza
risultante e \emph{inversamente proporzionale} alla sua massa:
\[
    \vec{a} = \vec{R} \;\Longrightarrow\;
    \vec{a} = \frac{\sum\vec{F}}{m} \;\Longrightarrow\;
    \underbrace{\frac{\sum\vec{F}(\vec{r},\dot{\vec{r}})}{m}
               = \ddot{\vec{r}}}_{\text{equazione del moto}}
\]
\end{definizione}

\paragraph{Equazioni Scalari.}
Essendo un'equazione vettoriale, equivale a tre equazioni scalari:
\[
\begin{cases}
    \sum F_x = m\, a_x \\
    \sum F_y = m\, a_y \\
    \sum F_z = m\, a_z
\end{cases}
\qquad\text{con }\; a_i = \dfrac{d^2 x_i}{dt^2}
\]

\paragraph{Problemi della Dinamica.}
\begin{itemize}
    \item \textbf{Diretto}: note forze e c.i.\ ($\vec{r}_0,\vec{v}_0$)
          $\;\ \to\;$ trovare $\vec{r}(t)$.
    \item \textbf{Inverso}: noto $\vec{r}(t)$ (e quindi $\vec{a}$)
          $\;\ \to\;$ trovare le forze.
\end{itemize}

% ---------------------------------------------------------------
\subsection{Formulazione Generale di Newton (Quantità di Moto)}
% ---------------------------------------------------------------
Newton fornì una definizione più generale introducendo la
\textbf{quantità di moto} $\vec{p} = m\vec{v}$.
Per $m = \text{costante}$:
\[
    \sum \vec{F} = \frac{d\vec{p}}{dt}
    \qquad\text{poiché}\qquad
    \frac{d\vec{p}}{dt} = m\frac{d\vec{v}}{dt} = m\vec{a}
\]

\begin{teorema}{Secondo Principio — Forma Generale}
In un sistema inerziale, per un punto materiale con
$\vec{p} = m\vec{v}$:
\begin{equation}\label{eq:2nd_general}
    \sum \vec{F} = \frac{d\vec{p}}{dt}
\end{equation}
Questa forma è valida anche per sistemi a \textbf{massa variabile}
(es.\ un razzo che espelle carburante).
\end{teorema}

\smallskip
Sfruttare $\vec{p}$ risulta efficace in: studio di sistemi multi-corpo,
moto rapidi (urti fra particelle), sistemi a massa variabile.

% ---------------------------------------------------------------
\subsection{Teorema dell'Impulso}
% ---------------------------------------------------------------
\begin{teorema}{Teorema dell'Impulso}
L'impulso della forza/sistema di forze è uguale alla variazione della
quantità di moto del corpo su cui agisce.
Integrando $\sum\vec{F}$ tra $t_1$ e $t_2$:
\begin{align*}
    \vec{I}_{t_1 \to t_2}
        &= \int_{t_1}^{t_2}\sum\vec{F}\,dt
         = \int_{t_1}^{t_2}\frac{d\vec{p}}{dt}\,dt
         = \int_{t_1}^{t_2}d\vec{p} 
        = \bigl[\vec{p}\bigr]_{t_1}^{t_2}
         = \vec{p}(t_2)-\vec{p}(t_1)
         = \Delta\vec{p}
\end{align*}
\end{teorema}

% ---------------------------------------------------------------
\subsubsection{Forze Impulsive}
% ---------------------------------------------------------------
\noindent
Le forze ``impulsive'' (es.\ urto) sono \textbf{molto intense} e
agiscono per un \textbf{breve intervallo} $\Delta t\approx 1\,\mathrm{ms}$.\medskip

\noindent
Sfruttando il teorema dell'impulso, la \textbf{forza impulsiva media}
è definita come:
\begin{align*}
    \int_{0^-}^{0^+}\vec{F}\,dt
        &= \langle\vec{F}\rangle_{\Delta t}\cdot\Delta t
         = \vec{p}(0^+)-\vec{p}(0^-)
    \Longrightarrow\quad
    \boxed{\langle\vec{F}\rangle_{\Delta t} = \frac{\Delta\vec{p}}{\Delta t}}
    \qquad\text{(media integrale)}
\end{align*}

\noindent
Analizziamo gli stati immediatamente prima ($t=0^-$) e dopo ($t=0^+$) dell'urto:

\begin{figure}[ht]
    \centering
    \begin{tikzpicture}[scale=1.05, >=stealth]
        % Muro verticale con tratteggio
        \draw[line width=3pt, color=darkgray] (0,-1.8) -- (0,1.8);
        \fill[pattern=north east lines, pattern color=gray!70]
            (-0.35,-1.8) rectangle (0,1.8);
        \node[above, color=darkgray, font=\small\itshape] at (0,1.85)
            {Parete Rigida};

        % --- PRIMA DELL'URTO (t = 0-) ---
        \begin{scope}[yshift=0.9cm]
            \node[circle, draw=gray!70, dashed, minimum size=1.1cm,
                  thick, fill=gray!8] (ball_in) at (2.2, 0) {};
            \node[gray, font=\small] at (2.2, 0) {$m$};
            \draw[->, thick, gray] (ball_in.west) -- ++(-1.2,0)
                node[midway, above, font=\small] {$\vec{v}_{in}$};
            \node[right, gray!80, font=\footnotesize] at (2.9,0) {$t=0^-$};
        \end{scope}

        % --- DURANTE L'URTO (t = 0): forza impulsiva ---
        \fill[red!80] (0,0) circle (2.5pt);
        \draw[->, ultra thick, red!80] (0,0) -- (1.8,0)
            node[midway, below, font=\small] {$\vec{F}_{imp}$};
        \node[left, red!80, font=\footnotesize] at (-0.05,0) {$t=0$};

        % --- DOPO L'URTO (t = 0+) ---
        \begin{scope}[yshift=-0.9cm]
            \node[circle, draw=black, fill=white, minimum size=1.1cm,
                  thick] (ball_out) at (2.2, 0) {};
            \node[font=\small] at (2.2, 0) {$m$};
            \draw[->, thick, black] (ball_out.east) -- ++(1.2,0)
                node[midway, above, font=\small] {$\vec{v}_{out}$};
            \node[right, black, font=\footnotesize] at (2.9,-0.4) {$t=0^+$};
        \end{scope}
    \end{tikzpicture}
    \caption{Urto elastico contro una parete rigida. La forza impulsiva
        $\vec{F}_{imp}$ agisce per $\Delta t\to 0$, causando la variazione
        di quantità di moto $\Delta\vec{p} = \vec{p}(0^+)-\vec{p}(0^-)$.}
    \label{fig:impulso}
\end{figure}

\section{Terzo Principio: Azione e Reazione}
\begin{definizione}{Terzo Principio: Azione e Reazione}
Se un corpo $A$ esercita una forza $\vec{F}_{AB}$ su un corpo $B$, il corpo $B$ esercita contemporaneamente una forza $\vec{F}_{BA}$ su $A$, tale che:
\begin{equation}
    \vec{F}_{AB} = -\vec{F}_{BA}
\end{equation}
Le due forze hanno stessa retta d'azione, stesso modulo, ma versi opposti. Si applicano a \textbf{punti materiali diversi}.
\end{definizione}
Questo significa che le forze di interazione ("Azione" e "Reazione") compaiono sempre in \textbf{coppia} e condividono le seguenti caratteristiche:
\begin{itemize}
    \item Hanno lo stesso \textbf{modulo}: $|\vec{F}_{AB}| = |\vec{F}_{BA}|$.
    \item Hanno la stessa \textbf{direzione} (la retta congiungente i due corpi).
    \item Hanno \textbf{verso opposto}.
    \item \textbf{Agiscono su corpi DIVERSI}:
    \begin{itemize}
        \item $\vec{F}_{AB}$ agisce sul corpo $B$ (dovuta ad A).
        \item $\vec{F}_{BA}$ agisce sul corpo $A$ (dovuta a B).
    \end{itemize}
    \textit{Pertanto, azine e reazione non si elidono mai a vicenda perché applicate a oggetti distinti.}
\end{itemize}

\begin{figure}[ht]
    \centering
    \begin{tikzpicture}[scale=1.2, >=stealth]
        % Corpo A
        \node[circle, draw=black, fill=mypen!10, minimum size=1.2cm, thick] (A) at (0,0) {\textbf{A}};
        
        % Corpo B
        \node[circle, draw=black, fill=mypen!10, minimum size=1.2cm, thick] (B) at (4,0) {\textbf{B}};
        
        % Linea tratteggiata di interazione
        \draw[dashed, gray] (A) -- (B);
        
        % Forza F_BA (Su A dovuta a B) <-- verso sinistra
        \draw[->, ultra thick, mypen] (A.west) -- ++(-1.5, 0) node[above] {$\vec{F}_{BA}$};
        \node[below, mypen, font=\scriptsize] at ($(A.west)+(-0.75,0)$) {su A da B};

        % Forza F_AB (Su B dovuta ad A) --> verso destra
        \draw[->, ultra thick, mypen] (B.east) -- ++(1.5, 0) node[above] {$\vec{F}_{AB}$};
        \node[below, mypen, font=\scriptsize] at ($(B.east)+(0.75,0)$) {su B da A};
        
        % Note visive
        \node[align=center, font=\small] at (2, 0.5) {\textit{Interazione}};
    \end{tikzpicture}
    \caption{Illustrazione del Terzo Principio. Le forze $\vec{F}_{AB}$ e $\vec{F}_{BA}$ sono uguali e contrarie, ma applicate a corpi diversi.}
    \label{fig:terzo_principio}
\end{figure}

\section{Massa e Forza Peso}
\label{sec:massa_peso}

La \textbf{massa} è definita come una proprietà intrinseca del corpo.
Dalla seconda legge della dinamica, possiamo definirla operativamente come:
\begin{equation}
    m = \frac{|\sum \vec{F}|}{|\vec{a}|}
\end{equation}

La massa gode della \textbf{proprietà di additività} (tranne a livello subnucleare, dove la massa può trasformarsi in energia $E=mc^2$). Ovvero, considerando due corpi con masse $m_1$ e $m_2$:
\[ m_{(1+2)} = m_1 + m_2 = m_{tot} \]

\subsection{Misurazione della Massa}
La massa si misura attraverso la \textbf{bilancia}, sfruttando la forza peso. Se la bilancia si trova in una situazione di \textbf{equilibrio}, significa che la risultante delle forze è nulla ($\sum \vec{F} = 0$).
Se ci troviamo sullo stesso punto della Terra (dove $\vec{g}$ è costante):
\[ \vec{P}_1 = \vec{P}_2 \implies m_1 \vec{g} = m_2 \vec{g} \implies m_1 = m_2 \]

Conoscendo le forze, misuriamo le accelerazioni e di conseguenza le masse (usando strumenti come la bilancia o lo spettrometro di massa).

\begin{figure}[h!]
    \centering
    \begin{tikzpicture}[scale=0.8]
        % Base
        \fill[gray!30] (-0.5,0) rectangle (0.5,-0.2);
        \draw[thick] (0,0) -- (0,2); % Asta verticale
        \filldraw (0,2) circle (2pt); % Fulcro
        
        % Braccio orizzontale
        \draw[thick] (-2.5,2) -- (2.5,2);
        
        % Piatti
        \draw (-2.5,2) -- (-2.5,1);
        \draw (-3,1) -- (-2,1); % Piatto sx
        \node[above] at (-2.5,1) {$m_1$};
        
        \draw (2.5,2) -- (2.5,1);
        \draw (2,1) -- (3,1); % Piatto dx
        \node[above] at (2.5,1) {$m_2$};
        
        % Indicatore equilibrio
        \draw[->, thick] (0,1.5) -- (0,0.5);
    \end{tikzpicture}
    \caption{Schema di principio della bilancia a bracci uguali.}
    \label{fig:bilancia_schema}
\end{figure}

\subsection{Massa Inerziale e Gravitazionale}
In realtà, la massa può essere distinta in due diverse tipologie concettuali:
\begin{itemize}
    \item \textbf{Massa Inerziale} ($m_i$): rappresenta la capacità di un corpo di opporsi al moto (resistenza all'accelerazione).
    \item \textbf{Massa Gravitazionale} ($m_g$): rappresenta la capacità di un corpo di attrarre altri corpi (carica gravitazionale).
\end{itemize}

\begin{teorema}{Equivalenza Massa Inerziale e Gravitazionale}
Sebbene concettualmente diverse (\textbf{$m_i$}: resistenza al moto; \textbf{$m_g$}: capacità di attrarre/essere attratti), gli esperimenti di \textbf{Eötvös} hanno dimostrato con altissima precisione che il loro rapporto è costante per tutti i corpi. Nel SI si assume:
\[ m_i = m_g = m \]
\end{teorema}

\begin{osservazione}{Principio di Equivalenza Debole}
Pertanto si può costatare l'uguaglianza numerica tra le due grandezze:
\[ m_i = m_g = m \]
\end{osservazione}

\subsection{La Forza Peso}
La \textbf{Forza Peso} è la forza esercitata dalla Terra su un corpo.
\begin{equation}
    \vec{P} = m \vec{g} \quad \text{con } |\vec{g}| \approx \SI{9.81}{m/s^2}
\end{equation}

L'accelerazione di gravità $\vec{g}$ \textbf{non è costante} su tutta la superficie terrestre. In effetti, essa è una grandezza puntuale e locale che varia in base a:
\begin{enumerate}
    \item \textbf{Rotazione Terrestre}: $\vec{g}$ è maggiore ai poli e minore all'equatore (a causa della forza centrifuga).
    \item \textbf{Geoide}: La Terra è una sfera imperfetta.
    \begin{itemize}
        \item Più in alto $\rightarrow$ $\vec{g}$ diminuisce (lontananza dal centro).
        \item Più in basso $\rightarrow$ $\vec{g}$ aumenta.
    \end{itemize}
    \item \textbf{Disomogeneità Locali}: Presenza di materiali a diversa densità nel sottosuolo.
\end{enumerate}

La misurazione precisa di $\vec{g}$ avviene attraverso particolari pendoli con una risoluzione dell'ordine di $10^{-5}$ (Gravimetria).

\subsection{Differenza tra Massa e Peso}

\begin{figure}[h!]
    \centering
    \begin{tikzpicture}[
        node distance=2cm,
        box/.style={draw, rounded corners, align=center, minimum width=2.5cm, minimum height=1cm, font=\small},
        arrow/.style={->, >=stealth, thick}
    ]
        % Nodo radice
        \node (diff) [box, fill=gray!10] {\textbf{Differenza Massa $\neq$ Peso}};
        
        % Nodi figli
        \node (massa) [box, below left=1cm and 0.2cm of diff, fill=mypen!10] {\textbf{Massa}};
        \node (peso) [box, below right=1cm and 0.2cm of diff, fill=mypen!10] {\textbf{Peso}};
        
        % Caratteristiche Massa
        \node (m_desc) [below=0.2cm of massa, align=center, font=\small] {
            $\bullet$ Grandezza \textbf{Scalare}\\
            $\bullet$ Proprietà \textbf{Intrinseca}
        };
        
        % Caratteristiche Peso
        \node (p_desc) [below=0.2cm of peso, align=center, font=\small] {
            $\bullet$ Grandezza \textbf{Vettoriale}\\
            $\bullet$ Dipendente da $\vec{g}$ e $\vec{a}$
        };
        
        % Frecce
        \draw[arrow] (diff) -- (massa);
        \draw[arrow] (diff) -- (peso);
        
    \end{tikzpicture}
    \caption{Schema delle differenze fondamentali tra Massa e Peso.}
    \label{fig:diff_massa_peso}
\end{figure}

\begin{esercizio}{L'Ascensore e il Peso Apparente}
Un corpo di massa $m$, posto su una bilancia all'interno di un ascensore, è soggetto alla forza peso $\vec{P} = m\vec{g}$ e alla reazione vincolare $\vec{N}$ della bilancia. La bilancia misura la reazione vincolare $\vec{N}$ (il "peso apparente").

Dalla seconda legge della dinamica: $N = m(g + a)$. Analizziamo i casi:
\begin{enumerate}
    \item \textbf{Ascensore fermo ($a=0$)}: $N = mg$. Peso reale.
    \item \textbf{Ascensore accelera verso l'alto ($a > 0$)}: $N = m(g + a) > mg$. Ci sentiamo più pesanti.
    \item \textbf{Ascensore accelera verso il basso ($a < 0$)}: $N = m(g - |a|) < mg$. Ci sentiamo più leggeri.
    \item \textbf{Caduta Libera ($a = -g$)}: $N = 0$. \textbf{Imponderabilità} (assenza di peso).
\end{enumerate}

    \begin{center}
    \begin{tikzpicture}[scale=1.0, >=stealth]
        \tikzset{
            cabin/.style={draw=black!70, thick, rounded corners=3pt, fill=gray!6, minimum width=3.2cm, minimum height=3.8cm},
            person/.style={circle, draw=mypen!60, fill=mypen!25, minimum size=0.7cm, thick},
            fvec/.style={->, line width=1.6pt},
            avec/.style={->, line width=2.2pt, violet}
        }
        % Caso a=0
        \begin{scope}[xshift=0cm]
            \node[cabin] at (0, 1.9) {};
            \node[person] (Hd1) at (0, 1.9) {\scriptsize $m$};
            \draw[fvec, mygreen!80] (Hd1.south) -- ++(0,-1.1) node[right=3pt] {$\vec{P}$};
            \draw[fvec, mypen!80] (Hd1.north) -- ++(0,+0.85) node[right=3pt] {$\vec{N}$};
            \node[below=6pt] at (0,0) {\textbf{$a=0$}};
        \end{scope}
        % Caso a > 0
        \begin{scope}[xshift=5.5cm]
            \node[cabin] at (0, 1.9) {};
            \node[person] (Hd2) at (0, 1.9) {\scriptsize $m$};
            \draw[fvec, mygreen!80]  (Hd2.south) -- ++(0,-1.1) node[right=3pt] {$\vec{P}$};
            \draw[fvec, line width=2.4pt, mypen!80] (Hd2.north) -- ++(0,+1.55) node[right=3pt] {$\vec{N}$};
            \draw[avec] (-2.2, 1.2) -- (-2.2, 2.8) node[above, violet] {$\vec{a}$};
            \node[below=6pt] at (0,0) {\textbf{$a > 0$}};
        \end{scope}
        % Caso a < 0
        \begin{scope}[xshift=11cm]
            \node[cabin] at (0, 1.9) {};
            \node[person] (Hd3) at (0, 1.9) {\scriptsize $m$};
            \draw[fvec, mygreen!80]  (Hd3.south) -- ++(0,-1.1) node[right=3pt] {$\vec{P}$};
            \draw[fvec, mypen!80] (Hd3.north) -- ++(0,+0.45) node[right=3pt] {$\vec{N}$};
            \draw[avec] (-2.2, 2.8) -- (-2.2, 1.2) node[below, violet] {$\vec{a}$};
            \node[below=6pt] at (0,0) {\textbf{$a < 0$}};
        \end{scope}
    \end{tikzpicture}
    \captionof{figure}{Variazione del peso apparente $\vec{N}$ nell'ascensore.}
    \end{center}
\end{esercizio}

\section{Forze di Contatto}
\label{sec:forze_contatto}

Le \textbf{Forze Vincolari (o Normali)} sono forze che uno o più corpi esercitano su un altro limitandone il moto.

\subsection{Equazioni del Vincolo}
Un vincolo impone delle condizioni geometriche alle coordinate del corpo. Ecco i tre casi principali (Fig. \ref{fig:vincoli_rappresentazione}):

\begin{figure}[h!]
    \centering
    % Use minipage for better alignment
    \begin{minipage}{0.31\textwidth}
        \centering
        \textbf{Piano (Tavolo)}
        \vspace{0.2cm}
        \begin{tikzpicture}[scale=0.55, x={(-0.4cm,-0.3cm)}, y={(1cm,0cm)}, z={(0cm,1cm)}]
            % Assi
            \draw[->] (0,0,0) -- (3,0,0) node[left] {\tiny $x$};
            \draw[->] (0,0,0) -- (0,3,0) node[right] {\tiny $y$};
            \draw[->] (0,0,0) -- (0,0,3) node[above] {\tiny $z$};
            % Piano
            \fill[blue!10, opacity=0.7] (0,0,0) -- (2.5,0,0) -- (2.5,2.5,0) -- (0,2.5,0) -- cycle;
            \draw[dashed] (0,0,0) -- (2.5,0,0) -- (2.5,2.5,0) -- (0,2.5,0) -- cycle;
        \end{tikzpicture}
        \vspace{0.1cm}
        {\scriptsize $z(t) = 0$}
    \end{minipage}
    \hfill
    \begin{minipage}{0.31\textwidth}
        \centering
        \textbf{Circolare}
        \vspace{0.2cm}
        \begin{tikzpicture}[scale=0.55, x={(-0.4cm,-0.3cm)}, y={(1cm,0cm)}, z={(0cm,1cm)}]
            % Assi
            \draw[->] (0,0,0) -- (3,0,0);
            \draw[->] (0,0,0) -- (0,3,0);
            \draw[->] (0,0,0) -- (0,0,3);
            % Cerchio
            \draw[thick, red] (2,0,0) arc (0:90:2);
            \draw[dashed, red] (2,0,0) -- (0,0,0) -- (0,2,0);
            \node[red, right] at (1,0,0) {\tiny $R$};
        \end{tikzpicture}
        \vspace{0.1cm}
        {\scriptsize $x^2 + y^2 = R^2$}
    \end{minipage}
    \hfill
    \begin{minipage}{0.31\textwidth}
        \centering
        \textbf{Piano Inclinato}
        \vspace{0.2cm}
        \begin{tikzpicture}[scale=0.7]
            % Cuneo
            \draw[thick] (0,0) -- (2.5,0) -- (2.5,1.2) -- cycle;
            \draw (0.5,0) arc (0:26:0.5);
            \node at (0.8,0.2) {\tiny $\theta$};
            % Sistema ruotato
            \begin{scope}[shift={(1.2, 0.6)}, rotate=26.56]
                \draw[->, blue] (0,0) -- (0.8,0) node[right] {\tiny $x'$};
                \draw[->, blue] (0,0) -- (0,0.8) node[above] {\tiny $y'$};
            \end{scope}
        \end{tikzpicture}
        \vspace{0.1cm}
        {\scriptsize $y'(t) = 0$}
    \end{minipage}
    
    \caption{Rappresentazione grafica delle equazioni dei vincoli per diverse geometrie.}
    \label{fig:vincoli_rappresentazione}
\end{figure}

\begin{osservazione}{La Reazione Vincolare}
Il fatto che sul piano orizzontale statico risulti $\vec{N} + \vec{P} = 0$ (ovvero $\vec{N} = -\vec{P}$) è \textbf{del tutto casuale}. Infatti $|\vec{N}|$ non corrisponde a $|-\vec{P}|$ in un piano inclinato, né in un ascensore accelerato.

$\vec{P} = m\vec{g}$ esiste anche se non ci fosse la $\vec{N}$ del tavolo. Questo perché il $3^{\circ}$ principio non si applica al sistema corpo-tavolo, ma al sistema \textbf{Corpo-Terra}. Il tavolo è solo il vincolo che impedisce il naturale moto dei gravi verso la Terra.
\end{osservazione}

\subsection{Tipologie di Vincoli}

\subsubsection{Cerniere e Tensioni}

\begin{figure}[h!]
    \centering
    \begin{tikzpicture}[>=stealth, scale=0.9]

        % 1. CERNIERE
        \begin{scope}[xshift=-4cm]
            \node[align=center, anchor=south, font=\bfseries] at (0, 2.3) {CERNIERE};
            % Cilindro/Cerniera
            \fill[gray!20] (-0.2,0) rectangle (0.2,2);
            \draw[thick] (-0.2,0) -- (-0.2,2); \draw[thick] (0.2,0) -- (0.2,2);
            \draw (0,2) ellipse (0.2 and 0.05); \draw (0,0) ellipse (0.2 and 0.05);
            % Staffa
            \draw[thick] (0.2, 1.5) -- (1, 1.8) -- (1, 0.8) -- (0.2, 0.5);
            % Vettore
            \draw[->, thick] (0.5, 1.6) -- (1.5, 2.2) node[right] {$\vec{F}_{\perp}$};
        \end{scope}

        % 2. TENSIONE (BLOCCHI)
        \begin{scope}[xshift=2cm, yshift=1cm]
            \node[anchor=west, font=\bfseries] at (-1.5, 1.2) {FILI INESTENSIBILI};
            \draw[thick, fill=gray!30] (-1.5, -0.3) rectangle (-0.5, 0.3); % Blocco 1
            \draw[thick, fill=gray!30] (1.5, -0.3) rectangle (2.5, 0.3); % Blocco 2
            \draw[thick] (-0.5, 0) -- (1.5, 0); % Filo
            % Vettori Tensione
            \draw[->, thick, red] (-0.5, 0) -- (0, 0) node[above, font=\tiny] {$\vec{T}$};
            \draw[->, thick, red] (1.5, 0) -- (1, 0) node[above, font=\tiny] {$\vec{T}$};
            \node[below, font=\scriptsize] at (0.5, -0.3) {Tensione Costante};
        \end{scope}

        % 3. CARICO DI ROTTURA (FUNE)
        \begin{scope}[xshift=2cm, yshift=-1.5cm]
            % Fune spessa
            \draw[thick] (-1, 0.2) -- (0.5, 0.2); 
            \draw[thick] (-1, -0.2) -- (0.5, -0.2);
            % Rottura
            \draw[decorate, decoration={zigzag, segment length=3mm, amplitude=1mm}] (0.5,0.2) -- (0.5,-0.2);
            
            \draw[->, thick] (0.5, 0) -- (1.5, 0) node[above] {$\vec{F}_{max}$};
            \node[right, font=\scriptsize, align=left] at (1.5, -0.2) {Carico di\\Rottura};
        \end{scope}

    \end{tikzpicture}
    \caption{Esempi di forze vincolari: cerniere che esercitano forze perpendicolari all'asse e fili che trasmettono tensioni.}
    \label{fig:vincoli_tipi}
\end{figure}

\subsubsection{Vincoli Lisci e Scabri}

\begin{itemize}
    \item \textbf{VINCOLO LISCIO}: Esercita solo forze vincolari perpendicolari alla superficie di contatto ($\vec{N}$). Non vi è alcuna forza che si oppone al moto laterale.
    \item \textbf{VINCOLO SCABRO}: Esercita una reazione vincolare $\vec{R}$ che ha due componenti:
    \begin{enumerate}
        \item \textbf{Normale} ($\vec{N}$): perpendicolare alla superficie.
        \item \textbf{Attrito} ($\vec{F}_{att}$): parallela alla superficie e opposta al verso del moto (o dell'imminente moto).
    \end{enumerate}
\end{itemize}

\begin{figure}[h!]
    \centering
    \begin{tikzpicture}[scale=1, >=stealth]
        % --- Liscio ---
        \begin{scope}[xshift=-3.5cm]
            \draw[thick] (-1.5,0) -- (1.5,0);
            \draw[thick, fill=blue!5] (-0.5,0) rectangle (0.5,0.8);
            
            \draw[->, thick, blue] (0,0.8) -- (0,1.8) node[above] {$\vec{N}$};
            \draw[->, thick, red] (0,0.4) -- (0,-0.6) node[right] {$\vec{P}$};
            \draw[->, dashed, black] (0.6,0.4) -- (1.4,0.4) node[right] {$\vec{v}$};
            
            \node[below, font=\bfseries] at (0,-0.8) {LISCIO};
            \node[below, font=\scriptsize] at (0,-1.2) {Reazione solo $\perp$};
        \end{scope}

        % --- Scabro ---
        \begin{scope}[xshift=3.5cm]
            \draw[thick] (-1.5,0) -- (1.5,0);
            \foreach \x in {-1.5,-1.3,...,1.5} \draw (\x,0) -- (\x-0.1,-0.15); % Tratteggio terreno
            \draw[thick, fill=orange!5] (-0.5,0) rectangle (0.5,0.8);
            
            % Reazione Totale R
            \draw[->, thick, violet, dashed] (0,0) -- (-0.8, 1.5) node[left] {$\vec{R}$};
            
            % Componenti
            \draw[->, thick, blue] (0,0.8) -- (0,1.8) node[above] {$\vec{N}$};
            \draw[->, thick, red] (0,0.4) -- (0,-0.6) node[right] {$\vec{P}$};
            \draw[->, thick, violet] (0,0) -- (-1.2,0) node[below] {$\vec{F}_{att}$};
            
            \draw[->, dashed, black] (0.6,0.4) -- (1.4,0.4) node[right] {$\vec{v}$};
            
            \node[below, font=\bfseries] at (0,-0.8) {SCABRO};
            \node[below, font=\scriptsize] at (0,-1.2) {$\vec{R} = \vec{N} + \vec{F}_{att}$};
        \end{scope}
    \end{tikzpicture}
    \caption{Differenza tra vincolo liscio e scabro. Nel vincolo scabro la reazione vincolare $\vec{R}$ non è perpendicolare alla superficie.}
    \label{fig:liscio_vs_scabro}
\end{figure}

Dove $\mu_s$ è il \textbf{Coefficiente di Attrito Statico} (adimensionale), che dipende dalla natura dei materiali e dalla levigatezza delle superfici.

\begin{figure}[h!]
    \centering
    \begin{tikzpicture}[>=stealth, scale=1]
        % Piano
        \draw[thick] (-2,0) -- (2,0);
        \foreach \x in {-2,-1.8,...,2} \draw (\x,0) -- (\x-0.1,-0.15);
        
        % Corpo
        \draw[thick, fill=gray!10] (-0.75,0) rectangle (0.75,1);
        \node at (0,0.5) {$m$};
        
        % Forze
        \draw[->, thick, blue] (0,1) -- (0,1.8) node[above] {$\vec{N}$};
        \draw[->, thick, red] (0,0.5) -- (0,-0.8) node[right] {$\vec{P}$};
        
        \draw[->, thick, black] (0.75,0.5) -- (1.8,0.5) node[right] {$\vec{F}_{ext}$};
        \draw[->, thick, violet] (-0.75,0.5) -- (-1.8,0.5) node[left] {$\vec{F}_s$};
        
        \node[below, align=center, font=\small] at (0,-1) {$\vec{v} = 0$\\Equilibrio: $\vec{F}_s = - \vec{F}_{ext}$};
    \end{tikzpicture}
    \caption{Attrito Statico: si oppone alla forza motrice finché il corpo è fermo.}
\end{figure}

\subsubsection{Attrito Dinamico ($F_d$)}
È la forza parallela alla superficie presente quando il corpo è \textbf{in moto} relativo su una superficie scabra.
Si oppone sempre al versore della velocità istantanea $\hat{v}$.

A differenza dell'attrito statico, l'attrito dinamico ha un modulo (approssimativamente) \textbf{costante} durante il moto, indipendentemente dalla velocità:

\begin{equation}
    |\vec{F}_d| = \mu_d |\vec{N}|
\end{equation}

In forma vettoriale:
\begin{equation}
    \vec{F}_d = - \mu_d |\vec{N}| \hat{v}
\end{equation}

Dove:
\begin{itemize}
    \item $\mu_d$ è il \textbf{Coefficiente di Attrito Dinamico} (solitamente $\mu_d < \mu_s$).
    \item Il segno meno unito al versore $\hat{v}$ indica che il verso di $\vec{F}_d$ è sempre \textbf{opposto} a quello del moto (della velocità e dello spostamento elementare).
    \item \textbf{Nota Bene}: È SBAGLIATO scrivere $\vec{F}_d = -\mu_d \vec{N}$, poiché $\vec{F}_d$ (parallela al piano) e $\vec{N}$ (perpendicolare) sono vettori ortogonali! La relazione vettoriale corretta lega attrito e velocità.
\end{itemize}

\begin{figure}[h!]
    \centering
    \begin{tikzpicture}[>=stealth, scale=1]
        % Piano
        \draw[thick] (-2,0) -- (2,0);
        \foreach \x in {-2,-1.8,...,2} \draw (\x,0) -- (\x-0.1,-0.15);
        
        % Corpo
        \draw[thick, fill=gray!10] (-0.75,0) rectangle (0.75,1);
        \node at (0,0.5) {$m$};
        
        % Velocità
        \draw[->, dashed, thick, black] (0.8, 1.2) -- (1.8, 1.2) node[above] {$\vec{v}$};
        
        % Forze
        \draw[->, thick, blue] (0,1) -- (0,1.8) node[above] {$\vec{N}$};
        \draw[->, thick, red] (0,0.5) -- (0,-0.8) node[right] {$\vec{P}$};
        
        \draw[->, thick, violet] (-0.75,0.5) -- (-1.8,0.5) node[left] {$\vec{F}_d$};
        
        \node[below, align=center, font=\small] at (0,-1) {Moto in corso ($v \neq 0$)\\$F_d = \text{costante}$};
    \end{tikzpicture}
    \caption{Attrito Dinamico: si oppone alla velocità del corpo.}
    \label{fig:attrito_dinamico}
\end{figure}

\subsubsection{Confronto tra Attriti: Il Caso del Piano Inclinato}
In generale, la forza limite che separa lo stato di \textbf{QUIETE} (in cui agisce l'attrito statico $\vec{F}_s$) dallo stato di \textbf{MOTO} (in cui agisce l'attrito dinamico $\vec{F}_d$) corrisponde alla forza di attrito statico massima ($F_{s,max}$).

Consideriamo il caso generico di un corpo su un \textbf{piano inclinato} di un angolo $\theta$. Ci troviamo nella situazione limite in cui il corpo è sul punto di muoversi (imminente moto), quindi $F_s = F_{s,max} = \mu_s N$.

Scomponiamo le forze lungo gli assi:
\begin{itemize}
    \item Asse $x$ (parallelo al piano, verso il basso): componente parallela del peso $P_{\parallel} = mg \sin\theta$ e forza d'attrito statico $F_s$ (che si oppone allo scivolamento).
    \item Asse $y$ (perpendicolare al piano): reazione vincolare $N$ e componente normale del peso $P_{\perp} = mg \cos\theta$.
\end{itemize}

Il sistema all'equilibrio limite diventa:
\begin{equation}
\begin{cases}
    \sum F_x = m g \sin\theta - F_{s,max} = 0 \\
    \sum F_y = N - m g \cos\theta = 0
\end{cases}
\implies
\begin{cases}
    \mu_s N = m g \sin\theta \\
    N = m g \cos\theta
\end{cases}
\end{equation}

Sostituendo la seconda equazione nella prima:
\[ \mu_s (mg \cos\theta) = mg \sin\theta \]

Semplificando la massa $m$ e l'accelerazione $g$, otteniamo una relazione fondamentale:
\begin{equation}
    \mu_s = \frac{\sin\theta}{\cos\theta} = \tan\theta
\end{equation}

\paragraph{Angolo Critico ($\theta_{max}$)}
Ragionando all'opposto, possiamo calcolare l'angolo massimo di inclinazione $\theta_{max}$ per cui il corpo rimane fermo:
\[ \theta_{max} = \arctan(\mu_s) \]

\begin{itemize}
    \item Finché $\theta < \theta_{max}$: il corpo resta \textbf{fermo} ($F_s < F_{s,max}$).
    \item Quando $\theta > \theta_{max}$: la componente parallela del peso supera l'attrito massimo e il corpo \textbf{scivola}.
\end{itemize}

Geometricamente, questo definisce il cosiddetto \textbf{Cono di Attrito Statico}: se la risultante delle forze di contatto cade all'interno di questo cono di apertura $2\theta_{max}$, il corpo rimane in equilibrio.

\begin{figure}[h!]
    \centering
    \begin{minipage}{0.48\textwidth}
        \centering
        \begin{tikzpicture}[scale=1.2, >=stealth]
            % Piano Inclinato
            \coordinate (A) at (0,0);
            \coordinate (B) at (3,0);
            \coordinate (C) at (3,1.5); % tan(theta) = 0.5 -> theta approx 26.5 deg
            \draw[thick] (A) -- (B) -- (C) -- cycle;
            \draw (0.5,0) arc (0:26.5:0.5);
            \node at (0.8,0.2) {$\theta$};

            % Corpo
            \begin{scope}[shift={(1.5, 0.75)}, rotate=26.56]
                \draw[thick, fill=gray!10] (-0.4, 0) rectangle (0.4, 0.5);
                \fill (0, 0.25) circle (1.5pt); % Centro di massa

                % Sistema di riferimento locale
                \draw[gray, dashed] (-1,0) -- (1,0) node[right] {$x$};
                \draw[gray, dashed] (0,-1) -- (0,1) node[above] {$y$};

                % Forze
                \draw[->, thick, blue] (0, 0) -- (0, 1) node[right] {$\vec{N}$};
                \draw[->, thick, violet] (0, 0) -- (-0.8, 0) node[above] {$\vec{F}_{s}$};
                
                % Peso (non ruotato rispetto alla terra, SCOMPOSTO)
                \draw[->, thick, red] (0, 0) -- (0.8, 0) node[above, xshift=2pt] {$P_{\parallel}$};
                \draw[->, thick, red] (0, 0) -- (0, -0.8) node[right] {$P_{\perp}$};
                
            \end{scope}
            
        \end{tikzpicture}
        \caption{Equilibrio limite su piano inclinato.}
        \label{fig:piano_inclinato_attrito}
    \end{minipage}
    \hfill
    \begin{minipage}{0.48\textwidth}
        \centering
        \begin{tikzpicture}[scale=0.9]
            % Piano
            \draw[thick] (-2,0) -- (2,0);
            % Cono
            \fill[gray!10] (0,0) -- (1.5, 2.5) -- (-1.5, 2.5) -- cycle;
            \draw[dashed] (0,0) -- (0,3);
            
            \draw[thick] (0,0) -- (1.5, 2.5);
            \draw[thick] (0,0) -- (-1.5, 2.5);
            
            % Angolo
            \draw (0,0.8) arc (90:59:0.8);
            \node at (0.4, 1.1) {$\theta_{max}$};
            
            \node[align=center, below] at (0,0) {\textbf{CONO DI}\\\textbf{ATTRITO STATICO}};
        \end{tikzpicture}
        \caption{Cono di Attrito.}
        \label{fig:cono_attrito}
    \end{minipage}
\end{figure}

\paragraph{Nota sul piano orizzontale ($\theta = 0$):}
Su un piano orizzontale, in assenza di altre forze esterne, $P_{\parallel} = m g \sin(0) = 0$, quindi la forza di attrito necessaria per l'equilibrio è nulla.
Per misurare $\mu_s$ su un piano orizzontale, bisogna applicare una forza esterna $F_{min}$ finché il corpo non si stacca:
\[ F_{min} = |F_{s,max}| = \mu_s m g \implies \mu_s = \frac{F_{min}}{m g} \]
Dopo lo stacco, se vogliamo mantenere il corpo a velocità costante ($a=0$), dovremo applicare una forza pari all'attrito dinamico: 
\[ F_{motrice} = F_d = \mu_d m g \]
Essendo solitamente $\mu_d < \mu_s$, la forza necessaria per mantenere il moto uniforme è minore di quella necessaria per avviarlo (forza di primo distacco).

\begin{esercizio}{Dinamica del Moto sul Piano Inclinato con Attrito}
Analizziamo il moto di un corpo di massa $m$ su un piano inclinato di un angolo $\theta$.

    \begin{center}
    \begin{tikzpicture}[scale=1.2, >=stealth]
        % Piano Inclinato
        \draw[thick] (0,0) -- (4,0) -- (4,2.3) -- cycle;
        \draw (0.8,0) arc (0:30:0.8);
        \node at (1.1, 0.25) {$\theta$};
        
        % Assi ruotati
        \begin{scope}[rotate=30, shift={(2.5, 0.5)}]
             \draw[->, gray] (0,0) -- (1.5,0) node[below] {$x$};
             \draw[->, gray] (0,0) -- (0,1.5) node[left] {$y$};
             
             % Corpo
             \draw[thick, fill=gray!10] (-0.4,0) rectangle (0.4,0.6);
             \coordinate (C) at (0,0.3);
             
             % Forze
             \draw[->, thick, blue] (C) -- (0, 1.2) node[above] {$\vec{N}$};
             \draw[->, thick, violet] (C) -- (-1.2, 0) node[left] {$\vec{F}_{att}$};
             
             % Peso
             \draw[->, thick, red] (C) -- ++(0.5, -0.866) node[below right] {$\vec{P} = m\vec{g}$};
             
             % Scomposizione Peso
             \draw[->, dashed, red!60] (C) -- (0.5, 0.3) node[right] {$P_{\parallel}$};
             \draw[->, dashed, red!60] (C) -- (0, -0.866) node[left] {$P_{\perp}$};
        \end{scope}
    \end{tikzpicture}
    \end{center}

\textbf{Equazioni del moto}:
Proiettando le forze sugli assi $x$ (parallelo al piano) e $y$ (normale):
\begin{itemize}
    \item \textbf{Asse $y$}: $N - mg \cos\theta = 0 \implies N = mg \cos\theta$
    \item \textbf{Asse $x$ (Corpo fermo)}: $mg \sin\theta - F_s = 0 \implies F_s = mg \sin\theta$
    \item \textbf{Asse $x$ (Corpo in moto)}: $mg \sin\theta \mp F_D = ma$
\end{itemize}

\textbf{Accelerazione in discesa}:
\[ a = g(\sin\theta - \mu_D \cos\theta) \]
Velocità finale dopo aver percorso una distanza $L$ partendo da fermo:
\[ v_f = \sqrt{2 L g (\sin\theta - \mu_D \cos\theta)} \]

\textbf{Accelerazione in salita}:
\[ a = g(\sin\theta + \mu_D \cos\theta) \]
Se lanciato con velocità $v_0$, lo spazio $S$ percorso prima di fermarsi è:
\[ S = \frac{v_0^2}{2g(\sin\theta + \mu_D \cos\theta)} \]
\end{esercizio}

\section{Sistemi con Forze di Tensione}
Le forze di tensione si manifestano in presenza di fili o funi tese. Esiste una differenza sostanziale tra \textbf{funi} (con massa non trascurabile) e \textbf{fili} (ideali, con massa trascurabile).

\subsection{Funi e Fili}
Si definisce \textbf{densità lineare di massa} $\mu = \frac{M}{L} [kg/m]$.

\begin{itemize}
    \item \textbf{FUNI}: Considerando un piccolo elemento di corda $\Delta l$ di massa $\Delta m = \mu \Delta l$:
    
    \begin{center}
    \begin{tikzpicture}
        \draw[thick, fill=gray!10] (0,-0.2) rectangle (3,0.2);
        \draw[->, thick, blue] (3,0) -- (4,0) node[right] {$\vec{T}_P$};
        \draw[->, thick, blue] (0,0) -- (-1,0) node[left] {$\vec{T}_Q$};
        \draw[<->] (0,-0.4) -- (3,-0.4) node[midway, below] {$\Delta l$};
        \node at (1.5, 0) {\small $\Delta m$};
    \end{tikzpicture}
    \end{center}
    
    Se il sistema si muove con accelerazione $\vec{a}$, per il II principio:
    \[ \vec{T}_P - \vec{T}_Q = \Delta m \cdot \vec{a} = (\mu \Delta l) \vec{a} \]
    La tensione non è costante lungo la fune.
    
    \item \textbf{FILI (Ideali)}: Se la massa della fune è trascurabile ($M \approx 0 \implies \mu \approx 0$), allora:
    \[ \vec{T}_P - \vec{T}_Q \approx \vec{0} \implies T_P = T_Q = T \]
    In un filo ideale la tensione è \textbf{costante} in ogni punto.
\end{itemize}

\subsection{Esempio: Sistema Uomo-Cassa}
Consideriamo un uomo che tira una cassa tramite una fune di massa $M$.
\begin{itemize}
    \item Eq. moto fune: $M \vec{a} = \vec{T}_D - \vec{T}_P$
    \item Eq. moto cassa: $m \vec{a} = -\vec{T}_P$
\end{itemize}
Se la fune è un filo ideale ($M=0$), allora $|\vec{F}_{uomo}| = |\vec{T}_P|$, ovvero l'uomo applica la forza direttamente alla cassa.

\subsection{Macchina di Atwood}
La macchina di Atwood consiste in due masse $m_1$ e $m_2$ collegate da un filo ideale che passa sopra una carrucola liscia di massa trascurabile.

\begin{figure}[ht]
    \centering
    \begin{tikzpicture}[scale=0.8]
        % Carrucola
        \draw[thick] (0,0) circle (0.5cm);
        \filldraw (0,0) circle (2pt);
        \draw[thick] (-2,2) -- (2,2);
        \fill[pattern=north east lines] (-2,2) rectangle (2,2.2);
        \draw[thick] (0,2) -- (0,0.5); % Sostegno carrucola
        
        % Filo e masse
        \draw[thick] (-0.5,0) -- (-0.5,-2);
        \draw[thick] (0.5,0) -- (0.5,-3);
        
        \draw[fill=gray!20] (-0.9,-2) rectangle (-0.1,-2.8) node[midway] {$m_1$};
        \draw[fill=gray!20] (0.1,-3) rectangle (0.9,-3.8) node[midway] {$m_2$};
        
        % Forze su m1
        \draw[->, thick, blue] (-0.5,-2) -- (-0.5,-1) node[left] {$\vec{T}$};
        \draw[->, thick, red] (-0.5,-2.4) -- (-0.5,-3.4) node[left] {$m_1\vec{g}$};
        
        % Forze su m2
        \draw[->, thick, blue] (0.5,-3) -- (0.5,-2) node[right] {$\vec{T}$};
        \draw[->, thick, red] (0.5,-3.4) -- (0.5,-4.4) node[right] {$m_2\vec{g}$};
        
        % Assi
        \draw[->, dashed, gray] (1.5, -4) -- (1.5, 0) node[above] {$y$};
    \end{tikzpicture}
    \caption{Schema della Macchina di Atwood.}
\end{figure}

Assumendo $m_1 > m_2$, impostiamo il sistema delle equazioni del moto:
\[
\begin{cases}
    m_1 g - T = m_1 a \\
    T - m_2 g = m_2 a
\end{cases}
\]
Sommando le due equazioni eliminiamo la tensione $T$:
\[ (m_1 - m_2)g = (m_1 + m_2)a \implies \boxed{a = g \frac{m_1 - m_2}{m_1 + m_2}} \]
Sostituendo $a$ in una delle equazioni, troviamo la tensione del filo:
\[ \boxed{T = \frac{2 m_1 m_2}{m_1 + m_2} g} \]

\section{Forze di Attrito Viscoso}
L'attrito viscoso si manifesta quando un corpo solido si muove all'interno di un \textbf{fluido} (liquido o gas). Si tratta di forze frenanti che dipendono dalla velocità del corpo.

\subsection{Legge Empirica e Regimi di Velocità}
In generale, la forza di attrito viscoso può essere espressa come:
\[ \vec{F}_v = -\vec{v} \cdot f(v) \]
Espandendo $f(v)$ in serie di potenze: $f(v) = \alpha + \beta v + \gamma v^2 + \dots$
\begin{itemize}
    \item \textbf{Regime di basse velocità}: $f(v) \approx \alpha$ (costante). La forza è proporzionale alla velocità: $\vec{F}_v = -\alpha \vec{v}$.
    \item \textbf{Regime di alte velocità}: La dipendenza diventa quadratica: $\vec{F}_v \approx -\gamma v^2 \hat{v}$.
\end{itemize}

Per una sfera di raggio $r$, il coefficiente $\alpha$ è dato dalla \textbf{Legge di Stokes}:
\[ \alpha = 6 \pi \eta r \]
dove $\eta$ è la viscosità del mezzo.

\subsection{Equazione del Moto e Velocità Limite}
Se applichiamo una forza costante $\vec{F}$ (es. la gravità) a un corpo in un mezzo viscoso:
\[ m \frac{dv}{dt} = F - \alpha v \implies \frac{dv}{dt} = \frac{F}{m} - \frac{\alpha}{m} v \]
Questa è un'equazione differenziale del I ordine non omogenea. La soluzione è:
\[ v(t) = \frac{F}{\alpha} + \left( v_0 - \frac{F}{\alpha} \right) e^{-t/\tau} \]
dove $\tau = \frac{m}{\alpha}$ è la \textbf{costante di tempo} del sistema.

\paragraph{Velocità di Saturazione (Limite):}
Per $t \to \infty$, il termine esponenziale decade e la velocità tende asintoticamente a un valore costante:
\begin{equation}
    v_L = \lim_{t \to \infty} v(t) = \frac{F}{\alpha}
\end{equation}
In questo stato la forza viscosa bilancia esattamente la forza applicata ($a=0$).

\begin{figure}[ht]
    \centering
    \begin{tikzpicture}
        \draw[->] (0,0) -- (5,0) node[right] {$t$};
        \draw[->] (0,0) -- (0,3) node[above] {$v(t)$};
        
        \draw[dashed, gray] (0,2) -- (5,2) node[right] {$v_L = F/\alpha$};
        
        % Caso v0 < vL
        \draw[thick, blue] plot[domain=0:4.5, samples=100] (\x, {2*(1 - exp(-\x))});
        \node[blue] at (1.5,1) {\small $v_0 < v_L$};
        
        % Caso v0 > vL
        \draw[thick, red] plot[domain=0:4.5, samples=100] (\x, {2 + 0.8*exp(-\x)});
        \node[red] at (1.5,2.6) {\small $v_0 > v_L$};
        
        \node[below left] at (0,0) {$0$};
    \end{tikzpicture}
    \caption{Andamento della velocità verso il valore di saturazione $v_L$.}
\end{figure}

\subsection{Calcolo dello Spazio Percorso}
Consideriamo il caso in cui $F=0$ (corpo lanciato con $v_0$ che viene frenato solo dal fluido):
\[ v(t) = v_0 e^{-t/\tau} \]
Integrando per trovare la posizione $x(t)$:
\[ x(t) = \int_0^t v(t') dt' = \int_0^t v_0 e^{-t'/\tau} dt' = v_0 \tau \left( 1 - e^{-t/\tau} \right) \]
Il corpo percorre una distanza finita prima di fermarsi idealmente (per $t \to \infty$):
\begin{equation}
    \Delta x_{total} = v_0 \tau = \frac{m v_0}{\alpha}
\end{equation}
\subsection{Analisi per Tempi Brevi (\texorpdfstring{$t \ll \tau$}{t << tau})}
Per tempi molto piccoli rispetto alla costante di tempo, possiamo approssimare l'esponenziale con lo sviluppo di Taylor al primo ordine: $e^{-t/\tau} \approx 1 - \frac{t}{\tau}$.
Sostituendo nella legge oraria della velocità:
\[ v(t) = \frac{F}{\alpha} + \left( v_0 - \frac{F}{\alpha} \right) \left( 1 - \frac{t}{\tau} \right) = \frac{F}{\alpha} + v_0 - \frac{F}{\alpha} - v_0\frac{t}{\tau} + \frac{F}{\alpha}\frac{t}{\tau} \]
Dunque:
\begin{equation}
    v(t) \approx v_0 + \left( \frac{F}{\alpha} - v_0 \right) \frac{t}{\tau}
\end{equation}
Questa espressione mostra che per tempi brevi la velocità varia linearmente, come in un moto uniformemente accelerato con accelerazione $a = \frac{F - \alpha v_0}{m}$.

\begin{tcolorbox}[colback=blue!5!white,colframe=blue!75!black,title=\textbf{Effetti della Forza sulla Velocità}]
Le modalità con cui una forza modifica il moto dipendono dalla sua orientazione rispetto al vettore velocità $\vec{v}_0$ istantaneo:
\begin{enumerate}
    \item \textbf{Forza Parallela} ($\vec{F} \parallel \vec{v}_0$): La forza modifica solo il \textbf{modulo} della velocità, non la sua direzione (moto unidimensionale).
    \item \textbf{Forza Perpendicolare} ($\vec{F} \perp \vec{v}_0$): La forza modifica solo la \textbf{direzione} della velocità, mantenendo il modulo costante (es. moto circolare uniforme).
\end{enumerate}
\end{tcolorbox}

\chapter{Lavoro ed Energia}
\section{Definizione di Lavoro}
L'attrito statico è essenziale affinché un veicolo o un uomo/animale possa spostarsi, poiché permette di sfruttare l'aderenza del piano. Tuttavia, da solo non basta; senza una "forza interna" da parte dell'uomo (muscoli) il movimento non potrebbe avvenire.

\textbf{LAVORO ($L$)}: Quando un punto materiale subisce uno spostamento sotto l'azione di una forza, si dice che quest'ultima compie un lavoro sul punto. 
Il lavoro può essere prodotto da una forza:
\begin{itemize}
    \item \textbf{UNIFORME}: Indipendente dalla posizione.
    \item \textbf{VARIA}: Dipendente dalla posizione.
\end{itemize}

\subsection{Caso di Forza Uniforme}
Il lavoro $L$ compiuto da una forza costante $\vec{F}$ durante uno spostamento rettilineo $\vec{s}$ è definito come il prodotto scalare:
\begin{equation}
    L = \vec{F} \cdot \vec{s} = F \cdot s \cdot \cos\theta
\end{equation}

Si distinguono tre casi fondamentali:
\begin{itemize}
    \item $L > 0$ se $0 \leq \theta < \pi/2$: \textbf{LAVORO MOTORE}.
    \item $L < 0$ se $\pi/2 < \theta \leq \pi$: \textbf{LAVORO RESISTENTE (DISSIPATIVO)}.
    \item $L = 0$ se $\theta = \pi/2$: La forza è perpendicolare allo spostamento ($\vec{F} \perp \vec{s}$).
\end{itemize}

\subsection{Caso di Forza Varia}
Per una forza che varia lungo la traiettoria $\gamma$, definiamo il lavoro infinitesimo $dL$ per uno spostamento $d\vec{s}$:
\begin{equation}
    dL = \vec{F} \cdot d\vec{s} = F \cdot ds \cdot \cos\theta
\end{equation}
Il lavoro totale è dato dall'integrale di linea lungo la curva $\gamma$:
\begin{equation}
    L^{(\gamma)} = \int_{R_{in}}^{R_F} \vec{F} \cdot d\vec{s} = \sum_{i=1}^n \Delta L_i \quad \text{per } n \to \infty
\end{equation}

In coordinate cartesiane:
\begin{equation}
    L^{(\gamma)} = \int_{R_{in}}^{R_F} (F_x dx + F_y dy + F_z dz) = L_x + L_y + L_z
\end{equation}
In caso unidimensionale lungo $x$: $L = \int_{x_{in}}^{x_F} F_x dx$, dove $F_x = F \cos\theta$.

\textbf{Unità di misura}: $[L] = [F] \cdot [s] = \text{Newton} \cdot \text{metro} = \text{Joule (J)}$.
$1\text{ J} = 1\text{ kg} \cdot \text{m}^2/\text{s}^2 = 1\text{ N} \cdot \text{m}$.

\section{Lavoro della Forza Peso}
La forza peso è $\vec{P} = m\vec{g} = mg(-\hat{y})$. Il lavoro compiuto dalla forza peso per uno spostamento da $\vec{R}_{in}$ a $\vec{R}_F$ su una traiettoria $\gamma$ è:
\begin{equation}
    L_P^{(\gamma)} = \int_{R_{in}}^{R_F} m\vec{g} \cdot d\vec{s} = m\vec{g} \cdot (\vec{R}_F - \vec{R}_{in}) = m g (y_{in} - y_F) = -mg\Delta h
\end{equation}
Notiamo che $\Delta h = y_F - y_{in}$ deve essere considerato con il suo segno. 
$L_P$ dipende solo da $\Delta h$, poiché $m$ e $g$ (vicino al suolo) rimangono costanti. In particolare:
\begin{itemize}
    \item $L_P > 0$ se $\Delta h < 0$: il corpo scende, $y_F < y_{in}$.
    \item $L_P < 0$ se $\Delta h > 0$: il corpo sale, $y_F > y_{in}$.
\end{itemize}
$\Delta h$ identifica lo \textbf{spostamento globale} lungo la verticale, non lo spazio effettivo percorso seguendo $\gamma$. Infatti $L_P$ non dipende dal percorso seguito ($L_{Pa} = L_{Pb}$).

\begin{figure}[ht]
    \centering
    \begin{tikzpicture}[scale=0.8]
        \draw[->] (0,0) -- (0,2.5) node[above] {$\hat{y}$};
        \draw[->] (0,0) -- (2.5,0) node[right] {$\hat{x}$};
        \node at (0.5,2) (A) [circle, fill, inner sep=1.5pt] {};
        \node at (2,0.5) (B) [circle, fill, inner sep=1.5pt] {};
        \draw[->, thick, blue] (A) .. controls (1,1.5) and (0.5,0.5) .. (B) node[midway, right] {$\gamma_a$};
        \draw[->, thick, red] (A) .. controls (1.5,2) and (2.5,1.5) .. (B) node[midway, right] {$\gamma_b$};
        \node[above] at (A) {$\vec{R}_{in}$};
        \node[right] at (B) {$\vec{R}_F$};
    \end{tikzpicture}
    \caption{Indipendenza del lavoro della forza peso dal cammino.}
\end{figure}

\subsection{Esempi di calcolo della velocità finale}
Vediamo come calcolare $v_f$ o $v_0$ in vari scenari:
\begin{enumerate}
    \item \textbf{Caduta libera}: Partendo da fermo ($v_{in}=0$), $v_f = \sqrt{2gh}$.
    \item \textbf{Lancio verso l'alto}: Per raggiungere quota $h$ con $v_f=0$, serve $v_0 = \sqrt{2gh}$.
    \item \textbf{Relazione tra $v^2$ ed $h$}: Come visto dagli esempi, $h = \frac{v^2}{2g}$ e $v^2 = 2gh$. Moltiplicando per $m/2$:
    \begin{equation}
        mgh = \frac{1}{2} m v^2
    \end{equation}
    Dove $\frac{1}{2}mv^2 = E_c$ è l'\textbf{Energia Cinetica} di un punto di massa $m$.
\end{enumerate}

\section{Teorema delle Forze Vive (Lavoro-Energia)}
Si definisce \textbf{Energia Cinetica} ($K$ o $E_c$) di un punto materiale:
\begin{equation}
    K = \frac{1}{2} m v^2
\end{equation}
Relazione fondamentale: $mgh = \frac{1}{2} m v^2 \implies v = \sqrt{2gh}$.

\begin{tcolorbox}[colback=green!5!white,colframe=green!75!black,title=\textbf{Teorema delle Forze Vive}]
Il lavoro compiuto dalla risultante delle forze agenti su un punto materiale è pari alla variazione della sua energia cinetica:
\begin{equation}
    L = \Delta K = K_F - K_{in}
\end{equation}
\end{tcolorbox}

\textit{Dimostrazione}:
\[ L = \int \vec{F} \cdot d\vec{s} = \int m \frac{d\vec{v}}{dt} \cdot \vec{v} dt = \int m \left( \frac{1}{2} \frac{d v^2}{dt} \right) dt = \left[ \frac{1}{2} m v^2 \right]_{in}^F \]

\section{Caso del Pendolo Semplice}
Consideriamo una massa $m$ appesa a un filo di lunghezza $R$. Indicando con $\phi$ l'angolo tra il filo e la verticale (positivo a destra):
\begin{equation}
    m\vec{a} = m\vec{g} + \vec{T} \implies m(\vec{a}_t + \vec{a}_n) = m\vec{g} + \vec{T}
\end{equation}
L'accelerazione in componenti polari ($R$ costante) è:
$\vec{a} = -R\dot{\phi}^2 \hat{n} + R\ddot{\phi} \hat{\tau}$. 
Scriviamo le equazioni del moto per $\hat{\tau}$ ed $\hat{n}$:
\begin{equation}
    \begin{cases}
        \hat{\tau}: m R \ddot{\phi} = -mg \sin\phi \implies \ddot{\phi} = -\frac{g}{R} \sin\phi \\
        \hat{n}: m \frac{v^2}{R} = T - mg \cos\phi \implies T = m \frac{v^2}{R} + mg \cos\phi
    \end{cases}
\end{equation}
Il punto $\phi = 0$ è detto \textbf{centro di oscillazione}.

\subsection{Analisi della Tensione e delle Piccole Oscillazioni}
Per la componente normale, deve essere $T \geq 0$ affinché il filo rimanga teso, altrimenti il corpo cadrebbe. Questo implica $v^2 \geq -R g \cos\phi$. Questo è sempre vero se $-\pi/2 \leq \phi \leq \pi/2$.

\textbf{Piccole oscillazioni}: Per $\phi \ll 1$, usiamo $\sin\phi \approx \phi$. L'equazione diventa $\ddot{\phi} + \frac{g}{R} \phi = 0$.
Definiamo $\omega_0^2 = \frac{g}{R}$, dove $\omega_0$ è la \textbf{pulsazione} (costante).
La soluzione è $\phi(t) = A \cos(\omega_0 t) + B \sin(\omega_0 t)$.
Se il pendolo parte da $\phi_0$ con velocità nulla ($\dot{\phi}=0$):
\begin{equation}
    \phi(t) = \phi_0 \cos(\omega_0 t)
\end{equation}
I punti $\phi = \pm \phi_0$ sono i \textbf{punti di inversione} del moto, dove la velocità $\dot{\phi}(t)$ cambia verso dopo essersi annullata e l'accelerazione $\ddot{\phi}(t)$ è massima e opposta allo spostamento ($\ddot{\phi} \propto -\phi$).

\begin{figure}[ht]
    \centering
    \begin{tikzpicture}[scale=1.2]
        \draw[thick] (-1.5,0) -- (1.5,0);
        \draw[dashed, gray] (0,0) -- (0,-2.5);
        \draw[thick] (0,0) -- (-0.8,-2.2) node[midway, left] {$R$};
        \filldraw (-0.8,-2.2) circle (2.5pt) node[below left] {$m$};
        \draw[thick] (0,0) -- (0.8,-2.2);
        \filldraw[gray!30] (0.8,-2.2) circle (2pt);
        \draw (0,-0.6) arc (-90:-110:0.6);
        \node at (-0.2,-0.8) {$\phi$};
        \draw[->, brown] (-0.8,-2.2) -- (-0.8,-3.2) node[below] {$m\vec{g}$};
        \draw[->, purple] (-0.8,-2.2) -- (-0.4,-1.1) node[right] {$\vec{T}$};
    \end{tikzpicture}
    \caption{Diagramma del pendolo semplice e forze agenti.}
\end{figure}

\section{Forze Elastiche}
L'elasticità è la capacità di tornare allo stato di equilibrio dopo una deformazione. La \textbf{Legge di Hooke} recita: $\vec{F}_{el} = -k \Delta \vec{R}$.
In un piano senza attrito, il moto è descritto da $m \ddot{x} = -k(x - x_0)$, da cui $\ddot{x} + \frac{k}{m}(x - x_0) = 0$.
Questa è un'equazione differenziale di II ordine lineare a coefficienti costanti non omogenea (ma riconducibile a omogenea spostando l'origine). La soluzione è:
\begin{equation}
    x(t) = x_0 + A \cos(\omega_0 t) + B \sin(\omega_0 t) \quad \text{con } \omega_0 = \sqrt{k/m}
\end{equation}

\subsection{Caso della Molla con Forza Costante}
Se la molla è verticale (soggetta a gravità), l'equazione del moto è $m \ddot{x} = mg - k (x - x_0)$.
Ricidendo l'equilibrio ($\ddot{x}=0$), troviamo il punto di equilibrio statico:
$x_{eq} = x_0 + \frac{mg}{k}$.
L'equazione diventa $\ddot{x} = -\frac{k}{m} (x - x_{eq})$.
La soluzione generale per il problema di Cauchy con condizioni iniziali a $t=0$ è:
\begin{equation}
    x(t) = x_{eq} + (x(0) - x_{eq}) \cos(\omega_0 t) + \frac{v(0)}{\omega_0} \sin(\omega_0 t)
\end{equation}
Per risolvere la situazione particolare dobbiamo considerare una situazione specifica (come quella di equilibrio) mentre per l'intera soluzione dobbiamo sfruttare le condizioni iniziali (\textbf{Teorema di Cauchy}).

\begin{definizione}{Energia Potenziale}
Si definisce \textbf{Energia Potenziale} ($U$) la capacità di un sistema di compiere lavoro in virtù della sua configurazione o posizione. La relazione fondamentale tra una forza conservativa e l'energia potenziale associata è:
\begin{equation}
    \vec{F} = -\nabla U
\end{equation}
\end{definizione}

\subsection{Relazione Differenziale e Gradiente}
Considerando un moto unidimensionale lungo l'asse $x$, vogliamo ricavare $F_x(x)$ conoscendo $U(x)$. Partendo dalla definizione di $\Delta U$ per uno spostamento da $x_0$ a $x$:
\begin{equation}
    U(x) - U(x_0) = -\int_{x_0}^x F_{x'}(x') dx'
\end{equation}
Derivando entrambi i membri rispetto a $x$ e applicando il \textbf{Teorema Fondamentale del Calcolo Integrale}:
\begin{equation}
    \frac{d}{dx} [U(x) - U(x_0)] = \frac{d}{dx} \left[ -\int_{x_0}^x F_{x'} dx' \right] \implies \frac{dU}{dx} = -F_x(x)
\end{equation}
Pertanto: $\boxed{F_x(x) = -\frac{dU}{dx}}$. In tre dimensioni, considerando le variazioni lungo ogni asse ($x, y, z$):
\begin{equation}
    F_x = -\frac{\partial U}{\partial x}; \quad F_y = -\frac{\partial U}{\partial y}; \quad F_z = -\frac{\partial U}{\partial z}
\end{equation}
Queste componenti possono essere riunite dal concetto di \textbf{Gradiente} ($\nabla$):
\begin{equation}
    \vec{F} = -\nabla U = -\text{grad } U = -\left( \frac{\partial U}{\partial x} \hat{x} + \frac{\partial U}{\partial y} \hat{y} + \frac{\partial U}{\partial z} \hat{z} \right)
\end{equation}

\subsection{Superfici Equipotenziali}
Una \textbf{superficie equipotenziale} è l'insieme dei punti dello spazio in cui l'energia potenziale assume lo stesso valore ($U(x,y,z) = \text{cost}$). 
Poiché su tale superficie $\Delta U = 0$, il lavoro compiuto dalla forza per spostamenti tangenti alla superficie è nullo. Ne consegue che:
\begin{itemize}
    \item La forza $\vec{F}$ è sempre \textbf{perpendicolare} alle superfici equipotenziali.
    \item La forza punta nella direzione di massima decrescita di $U$.
\end{itemize}

\begin{center}
\begin{tikzpicture}
    % Surface lines
    \draw[thick, blue!50] (0,0) ellipse (2cm and 1cm);
    \draw[thick, blue!50] (0,0) ellipse (2.5cm and 1.25cm);
    \draw[thick, blue!10] (0,0) ellipse (3cm and 1.5cm);
    \node[blue] at (2.2,0.8) {\small $U_1$};
    \node[blue] at (2.7,1.1) {\small $U_2$};
    
    % Force vector
    \draw[->, thick, red] (1.8,0.44) -- (1.2,0.3) node[left] {$\vec{F}$};
    \node at (2,0.2) {\small $U = \text{const}$};
\end{tikzpicture}
\end{center}

\subsection{Punti di Equilibrio e Buche di Potenziale}
L'analisi dell'andamento di $U(x)$ permette di identificare i punti di equilibrio ($F_x = -dU/dx = 0$):
\begin{enumerate}
    \item \textbf{Equilibrio Stabile}: Corrisponde a un \textbf{minimo} di $U$. Una piccola perturbazione genera una forza di richiamo verso il punto di equilibrio.
    \item \textbf{Equilibrio Instabile}: Corrisponde a un \textbf{massimo} di $U$. Una perturbazione allontana definitivamente il corpo.
    \item \textbf{Equilibrio Indifferente}: Corrisponde a una regione dove $U$ è costante.
\end{enumerate}

\begin{center}
\begin{tikzpicture}[scale=0.8]
    \draw[->] (0,0) -- (6,0) node[right] {$x$};
    \draw[->] (0,0) -- (0,4) node[above] {$U(x)$};
    \draw[thick, teal] plot [smooth, tension=1] coordinates { (0.5,3.5) (2,1) (4,3) (5.5,1.5)};
    
    \node at (2,1) [circle, fill, inner sep=1.5pt] {};
    \node[below] at (2,1) {\small Stabile (Min)};
    \node at (4,3) [circle, fill, inner sep=1.5pt] {};
    \node[above] at (4,3) {\small Instabile (Max)};
    
    % Ball illustrations
    \draw (2,1.2) circle (0.15);
    \draw (4,3.2) circle (0.15);
\end{tikzpicture}
\end{center}

\section{Campi di Forza ed Energia Meccanica}
Un \textbf{Campo} è l'insieme dei valori che una grandezza fisica assume in certe regioni dello spazio. Le forze conservative generano un campo conservativo dove la forza dipende dalla posizione del corpo.
Unendo i punti di applicazione delle forze si ottiene un \textbf{inviluppo} che prende il nome di \textbf{linea di forza} o di campo, a cui $\vec{F}$ è tangente.
\textbf{Proprietà delle linee di forza}:
\begin{enumerate}[label=\alph*)]
    \item Non possono incrociarsi.
    \item Non possono essere chiuse se il campo è conservativo (altrimenti $L = \oint \vec{F} \cdot d\vec{s} \neq 0$).
\end{enumerate}

\subsection{Energia Meccanica}
L'unione dell'energia cinetica $K$ e potenziale $U$ definisce l'\textbf{Energia Meccanica} ($E_{mecc}$):
\begin{equation}
    E_{mecc} = K + U = \frac{1}{2}mv^2 + U(\vec{r}) = \frac{p^2}{2m} + U(\vec{r})
\end{equation}

Dato che in presenza di sole forze conservative ($F_{con}$) l'energia non viene persa, mentre in presenza di forze non conservative ($F_{diss}$) sì, capiamo che la variazione $\Delta E_{mecc}$ è dovuta esclusivamente alle forze dissipative.

\begin{teorema}{Teorema delle Forze Vive (Energia Meccanica)}
Il lavoro totale è la somma del lavoro delle forze conservative e dissipative. $\Delta K + \Delta U = L_{diss}$, ovvero:
\[ \Delta E_{mecc} = L_{diss} \]
\end{teorema}

\subsection{Principio di Conservazione dell'Energia}
"L'energia non si crea né si distrugge, ma si trasforma da una forma all'altra":
\begin{itemize}
    \item Le forze conservative trasformano $K$ in $U$ e viceversa.
    \item Le forze dissipative trasformano $E_{mecc}$ in calore, luce, ecc.
\end{itemize}

\subsection{Teorema della Conservazione dell'Energia Meccanica}
In presenza di sole forze conservative, l'energia meccanica del sistema resta costante durante il moto:
\begin{equation}
    \Delta E_{mecc} = 0 \implies K_i + U_i = K_f + U_f \implies \Delta K = -\Delta U
\end{equation}

\subsection{Regioni Accessibili al Moto (1-D)}
Nei moti 1-D lungo $x$, la conservazione dell'energia implica $E_0 = \frac{1}{2}m\dot{x}^2 + U(x)$. Poiché $K \ge 0$, il moto è possibile solo nelle regioni dove $E_0 \ge U(x)$:
\begin{itemize}
    \item \textbf{Punti di inversione}: $\dot{x} = 0 \implies E_0 = U(x)$.
    \item \textbf{Regioni di moto}: $\dot{x} \neq 0 \implies E_0 > U(x)$.
    \item \textbf{Regioni proibite}: $E_0 < U(x)$ (richiederebbero $K < 0$, impossibile).
\end{itemize}

\begin{center}
\begin{tikzpicture}[scale=0.9]
    \draw[->] (0,0) -- (6,0) node[right] {$x$};
    \draw[->] (0,0) -- (0,4) node[above] {$U(x)$};
    \draw[thick, teal] plot [smooth, tension=1] coordinates { (0.5,3.5) (2,1) (4,3) (5.5,1.5)};
    
    % Energy level
    \draw[dashed, red, thick] (0,2.2) -- (6,2.2) node[right] {$E_0$};
    
    % Intersection points
    \filldraw[red] (1,2.2) circle (2pt) node[above left] {$x_{inv}$};
    \filldraw[red] (2.9,2.2) circle (2pt);
    \filldraw[red] (5.1,2.2) circle (2pt);
    
    \node at (2,0.5) {\small Moto permesso};
    \node at (4,3.5) {\small Proibita};
\end{tikzpicture}
\end{center}

\section{Dinamica dei Sistemi di Punti Materiali}
Un \textbf{sistema materiale} è un insieme di $N$ punti materiali. Può essere:
\begin{itemize}
    \item \textbf{Discreto}: Punti separati con masse $m_i$.
    \item \textbf{Continuo}: Distribuzione continua di massa con densità $\rho$.
\end{itemize}

\begin{definizione}{Centro di Massa (CM)}
Il centro di massa è il punto in cui immaginiamo concentrata tutta la massa del sistema.
\begin{itemize}
    \item \textbf{Sistema Discreto}: $\vec{R}_{CM} = \frac{\sum m_i \vec{r}_i}{M}$ con $M = \sum m_i$.
    \item \textbf{Sistema Continuo}: $\vec{R}_{CM} = \frac{1}{M} \int \vec{r} dm$, dove $dm = \rho dV$.
\end{itemize}
\end{definizione}

\subsection{Sistema di Riferimento del Centro di Massa}
Il CM può essere usato come origine di un sistema di riferimento rispetto al quale i punti si esprimono con \textbf{coordinate relative} $\vec{r}_i'$:
\begin{equation}
     S \to S' : \vec{r}_i' = \vec{r}_i - \vec{R}_{CM} \iff \vec{r}_i = \vec{r}_i' + \vec{R}_{CM}
\end{equation}
Rispetto a questo sistema, le coordinate del CM sono banalmente $(0,0,0)$. Una proprietà fondamentale è che la sommatoria del momento del primo ordine rispetto al CM è nulla:
\begin{equation}
    \sum m_i \vec{r}_i' = \sum m_i (\vec{r}_i - \vec{R}_{CM}) = \underbrace{\sum m_i \vec{r}_i}_{M \vec{R}_{CM}} - \underbrace{\sum m_i}_{M} \vec{R}_{CM} = \vec{0}
\end{equation}

\subsection{Grandezze Cinetiche del CM}
Seguendo la stessa logica di trasformazione delle coordinate, definiamo le grandezze derivate:
\begin{itemize}
    \item \textbf{Velocità}: $\vec{V}_{CM} = \frac{d\vec{R}_{CM}}{dt} = \frac{1}{M} \sum m_i \vec{v}_i$
    \item \textbf{Quantità di Moto}: $\vec{P}_{tot} = M \vec{V}_{CM} = \sum m_i \vec{v}_i$
    \item \textbf{Accelerazione}: $\vec{A}_{CM} = \frac{d\vec{V}_{CM}}{dt} = \frac{1}{M} \sum m_i \vec{a}_i$
\end{itemize}
Si noti che $\vec{v}_i = \vec{v}_i' + \vec{V}_{CM}$ e $\vec{a}_i = \vec{a}_i' + \vec{A}_{CM}$.

\begin{teorema}{Prima Equazione Cardinale}
Per ogni punto del sistema vale $\vec{F}_i^{ext} + \vec{F}_i^{int} = m_i \vec{a}_i$. Sommando su tutti gli $N$ punti:
\begin{equation}
    \sum_{i=1}^N \vec{F}_i^{ext} + \underbrace{\sum_{i=1}^N \vec{F}_i^{int}}_{\vec{0}} = \sum_{i=1}^N m_i \vec{a}_i = M \vec{A}_{CM}
\end{equation}
Risulta quindi: 
\begin{equation}
    \vec{F}_{tot}^{ext} = \frac{d\vec{P}_{tot}}{dt} = M \vec{A}_{CM}
\end{equation}
\end{teorema}
\textbf{Conseguenza}: Se $\sum \vec{F}^{ext} = \vec{0}$, allora $\vec{V}_{CM} = \text{costante}$. Il CM si muove di moto rettilineo uniforme o è in quiete.

\subsection{Massa Ridotta ($\mu$)}
La \textbf{Massa Ridotta} è definita come la sommatoria del reciproco dei reciproci delle masse del sistema. Per due corpi $m_1$ e $m_2$:
\begin{equation}
    \mu = \left( \frac{1}{m_1} + \frac{1}{m_2} \right)^{-1} = \frac{m_1 m_2}{m_1 + m_2}
\end{equation}
\textbf{Casi Particolari}:
\begin{itemize}
    \item Se $m_1 = m_2 = m \implies \mu = \frac{m^2}{2m} = \frac{m}{2}$ (es: sistemi $e^+e^-$, stelle doppie).
    \item Se $m_1 \gg m_2 \implies \mu = \frac{m_1 m_2}{m_1 + m_2} \approx \frac{m_1 m_2}{m_1} = m_2$ (es: sistema Terra-Sole).
\end{itemize}

\begin{tcolorbox}[colback=yellow!5!white,colframe=yellow!75!black,title=\textbf{Esempio: Centro di Massa di 3 punti}]
Siano tre punti di massa $m$ posti in $A(a,0,0)$, $B(0,a,0)$ e $C(0,0,0)$.
Le coordinate del CM sono:
\[ \vec{R}_{CM} = \frac{m(a,0,0) + m(0,a,0) + m(0,0,0)}{3m} = \frac{m(a, a, 0)}{3m} = \left( \frac{a}{3}, \frac{a}{3}, 0 \right) \]
\end{tcolorbox}

\subsection{Sistema di Due Punti}
Consideriamo un sistema composto da due punti $A$ e $B$ con masse $m_A, m_B$. Definiamo la \textbf{coordinata relativa} $\vec{r}_r = \vec{r}_A - \vec{r}_B$. Possiamo ricavare le posizioni assolute partendo dalla definizione di centro di massa:
\[ \vec{R}_{CM} = \frac{m_A \vec{r}_A + m_B \vec{r}_B}{m_A + m_B} = \frac{m_A (\vec{r}_r + \vec{r}_B) + m_B \vec{r}_B}{M} = \frac{m_A \vec{r}_r + M \vec{r}_B}{M} \]
Da cui otteniamo le relazioni fondamentali:
\begin{equation}
    \vec{r}_B = \vec{R}_{CM} - \frac{m_A}{M} \vec{r}_r, \quad \vec{r}_A = \vec{R}_{CM} + \frac{m_B}{M} \vec{r}_r
\end{equation}

\subsubsection{Caso 1: Sistema ISOLATO}
Se il sistema è isolato ($\sum \vec{F}_i^{ext} = \vec{0}$), agiscono solo le forze interne $\vec{F}_{AB} = -\vec{F}_{BA}$. Considerando il moto di $A$:
\begin{equation}
    \vec{F}_{BA} = m_A \ddot{\vec{r}}_A = m_A \left( \ddot{\vec{R}}_{CM} + \frac{m_B}{M} \ddot{\vec{r}}_r \right)
\end{equation}
Poiché $\ddot{\vec{R}}_{CM} = \vec{0}$ in un sistema isolato (per la I Eq. Cardinale), risulta:
\begin{equation}
    \vec{F}_{BA} = \frac{m_A m_B}{M} \ddot{\vec{r}}_r = \mu \ddot{\vec{r}}_r, \quad \vec{F}_{AB} = -\mu \ddot{\vec{r}}_r
\end{equation}
La dinamica interna dipende quindi solo dalla massa ridotta e dall'accelerazione relativa.

\subsubsection{Caso 2: Sistema NON ISOLATO}
In presenza di forze esterne, sommiamo le equazioni del moto dei due punti:
\begin{equation}
    m_A \ddot{\vec{r}}_A + m_B \ddot{\vec{r}}_B = (\vec{F}_{BA} + \vec{F}_A^{ext}) + (\vec{F}_{AB} + \vec{F}_B^{ext}) = \vec{F}_A^{ext} + \vec{F}_B^{ext}
\end{equation}
Ricordando che la somma delle masse moltiplicata per l'accelerazione del CM è $M \ddot{\vec{R}}_{CM}$, ritroviamo la prima equazione cardinale: $\sum \vec{F}^{ext} = M \vec{A}_{CM}$.

\begin{teorema}{I Teorema di Koenig}
Il teorema di Koenig permette di separare il contributo del moto del centro di massa dal moto relativo delle parti interne del sistema.
\begin{equation}
    \boxed{K = \frac{1}{2} M V_{CM}^2 + K'} \quad \text{con } K' = \sum_{i=1}^N \frac{1}{2} m_i (v_i')^2
\end{equation}
\end{teorema}
Dove $K_{CM} = \frac{1}{2} M V_{CM}^2$ è l'energia del CM e $K'$ è l'energia dei corpi \textbf{nel sistema del centro di massa}.
Dove:
\begin{itemize}
    \item $K_{CM} = \frac{1}{2} M V_{CM}^2$ è l'energia cinetica associata al moto del centro di massa.
    \item $K'$ è l'energia cinetica interna (rispetto al CM).
\end{itemize}

\textbf{Nota}: Se il sistema è isolato, $\vec{V}_{CM}$ è costante e quindi $K_{CM}$ è una costante del moto. Il bilancio energetico dipende esclusivamente dalle variazioni di $K'$.

\subsection{Caso del Sistema a Due Corpi}
Per un sistema con due corpi, l'energia cinetica relativa $K'$ può essere espressa in funzione della velocità relativa $\vec{v}_r = \vec{v}_A - \vec{v}_B$:
\begin{equation}
    K' = \frac{1}{2} (m_A (\vec{v}_A')^2 + m_B (\vec{v}_B')^2) = \dots = \frac{1}{2} \mu v_r^2
\end{equation}
Sostituendo tutto nel teorema di Koenig:
\begin{equation}
    K = \frac{1}{2} M V_{CM}^2 + \frac{1}{2} \mu v_r^2
\end{equation}

\section{Teorema delle Forze Vive (Sistemi)}
"La variazione di energia cinetica totale di un sistema materiale è pari al lavoro compiuto da tutte le forze agenti: interne, esterne ed eventualmente apparenti."
\begin{equation}
    \Delta K = L_{tot} = L_{con} + L_{diss}
\end{equation}

\section{Teorema della Conservazione dell'Energia Meccanica}
"Se tutte le forze agenti sul sistema materiale sono conservative, l'energia meccanica totale è costante."
In presenza di forze dissipative:
\begin{equation}
    \Delta K + \Delta U = \Delta E_{mecc} = L_{diss}
\end{equation}

\chapter{Dinamica degli Urti}

Lo studio degli urti è il primo caso in cui la dinamica dei sistemi di punti materiali mostra tutta la sua potenza: durante l'interazione le forze in gioco sono così intense e così brevi da risultare impossibili da misurare istante per istante, eppure lo stato finale del sistema resta completamente determinato dalle sole \textbf{leggi di conservazione}. È proprio questa la lezione del capitolo: quando la descrizione dettagliata della forza è inaccessibile, i teoremi di conservazione forniscono ugualmente la risposta.

\begin{definizione}{Urto}
In fisica, un \textbf{urto} è un'interazione fra corpi (generalmente due) che si compie in un intervallo di tempo brevissimo, $\Delta t \to 0$. Durante tale intervallo avviene uno scambio significativo di energia cinetica $K$ e di quantità di moto $\vec{P}$.
\end{definizione}

La brevità dell'interazione non è un dettaglio accessorio ma l'ipotesi centrale di tutta la trattazione: è da essa che discenderanno sia l'approssimazione di sistema isolato, sia la possibilità di separare nettamente uno ``stato prima'' e uno ``stato dopo'' senza doversi occupare di ciò che accade nel mezzo.

\begin{figure}[ht]
    \centering
    \begin{tikzpicture}[>=stealth, scale=1.0,
        corpoA/.style={fill=boxoss!35, draw=boxoss, thick},
        corpoB/.style={fill=notered!35, draw=notered, thick}]

        % ---------------- PRIMA ----------------
        \begin{scope}
            \node[font=\small\bfseries, mypen] at (2.1,1.75) {PRIMA};
            \draw[thin, mypen!35] (0.0,0) -- (4.2,0);
            \filldraw[corpoA] (0.6,0) circle (0.24);
            \node[font=\small, mypen, below] at (0.6,-0.42) {$m_A$};
            \draw[->, very thick, boxoss] (1.0,0) -- (1.9,0)
                node[midway, above, font=\small, boxoss] {$\vec{v}_{Ai}$};
            \filldraw[corpoB] (3.6,0) circle (0.28);
            \node[font=\small, mypen, below] at (3.6,-0.46) {$m_B$};
            \draw[->, very thick, notered] (3.2,0) -- (2.3,0)
                node[midway, above, font=\small, notered] {$\vec{v}_{Bi}$};
        \end{scope}

        % ---------------- ZONA D'URTO ----------------
        \draw[decorate, decoration={zigzag, amplitude=2.6pt, segment length=5pt},
              very thick, boxatt] (4.5,0) -- (5.7,0);
        \node[font=\small\bfseries, boxatt, align=center] at (5.1,1.35)
            {urto\\[-2pt] $\Delta t \to 0$};

        % ---------------- DOPO ----------------
        \begin{scope}[xshift=6.9cm]
            \node[font=\small\bfseries, mypen] at (2.1,1.75) {DOPO};
            \draw[thin, mypen!35] (0.0,0) -- (4.2,0);
            \filldraw[corpoA] (1.1,0) circle (0.24);
            \node[font=\small, mypen, below] at (1.1,-0.42) {$m_A$};
            \draw[->, very thick, boxoss] (0.7,0) -- (-0.2,0)
                node[midway, above, font=\small, boxoss] {$\vec{v}_{Af}$};
            \filldraw[corpoB] (3.2,0) circle (0.28);
            \node[font=\small, mypen, below] at (3.2,-0.46) {$m_B$};
            \draw[->, very thick, notered] (3.6,0) -- (4.5,0)
                node[midway, above, font=\small, notered] {$\vec{v}_{Bf}$};
        \end{scope}
    \end{tikzpicture}
    \caption{Struttura logica di un urto. L'interazione vera e propria, confinata in un intervallo $\Delta t \to 0$, viene trattata come una scatola nera: si confrontano soltanto lo stato iniziale e lo stato finale. La quantità di moto totale si conserva in ogni caso, $\vec{P}_f = \vec{P}_i$; l'energia cinetica soltanto se l'urto è elastico.}
    \label{fig:schema_urto}
\end{figure}

\section{Forze Impulsive}

Lo scambio di $K$ e $\vec{P}$ è dovuto all'azione di forze estremamente intense che agiscono per un tempo limitatissimo, dette \textbf{forze impulsive}. Di una forza siffatta non conosciamo quasi mai l'andamento istantaneo $\vec{F}(t)$ --- che dipende dai dettagli microscopici della deformazione dei corpi a contatto --- ma non ne abbiamo bisogno: l'unica informazione dinamicamente rilevante è il suo \textbf{integrale} nel tempo.

\begin{definizione}{Impulso di una Forza}
Si definisce \textbf{impulso} $\vec{I}$ della forza $\vec{F}$ agente nell'intervallo dell'urto la quantità
\begin{equation}
    \vec{I} = \int_{0^-}^{0^+} \vec{F}(t) \, dt = \vec{P}_f - \vec{P}_i = \Delta \vec{P}
\end{equation}
Geometricamente, $\vec{I}$ è l'\textbf{area sottesa} dalla curva $F(t)$ nell'intervallo di contatto.
\end{definizione}

L'uguaglianza $\vec{I} = \Delta\vec{P}$ non è altro che il secondo principio della dinamica riscritto in forma integrale, ed è il vero contenuto operativo del concetto di forza impulsiva: due forze molto diverse fra loro, purché sottendano la stessa area, producono sullo stesso corpo la medesima variazione di quantità di moto.

\begin{figure}[ht]
    \centering
    \begin{tikzpicture}[>=stealth, scale=1.0]
        % Assi
        \draw[->, thick, mypen] (-0.3,0) -- (6.6,0) node[right, font=\small] {$t$};
        \draw[->, thick, mypen] (0,-0.1) -- (0,3.9) node[above, font=\small] {$F(t)$};

        % Area sottesa
        \fill[notered!18, domain=0.5:5.5, samples=140]
            plot (\x, {3.1*exp(-2.0*(\x-3)^2)}) -- (5.5,0) -- (0.5,0) -- cycle;
        % Curva impulsiva
        \draw[ultra thick, notered, domain=0.5:5.5, samples=140]
            plot (\x, {3.1*exp(-2.0*(\x-3)^2)});

        % Etichetta della curva
        \node[font=\small\bfseries, notered, right] at (4.6,3.45) {forza impulsiva};
        \draw[thin, notered!70] (4.55,3.35) -- (3.55,2.55);

        % Etichetta dell'area
        \node[font=\small, mypen, left] at (1.55,1.55) {area $= |\vec{I}| = |\Delta \vec{P}|$};
        \draw[thin, mypen!70] (1.6,1.5) -- (2.55,0.7);

        % Intervallo di contatto
        \draw[dashed, thin, mypen!55] (2.1,0) -- (2.1,-0.95);
        \draw[dashed, thin, mypen!55] (3.9,0) -- (3.9,-0.95);
        \draw[<->, thick, mypen] (2.1,-0.8) -- (3.9,-0.8);
        \node[font=\small, mypen, below] at (3.0,-0.88) {$\Delta t \to 0$};

        % Forza esterna ordinaria, per confronto
        \draw[thick, boxese, dashed] (0.3,0.22) -- (5.2,0.22);
        \node[font=\scriptsize, boxese, above right] at (5.2,0.2) {$F_{\text{ext}}$};
    \end{tikzpicture}
    \caption{Andamento temporale di una forza impulsiva. L'area sottesa dalla curva è l'impulso $\vec{I} = \Delta\vec{P}$. In tratteggio verde, per confronto, una forza esterna ordinaria: il suo modulo è incomparabilmente più piccolo e, sull'intervallo $\Delta t \to 0$, l'area che essa sottende è trascurabile.}
    \label{fig:forza_impulsiva}
\end{figure}

Per il bilancio energetico dell'urto useremo invece i due teoremi generali già stabiliti per i sistemi di punti materiali, che qui richiamiamo perché saranno lo strumento con cui distingueremo i vari tipi di urto.

\begin{teorema}{Teorema delle Forze Vive (Sistemi)}
La variazione di energia cinetica totale di un sistema materiale è pari al lavoro compiuto da tutte le forze agenti: interne, esterne ed eventualmente apparenti,
\begin{equation}
    \Delta K = L_{\text{tot}} = L_{\text{ext}} + L_{\text{int}} \; (+ \, L_{\text{app}})
\end{equation}
\end{teorema}

\begin{teorema}{Conservazione dell'Energia Meccanica}
Se tutte le forze agenti sul sistema materiale sono conservative, l'energia meccanica totale è costante. In presenza di forze dissipative,
\begin{equation}
    \Delta E_{\text{mecc}} = \Delta (K + U) = L_{\text{diss}}
\end{equation}
\end{teorema}

\begin{osservazione}{Perché l'Energia Cinetica Può Non Conservarsi}
Nel caso degli urti il termine decisivo è $L_{\text{int}}$, il lavoro delle forze interne. Se le forze di contatto sono \emph{conservative} (deformazione puramente elastica), esse restituiscono integralmente il lavoro assorbito durante la compressione e $L_{\text{int}} = 0$ sull'intero urto: l'energia cinetica si conserva. Se invece una parte del lavoro viene spesa in deformazioni permanenti, calore o vibrazioni, allora $L_{\text{int}} < 0$ e l'energia cinetica finale è minore di quella iniziale. La quantità di moto, al contrario, è \textbf{indifferente} a questa distinzione: le forze interne sono a due a due uguali e opposte per il terzo principio, e il loro contributo a $\Delta\vec{P}$ si annulla sempre.
\end{osservazione}

\subsection{Approssimazione di Sistema Isolato}

Durante l'urto possono essere presenti anche forze esterne --- gravità, attrito, reazioni vincolari --- e in linea di principio esse violerebbero la conservazione della quantità di moto. Il loro modulo è però trascurabile rispetto a quello delle forze impulsive, $F_{\text{ext}} \ll F_{\text{imp}}$, e soprattutto esse agiscono su un intervallo che stiamo facendo tendere a zero. Ne consegue che il loro impulso è nullo nel limite $\Delta t \to 0$:
\begin{equation}
    \vec{I}_{\text{ext}} = \int_{\Delta t} \sum \vec{F}^{\,\text{ext}} \, dt \approx \vec{0}
\end{equation}

\begin{attenzione}{È l'Impulso a Essere Trascurabile, non la Forza}
La gravità non ``sparisce'' durante l'urto: continua ad agire con il suo valore abituale. Ciò che si annulla è il suo \emph{impulso}, perché una forza di modulo finito integrata su un tempo infinitesimo dà un contributo infinitesimo. Subito prima e subito dopo l'urto le forze esterne tornano invece a governare il moto in modo del tutto ordinario.
\end{attenzione}

Pertanto il sistema può essere considerato \textbf{isolato} nell'istante dell'urto, e ciò implica la conservazione della quantità di moto totale:
\begin{equation}
    \sum \vec{F}^{\,\text{ext}} \approx \vec{0}
    \implies \Delta \vec{P}_{\text{tot}} = \vec{0}
    \implies \vec{P}_{\text{tot}} = \text{costante}
\end{equation}

\begin{formulario}{Le Due Leggi che Governano Ogni Urto}
\begin{minipage}[t]{0.47\textwidth}
\centering \textbf{Sempre valida}
\begin{equation*}
    \vec{P}_i = \vec{P}_f
\end{equation*}
\centering \footnotesize (le forze interne non alterano $\vec{P}_{\text{tot}}$)
\end{minipage}
\hfill
\begin{minipage}[t]{0.47\textwidth}
\centering \textbf{Valida solo se elastico}
\begin{equation*}
    K_i = K_f
\end{equation*}
\centering \footnotesize (nessuna dissipazione interna)
\end{minipage}
\vspace{4pt}
\end{formulario}

\section{Sistemi di Riferimento per lo Studio degli Urti}

Per analizzare un urto è utile sfruttare due diversi sistemi di riferimento inerziali, scegliendo di volta in volta quello in cui i conti risultano più semplici.

\begin{itemize}
    \item \textbf{Sistema del Laboratorio}: il sistema inerziale ``fermo'' in cui osserviamo le velocità iniziali e finali dei corpi nel nostro apparato sperimentale. È il sistema in cui i dati vengono effettivamente misurati.
    \item \textbf{Sistema del Centro di Massa (CM)}: il sistema solidale con il centro di massa, in moto con velocità $\vec{V}_{CM}$ rispetto al laboratorio. Poiché il sistema è isolato durante l'urto, $\vec{V}_{CM}$ è costante e anche questo è un sistema di riferimento inerziale.
\end{itemize}

Nel sistema del CM la quantità di moto totale è nulla per definizione, in ogni istante, e dunque
\begin{equation}
    \vec{p}_{Ai}\,' = -\,\vec{p}_{Bi}\,'
    \qquad \text{e} \qquad
    \vec{p}_{Af}\,' = -\,\vec{p}_{Bf}\,'
\end{equation}
I due corpi si muovono cioè sempre con quantità di moto uguali e opposte: il numero di incognite indipendenti si dimezza e il problema, come vedremo, si riduce a una semplice questione di moduli.

\begin{figure}[ht]
    \centering
    \begin{tikzpicture}[>=stealth, scale=1.0,
        corpoA/.style={fill=boxoss!35, draw=boxoss, thick},
        corpoB/.style={fill=notered!35, draw=notered, thick}]

        % ---------------- LAB ----------------
        \begin{scope}
            \node[font=\small\bfseries, mypen, right] at (0.0,1.35) {LAB};
            \draw[thin, mypen!35] (0.2,0) -- (6.4,0);
            \filldraw[corpoA] (1.0,0) circle (0.23);
            \node[font=\small, mypen, above] at (1.0,0.4) {$m_A$};
            \draw[->, very thick, boxoss] (1.35,0) -- (2.45,0)
                node[midway, below, font=\small, boxoss] {$\vec{v}_{Ai}$};
            \filldraw[corpoB] (4.6,0) circle (0.27);
            \node[font=\small, mypen, above] at (4.6,0.44) {$m_B$};
            \draw[->, very thick, notered] (4.95,0) -- (5.65,0)
                node[midway, below, font=\small, notered] {$\vec{v}_{Bi}$};
            \draw[->, thick, boxese] (2.9,-1.0) -- (3.9,-1.0);
            \node[font=\small, boxese, right] at (3.98,-1.0) {$\vec{V}_{CM}$};
        \end{scope}

        % ---------------- CM ----------------
        \begin{scope}[yshift=-3.3cm]
            \node[font=\small\bfseries, mypen, right] at (0.0,1.35) {CM};
            \draw[thin, mypen!35] (0.2,0) -- (6.4,0);
            \filldraw[corpoA] (1.6,0) circle (0.23);
            \node[font=\small, mypen, above] at (1.6,0.4) {$m_A$};
            \draw[->, very thick, boxoss] (1.95,0) -- (2.85,0)
                node[midway, below, font=\small, boxoss] {$\vec{p}_A\,'$};
            \filldraw[corpoB] (5.0,0) circle (0.27);
            \node[font=\small, mypen, above] at (5.0,0.44) {$m_B$};
            \draw[->, very thick, notered] (4.65,0) -- (3.75,0)
                node[midway, below, font=\small, notered] {$\vec{p}_B\,'$};
            \fill[boxese] (3.3,0) circle (2.6pt);
            \node[font=\scriptsize, boxese, above] at (3.3,0.14) {CM fermo};
            \node[font=\small, mypen] at (3.3,-1.0) {$\vec{p}_A\,' + \vec{p}_B\,' = \vec{0}$};
        \end{scope}
    \end{tikzpicture}
    \caption{Lo stesso urto visto dal laboratorio e dal centro di massa. Nel LAB le due quantità di moto sono indipendenti; nel CM sono per costruzione uguali e opposte, e il centro di massa resta immobile prima, durante e dopo l'urto.}
    \label{fig:lab_vs_cm}
\end{figure}

\section{Tipologie di Urto e Coefficiente di Restituzione}

Gli urti si classificano in base alla conservazione dell'\textbf{energia cinetica interna} $K'$, cioè dell'energia cinetica valutata nel sistema del centro di massa. La scelta di $K'$ e non di $K$ non è arbitraria: come mostreremo fra poco, $K'$ è l'unica aliquota di energia che un urto può effettivamente dissipare.

\begin{enumerate}
    \item \textbf{Urto elastico}: l'energia cinetica interna si conserva perfettamente, $\Delta K' = 0$.
    \item \textbf{Urto totalmente anelastico}: l'energia cinetica interna si dissipa completamente e i corpi restano uniti dopo l'urto.
    \item \textbf{Urto parzialmente anelastico}: l'energia cinetica interna si dissipa solo in parte, $0 < K'_f < K'_i$.
\end{enumerate}

\begin{definizione}{Coefficiente di Restituzione ($e$)}
Il \textbf{coefficiente di restituzione} $e$ è il rapporto fra i moduli della velocità relativa dopo e prima dell'urto, ovvero la radice quadrata del rapporto fra le energie cinetiche interne:
\begin{equation}
    e = \frac{|\vec{v}_{rf}|}{|\vec{v}_{ri}|}
    \qquad \Longleftrightarrow \qquad
    e^2 = \frac{K'_{f}}{K'_{i}}
\end{equation}
Esso è un \textbf{parametro sperimentale} che riassume in un solo numero il comportamento dei materiali a contatto, senza richiedere alcuna conoscenza della forza impulsiva.
\end{definizione}

In base al valore numerico di $e$ si distinguono i tre casi:
\begin{itemize}
    \item $e = 1$: urto \textbf{elastico}, con $K'_{f} = K'_{i}$;
    \item $e = 0$: urto \textbf{totalmente anelastico}, con $K'_{f} = 0$;
    \item $0 < e < 1$: urto \textbf{parzialmente anelastico}, con $K'_{f} < K'_{i}$.
\end{itemize}

\begin{figure}[ht]
    \centering
    \begin{tikzpicture}[>=stealth, scale=1.0,
        pallina/.style={fill=boxoss!35, draw=boxoss, thick}]

        % ---------------- (a) elastico ----------------
        \begin{scope}
            \node[font=\small\bfseries, boxese] at (2.0,3.15) {elastico};
            \node[font=\small, boxese] at (2.0,2.75) {$e = 1$};
            \draw[line width=1.1pt, mypen!70] (0.35,0) -- (3.65,0);
            \fill[pattern=north east lines, pattern color=mypen!35]
                (0.35,-0.24) rectangle (3.65,0);
            \draw[dashed, thin, mypen!50] (1.25,2.05) -- (1.25,0.22);
            \filldraw[pallina] (1.25,2.05) circle (0.2);
            \draw[->, very thick, boxoss] (0.85,1.85) -- (0.85,0.3)
                node[midway, left, font=\scriptsize] {$v_{ri}$};
            \draw[->, very thick, boxese] (2.75,0.3) -- (2.75,1.85)
                node[midway, right, font=\scriptsize] {$v_{rf} = v_{ri}$};
        \end{scope}

        % ---------------- (b) parzialmente anelastico ----------------
        \begin{scope}[xshift=4.9cm]
            \node[font=\small\bfseries, boxatt] at (2.0,3.15) {parz. anelastico};
            \node[font=\small, boxatt] at (2.0,2.75) {$0 < e < 1$};
            \draw[line width=1.1pt, mypen!70] (0.35,0) -- (3.65,0);
            \fill[pattern=north east lines, pattern color=mypen!35]
                (0.35,-0.24) rectangle (3.65,0);
            \draw[dashed, thin, mypen!50] (1.25,2.05) -- (1.25,0.22);
            \filldraw[pallina] (1.25,2.05) circle (0.2);
            \draw[->, very thick, boxoss] (0.85,1.85) -- (0.85,0.3)
                node[midway, left, font=\scriptsize] {$v_{ri}$};
            \draw[->, very thick, boxatt] (2.75,0.3) -- (2.75,1.05)
                node[midway, right, font=\scriptsize] {$v_{rf} < v_{ri}$};
        \end{scope}

        % ---------------- (c) totalmente anelastico ----------------
        \begin{scope}[xshift=9.8cm]
            \node[font=\small\bfseries, boxdef] at (2.0,3.15) {tot. anelastico};
            \node[font=\small, boxdef] at (2.0,2.75) {$e = 0$};
            \draw[line width=1.1pt, mypen!70] (0.35,0) -- (3.65,0);
            \fill[pattern=north east lines, pattern color=mypen!35]
                (0.35,-0.24) rectangle (3.65,0);
            \draw[dashed, thin, mypen!50] (1.25,2.05) -- (1.25,0.22);
            \filldraw[pallina] (1.25,2.05) circle (0.2);
            \draw[->, very thick, boxoss] (0.85,1.85) -- (0.85,0.3)
                node[midway, left, font=\scriptsize] {$v_{ri}$};
            \filldraw[pallina] (2.75,0.2) circle (0.2);
            \node[font=\scriptsize, boxdef, right] at (3.0,0.42) {$v_{rf} = 0$};
        \end{scope}
    \end{tikzpicture}
    \caption{Lettura sperimentale del coefficiente di restituzione tramite il rimbalzo di un corpo su un piano fisso. Il rapporto fra la velocità relativa di allontanamento e quella di avvicinamento vale $1$ nell'urto elastico, è compreso fra $0$ e $1$ in quello parzialmente anelastico e si annulla in quello totalmente anelastico, in cui il corpo resta aderente al piano.}
    \label{fig:tipi_urto}
\end{figure}

\begin{osservazione}{Frazione di Energia Persa}
Poiché l'energia cinetica associata al moto del centro di massa, $K_{CM}$, resta rigorosamente costante in un sistema isolato, la variazione di energia cinetica \emph{totale} coincide con la variazione della sola energia cinetica \emph{interna}:
\begin{equation}
    \frac{\Delta K}{K'_i} = \frac{K'_f - K'_i}{K'_i} = e^2 - 1
\end{equation}
La frazione di energia persa dipende dunque soltanto da $e$, e non dalle masse né dalla velocità con cui i corpi si avvicinano.
\end{osservazione}

\section{Urto Elastico ($e=1$)}

In un urto elastico i corpi si scontrano e tornano separati conservando rigorosamente sia la quantità di moto totale $\vec{P}_{\text{tot}}$ sia l'energia cinetica totale $K_{\text{tot}}$. Il problema consiste allora nel determinare le velocità finali a partire da quelle iniziali, e la sua risolubilità dipende in modo essenziale dal numero di dimensioni spaziali coinvolte.

\subsection{Sistema di Laboratorio --- 3 Dimensioni}

In tre dimensioni le leggi di conservazione impongono il sistema
\begin{equation}
\begin{cases}
    \tfrac{1}{2} m_A v_{Ai}^2 + \tfrac{1}{2} m_B v_{Bi}^2
        = \tfrac{1}{2} m_A v_{Af}^2 + \tfrac{1}{2} m_B v_{Bf}^2 \\[4pt]
    \hat{x}: \; p_{A,xi} + p_{B,xi} = p_{A,xf} + p_{B,xf} \\[2pt]
    \hat{y}: \; p_{A,yi} + p_{B,yi} = p_{A,yf} + p_{B,yf} \\[2pt]
    \hat{z}: \; p_{A,zi} + p_{B,zi} = p_{A,zf} + p_{B,zf}
\end{cases}
\end{equation}

Disponiamo di \textbf{4 equazioni scalari} e di \textbf{6 incognite}: le tre componenti cartesiane di $\vec{v}_{Af}$ e le tre di $\vec{v}_{Bf}$. Il sistema è dunque indeterminato, e nessuna legge di conservazione può colmare il divario: le due informazioni mancanti riguardano la \emph{geometria} dell'interazione, cioè quanto i due corpi si siano presentati centrati l'uno sull'altro. Occorrono perciò dati sperimentali aggiuntivi, tipicamente gli angoli di deflessione misurati dai rivelatori (analisi di \textbf{scattering}).

\subsection{Sistema di Laboratorio --- 2 Dimensioni}

Se il moto si svolge su un piano, una delle equazioni della quantità di moto cade e restano
\begin{equation}
\begin{cases}
    K_f = K_i \\[3pt]
    \hat{x}: \; p_{xi} = p_{xf} \\[2pt]
    \hat{y}: \; p_{yi} = p_{yf}
\end{cases}
\end{equation}
cioè \textbf{3 equazioni} e \textbf{4 incognite}. L'indeterminazione si riduce a uno solo parametro: è sufficiente misurare \emph{un} angolo di deflessione perché tutto il resto della cinematica finale risulti determinato.

\begin{figure}[ht]
    \centering
    \begin{tikzpicture}[>=stealth, scale=1.15]
        % Asse di riferimento
        \draw[dashed, thin, mypen!55] (-3.4,0) -- (3.3,0);

        % Proiettile in arrivo
        \draw[->, very thick, boxoss] (-3.0,0) -- (-0.55,0)
            node[midway, above, font=\small] {$\vec{v}_{Ai}$};
        \filldraw[boxoss!40, draw=boxoss, thick] (-3.0,0) circle (0.19)
            node[above=0.26cm, font=\small, mypen] {$m_A$};

        % Bersaglio fermo
        \filldraw[mypen!18, draw=mypen, thick] (0,0) circle (0.24);
        \node[font=\small, mypen] at (-0.05,-0.62) {$m_B$ (fermo)};

        % Deflessione di A
        \draw[->, very thick, boxoss] (0.18,0.12) -- (2.5,1.42)
            node[right, font=\small] {$\vec{v}_{Af}$};
        \draw[->, thin, mypen] (1.25,0) arc (0:29.3:1.25);
        \node[font=\small, boxoss] at (1.15,0.42) {$\theta_A$};

        % Rinculo di B
        \draw[->, very thick, notered] (0.18,-0.12) -- (2.5,-1.28)
            node[right, font=\small] {$\vec{v}_{Bf}$};
        \draw[->, thin, mypen] (1.55,0) arc (0:-26.8:1.55);
        \node[font=\small, notered] at (1.48,-0.45) {$\theta_B$};
    \end{tikzpicture}
    \caption{Scattering elastico in due dimensioni nel sistema del laboratorio, con il bersaglio $m_B$ inizialmente fermo. Misurato uno dei due angoli di deflessione, le leggi di conservazione determinano univocamente l'altro angolo e i moduli delle due velocità finali.}
    \label{fig:scattering_2d}
\end{figure}

\subsection{Sistema di Laboratorio --- 1 Dimensione}

In una sola dimensione spaziale il conteggio torna finalmente in pareggio: \textbf{2 equazioni} e \textbf{2 incognite}.
\begin{equation}
\begin{cases}
    \tfrac{1}{2} m_A v_{Ai}^2 + \tfrac{1}{2} m_B v_{Bi}^2
        = \tfrac{1}{2} m_A v_{Af}^2 + \tfrac{1}{2} m_B v_{Bf}^2 \\[4pt]
    m_A v_{Ai} + m_B v_{Bi} = m_A v_{Af} + m_B v_{Bf}
\end{cases}
\end{equation}
Le sole condizioni iniziali bastano quindi a determinare univocamente lo stato finale del moto, e la soluzione analitica generale è
\begin{equation}
    v_{Af} = \frac{(m_A - m_B)\,v_{Ai} + 2m_B \,v_{Bi}}{m_A + m_B}
    \, , \qquad
    v_{Bf} = \frac{(m_B - m_A)\,v_{Bi} + 2m_A \,v_{Ai}}{m_A + m_B}
\end{equation}
Si noti la simmetria delle due espressioni: l'una si ottiene dall'altra scambiando fra loro le etichette $A$ e $B$, come dev'essere, dato che nessuno dei due corpi gioca un ruolo privilegiato.

\subsection{Casi Particolari Notevoli}

Le formule generali diventano molto più espressive in tre situazioni limite, che conviene tenere a mente perché ricorrono continuamente nei problemi.

\begin{itemize}
    \item \textbf{Masse uguali} ($m_A = m_B$). Le velocità dei due corpi si \textbf{scambiano} semplicemente:
    \begin{equation*}
        v_{Af} = v_{Bi} \qquad \text{e} \qquad v_{Bf} = v_{Ai}
    \end{equation*}
    È il comportamento delle biglie da biliardo: la prima si arresta di colpo e la seconda riparte con la velocità che aveva la prima.

    \item \textbf{Bersaglio molto pesante} ($m_B \gg m_A$, con $v_{Bi} = 0$). Il proiettile leggero rimbalza indietro con velocità praticamente opposta, mentre il bersaglio resta quasi immobile:
    \begin{equation*}
        v_{Af} \approx -\,v_{Ai} \qquad \text{e} \qquad v_{Bf} \approx 0
    \end{equation*}
    È il caso di una pallina che rimbalza su un muro: il muro riceve quantità di moto $2m_Av_{Ai}$, ma essendo enormemente più massivo ne ricava una velocità trascurabile.

    \item \textbf{Proiettile molto pesante} ($m_A \gg m_B$, con il bersaglio fermo). Il proiettile prosegue quasi indisturbato, $v_{Af} \approx v_{Ai}$, mentre il bersaglio leggero viene ``sparato'' a velocità circa doppia:
    \begin{equation*}
        v_{Bf} \approx 2\,v_{Ai}
    \end{equation*}
\end{itemize}

\begin{osservazione}{Il Fattore 2 nell'Ultimo Caso}
Il risultato $v_{Bf} \approx 2 v_{Ai}$ si comprende immediatamente passando al sistema del centro di massa. Se $m_A \gg m_B$, il centro di massa praticamente coincide con $A$ e viaggia con $V_{CM} \approx v_{Ai}$. Nel CM il bersaglio arriva con velocità $-v_{Ai}$ e, per quanto vedremo nel prossimo paragrafo, riparte con $+v_{Ai}$. Tornando al laboratorio si somma $V_{CM}$ e si ottiene $v_{Ai} + v_{Ai} = 2v_{Ai}$.
\end{osservazione}

\subsection{Analisi nel Sistema del CM (1D)}

Il vantaggio del sistema del centro di massa è che una delle due leggi di conservazione diventa un'identità, e il problema si risolve quasi senza calcoli. Nel CM la quantità di moto totale è identicamente nulla in ogni istante, $\vec{p}_{\text{tot}}\,' = \vec{p}_A\,' + \vec{p}_B\,' = \vec{0}$, cioè
\begin{equation}
    m_A v_{Ai}\,' = -\,m_B v_{Bi}\,'
    \qquad \text{e} \qquad
    m_A v_{Af}\,' = -\,m_B v_{Bf}\,'
\end{equation}

Imponiamo ora anche la conservazione dell'energia cinetica:
\begin{equation}
    \tfrac{1}{2} m_A (v_{Ai}\,')^2 + \tfrac{1}{2} m_B (v_{Bi}\,')^2
    = \tfrac{1}{2} m_A (v_{Af}\,')^2 + \tfrac{1}{2} m_B (v_{Bf}\,')^2
\end{equation}
e raggruppiamo i termini relativi a ciascuna massa:
\begin{equation}
    m_A \left[ (v_{Ai}\,')^2 - (v_{Af}\,')^2 \right]
    = m_B \left[ (v_{Bf}\,')^2 - (v_{Bi}\,')^2 \right]
\end{equation}
Applicando il prodotto notevole $a^2 - b^2 = (a-b)(a+b)$:
\begin{equation}
    m_A (v_{Ai}\,' - v_{Af}\,')(v_{Ai}\,' + v_{Af}\,')
    = m_B (v_{Bf}\,' - v_{Bi}\,')(v_{Bf}\,' + v_{Bi}\,')
\end{equation}
Dividiamo infine membro a membro per la conservazione della quantità di moto nella forma $m_A (v_{Ai}\,' - v_{Af}\,') = m_B (v_{Bf}\,' - v_{Bi}\,')$, ottenendo la relazione puramente lineare
\begin{equation}
    v_{Ai}\,' + v_{Af}\,' = v_{Bf}\,' + v_{Bi}\,'
\end{equation}

Sostituendo $v_{Bi}\,' = -\frac{m_A}{m_B}v_{Ai}\,'$ e $v_{Bf}\,' = -\frac{m_A}{m_B}v_{Af}\,'$:
\begin{equation}
    (v_{Ai}\,' + v_{Af}\,') = -\frac{m_A}{m_B}\,(v_{Ai}\,' + v_{Af}\,')
    \implies \left(1 + \frac{m_A}{m_B}\right)(v_{Ai}\,' + v_{Af}\,') = 0
\end{equation}
Poiché $1 + m_A/m_B \neq 0$, l'unica soluzione fisicamente accettabile è

\begin{formulario}{Urto Elastico 1D nel Centro di Massa}
\begin{equation}
    v_{Af}\,' = -\,v_{Ai}\,'
    \qquad\qquad
    v_{Bf}\,' = -\,v_{Bi}\,'
\end{equation}
Nel sistema del CM un urto elastico unidimensionale consiste semplicemente nell'\textbf{inversione del verso delle velocità}, con i moduli inalterati.
\end{formulario}

\begin{figure}[ht]
    \centering
    \begin{tikzpicture}[>=stealth, scale=1.0,
        corpoA/.style={fill=boxoss!35, draw=boxoss, thick},
        corpoB/.style={fill=notered!35, draw=notered, thick}]

        % Asse e centro di massa
        \draw[thin, mypen!35] (-3.9,0) -- (3.9,0);
        \filldraw[boxese] (0,0) circle (2.6pt);
        \node[font=\small\bfseries, boxese, above] at (0,0.10) {CM};

        % ----- PRIMA: in avvicinamento, sopra l'asse
        \node[font=\small\bfseries, mypen, left] at (-4.15,1.0) {prima};
        \filldraw[corpoA] (-3.0,1.0) circle (0.2);
        \node[font=\small, mypen, above] at (-3.0,1.22) {$m_A$};
        \draw[->, very thick, boxoss] (-2.65,1.0) -- (-1.1,1.0)
            node[midway, above, font=\small, boxoss] {$\vec{v}_{Ai}\,'$};
        \filldraw[corpoB] (3.0,1.0) circle (0.24);
        \node[font=\small, mypen, above] at (3.0,1.26) {$m_B$};
        \draw[->, very thick, notered] (2.65,1.0) -- (1.1,1.0)
            node[midway, above, font=\small, notered] {$\vec{v}_{Bi}\,'$};

        % ----- DOPO: in allontanamento, sotto l'asse
        \node[font=\small\bfseries, mypen, left] at (-4.15,-1.1) {dopo};
        \draw[->, very thick, boxoss, dashed] (-1.1,-1.1) -- (-2.65,-1.1);
        \node[font=\small, boxoss, below] at (-1.9,-1.22) {$\vec{v}_{Af}\,' = -\vec{v}_{Ai}\,'$};
        \draw[->, very thick, notered, dashed] (1.1,-1.1) -- (2.65,-1.1);
        \node[font=\small, notered, below] at (1.9,-1.22) {$\vec{v}_{Bf}\,' = -\vec{v}_{Bi}\,'$};
    \end{tikzpicture}
    \caption{Urto elastico unidimensionale nel sistema del centro di massa. I due corpi si avvicinano con quantità di moto uguali e opposte e si riallontanano con le stesse quantità di moto invertite: l'urto si riduce a un cambio di segno.}
    \label{fig:inversione_velocita_cm}
\end{figure}

\subsection{Passaggio dal CM al Laboratorio}

Il risultato appena ottenuto è tanto semplice quanto poco utile di per sé, perché le misure si eseguono nel laboratorio. Per tornarvi basta però applicare le trasformazioni galileiane delle velocità, $v = v' + V_{CM}$:
\begin{equation}
\begin{split}
    v_{Af} &= v_{Af}\,' + V_{CM} = -\,v_{Ai}\,' + V_{CM} \\
           &= -\,(v_{Ai} - V_{CM}) + V_{CM} = 2V_{CM} - v_{Ai}
\end{split}
\end{equation}

Sostituendo la definizione del centro di massa, $V_{CM} = \dfrac{m_A v_{Ai} + m_B v_{Bi}}{m_A + m_B} = \dfrac{m_A v_{Ai} + m_B v_{Bi}}{M}$:
\begin{equation}
\begin{split}
    v_{Af} &= 2 \left( \frac{m_A v_{Ai} + m_B v_{Bi}}{M} \right) - v_{Ai} \\
           &= \frac{2m_A v_{Ai} + 2m_B v_{Bi} - (m_A + m_B)v_{Ai}}{M} \\
           &= \frac{(m_A - m_B)\,v_{Ai} + 2m_B \,v_{Bi}}{m_A + m_B}
\end{split}
\end{equation}

Ritroviamo così, per una via molto più breve, le formule generali già scritte per il laboratorio:
\begin{equation}
    v_{Af} = \frac{(m_A - m_B)\,v_{Ai} + 2m_B \,v_{Bi}}{m_A + m_B}
    \, , \qquad
    v_{Bf} = \frac{(m_B - m_A)\,v_{Bi} + 2m_A \,v_{Ai}}{m_A + m_B}
\end{equation}

\begin{osservazione}{La Strategia Generale}
Lo schema appena seguito --- \textbf{passare al CM, risolvere, tornare al laboratorio} --- è il metodo standard per ogni problema d'urto, e verrà riutilizzato tale e quale nella dinamica degli urti relativistici. Il motivo è sempre lo stesso: nel CM una delle due leggi di conservazione si riduce a $\vec{p}_{\text{tot}}\,' = \vec{0}$, e il problema perde metà delle sue incognite.
\end{osservazione}

\section{Urto Totalmente Anelastico ($e=0$)}

In un urto totalmente anelastico i corpi rimangono uniti dopo l'impatto e proseguono come un unico corpo di massa $M = m_A + m_B$. Questa volta l'energia cinetica non si conserva e la sua legge ci è ignota; disponiamo però della sola conservazione della quantità di moto, che per fortuna basta, perché anche le incognite si sono ridotte a una (l'unica velocità finale $V_f$):
\begin{equation}
    P_i = m_A v_{Ai} + m_B v_{Bi} = (m_A + m_B)\,V_f
    \implies V_f = \frac{m_A v_{Ai} + m_B v_{Bi}}{m_A + m_B} = V_{CM}
\end{equation}

La velocità finale dell'aggregato coincide esattamente con la velocità del centro di massa del sistema --- risultato prevedibile, dato che dopo l'urto i due corpi \emph{sono} il loro centro di massa.

\subsection{Energia Dissipata e Teorema di Koenig}

Per calcolare l'energia meccanica persa sfruttiamo la scomposizione fornita dal teorema di Koenig, $K = K_{CM} + K'$, che separa il moto globale di traslazione dal moto interno. Poiché nell'urto totalmente anelastico i corpi restano uniti, la loro velocità relativa e la loro velocità rispetto al CM dopo l'urto sono identicamente nulle, $\vec{v}_f\,' = \vec{0}$, il che implica $K'_f = 0$. Quindi
\begin{align*}
    K_i &= K_{CM} + K'_i \\
    K_f &= K_{CM} + \cancelto{0}{K'_f} = K_{CM}
    \qquad \text{(poiché } V_f = V_{CM}\text{)}
\end{align*}
e la variazione di energia cinetica totale risulta
\begin{equation}
    \Delta K = K_f - K_i = K_{CM} - (K_{CM} + K'_i)
             = -\,K'_i = -\,\frac{1}{2}\mu \, v_{ri}^2
\end{equation}
dove $\mu = \dfrac{m_A m_B}{m_A + m_B}$ è la massa ridotta del sistema e $v_{ri} = |v_{Ai} - v_{Bi}|$ la velocità relativa iniziale. L'intera energia interna viene dissipata in calore e lavoro di deformazione permanente.

\begin{attenzione}{Non Tutta l'Energia Cinetica Può Essere Dissipata}
Anche nel caso più violento possibile --- quello totalmente anelastico --- il sistema \textbf{non} si ferma e non perde tutta la sua energia cinetica: sopravvive sempre l'aliquota $K_{CM}$ legata al moto del centro di massa. Dissiparla significherebbe portare a zero $\vec{P}_{\text{tot}}$, il che è vietato dalla conservazione della quantità di moto. L'urto totalmente anelastico è quindi il \emph{limite superiore} della dissipazione possibile, non una dissipazione totale.
\end{attenzione}

\section{Urto Parzialmente Anelastico ($0 < e < 1$)}

È il caso generale e realistico: i corpi si separano dopo il contatto, ma solo una frazione dell'energia cinetica interna viene restituita. Le due leggi a disposizione sono la conservazione della quantità di moto e la definizione stessa del coefficiente di restituzione, che riscriviamo in forma algebrica come
\begin{equation}
    v_{Bf} - v_{Af} = e\,(v_{Ai} - v_{Bi})
\end{equation}
cioè: la velocità di allontanamento è pari a $e$ volte la velocità di avvicinamento. Mettendo a sistema le due relazioni,
\begin{equation}
\begin{cases}
    m_A v_{Ai} + m_B v_{Bi} = m_A v_{Af} + m_B v_{Bf} \\[3pt]
    v_{Bf} - v_{Af} = e\,(v_{Ai} - v_{Bi})
\end{cases}
\end{equation}
e risolvendo rispetto alle due incognite si ottengono le formule generali

\begin{formulario}{Velocità Finali per un Urto di Restituzione $e$}
\begin{equation}
    v_{Af} = \frac{(m_A - e\,m_B)\,v_{Ai} + m_B(1+e)\,v_{Bi}}{m_A + m_B}
    \, , \qquad
    v_{Bf} = \frac{(m_B - e\,m_A)\,v_{Bi} + m_A(1+e)\,v_{Ai}}{m_A + m_B}
\end{equation}
Esse contengono come casi particolari tutti quelli visti finora: ponendo $e = 1$ si riottengono le formule dell'urto elastico, ponendo $e = 0$ si ottiene $v_{Af} = v_{Bf} = V_{CM}$, cioè l'urto totalmente anelastico.
\end{formulario}

\begin{esempio}{Analisi Energetica tramite il Teorema di Koenig}
Il teorema di Koenig, $K = K_{CM} + K'$, consente di separare in modo netto l'energia legata al moto traslatorio globale del sistema --- $K_{CM}$, che non varia in assenza di forze esterne --- dall'energia interna al centro di massa --- $K'$, che rappresenta l'unica aliquota effettivamente trasformabile o dissipabile negli urti.

\medskip
\textbf{Caso 1: masse uguali, entrambe mobili ($m_A = m_B = m$).}
Dalla scomposizione di Koenig si ha $K' = K - K_{CM}$. Per lo stato iniziale, essendo $V_{CM} = \frac{1}{2}(v_{Ai} + v_{Bi})$ e dunque $v_{Ai}\,' = \frac{1}{2}(v_{Ai} - v_{Bi})$:
\begin{equation}
    K'_i = \frac{1}{2}m \left[(v_{Ai}\,')^2 + (v_{Bi}\,')^2\right]
         = \frac{m}{4}\,(v_{Ai} - v_{Bi})^2 = \frac{m}{4}\,v_{ri}^2
\end{equation}
In modo del tutto analogo, per lo stato finale dopo l'urto si ottiene
\begin{equation}
    K'_f = \frac{m}{4}\,v_{rf}^2
\end{equation}
e dalla definizione di coefficiente di restituzione segue immediatamente la coerenza delle due scritture:
\begin{equation*}
    e^2 = \frac{K'_f}{K'_i} = \frac{v_{rf}^2}{v_{ri}^2}
    \implies e = \frac{v_{rf}}{v_{ri}}
\end{equation*}

\medskip
\textbf{Caso 2: proiettile contro bersaglio fermo ($v_{Bi} = 0$).}
Le due componenti dell'energia iniziale si ripartiscono come
\begin{equation}
    K_{CM} = \frac{1}{2}\,\frac{m_A^2 \,v_{Ai}^2}{M}
    \, , \qquad
    K'_i = \frac{1}{2}\mu \, v_{ri}^2 = \frac{1}{2}\,\frac{m_A m_B}{M}\,v_{Ai}^2
\end{equation}
In un urto contro un bersaglio fermo, la frazione di energia cinetica associata al centro di massa \textbf{non può mai essere dissipata}, in virtù della conservazione della quantità di moto totale. La massima energia dissipabile è esclusivamente l'aliquota interna $K'_i$, e la si dissipa tutta quando $e = 0$.

\medskip
Si osservi il rapporto fra le due aliquote:
\begin{equation*}
    \frac{K'_i}{K_{CM}} = \frac{m_B}{m_A}
\end{equation*}
Un proiettile leggero contro un bersaglio pesante ($m_B \gg m_A$) mette a disposizione dell'urto quasi tutta la propria energia; viceversa, un proiettile pesante contro un bersaglio leggero ($m_A \gg m_B$) ne conserva quasi tutta nel moto del centro di massa e ne può dissipare pochissima. È la ragione per cui, negli acceleratori, i fasci collidenti sono molto più efficienti dei bersagli fissi.
\end{esempio}

\chapter{Dinamica delle Rotazioni}

\section{Introduzione}

In questo capitolo studiamo il comportamento dei sistemi in rotazione, introducendo le tre grandezze vettoriali che ne governano la dinamica: il \textbf{momento di un vettore}, il \textbf{momento di una forza} e il \textbf{momento angolare}.

Tutte e tre hanno una caratteristica in comune che conviene mettere subito in chiaro: non sono proprietà intrinseche del corpo, ma dipendono dal \emph{punto} rispetto al quale le si calcola. Prima ancora di definirle occorre dunque fissare tale punto di riferimento.

\begin{definizione}{Polo della Rotazione}
Si chiama \textbf{polo della rotazione} il punto $O$ rispetto al quale vengono calcolati tutti i momenti. Per comodità lo si sceglie in genere coincidente con l'origine del sistema di riferimento cartesiano adottato.
\end{definizione}

\begin{attenzione}{Ogni Momento Va Riferito a un Polo}
Parlare del ``momento angolare di un corpo'' senza specificare il polo è privo di significato: lo stesso corpo, nello stesso istante, ha momenti angolari diversi rispetto a poli diversi. Per questo la notazione porta sempre il polo a pedice, $\vec{L}_O$ e $\vec{\tau}_O$.
\end{attenzione}

\section{Momento di un Vettore}

Per quantificare l'attitudine di una grandezza vettoriale a produrre una rotazione attorno a un punto si definisce il suo momento. Dato un generico vettore $\vec{B}$ applicato nel punto $A$, il suo momento rispetto al polo $O$ è il prodotto vettoriale fra il raggio vettore $\vec{r} = \vec{OA}$ e il vettore stesso:
\begin{equation}
    \vec{M}_B = \vec{OA} \times \vec{B}
\end{equation}

La retta su cui giace $\vec{B}$ prende il nome di \textbf{retta d'azione}. Il momento gode di una proprietà notevole: dipende soltanto dalla retta d'azione, non dal particolare punto di applicazione scelto lungo di essa. Possiamo cioè far scorrere il vettore lungo la propria retta d'azione senza alterarne l'effetto rotazionale.

\begin{osservazione}{Indipendenza dal Punto di Applicazione}
Scegliamo come punto di applicazione un altro punto $A'$ della stessa retta d'azione e calcoliamo il momento corrispondente:
\begin{equation}
    \vec{M}_B' = \vec{OA}' \times \vec{B}
    = (\vec{OA} + \vec{AA}') \times \vec{B}
    = \vec{OA} \times \vec{B} + \vec{AA}' \times \vec{B}
\end{equation}
Poiché $A$ e $A'$ appartengono entrambi alla retta d'azione di $\vec{B}$, lo spostamento $\vec{AA}'$ è parallelo a $\vec{B}$; il prodotto vettoriale fra vettori paralleli è nullo, perché l'angolo compreso vale zero e $\sin 0 = 0$. Il secondo termine si annulla e resta
\begin{equation}
    \vec{M}_B' = \vec{OA} \times \vec{B} = \vec{M}_B
\end{equation}
\end{osservazione}

La conseguenza pratica è che il modulo del momento si calcola comodamente come
\begin{equation}
    |\vec{M}_B| = |\vec{B}| \, b
\end{equation}
dove $b = |\vec{OH}|$ è il \textbf{braccio}, cioè la distanza fra il polo e la retta d'azione, misurata perpendicolarmente. Tutta l'informazione geometrica rilevante è contenuta in quella sola distanza.

\begin{figure}[ht]
    \centering
    \begin{tikzpicture}[x={(-0.55cm,-0.38cm)}, y={(1cm,0cm)}, z={(0cm,1cm)},
                        >=stealth, scale=1.15, font=\small]
        \coordinate (O) at (0,0,0);

        % Piano xy, per dare profondità
        \fill[mypen, opacity=0.05] (-1,-3,0) -- (3.2,-3,0) -- (3.2,4.3,0) -- (-1,4.3,0) -- cycle;

        % Assi
        \draw[->, mypen!55, thick] (0,0,0) -- (3.2,0,0) node[anchor=north east] {$x$};
        \draw[->, mypen!55, thick] (0,0,0) -- (0,4.3,0) node[anchor=north west] {$y$};
        \draw[->, mypen!55, thick] (0,0,0) -- (0,0,4.3) node[anchor=south] {$z$};
        \node[anchor=north west, mypen, xshift=2pt] at (O) {$O$};

        % Retta d'azione
        \draw[dashed, thick, mypen!70] (2,-3,0) -- (2,4.6,0)
            node[above, anchor=south west, mypen] {retta d'azione};

        % Braccio
        \draw[very thick, boxatt] (O) -- (2,0,0) node[pos=0.55, anchor=south, boxatt, yshift=2pt] {$b$};
        \fill[boxatt] (2,0,0) circle (1.6pt) node[anchor=north, boxatt, yshift=-2pt] {$H$};

        % Raggi vettori
        \coordinate (A)  at (2,-2.2,0);
        \coordinate (Ap) at (2,2.3,0);
        \draw[->, thick, boxese] (O) -- (A)  node[midway, anchor=south east, boxese] {$\vec{r}$};
        \draw[->, thick, boxese!70, dashed] (O) -- (Ap) node[pos=0.72, anchor=south, boxese, yshift=2pt] {$\vec{r}\,'$};

        % Vettori B, applicati in A e in A'
        \fill[mypen] (A) circle (1.6pt) node[anchor=east, mypen, xshift=-2pt] {$A$};
        \fill[mypen] (Ap) circle (1.6pt) node[anchor=north, mypen, yshift=-2pt] {$A'$};
        \draw[->, ultra thick, boxoss] (A) -- (2,-0.85,0) node[midway, anchor=north, boxoss, yshift=-1pt] {$\vec{B}$};
        \draw[->, thick, boxoss!45, dashed] (Ap) -- (2,4.1,0) node[midway, anchor=north west, boxoss] {$\vec{B}$};

        % Momento
        \draw[->, ultra thick, notered] (O) -- (0,0,3.5)
            node[anchor=west, notered] {$\vec{M}_B = \vec{r} \times \vec{B}$};

        % Simboli di ortogonalità
        \draw[thin, mypen] (0,0,0.3) -- (0.35,0,0.3) -- (0.35,0,0);
        \draw[thin, mypen] (1.75,0,0) -- (1.75,0.25,0) -- (2.0,0.25,0);
    \end{tikzpicture}
    \caption{Momento del vettore $\vec{B}$ rispetto al polo $O$. I punti $A$ e $A'$ sono due scelte arbitrarie del punto di applicazione lungo la stessa retta d'azione: il momento $\vec{M}_B$, ortogonale al piano individuato da $\vec{r}$ e $\vec{B}$, è il medesimo in entrambi i casi e il suo modulo vale $|\vec{B}|\,b$, dove $b$ è il braccio.}
    \label{fig:momento_vettore}
\end{figure}

\section{Grandezze della Dinamica Rotazionale}

Applicando la definizione appena data ai due vettori fondamentali della dinamica --- la forza e la quantità di moto --- si ottengono le due grandezze su cui si regge l'intero capitolo.

\begin{definizione}{Momento di una Forza}
Se il vettore considerato è una forza $\vec{F}$, il suo momento rispetto al polo $O$ si chiama \textbf{momento della forza} (o \emph{momento torcente}) e si indica con $\vec{M}_F$ o, più comunemente, con $\vec{\tau}$:
\begin{equation}
    \vec{\tau} = \vec{OA} \times \vec{F} = \vec{r} \times (m\vec{a})
\end{equation}
dove $\vec{r}$ è il raggio vettore del punto di applicazione. Il modulo è massimo quando la forza è perpendicolare al raggio vettore, e in tal caso vale semplicemente
\begin{equation*}
    \tau = m\,r\,a \, \sin\!\left(\frac{\pi}{2}\right) = m\,r\,a
\end{equation*}
\end{definizione}

\begin{definizione}{Momento Angolare}
Se il vettore considerato è la quantità di moto $\vec{p} = m\vec{v}$ di un punto materiale, il suo momento rispetto al polo $O$ si chiama \textbf{momento angolare} e si indica con $\vec{M}_p$ o, più comunemente, con $\vec{L}$:
\begin{equation}
    \vec{L} = \vec{OA} \times \vec{p} = \vec{r} \times (m\vec{v})
\end{equation}
\end{definizione}

Dalla definizione discende immediatamente una proprietà geometrica di grande portata. I vettori $\vec{r}$ e $\vec{v}$ individuano un piano, detto \textbf{piano istantaneo della traiettoria}, e il momento angolare, essendo un prodotto vettoriale, è per costruzione perpendicolare a tale piano.

\begin{osservazione}{Momento Angolare Costante Implica Moto Piano}
Se $\vec{L}$ mantiene direzione e verso costanti, allora anche il piano individuato da $\vec{r}$ e $\vec{v}$ resta fisso nello spazio, e la traiettoria è necessariamente \textbf{piana}. È il motivo per cui i pianeti, soggetti a una forza che conserva il momento angolare, percorrono orbite contenute in un piano anziché curve sghembe nello spazio.
\end{osservazione}

\subsection{Significato Fisico del Momento Angolare}

Per legare il momento angolare allo stato di rotazione conviene esprimere i vettori in coordinate polari, usando il versore radiale $\hat{r}$ e il versore trasverso $\hat{\phi}$:
\begin{itemize}
    \item il raggio vettore ha solo componente radiale, $\vec{r} = r\hat{r}$;
    \item la velocità, essendone la derivata temporale, contiene sia la variazione della distanza sia quella dell'angolo: $\vec{v} = \dot{\vec{r}} = \dot{r}\hat{r} + r\dot{\phi}\hat{\phi}$.
\end{itemize}

Sostituendo nella definizione di momento angolare:
\begin{equation}
    \begin{split}
        \vec{L} &= \vec{r} \times m\vec{v}
        = m \left[ r\hat{r} \times \big(\dot{r}\hat{r} + r\dot{\phi}\hat{\phi}\big) \right] \\
        &= m \left( r\dot{r}\,(\hat{r} \times \hat{r}) + r^2 \dot{\phi}\,(\hat{r} \times \hat{\phi}) \right)
    \end{split}
\end{equation}
Il primo termine si annulla perché $\hat{r} \times \hat{r} = \vec{0}$; il secondo si semplifica ricordando che $\hat{r} \times \hat{\phi} = \hat{z}$, versore uscente dal piano del moto. Indicando con $\omega = \dot{\phi}$ la velocità angolare si ottiene
\begin{equation}
    \vec{L} = m r^2 \dot{\phi} \, \hat{z}
\end{equation}

\begin{osservazione}{Il Momento Angolare Misura la Rotazione, non la Traslazione}
Il risultato è molto istruttivo: $\vec{L}$ si annulla se e solo se $\dot{\phi} = \omega = 0$, cioè quando il moto è puramente radiale. La componente $\dot{r}$ della velocità --- per quanto grande --- non contribuisce affatto. Il momento angolare è dunque una grandezza \textbf{associata esclusivamente alla rotazione} attorno al polo, e il suo valore dipende solo da $r$ e da $\dot{\phi}$.
\end{osservazione}

Il caso notevole è quello in cui si impedisce ogni variazione radiale, $\dot{r} = 0$: il raggio è fissato e il moto è \textbf{circolare}. All'interno di questa famiglia si distinguono due sottocasi:
\begin{itemize}
    \item se l'accelerazione angolare è non nulla ($a_\phi \neq 0$), e dunque $\dot{\phi}$ varia nel tempo, si ha un \textbf{moto circolare vario};
    \item se l'accelerazione angolare è nulla ($a_\phi = 0$), e dunque $\dot{\phi}$ è costante, si ha un \textbf{moto circolare uniforme}.
\end{itemize}

\section{Moto non Planare rispetto a \texorpdfstring{$O$}{O}}

Non sempre il polo $O$ appartiene al piano in cui avviene il moto. Quando il moto circolare si svolge su un piano che non contiene il polo, il momento angolare rispetto a $O$ acquista una struttura più ricca, e conviene calcolarlo scomponendo il raggio vettore.

L'idea è semplice: si proietta il polo $O$ sul piano dell'orbita, ottenendo il centro $O'$ della circonferenza, e si scrive $\vec{r}$ come somma di due contributi ortogonali --- il raggio $\vec{R}$ della circonferenza, che giace nel piano del moto, e la quota $z$ che separa i due poli lungo l'asse di rotazione.

\begin{figure}[ht]
    \centering
    \tdplotsetmaincoords{70}{120}
    \begin{tikzpicture}[tdplot_main_coords, scale=1.45, >=stealth,
                        line cap=round, line join=round, font=\small]
        \def\R{1.5}
        \def\H{2.5}
        \def\angP{45}

        % Assi
        \draw[->, thick, mypen!55] (0,0,0) -- (2.5,0,0) node[anchor=north east] {$x$};
        \draw[->, thick, mypen!55] (0,0,0) -- (0,2.5,0) node[anchor=north west] {$y$};
        \draw[->, thick, mypen!55] (0,0,0) -- (0,0,\H+1.4) node[anchor=south] {$z$};
        \node[anchor=north east, mypen] at (0,0,0) {$O$};

        % Circonferenza e centro proiettato
        \coordinate (Op) at (0,0,\H);
        \node[anchor=north east, mypen] at (Op) {$O'$};
        \fill[mypen] (Op) circle (1.4pt);
        \draw[thick, mypen!75] plot[domain=0:360, samples=60, smooth]
            ({\R*cos(\x)}, {\R*sin(\x)}, \H);

        % Quota z fra i due poli
        \draw[very thick, boxatt] (0,0,0) -- (Op);
        \node[anchor=east, boxatt] at (0,0,{\H/2}) {$z$};

        % Punto P
        \coordinate (P) at ({\R*cos(\angP)}, {\R*sin(\angP)}, \H);
        \fill[notered] (P) circle (1.7pt) node[anchor=south west, mypen] {$P$};

        % Vettori
        \draw[->, thick, boxese] (0,0,0) -- (P)
            node[midway, anchor=north west, boxese] {$\vec{r}$};
        \draw[->, thick, boxoss] (Op) -- (P)
            node[midway, anchor=south, boxoss] {$\vec{R}$};
        \coordinate (V) at ({\R*cos(\angP) - 1.2*sin(\angP)}, {\R*sin(\angP) + 1.2*cos(\angP)}, \H);
        \draw[->, thick, notered] (P) -- (V) node[anchor=west, notered] {$\vec{v}$};

        % Velocità angolare
        \draw[->, thick, mypen] plot[domain=150:-30, samples=30, smooth]
            ({0.42*cos(\x)}, {0.42*sin(\x)}, {\H+0.85});
        \node[mypen] at (0.95, 0.85, {\H+0.95}) {$\vec{\omega}$};

        % Versori locali, sul lato opposto della circonferenza
        \def\angL{170}
        \coordinate (L) at ({\R*cos(\angL)}, {\R*sin(\angL)}, \H);
        \draw[->, very thick, boxese]
            (L) -- ({(\R+0.75)*cos(\angL)}, {(\R+0.75)*sin(\angL)}, \H)
            node[anchor=north east, boxese] {$\hat{\rho}$};
        \draw[->, very thick, boxoss]
            (L) -- ({\R*cos(\angL) - 0.75*sin(\angL)}, {\R*sin(\angL) + 0.75*cos(\angL)}, \H)
            node[anchor=north west, boxoss] {$\hat{\phi}$};
    \end{tikzpicture}
    \caption{Moto circolare su un piano che non contiene il polo $O$. Il raggio vettore si scompone nel raggio $\vec{R} = R\hat{\rho}$ della circonferenza e nella quota $z$ che separa $O$ dal centro proiettato $O'$.}
    \label{fig:moto_non_planare}
\end{figure}

Fissato il piano del moto perpendicolarmente all'asse $\hat{z}$, il centro della circonferenza è $O'$ e la quota che lo separa da $O$ vale $z$. Adottiamo i versori $\hat{\rho}$ (radiale, uscente dal centro $O'$ verso il punto) e $\hat{\phi}$ (trasverso), osservando che il raggio della circonferenza è per costruzione parallelo al primo, $\vec{R} \parallel \hat{\rho}$.

Il raggio vettore rispetto a $O$ si scrive allora come somma dei due contributi ortogonali, $\vec{r} = \vec{R} + \vec{OO}'$, cioè
\begin{equation}
    \vec{r} = R\hat{\rho} + z\hat{z}
\end{equation}
mentre la velocità, essendo il moto circolare, è puramente tangenziale:
\begin{equation}
    \vec{v} = R\omega\hat{\phi}
\end{equation}

Sostituiamo nella definizione di momento angolare e distribuiamo il prodotto vettoriale:
\begin{align}
    \vec{L}_O &= m \, \vec{r} \times \vec{v} = m \, (R \hat{\rho} + z \hat{z}) \times (R \omega \hat{\phi}) \\
              &= m \big[ R \hat{\rho} \times R \omega \hat{\phi} + z \hat{z} \times R \omega \hat{\phi} \big] \\
              &= m \big[ R^2 \omega \, (\hat{\rho} \times \hat{\phi}) + R \omega z \, (\hat{z} \times \hat{\phi}) \big]
\end{align}

Ricordando le regole di moltiplicazione della terna destrorsa, $\hat{\rho} \times \hat{\phi} = \hat{z}$ e $\hat{z} \times \hat{\phi} = -\hat{\rho}$, si ottiene il risultato finale:
\begin{equation}
    \vec{L}_O = m R^2 \omega \, \hat{z} - m R \omega z \, \hat{\rho}
\end{equation}

\begin{osservazione}{Le Due Componenti del Momento Angolare}
Il risultato si legge come somma di due contributi distinti:
\begin{itemize}[leftmargin=1.4em]
    \item $\vec{L}_{O'} = m R^2 \omega \, \hat{z}$, diretto lungo l'asse di rotazione: è esattamente il momento angolare che si avrebbe scegliendo come polo il centro $O'$ della circonferenza, cioè in assenza di sfasamento fra i due poli;
    \item $-\, m R \omega z \, \hat{\rho}$, diretto radialmente verso l'asse: compare unicamente a causa della quota $z$ ed è proporzionale ad essa.
\end{itemize}
Di conseguenza $\vec{L}_O$ \textbf{non è parallelo a} $\vec{\omega}$, e per di più il termine radiale ruota insieme al punto: il momento angolare rispetto a $O$ non è costante in direzione, pur essendolo in modulo. Solo scegliendo il polo sul piano del moto ($z = 0$) le due grandezze tornano parallele.
\end{osservazione}

\section{Forze Centrali}

Si dicono \textbf{centrali} le forze dirette lungo il raggio vettore $\vec{r}$ che congiunge la particella a un punto fisso $O$, detto \textbf{centro di forza}.

\begin{definizione}{Forza Centrale}
Una forza è centrale se la sua retta d'azione passa in ogni istante per il centro di forza $O$, e se il suo modulo dipende soltanto dalla distanza da esso:
\begin{equation}
    \vec{F}_c = \vec{F}(r) = F(r)\,\hat{r}
\end{equation}
dove $F(r)$ è una funzione scalare della sola coordinata radiale $r$. Il segno di $F(r)$ distingue le forze attrattive ($F < 0$) da quelle repulsive ($F > 0$).
\end{definizione}

Le forze fondamentali della natura sono in larga parte di questo tipo: lo sono la forza \textbf{gravitazionale}, la forza \textbf{coulombiana} e la forza \textbf{elastica} isotropa. Lo studio che segue si applica dunque tanto al moto dei pianeti quanto a quello degli elettroni attorno al nucleo.

\subsection{Proprietà Fondamentali}

Le forze centrali possiedono due proprietà generali, dalle quali discende tutto il resto.

\paragraph{1. Sono forze conservative.}
Il lavoro compiuto da una forza centrale dipende unicamente dalle posizioni iniziale e finale. Per dimostrarlo scomponiamo lo spostamento elementare $d\vec{l}$ in una componente radiale $d\vec{r}$ e una trasversale $d\vec{s}$ (lungo il versore $\hat{t}$) e calcoliamo il lavoro elementare:
\begin{equation}
    d\mathcal{L}_{F_c} = \vec{F}_c \cdot d\vec{l}
    = \big(F(r)\hat{r}\big) \cdot \big(dr\,\hat{r} + ds\,\hat{t}\big) = F(r)\,dr
\end{equation}
Il contributo trasversale si annulla perché $\hat{r} \cdot \hat{t} = 0$. Integrando fra due posizioni $\vec{x}_1$ e $\vec{x}_2$:
\begin{equation}
    \mathcal{L}_{F_c} = \int_{\vec{x}_1}^{\vec{x}_2} \vec{F}_c \cdot d\vec{l}
    = \int_{r_1}^{r_2} F(r) \, dr
\end{equation}
L'integrale dipende solo dagli estremi radiali $r_1$ e $r_2$, e non dal cammino percorso per congiungerli: la forza è quindi conservativa e ammette un'energia potenziale.

\begin{figure}[ht]
    \centering
    \begin{tikzpicture}[>=stealth, scale=1.35, font=\small]
        \coordinate (O) at (0,0);
        \coordinate (P) at (2.5,2.0);

        % Traiettoria
        \draw[thick, mypen!45, dashed] (0.8,0.5) to[out=30, in=240] (P) to[out=60, in=190] (4.6,3.5);
        \node[mypen!70, right] at (4.6,3.5) {traiettoria};

        % Direzione radiale
        \draw[dashed, thin, mypen!55] (O) -- (4.1,3.28);

        % Forza centrale, verso O
        \draw[->, ultra thick, boxatt] (P) -- ($(P)!0.55!(O)$)
            node[midway, anchor=south east, boxatt] {$\vec{F}_c$};

        % Punti
        \fill[mypen] (O) circle (1.7pt) node[anchor=north east, mypen] {$O$};
        \fill[notered] (P) circle (1.7pt) node[anchor=north west, mypen] {$P$};

        % Scomposizione dello spostamento
        \coordinate (Pn) at ($(P) + (0.85,0.85)$);
        \coordinate (Pr) at ($(P) + (0.85,0.68)$);
        \draw[->, very thick, notered] (P) -- (Pn)
            node[anchor=south east, notered, yshift=2pt] {$d\vec{l}$};
        \draw[->, very thick, boxese] (P) -- (Pr)
            node[anchor=north west, boxese, xshift=-2pt] {$d\vec{r}$};
        \draw[->, very thick, boxoss] (Pr) -- (Pn)
            node[midway, anchor=west, boxoss] {$d\vec{s}$};
        \draw[dashed, thin, mypen!45] (Pr) -- (Pn);
    \end{tikzpicture}
    \caption{Scomposizione dello spostamento elementare $d\vec{l}$ in componente radiale $d\vec{r}$ e trasversale $d\vec{s}$. La forza centrale $\vec{F}_c$, parallela a $\hat{r}$, compie lavoro soltanto sulla prima.}
    \label{fig:lavoro_forza_centrale}
\end{figure}

\paragraph{2. Il loro momento rispetto a $O$ è nullo.}
Scegliendo come polo il centro di forza stesso, la verifica è immediata:
\begin{equation}
    \vec{\tau}_O = \vec{r} \times \vec{F}_c
    = (r \hat{r}) \times \big(F(r) \hat{r}\big)
    = r F(r) \, (\hat{r} \times \hat{r}) = \vec{0}
\end{equation}
perché $\vec{r}$ e $\vec{F}_c$ sono paralleli per definizione.

\begin{formulario}{Le Due Costanti del Moto in un Campo Centrale}
Da queste due proprietà discendono immediatamente due integrali primi del moto:
\begin{minipage}[t]{0.47\textwidth}
\centering \textbf{Energia totale}
\begin{equation*}
    E = K + U(r) = \text{costante}
\end{equation*}
\centering \footnotesize (la forza è conservativa)
\end{minipage}
\hfill
\begin{minipage}[t]{0.47\textwidth}
\centering \textbf{Momento angolare}
\begin{equation*}
    \frac{d\vec{L}}{dt} = \vec{\tau}_O = \vec{0} \implies \vec{L} = \text{costante}
\end{equation*}
\centering \footnotesize (purché $O \equiv$ centro di forza)
\end{minipage}
\vspace{4pt}
\end{formulario}

Sono queste due costanti a rendere risolubile il problema del moto in un campo centrale, e sarà da esse che ricaveremo, nella sezione conclusiva, la forma generale della traiettoria.

\subsection{Significato Geometrico del Momento Angolare}

La costanza di $\vec{L}$ ammette una lettura puramente geometrica, che storicamente precede di quasi un secolo la formulazione newtoniana. Mentre il punto si muove, il raggio vettore spazza una superficie; calcoliamone il ritmo di crescita.

\begin{figure}[ht]
    \centering
    \begin{tikzpicture}[>=stealth, scale=1.65, font=\small]
        \coordinate (O) at (0,0);
        \coordinate (P1) at (20:3);
        \coordinate (P2) at (50:3.2);

        % Settore spazzato
        \fill[boxatt!15] (O) -- (P1) to[bend right=10] (P2) -- cycle;

        % Traiettoria
        \draw[thick, mypen!45, dashed]
            (5:2.8) to[out=100, in=260] (P1) to[out=80, in=230] (P2) to[out=50, in=200] (66:3.5);
        \node[mypen!70, right] at (66:3.5) {traiettoria};

        % Raggi vettori
        \draw[->, very thick, boxese] (O) -- (P1)
            node[midway, anchor=north west, boxese] {$\vec{r}(t)$};
        \draw[->, very thick, boxese] (O) -- (P2)
            node[midway, anchor=south east, boxese] {$\vec{r}(t+dt)$};

        % Spostamento
        \draw[->, thick, notered] (P1) -- (P2)
            node[midway, anchor=west, notered, xshift=2pt] {$d\vec{r}$};

        % Angolo
        \draw[thick, mypen] (20:0.85) arc (20:50:0.85);
        \node[mypen] at (35:1.1) {$d\varphi$};

        % Punti
        \fill[mypen] (P1) circle (1.6pt) node[anchor=west, mypen, xshift=2pt] {$P(t)$};
        \fill[mypen] (P2) circle (1.6pt) node[anchor=south, mypen, yshift=2pt] {$P(t+dt)$};
        \fill[mypen] (O) circle (1.6pt) node[anchor=north east, mypen] {$O$};

        % Area
        \node[boxatt, font=\bfseries] at (35:1.85) {$dA$};
    \end{tikzpicture}
    \caption{Interpretazione geometrica della costanza del momento angolare: nell'intervallo $dt$ il raggio vettore spazza l'area infinitesima $dA$, assimilabile a un settore circolare di ampiezza $d\varphi$.}
    \label{fig:velocita_areolare}
\end{figure}

Nell'intervallo $dt$ il raggio vettore ruota di $d\varphi$ e spazza un triangolo mistilineo, assimilabile per $d\varphi \to 0$ a un settore circolare di area
\begin{equation}
    dA = \frac{r^2 \, d\varphi}{2}
\end{equation}
Dividendo per $dt$ si ottiene la \textbf{velocità areolare}, e moltiplicando e dividendo per la massa vi si riconosce il momento angolare:
\begin{equation}
    \frac{dA}{dt} = \frac{r^2}{2}\,\frac{d\varphi}{dt}
    = \frac{r^2 \dot{\varphi}}{2} \cdot \frac{m}{m}
    = \frac{|\vec{L}_O|}{2m} = \text{costante}
\end{equation}
avendo usato $|\vec{L}_O| = m r^2 \dot{\varphi}$. Poiché in un campo centrale $|\vec{L}_O|$ è costante, lo è anche la velocità areolare.

\begin{teorema}{Legge delle Aree (Seconda Legge di Keplero)}
In un campo di forze centrali il raggio vettore $\vec{r}$ spazza \textbf{aree uguali in tempi uguali}.
\end{teorema}

\begin{osservazione}{Conseguenza sulla Velocità Tangenziale}
Dalla costanza di $r^2\dot{\varphi}$ segue direttamente
\begin{equation}
    r^2 \dot{\varphi} = (r \dot{\varphi})\, r = v_{\text{tang}} \cdot r = \text{costante}
    \qquad \Longrightarrow \qquad
    v_{\text{tang}} \propto \frac{1}{r}
\end{equation}
La componente tangenziale della velocità decresce dunque in proporzione inversa alla distanza: un pianeta corre più veloce al perielio e più lento all'afelio, esattamente come richiede la legge delle aree.
\end{osservazione}

La costanza della velocità areolare permette infine di legare il \textbf{periodo} orbitale all'area racchiusa dall'orbita:
\begin{equation}
    A_{\text{orb}} = \int_0^T \frac{dA}{dt} \, dt
    = \frac{dA}{dt} \int_0^T dt
    = \frac{|\vec{L}_O|}{2m} \cdot T
    \qquad \Longrightarrow \qquad
    T = \frac{2 \, m \, A_{\text{orb}}}{|\vec{L}_O|}
\end{equation}

\subsection{Energia Potenziale di una Forza \texorpdfstring{$\vec{F}_c \propto r^{-2}$}{Fc proporzionale a r alla meno due}}

Applichiamo quanto visto al caso di gran lunga più importante, quello in cui l'intensità della forza decade come l'inverso del quadrato della distanza --- la legge comune alla gravitazione e all'elettrostatica:
\begin{equation}
    F(r) = \frac{\alpha}{r^2}
\end{equation}
dove la costante $\alpha$ è negativa per una forza attrattiva e positiva per una repulsiva.

Poiché la forza è conservativa, l'energia potenziale si ottiene integrando a partire da una posizione di riferimento $r_0$:
\begin{equation}
    \begin{split}
        U(r) - U(r_0) &= -\,\mathcal{L}_{r_0 \to r} = - \int_{r_0}^{r} \frac{\alpha}{r'^2} \, dr' \\
        &= \int_{r_0}^{r} \left(-\frac{\alpha}{r'^2}\right) dr'
        = \left[ \frac{\alpha}{r'} \right]_{r_0}^{r}
        = \frac{\alpha}{r} - \frac{\alpha}{r_0}
    \end{split}
\end{equation}

Resta da fissare il riferimento. La scelta naturale è porlo all'infinito: a distanza molto grande l'interazione si estingue e sia la forza sia l'energia potenziale tendono a zero. Ponendo dunque $r_0 \to +\infty$ il termine $\alpha/r_0$ si annulla e si ottiene l'espressione definitiva:
\begin{equation}
    U(r) - \cancelto{0}{U(+\infty)} = \frac{\alpha}{r}
    \qquad \Longrightarrow \qquad
    U(r) = \frac{\alpha}{r}
\end{equation}

\begin{attenzione}{La Scelta del Riferimento è Convenzionale}
Il valore assoluto di $U$ non ha significato fisico: solo le \emph{differenze} di energia potenziale sono misurabili, perché è da esse che si ricava il lavoro. Porre $U(+\infty) = 0$ è una convenzione comoda, non un risultato: cambiandola si aggiungerebbe a $U(r)$ una costante additiva senza alterare né le forze né il moto.
\end{attenzione}

\section{Seconda Equazione Cardinale}

La \textbf{seconda equazione cardinale} è l'equazione vettoriale che governa la dinamica rotazionale, e sta al momento angolare come il secondo principio di Newton sta alla quantità di moto. La ricaviamo in quattro casi distinti, combinando due alternative: il polo può essere \textit{fisso} o \textit{mobile}, e il sistema può essere un singolo \textit{punto materiale} oppure un \textit{sistema materiale}.

\subsection{Polo Fisso per un Punto Materiale}

Consideriamo un polo $O$ fisso, coincidente con il polo $O'$ rispetto a cui valutiamo il momento, per un singolo punto materiale. Deriviamo rispetto al tempo la definizione di momento angolare:
\begin{align}
    \left. \frac{d\vec{L}_O}{dt} \right|_{O} &= \left. \frac{d\vec{L}_O}{dt} \right|_{O'} = \frac{d}{dt}\left(\vec{r} \times m\vec{v}\right) \\
    &= m \, \frac{d\vec{r}}{dt} \times \vec{v} + m\,\vec{r} \times \frac{d\vec{v}}{dt} \\
    &= m\,(\vec{v} \times \vec{v}) + m\,(\vec{r} \times \vec{a})
\end{align}
Il primo termine è nullo, perché il prodotto vettoriale di un vettore con sé stesso si annulla. Nel secondo riconosciamo $m\vec{a} = \vec{F}$:
\begin{equation*}
    = \vec{r} \times m\vec{a} = \vec{r} \times \vec{F} = \vec{\tau}_O
\end{equation*}

\begin{teorema}{Seconda Equazione Cardinale --- Polo Fisso}
Per un punto materiale e un polo $O$ fisso vale, in \textbf{forma differenziale},
\begin{equation}
    \frac{d\vec{L}_O}{dt} = \vec{\tau}_O
\end{equation}
e, integrando fra due istanti $t_0$ e $t_1$, in \textbf{forma integrale},
\begin{equation}
    I_{\vec{\tau}_O} = \int_{t_0}^{t_1} \vec{\tau}_O(t') \, dt'
    = \int_{t_0}^{t_1} \frac{d\vec{L}_O}{dt'} \, dt'
    = \vec{L}_O(t_1) - \vec{L}_O(t_0) = \Delta\vec{L}_O
\end{equation}
L'impulso angolare del momento delle forze è dunque pari alla variazione del momento angolare.
\end{teorema}

\subsection{Polo Mobile per un Punto Materiale}

Supponiamo ora che $O$ resti fisso ma che il polo $O'$ rispetto a cui calcoliamo il momento sia \textbf{mobile}, con velocità $\vec{v}_{O'} = \frac{d\,\vec{OO'}}{dt} = -\frac{d\,\vec{O'O}}{dt}$. Partiamo dalla relazione fra i momenti rispetto ai due poli, $\vec{L}_{O'} = \vec{L}_O + \vec{O'O} \times m\vec{v}$, e deriviamo:
\begin{align}
    \left. \frac{d\vec{L}_{O'}}{dt} \right|_{O'}
    &= \frac{d}{dt}\left( \vec{L}_O + \vec{O'O} \times m\vec{v} \right) \\
    &= \frac{d\vec{L}_O}{dt} + \frac{d(\vec{O'O})}{dt} \times m\vec{v} + \vec{O'O} \times m\frac{d\vec{v}}{dt} \\
    &= \vec{\tau}_O + (-\vec{v}_{O'}) \times m\vec{v} + \vec{O'O} \times m\vec{a}
\end{align}
Riscriviamo ora $\vec{\tau}_O = \vec{OP} \times \vec{F}$ e raccogliamo i due termini che contengono la forza, ricomponendo il braccio rispetto al nuovo polo:
\begin{align}
    \left. \frac{d\vec{L}_{O'}}{dt} \right|_{O'}
    &= \big(\vec{OP} \times \vec{F} + \vec{O'O} \times \vec{F}\big) - \vec{v}_{O'} \times m\vec{v} \\
    &= (\vec{O'O} + \vec{OP}) \times \vec{F} - \vec{v}_{O'} \times m\vec{v} \\
    &= \vec{O'P} \times \vec{F} - \vec{v}_{O'} \times m\vec{v}
    = \vec{\tau}_{O'} - \vec{v}_{O'} \times m\vec{v}
\end{align}

\begin{teorema}{Seconda Equazione Cardinale --- Polo Mobile}
Per un punto materiale e un polo $O'$ in moto con velocità $\vec{v}_{O'}$ vale
\begin{equation}
    \frac{d\vec{L}_{O'}}{dt} + \vec{v}_{O'} \times m\vec{v} = \vec{\tau}_{O'}
\end{equation}
Rispetto al caso di polo fisso compare il termine aggiuntivo $\vec{v}_{O'} \times m\vec{v}$, che si annulla quando il polo è fermo ($\vec{v}_{O'} = \vec{0}$) oppure quando esso si muove parallelamente al punto ($\vec{v}_{O'} \parallel \vec{v}$).
\end{teorema}

\begin{osservazione}{Il Centro di Massa è un Polo Mobile Privilegiato}
Il secondo caso di annullamento è quello che rende il centro di massa un polo così comodo: scegliendo $O' \equiv CM$ si ha $\vec{v}_{O'} = \vec{V}_{CM}$ e, per un sistema, $\sum_i m_i\vec{v}_i = M\vec{V}_{CM}$, sicché il termine aggiuntivo diventa $\vec{V}_{CM} \times M\vec{V}_{CM} = \vec{0}$. L'equazione cardinale riprende allora la forma semplice valida per i poli fissi, anche se il centro di massa è in moto.
\end{osservazione}

\subsection{Polo Fisso per un Sistema Materiale}

Passiamo ora a un sistema di $N$ punti materiali, mantenendo il polo $O \equiv O'$ fisso. Il momento angolare totale è $\vec{L}_O^{\,\text{tot}} = \sum_i \vec{r}_i \times m_i \vec{v}_i$, e derivandolo termine a termine:
\begin{align}
    \left. \frac{d\vec{L}_O^{\,\text{tot}}}{dt} \right|_{O}
    &= \frac{d}{dt}\left( \sum_i \vec{r}_i \times m_i \vec{v}_i \right)
     = \sum_i \frac{d\vec{L}_{Oi}}{dt} \\
    &= \sum_i \vec{\tau}_{Oi}
     = \sum_i \vec{\tau}_{Oi}^{\,\text{INT}} + \sum_i \vec{\tau}_{Oi}^{\,\text{EXT}}
\end{align}
avendo applicato a ciascun punto il risultato del primo caso.

Il punto cruciale è che la sommatoria dei momenti delle forze \emph{interne} si annulla: per il terzo principio esse sono a due a due uguali e opposte, e agiscono lungo la congiungente dei due punti, sicché ogni coppia ha momento risultante nullo. Quindi $\sum_i \vec{\tau}_{Oi}^{\,\text{INT}} = \vec{0}$.

\begin{teorema}{Seconda Equazione Cardinale per i Sistemi}
Per un sistema materiale e un polo fisso $O$ vale, in \textbf{forma differenziale},
\begin{equation}
    \frac{d\vec{L}_O^{\,\text{tot}}}{dt} = \sum_i \vec{\tau}_{Oi}^{\,\text{EXT}}
\end{equation}
e, in \textbf{forma integrale},
\begin{equation}
    I^{\,\text{tot}}_{\vec{\tau}_O^{\,\text{EXT}}}
    = \int_{t_0}^{t_1} \sum_i \vec{\tau}_{Oi}^{\,\text{EXT}}(t') \, dt'
    = \Delta\vec{L}_O^{\,\text{tot}}
\end{equation}
Soltanto le forze \textbf{esterne} possono modificare il momento angolare complessivo di un sistema.
\end{teorema}

\subsection{Polo Mobile per un Sistema Materiale}

L'ultimo caso si ottiene sommando su tutte le particelle la relazione trovata per il punto materiale con polo mobile, $\vec{\tau}_{O'i} = \frac{d\vec{L}_{O'i}}{dt} + \vec{v}_{O'} \times m_i\vec{v}_i$:
\begin{equation}
    \sum_i \vec{\tau}_{O'i}^{\,\text{EXT}}
    = \frac{d\vec{L}_{O'}^{\,\text{tot}}}{dt} + \sum_i \vec{v}_{O'} \times m_i \vec{v}_i
\end{equation}
da cui l'espressione conclusiva
\begin{equation}
    \frac{d\vec{L}_{O'}^{\,\text{tot}}}{dt} + \vec{v}_{O'} \times M\vec{V}_{CM}
    = \sum_i \vec{\tau}_{O'i}^{\,\text{EXT}}
\end{equation}
avendo riconosciuto $\sum_i m_i\vec{v}_i = M\vec{V}_{CM}$, la quantità di moto totale del sistema.

\subsection{Conservazione del Momento Angolare}

Dalle equazioni cardinali discende immediatamente il teorema di conservazione, che è di gran lunga il loro impiego più frequente.

\begin{teorema}{Conservazione del Momento Angolare}
Se il sistema è \textbf{isolato}, o più in generale se la risultante dei momenti delle forze esterne è nulla o trascurabile,
\begin{equation*}
    \sum_i \vec{\tau}_{Oi}^{\,\text{EXT}} \approx \vec{0}
    \qquad \Longrightarrow \qquad
    \Delta\vec{L}_O^{\,\text{tot}} = \vec{0}
\end{equation*}
allora il momento angolare totale assume lo stesso valore in istanti diversi,
\begin{equation}
    \vec{L}_O^{\,\text{tot}}(t_1) = \vec{L}_O^{\,\text{tot}}(t_0)
\end{equation}
e costituisce una \textbf{costante del moto}.
\end{teorema}

\begin{attenzione}{Conservazione del Momento Angolare e della Quantità di Moto sono Indipendenti}
Le due leggi rispondono a condizioni diverse: $\vec{P}_{\text{tot}}$ si conserva se è nulla la \emph{risultante} delle forze esterne, $\vec{L}_O^{\,\text{tot}}$ se è nullo il loro \emph{momento risultante}. Una coppia di forze uguali e opposte applicate in punti diversi ha risultante nulla ma momento non nullo: conserva la quantità di moto e non il momento angolare. Viceversa, una forza centrale conserva il momento angolare rispetto al centro ma non la quantità di moto.
\end{attenzione}

\begin{osservazione}{Identità del Doppio Prodotto Vettoriale}
Nei calcoli che seguono comparirà ripetutamente il prodotto vettoriale triplo; conviene richiamarne l'identità fondamentale (regola \emph{BAC--CAB}):
\begin{equation*}
    \vec{A} \times (\vec{B} \times \vec{C})
    = \vec{B}\,(\vec{A} \cdot \vec{C}) - \vec{C}\,(\vec{A} \cdot \vec{B})
\end{equation*}
\end{osservazione}

\section{Teorema di Koenig per il Momento Angolare}

Come già per l'energia cinetica, anche il momento angolare di un sistema ammette una scomposizione che separa il moto d'insieme dal moto interno. Indichiamo con l'apice le grandezze riferite al sistema del centro di massa:
\begin{equation}
    \vec{r}_i = \vec{r}_i\,' + \vec{R}_{CM}
    \qquad \text{e} \qquad
    \vec{v}_i = \vec{v}_i\,' + \vec{V}_{CM}
\end{equation}

Calcoliamo il momento angolare totale rispetto al polo fisso $O$ e sviluppiamo il prodotto:
\begin{align}
    \vec{L}_O &= \sum_i \vec{r}_i \times m_i \vec{v}_i = \sum_i m_i (\vec{r}_i \times \vec{v}_i) \\
    &= \sum_i m_i \left[ (\vec{r}_i\,' + \vec{R}_{CM}) \times (\vec{v}_i\,' + \vec{V}_{CM}) \right] \\
    &= \sum_i m_i \left( \vec{r}_i\,' \times \vec{v}_i\,'
       + \vec{r}_i\,' \times \vec{V}_{CM}
       + \vec{R}_{CM} \times \vec{v}_i\,'
       + \vec{R}_{CM} \times \vec{V}_{CM} \right)
\end{align}

I due termini misti si annullano grazie alla definizione stessa di centro di massa, che impone $\sum_i m_i \vec{r}_i\,' = \vec{0}$ e, derivando, $\sum_i m_i \vec{v}_i\,' = \vec{0}$:
\begin{align}
    \sum_i m_i (\vec{r}_i\,' \times \vec{V}_{CM})
    &= \Big(\underbrace{\textstyle\sum_i m_i \vec{r}_i\,'}_{=\,\vec{0}}\Big) \times \vec{V}_{CM} = \vec{0} \\
    \sum_i m_i (\vec{R}_{CM} \times \vec{v}_i\,')
    &= \vec{R}_{CM} \times \Big(\underbrace{\textstyle\sum_i m_i \vec{v}_i\,'}_{=\,\vec{0}}\Big) = \vec{0}
\end{align}

Restano soltanto il primo e il quarto termine, e quest'ultimo si compatta osservando che $\sum_i m_i = M$:
\begin{equation}
    \vec{L}_O = \underbrace{\sum_i \vec{r}_i\,' \times m_i \vec{v}_i\,'}_{\vec{L}_{O'}}
    + \underbrace{M \,(\vec{R}_{CM} \times \vec{V}_{CM})}_{\vec{L}_{CM}}
\end{equation}

\begin{teorema}{Teorema di Koenig per il Momento Angolare}
Il momento angolare totale di un sistema rispetto a un polo fisso $O$ si decompone nella somma di due contributi indipendenti:
\begin{equation}
    \vec{L}_O = \vec{L}_{O'} + \vec{L}_{CM}
\end{equation}
dove
\begin{itemize}[leftmargin=1.4em]
    \item $\vec{L}_O$ è il momento angolare totale del sistema rispetto al polo fisso $O$;
    \item $\vec{L}_{O'}$ è il momento angolare \textbf{intrinseco} (o di spin), dovuto al moto delle particelle rispetto al centro di massa;
    \item $\vec{L}_{CM} = M(\vec{R}_{CM} \times \vec{V}_{CM})$ è il momento angolare \textbf{orbitale}, quello che avrebbe un unico punto materiale di massa $M$ posto nel centro di massa.
\end{itemize}
\end{teorema}

\begin{osservazione}{Orbitale e Intrinseco}
La distinzione ricalca quella fra il moto di rivoluzione della Terra attorno al Sole --- che contribuisce a $\vec{L}_{CM}$ --- e il suo moto di rotazione attorno al proprio asse, che contribuisce a $\vec{L}_{O'}$. I due termini sono indipendenti: si possono modificare l'uno senza toccare l'altro.
\end{osservazione}

\subsection{Caso del Sistema a Due Punti Materiali}

Applichiamo il teorema al caso più importante, quello di un sistema isolato di due soli punti materiali. Accanto alle coordinate del centro di massa introduciamo le coordinate \textbf{relative}, che descrivono la posizione di un punto rispetto all'altro:
\begin{equation}
    \vec{R}_{CM} = \frac{m_1 \vec{r}_1 + m_2 \vec{r}_2}{m_1 + m_2}
    \, ; \qquad
    \begin{cases}
        \vec{r}_r = \vec{r}_2 - \vec{r}_1 \\
        \vec{v}_r = \vec{v}_2 - \vec{v}_1
    \end{cases}
\end{equation}

Invertendo queste relazioni si esprimono le posizioni assolute in funzione delle due nuove variabili:
\begin{equation}
    \begin{cases}
        \vec{r}_1 = \vec{R}_{CM} - \dfrac{m_2}{m_1 + m_2}\,\vec{r}_r \\[8pt]
        \vec{r}_2 = \vec{R}_{CM} + \dfrac{m_1}{m_1 + m_2}\,\vec{r}_r
    \end{cases}
\end{equation}
e, sottraendo $\vec{R}_{CM}$, si ottengono le posizioni e le velocità nel riferimento del centro di massa:
\begin{equation}
    \begin{cases}
        \vec{r}_1\,' = - \dfrac{m_2}{m_1 + m_2}\,\vec{r}_r \\[8pt]
        \vec{r}_2\,' = + \dfrac{m_1}{m_1 + m_2}\,\vec{r}_r
    \end{cases}
    \qquad \Longrightarrow \qquad
    \begin{cases}
        \vec{v}_1\,' = - \dfrac{m_2}{m_1 + m_2}\,\vec{v}_r \\[8pt]
        \vec{v}_2\,' = + \dfrac{m_1}{m_1 + m_2}\,\vec{v}_r
    \end{cases}
\end{equation}

Sostituiamo nell'espressione del momento angolare intrinseco:
\begin{align}
    \vec{L}_{O'} &= m_1 (\vec{r}_1\,' \times \vec{v}_1\,') + m_2 (\vec{r}_2\,' \times \vec{v}_2\,') \\
    &= m_1 \left( \frac{m_2}{m_1+m_2} \right)^{\!2} (\vec{r}_r \times \vec{v}_r)
     + m_2 \left( \frac{m_1}{m_1+m_2} \right)^{\!2} (\vec{r}_r \times \vec{v}_r) \\
    &= \frac{m_1 m_2^2 + m_2 m_1^2}{(m_1+m_2)^2} \, (\vec{r}_r \times \vec{v}_r) \\
    &= \frac{m_1 m_2 \,\cancel{(m_1 + m_2)}}{(m_1+m_2)^{\cancel{2}}} \, (\vec{r}_r \times \vec{v}_r)
     = \frac{m_1 m_2}{m_1+m_2} \, (\vec{r}_r \times \vec{v}_r)
\end{align}
(entrambi i coefficienti compaiono elevati al quadrato, dunque il segno meno di $\vec{r}_1\,'$ e $\vec{v}_1\,'$ scompare e i due contributi si sommano).

\begin{definizione}{Massa Ridotta}
Il fattore scalare che compare a moltiplicare il prodotto vettoriale prende il nome di \textbf{massa ridotta} del sistema:
\begin{equation}
    \mu = \frac{m_1 m_2}{m_1+m_2}
\end{equation}
Essa è sempre minore della più piccola delle due masse, e tende a quest'ultima quando l'altra diventa molto grande.
\end{definizione}

Il momento angolare intrinseco assume dunque la forma compatta
\begin{equation}
    \vec{L}_{O'} = \mu \,(\vec{r}_r \times \vec{v}_r)
\end{equation}
che è formalmente identica al momento angolare di \emph{un solo} punto materiale, di massa $\mu$, posizione $\vec{r}_r$ e velocità $\vec{v}_r$. Il problema a due corpi si è ridotto a un problema a un corpo: chiameremo \textbf{particella equivalente} questo punto fittizio di parametri $(\mu, \vec{r}_r, \vec{v}_r)$.

\subsection{Equivalenza del Sistema a Due Corpi}

La riduzione appena ottenuta è un'identità algebrica; resta da verificare che essa si estenda anche alla \emph{dinamica}, cioè che la particella equivalente obbedisca effettivamente a un'equazione del moto del tipo previsto. Partiamo dalla seconda equazione cardinale per il sistema, scomponendo le posizioni rispetto al centro di massa:
\begin{align}
    \frac{d\vec{L}_O^{\,\text{tot}}}{dt}
    &= \sum_i \vec{\tau}_{Oi}^{\,\text{EXT}} = \sum_i (\vec{r}_i \times \vec{F}_i^{\,\text{EXT}})
     = \sum_i (\vec{r}_i\,' + \vec{R}_{CM}) \times \vec{F}_i^{\,\text{EXT}} \\
    &= \sum_i (\vec{r}_i\,' \times \vec{F}_i^{\,\text{EXT}}) + \sum_i (\vec{R}_{CM} \times \vec{F}_i^{\,\text{EXT}}) \\
    &= \sum_i \vec{\tau}_{O'i}^{\,\text{EXT}} + \vec{R}_{CM} \times \vec{F}_{\text{TOT}}^{\,\text{EXT}}
\end{align}
che riproduce la scomposizione di Koenig anche a livello di derivate:
\begin{equation*}
    \frac{d\vec{L}_O^{\,\text{tot}}}{dt}
    = \frac{d\vec{L}_{O'}^{\,\text{tot}}}{dt} + \vec{R}_{CM} \times \frac{d\vec{P}^{\,\text{TOT}}}{dt}
\end{equation*}

Calcoliamo ora esplicitamente il termine intrinseco per il sistema a due corpi:
\begin{align}
    \frac{d\vec{L}_{O'}}{dt}
    &= \frac{d}{dt}\left( m_1 (\vec{r}_1\,' \times \vec{v}_1\,') + m_2 (\vec{r}_2\,' \times \vec{v}_2\,') \right)
\end{align}
I termini $\vec{v}_i\,' \times \vec{v}_i\,'$ si annullano, e ricordando che $\vec{v}_i\,' = \vec{v}_i - \vec{V}_{CM}$ resta
\begin{align}
    &= m_1 \left( \vec{r}_1\,' \times \frac{d\vec{v}_1}{dt} \right)
     + m_2 \left( \vec{r}_2\,' \times \frac{d\vec{v}_2}{dt} \right)
     - \big(\underbrace{m_1 \vec{r}_1\,' + m_2 \vec{r}_2\,'}_{=\,\vec{0}}\big) \times \frac{d\vec{V}_{CM}}{dt}
\end{align}
dove l'ultimo termine si annulla per la definizione di centro di massa. Traducendo le accelerazioni in forze:
\begin{align}
    \frac{d\vec{L}_{O'}}{dt}
    &= \vec{r}_1\,' \times \vec{F}_1 + \vec{r}_2\,' \times \vec{F}_2 \\
    &= \vec{r}_1\,' \times (\vec{F}_1^{\,\text{EXT}} + \vec{F}_{21})
     + \vec{r}_2\,' \times (\vec{F}_2^{\,\text{EXT}} + \vec{F}_{12})
\end{align}
Il terzo principio impone $\vec{F}_{12} = -\vec{F}_{21}$, e raccogliendo:
\begin{align}
    &= \big(\vec{r}_1\,' \times \vec{F}_1^{\,\text{EXT}} + \vec{r}_2\,' \times \vec{F}_2^{\,\text{EXT}}\big)
       + (\vec{r}_1\,' - \vec{r}_2\,') \times \vec{F}_{21} \\
    &= \vec{\tau}_{O'}^{\,\text{EXT}} + (\vec{r}_1\,' - \vec{r}_2\,') \times \vec{F}_{21}
\end{align}

Resta da valutare il termine dovuto alle forze interne. Esprimiamolo nella coordinata relativa:
\begin{equation*}
    \vec{r}_1\,' - \vec{r}_2\,'
    = (\vec{r}_1 - \vec{R}_{CM}) - (\vec{r}_2 - \vec{R}_{CM})
    = \vec{r}_1 - \vec{r}_2 = -\,\vec{r}_r
\end{equation*}
Il contributo interno vale dunque $-\,\vec{r}_r \times \vec{F}_{21}$. Ma la forza che i due corpi si scambiano è diretta lungo la loro congiungente --- è cioè una forza centrale rispetto all'altro corpo --- e quindi $\vec{F}_{21} \parallel \vec{r}_r$. Il prodotto vettoriale fra vettori paralleli è nullo:
\begin{equation}
    -\,\vec{r}_r \times \vec{F}_{21} = \vec{0}
\end{equation}

\begin{teorema}{Riduzione del Problema a Due Corpi}
Un sistema isolato di due corpi interagenti equivale in tutto e per tutto a una singola \textbf{particella equivalente} di parametri $(\mu, \vec{r}_r, \vec{v}_r)$, soggetta alla sola forza centrale diretta lungo la congiungente $\vec{r}_r$.
\end{teorema}

\begin{osservazione}{Perché la Riduzione è Così Utile}
Il teorema trasforma un problema a sei gradi di libertà (due punti nello spazio) in due problemi indipendenti e molto più semplici: il moto rettilineo uniforme del centro di massa, che è banale, e il moto di un singolo punto in un campo centrale, per il quale disponiamo già delle due costanti $E$ e $\vec{L}$. È su quest'ultimo che si concentra tutto il resto del capitolo.
\end{osservazione}

\section{Energia e Potenziale Efficace}

Studiamo ora l'energia del sistema dal punto di vista della particella equivalente. Il suo momento angolare, che sappiamo costante, si scrive in coordinate polari come
\begin{equation}
    \vec{L}_{O'} = \mu \,(\vec{r}_r \times \vec{v}_r)
    = \mu \, \big(r_r \hat{r} \times (\dot{r}_r \hat{r} + r_r \dot{\phi} \hat{\phi})\big)
    = \mu \, r_r^2 \, \omega \, \hat{z}
\end{equation}

L'energia totale è la somma di energia cinetica ed energia potenziale; esplicitando la velocità in coordinate polari:
\begin{align}
    E_{\text{tot}} &= K + U = \frac{1}{2} \mu \, \vec{v}_r^{\;2} + U(r) \\
    &= \frac{1}{2} \mu \left( \dot{r}_r^{\,2} + r_r^2 \dot{\phi}^2 \right) + U(r)
\end{align}

Il passaggio decisivo consiste nell'eliminare $\dot{\phi}$ servendosi della conservazione del momento angolare. Da $L = \mu r_r^2 \dot{\phi}$ segue $\dot{\phi} = \dfrac{L}{\mu r_r^2}$, e sostituendo:
\begin{equation}
    E_{\text{tot}} = \frac{1}{2} \mu \dot{r}_r^{\,2} + \frac{1}{2}\,\frac{L^2}{\mu r_r^2} + U(r)
\end{equation}
L'energia dipende ora dalla sola coordinata radiale: il problema bidimensionale è diventato \textbf{unidimensionale}.

\subsection{Definizione di Potenziale Efficace}

Raggruppiamo tutti i termini che dipendono esclusivamente da $r$:
\begin{equation}
    E_{\text{tot}} = \underbrace{\frac{1}{2} \mu \dot{r}_r^{\,2}}_{K_{\text{rad}}}
    + \underbrace{\frac{1}{2}\,\frac{L^2}{\mu r_r^2} + U(r)}_{U_{\text{eff}}(r)}
\end{equation}
dove si riconoscono:
\begin{itemize}
    \item $\frac{1}{2}\mu \dot{r}_r^{\,2}$: l'\textbf{energia cinetica radiale};
    \item $\frac{L^2}{2\mu r^2}$: il \textbf{potenziale centrifugo}, che non è una vera energia potenziale ma il residuo dell'energia cinetica di rotazione, riscritto come funzione di $r$;
    \item $U(r)$: il potenziale dell'interazione centrale vera e propria.
\end{itemize}

\begin{definizione}{Potenziale Efficace}
Si definisce \textbf{potenziale efficace} la funzione
\begin{equation}
    U_{\text{eff}}(r) = U(r) + \frac{L^2}{2\mu r^2}
\end{equation}
Grazie a essa il moto radiale della particella equivalente si tratta esattamente come un moto unidimensionale in un campo di energia potenziale $U_{\text{eff}}$.
\end{definizione}

\begin{attenzione}{Il Potenziale Centrifugo Dipende dalle Condizioni Iniziali}
A differenza di $U(r)$, che è fissato una volta per tutte dalla natura dell'interazione, il termine centrifugo contiene $L$, che dipende da \emph{come} il moto è stato avviato. Cambiando il momento angolare iniziale cambia la forma di $U_{\text{eff}}$, e con essa il tipo di orbita: lo stesso campo di forze produce traiettorie qualitativamente diverse a seconda di $L$.
\end{attenzione}

\subsection{Condizione di Esistenza del Moto}

Perché il moto sia fisicamente possibile l'energia cinetica radiale deve essere non negativa. Questa richiesta si traduce in un vincolo fra l'energia totale e il potenziale efficace:
\begin{equation}
    \frac{1}{2} \mu \dot{r}^2 \geq 0
    \quad \Longrightarrow \quad E - U_{\text{eff}}(r) \geq 0
    \quad \Longrightarrow \quad E \geq U_{\text{eff}}(r)
\end{equation}
La condizione delimita i valori di $r$ che la particella può effettivamente assumere, per una data energia $E$.

\subsection{Regioni Accessibili e Punti di Inversione}

Il moto è quindi consentito solo dove la curva $U_{\text{eff}}(r)$ sta al di sotto del livello $E$; le regioni in cui ciò accade si dicono \textbf{accessibili}, le altre proibite. Il segno della velocità radiale distingue le due fasi del moto:
\begin{itemize}
    \item $\dot{r} > 0$: la particella si \textbf{allontana} dal centro di forza;
    \item $\dot{r} < 0$: la particella si \textbf{avvicina} al centro di forza;
    \item $\dot{r} = 0$: la particella ha esaurito l'energia disponibile per il moto radiale ed è costretta a invertire il verso. Il punto corrispondente si dice \textbf{punto di inversione} (o punto nodale).
\end{itemize}

Nei punti di inversione la velocità è puramente tangenziale, poiché la componente radiale si annulla:
\begin{equation}
    \vec{v} = \cancelto{0}{\dot{r}\hat{r}} + r\dot{\phi}\hat{\phi} = r\omega\hat{\phi} = \vec{v}_{\text{tang}}
\end{equation}

\begin{figure}[ht]
    \centering
    \begin{tikzpicture}[x=1.55cm, y=1.15cm, >=stealth, font=\small]
        % Assi
        \draw[->, thick, mypen] (0,-2.6) -- (0,3.6) node[above] {$U_{\text{eff}}(r)$};
        \draw[->, thick, mypen] (-0.2,0) -- (6.3,0) node[right] {$r$};

        % Livelli energetici
        \draw[dashed, boxoss]  (0,2.5)  node[left, boxoss]  {$E_1$} -- (5.8,2.5);
        \draw[dashed, boxatt]  (0,0.8)  node[left, boxatt]  {$E_2$} -- (5.8,0.8);
        \draw[dashed, boxese]  (0,-0.5) node[left, boxese]  {$E_3$} -- (5.8,-0.5);
        \draw[dashed, boxnota] (0,-1.2) node[left, boxnota] {$E_4$} -- (3.5,-1.2);
        \draw[dashed, notered] (0,-1.8) node[left, notered] {$E_5$} -- (2.6,-1.8);

        % Curva del potenziale efficace
        \draw[ultra thick, mypen, smooth, tension=0.7] plot coordinates {
            (0.35,3.4) (0.6,1.0) (1.5,-1.8) (2.8,0.8) (4.5,0.2) (5.8,0.05)
        };
        \node[mypen, right] at (5.0,0.62) {$U_{\text{eff}}(r)$};

        % Punti di inversione
        \fill[boxoss]  (0.42,2.5)  circle (2pt) node[above right, boxoss, xshift=-2pt]  {$r_0$};
        \fill[boxatt]  (0.65,0.8)  circle (2pt) node[below left, boxatt, xshift=4pt]    {$r_1$};
        \fill[boxatt]  (2.8,0.8)   circle (2pt) node[above, boxatt]                     {$r_7$};
        \fill[boxese]  (0.90,-0.5) circle (2pt) node[below left, boxese, xshift=4pt]    {$r_2$};
        \fill[boxese]  (2.13,-0.5) circle (2pt) node[below right, boxese, xshift=-3pt]  {$r_6$};
        \fill[boxese]  (3.80,-0.5) circle (2pt) node[above right, boxese, xshift=-3pt]  {$r_8$};
        \fill[boxnota] (1.10,-1.2) circle (2pt) node[below left, boxnota, xshift=5pt]   {$r_3$};
        \fill[boxnota] (1.85,-1.2) circle (2pt) node[below right, boxnota, xshift=-4pt] {$r_5$};
        \fill[notered] (1.50,-1.8) circle (2pt) node[below, notered]                    {$r_4$};
    \end{tikzpicture}
    \caption{Andamento qualitativo del potenziale efficace e classificazione delle orbite. Ogni livello orizzontale rappresenta un valore dell'energia totale; le sue intersezioni con la curva sono i punti di inversione $r_i$, che delimitano le regioni accessibili al moto.}
    \label{fig:potenziale_efficace}
\end{figure}

Il tipo di orbita dipende dunque interamente dalle \textbf{condizioni iniziali}, riassunte nella coppia $(L_0, E_0)$. Scorrendo la figura dall'alto verso il basso:

\begin{itemize}
    \item $\mathbf{E_0 = E_1}$: il moto è accessibile per $r \geq r_0$. La particella giunge da grande distanza, raggiunge la minima distanza $r_0$ e torna verso l'infinito: si tratta di un \textbf{moto aperto} (illimitato).
    \item $\mathbf{E_0 = E_2}$: il moto è accessibile per $r \geq r_1$. Oltre alla traiettoria aperta di avvicinamento e allontanamento, in $r = r_7$ --- dove il livello è tangente al \emph{massimo} della curva --- esiste un'orbita circolare, ma \textbf{instabile}: una perturbazione infinitesima allontana definitivamente la particella da essa.
    \item $\mathbf{E_0 = E_3}$: la curva presenta ora due regioni accessibili separate. La particella può restare intrappolata nella buca, $r_2 \leq r \leq r_6$, oppure muoversi liberamente all'esterno per $r \geq r_8$. Quale dei due destini si realizzi dipende da dove essa si trovi inizialmente: le due regioni non comunicano.
    \item $\mathbf{E_0 = E_4}$: il moto è accessibile solo per $r_3 \leq r \leq r_5$. La particella oscilla confinata fra due \textbf{absidi}, alternando avvicinamenti e allontanamenti: il \textbf{moto è limitato}.
    \item $\mathbf{E_0 = E_5}$: il livello è tangente al \emph{minimo} della curva, dunque $K_{\text{rad}} = 0$ in permanenza e l'unico valore possibile è $r = r_4$. Il raggio non varia mai: l'orbita è \textbf{circolare e stabile}, e corrisponde alla minima energia compatibile con il momento angolare assegnato.
    \item $\mathbf{E_0 < E_5}$: si avrebbe $K_{\text{rad}} < 0$, che è impossibile. Nessun moto è consentito.
\end{itemize}

\begin{osservazione}{Stabilità e Concavità}
La differenza fra l'orbita circolare in $r_7$ e quella in $r_4$ non sta nell'energia ma nella \emph{concavità} della curva: entrambe annullano $dU_{\text{eff}}/dr$, ma solo il minimo ($d^2U_{\text{eff}}/dr^2 > 0$) richiama la particella verso l'orbita dopo una perturbazione. Il massimo la respinge, ed è per questo instabile.
\end{osservazione}

\subsection{Orbite Circolari nei Sistemi a Due Corpi}

Il caso del minimo di $U_{\text{eff}}$ merita un esame a parte, perché descrive le orbite circolari --- la configurazione di gran lunga più comune nei sistemi astronomici. Per una singola particella la condizione è semplicemente $r = \text{costante}$; vediamo che cosa significhi per un sistema a due corpi.

Nel sistema isolato la particella equivalente ha raggio costante, e dunque resta costante la \textbf{distanza relativa} fra i due corpi:
\begin{equation}
    r_r = |\vec{r}_2 - \vec{r}_1| = |\vec{r}_2\,' - \vec{r}_1\,'| = \text{costante}
\end{equation}

Richiamando le posizioni rispetto al centro di massa già ricavate,
\begin{equation}
    \begin{cases}
        \vec{r}_1\,' = - \dfrac{m_2}{m_1+m_2}\,\vec{r}_r \\[8pt]
        \vec{r}_2\,' = + \dfrac{m_1}{m_1+m_2}\,\vec{r}_r
    \end{cases}
\end{equation}
si vede subito che, essendo $|\vec{r}_r|$ costante, sono costanti anche $|\vec{r}_1\,'|$ e $|\vec{r}_2\,'|$. Da questa sola osservazione discendono tre conclusioni.

\begin{enumerate}
    \item \textbf{Entrambi i corpi percorrono un'orbita circolare} attorno al centro di massa, con raggi $|\vec{r}_1\,'|$ e $|\vec{r}_2\,'|$ inversamente proporzionali alle rispettive masse: il corpo più pesante descrive il cerchio più piccolo.
    \item \textbf{La velocità angolare è la stessa per entrambi} ed è costante. I due corpi restano infatti sempre diametralmente opposti rispetto al centro di massa --- i vettori $\vec{r}_1\,'$ e $\vec{r}_2\,'$ sono antiparalleli per costruzione --- e quindi spazzano lo stesso angolo nello stesso tempo.
    \item \textbf{Il centro di forza di ciascun corpo è l'altro corpo}, non il centro di massa. Quest'ultimo è solo il punto attorno al quale le due orbite appaiono descritte: nessuna massa vi è collocata e nessuna forza vi è applicata.
\end{enumerate}

\begin{figure}[ht]
    \centering
    \begin{tikzpicture}[>=stealth, scale=1.5, font=\small]
        \def\rA{1.15}   % raggio del corpo pesante (orbita interna)
        \def\rB{2.50}   % raggio del corpo leggero (orbita esterna)
        \def\angT{250}  % posizione di m_1 all'istante t
        \def\angP{210}  % posizione di m_1 all'istante t'

        % Orbite
        \draw[dashed, thick, mypen!40] (0,0) circle (\rA);
        \draw[dashed, thick, mypen!40] (0,0) circle (\rB);

        % Verso di rotazione
        \draw[->, thick, mypen!65] (100:\rB+0.30) arc (100:140:\rB+0.30);
        \node[mypen!80] at (120:\rB+0.58) {$\vec{\omega}$};

        % ---- istante t' (disegnato prima, in secondo piano) ----
        \coordinate (Ap) at (\angP:\rA);
        \coordinate (Bp) at ({\angP+180}:\rB);
        \draw[thick, mypen!40, dashed] (Ap) -- (Bp);
        \filldraw[notered!18, draw=notered!55, thick] (Ap) circle (0.16);
        \filldraw[boxoss!18, draw=boxoss!55, thick] (Bp) circle (0.12);
        \node[notered!75] at (\angP:\rA+0.48) {$m_1(t')$};
        \node[boxoss!75] at ({\angP+180}:\rB+0.42) {$m_2(t')$};

        % ---- istante t ----
        \coordinate (A) at (\angT:\rA);
        \coordinate (B) at ({\angT+180}:\rB);
        \draw[very thick, mypen] (A) -- (B);
        \filldraw[notered!45, draw=notered, very thick] (A) circle (0.16);
        \filldraw[boxoss!45, draw=boxoss, very thick] (B) circle (0.12);
        \node[notered] at (\angT:\rA+0.48) {$m_1(t)$};
        \node[boxoss] at ({\angT+180}:\rB+0.40) {$m_2(t)$};

        % Raggi rispetto al centro di massa
        \draw[->, very thick, notered] (0,0) -- (A)
            node[pos=0.62, anchor=west, notered, xshift=2pt] {$\vec{r}_1\,'$};
        \draw[->, very thick, boxoss] (0,0) -- (B)
            node[pos=0.55, anchor=east, boxoss, xshift=-2pt] {$\vec{r}_2\,'$};

        % Angoli spazzati, uguali sui due lati
        \draw[thick, boxese] (\angT:0.58) arc (\angT:\angP:0.58);
        \node[boxese] at ({(\angT+\angP)/2}:0.82) {$\theta$};
        \draw[thick, boxese] ({\angT+180}:1.62) arc ({\angT+180}:{\angP+180}:1.62);
        \node[boxese] at ({(\angT+\angP)/2+180}:1.86) {$\theta$};

        % Centro di massa
        \filldraw[boxese] (0,0) circle (2.0pt);
        \node[boxese, anchor=south west] at (0.10,0.10) {$CM$};
    \end{tikzpicture}
    \caption{Orbite circolari in un sistema isolato a due corpi, ritratte in due istanti successivi $t$ e $t'$. I due corpi restano costantemente allineati con il centro di massa e su lati opposti rispetto a esso, e spazzano nello stesso tempo il medesimo angolo $\theta$: la velocità angolare è comune. Il corpo più pesante ($m_1$) percorre l'orbita di raggio minore.}
    \label{fig:orbite_circolari_dettagliate}
\end{figure}

\subsection{Orbite Circolari e Punti Stazionari del Potenziale Efficace}

Verifichiamo ora per via analitica quanto la figura mostra geometricamente, cercando i punti stazionari del potenziale efficace:
\begin{equation}
    U_{\text{eff}}(r) = \frac{L_0^2}{2\mu r^2} + U(r)
\end{equation}
Deriviamo rispetto al raggio:
\begin{align}
    \frac{dU_{\text{eff}}}{dr}
    &= \frac{dU}{dr} + \frac{L_0^2}{2\mu} \left( -\frac{2}{r^3} \right)
     = \frac{dU}{dr} - \frac{L_0^2}{\mu r^3} \\
    &= \frac{dU}{dr} - \frac{(\mu r^2 \omega)^2}{\mu r^3}
     = \frac{dU}{dr} - \frac{\mu^{\cancel{2}} r^{\cancel{4}} \omega^2}{\cancel{\mu} \, r^{\cancel{3}}}
     = \frac{dU}{dr} - \mu r \omega^2
\end{align}
avendo sostituito $L_0 = \mu r^2 \omega$. Nel caso del moto circolare il raggio si colloca in un punto stazionario, dove la derivata si annulla:
\begin{equation}
    \frac{dU_{\text{eff}}}{dr} = 0
    \quad \Longrightarrow \quad
    \frac{dU}{dr} - \mu r \omega^2 = 0
\end{equation}

Ricordiamo infine che una forza conservativa è l'opposto del gradiente del potenziale; proiettando lungo $\hat{r}$:
\begin{equation}
    -\frac{dU}{dr}\,\hat{r} = -\nabla U = \vec{F}
    \quad \Longrightarrow \quad
    \left(-\frac{dU}{dr}\hat{r}\right) \cdot \hat{r} = \vec{F} \cdot \hat{r} = F(r)
\end{equation}
cioè $-\dfrac{dU}{dr} = F(r)$. Sostituendo nella condizione di stazionarietà si ottiene
\begin{equation}
    -F(r) - \mu r \omega^2 = 0
    \quad \Longrightarrow \quad
    F(r) = - \mu r \omega^2
\end{equation}

\begin{osservazione}{Una Verifica di Coerenza}
Il risultato non è una nuova legge, ma la conferma che tutto il formalismo è consistente: $-\mu r\omega^2$ è esattamente la forza centripeta necessaria a mantenere la particella equivalente su una circonferenza di raggio $r$ alla velocità angolare $\omega$. Il segno negativo indica che la forza è diretta verso il centro, come dev'essere. Ritrovare per questa via un risultato elementare della cinematica circolare è la miglior verifica formale dei passaggi svolti.
\end{osservazione}

\subsection{Soluzione Formale per il Moto Radiale}

Concludiamo ricavando la legge oraria del moto radiale. Dalla conservazione dell'energia si esplicita la velocità radiale:
\begin{equation}
    \frac{1}{2}\mu \dot{r}^2 = E - U_{\text{eff}}(r)
    \quad \Longrightarrow \quad
    \dot{r} = \pm \sqrt{\frac{2}{\mu}\big(E - U_{\text{eff}}(r)\big)}
\end{equation}
Il doppio segno corrisponde alle due fasi del moto:
\begin{itemize}
    \item segno $(+)$, con $\dot{r} > 0$: fase di \textbf{allontanamento} dal centro di forza;
    \item segno $(-)$, con $\dot{r} < 0$: fase di \textbf{avvicinamento} al centro di forza.
\end{itemize}

Scrivendo $\dot{r} = dr/dt$ e separando le variabili:
\begin{equation}
    \frac{dr}{\sqrt{\dfrac{2}{\mu}\big(E - U_{\text{eff}}\big)}} = \pm \, dt
\end{equation}
Integrando fra l'istante iniziale $t_0$, in cui il raggio vale $r_0$, e l'istante generico $t$:
\begin{equation}
    \int_{r_0}^{r(t)} \frac{dr'}{\sqrt{\dfrac{2}{\mu}\big(E - U_{\text{eff}}(r')\big)}} = \pm \,(t - t_0)
\end{equation}

\begin{formulario}{Soluzione Formale del Moto Radiale}
\begin{equation}
    t = t_0 \pm \int_{r_0}^{r(t)} \frac{dr'}{\sqrt{\dfrac{2}{\mu}\big(E - U_{\text{eff}}(r')\big)}}
\end{equation}
\end{formulario}

\begin{osservazione}{Che Cosa Significa ``Formale''}
La soluzione è detta formale perché fornisce $t$ in funzione di $r$, cioè la funzione \emph{inversa} di quella cercata: per ottenere $r(t)$ occorre risolvere l'integrale --- il che è possibile in forma chiusa solo per pochi potenziali notevoli, fra cui quello kepleriano --- e poi invertire il risultato.

Ciò nonostante l'espressione è tutt'altro che inutile. Il doppio segno permette di ricostruire il moto a tratti, raccordando gli intervalli in cui $r(t)$ cresce ($\dot{r} > 0$) con quelli in cui decresce ($\dot{r} < 0$), separati dai punti di inversione. E anche senza risolverla si può usarla per calcolare quantità globali: valutata fra due punti di inversione consecutivi, essa fornisce direttamente il \textbf{periodo} dell'oscillazione radiale.
\end{osservazione}

\chapter{La Gravitazione}

\section{Storia e Scoperte in Campo Astronomico}

La gravitazione è il primo caso nella storia della fisica in cui un'unica legge matematica ha reso conto di fenomeni fino ad allora ritenuti di natura completamente diversa --- la caduta dei gravi sulla Terra e il moto dei pianeti nel cielo. Vale la pena ripercorrere brevemente la strada che vi ha condotto, perché il capitolo precedente ne contiene già tutti gli strumenti matematici: la gravitazione non è che il più importante esempio di forza centrale.

Il Sistema Solare è costituito dal Sole e dagli otto pianeti che gli orbitano attorno: Mercurio, Venere, Terra, Marte, Giove, Saturno, Urano e Nettuno. I Greci conoscevano soltanto quelli visibili a occhio nudo, fino a Saturno; Urano fu scoperto da Herschel nel 1781 e Nettuno nel 1846, mentre Plutone --- individuato nel 1930 --- è stato riclassificato come \emph{pianeta nano} nel 2006.

La concezione greca del cosmo era \textbf{geocentrica}: formalizzata da Tolomeo, ma già presente in Aristotele, essa prevedeva
\begin{itemize}
    \item la \textbf{Terra immobile al centro} dell'universo;
    \item il Sole, la Luna e i pianeti in moto su \textbf{orbite circolari}, il cerchio essendo ritenuto la sola figura perfetta;
    \item una separazione netta fra il mondo \textbf{sublunare}, imperfetto, dominato dalla corruzione e dal moto rettilineo, e quello \textbf{sovralunare}, perfetto, incorruttibile e sede del solo moto circolare, senza principio né fine.
\end{itemize}

\subsection{L'Evoluzione verso il Modello Newtoniano}

Il modello geocentrico si scontrava però con un'osservazione imbarazzante: il \textbf{moto retrogrado}. Osservati da Terra, i pianeti esterni come Marte non avanzano uniformemente fra le stelle fisse, ma di tanto in tanto rallentano, tornano indietro per qualche settimana e poi riprendono il cammino, descrivendo un cappio.

Per salvare i moti circolari Tolomeo introdusse il sistema degli \textbf{epicicli} e dei \textbf{deferenti}: il pianeta percorre un cerchio minore, l'epiciclo, il cui centro a sua volta ruota attorno alla Terra lungo un cerchio maggiore, il deferente. La composizione dei due moti circolari produce una traiettoria a cappi che riproduce il moto retrogrado osservato.

\begin{figure}[ht]
    \centering
    \begin{tikzpicture}[>=stealth, scale=1.0, font=\small]
        \def\Rd{2.25}   % raggio del deferente
        \def\Re{0.78}   % raggio dell'epiciclo
        \def\kk{5}      % rapporto fra le velocità angolari

        % Traiettoria risultante (epitrocoide, con i cappi)
        \draw[very thick, notered, domain=0:360, samples=400, smooth]
            plot ({\Rd*cos(\x) + \Re*cos(\kk*\x)}, {\Rd*sin(\x) + \Re*sin(\kk*\x)});

        % Deferente
        \draw[dashed, thick, mypen!50] (0,0) circle (\Rd);

        % Terra al centro
        \filldraw[boxoss!45, draw=boxoss, very thick] (0,0) circle (0.17);
        \node[mypen, anchor=north east] at (-0.1,-0.1) {Terra};

        % Centro dell'epiciclo a un istante dato
        \def\ang{58}
        \coordinate (C) at ({\Rd*cos(\ang)}, {\Rd*sin(\ang)});
        \draw[thick, mypen!70] (0,0) -- (C);
        \fill[mypen] (C) circle (1.6pt);
        \node[mypen, anchor=south east, xshift=2pt] at (C) {$C$};

        % Epiciclo
        \draw[thick, boxese] (C) circle (\Re);
        \coordinate (P) at ({\Rd*cos(\ang) + \Re*cos(\kk*\ang)},
                            {\Rd*sin(\ang) + \Re*sin(\kk*\ang)});
        \draw[thick, boxese] (C) -- (P);
        \filldraw[notered!45, draw=notered, very thick] (P) circle (0.115);
        \node[notered, anchor=west, xshift=3pt] at (P) {pianeta};

        % Etichette
        \node[mypen!75, anchor=north] at (-1.55,-2.35) {deferente};
        \draw[thin, mypen!60] (-1.55,-2.3) -- (-1.75,-1.95);
        \node[boxese, anchor=west] at (2.25,-1.75) {epiciclo};
        \draw[thin, boxese!70] (2.2,-1.72) -- (1.55,-1.42);
    \end{tikzpicture}
    \caption{Il sistema degli epicicli e dei deferenti. Il pianeta percorre l'epiciclo (in verde) mentre il centro $C$ di quest'ultimo percorre il deferente attorno alla Terra: la traiettoria risultante (in rosso) presenta i cappi che spiegano il moto retrogrado osservato.}
    \label{fig:epicicli}
\end{figure}

Nel corso dei secoli la visione dell'universo cambiò radicalmente, grazie a quattro contributi decisivi:
\begin{itemize}
    \item \textbf{Copernico} sposta il centro del sistema dalla Terra al Sole (\textbf{eliocentrismo}), riducendo drasticamente il numero di epicicli necessari;
    \item \textbf{Brahe} raccoglie, senza telescopio, osservazioni di precisione senza precedenti, e \textbf{Keplero} le riassume in tre leggi puramente empiriche, abbandonando per primo il dogma della circolarità in favore dell'ellisse;
    \item \textbf{Galileo} introduce il telescopio e osserva le fasi di Venere --- incompatibili col geocentrismo --- i satelliti di Giove e le macchie solari, che segnano la fine dell'incorruttibilità celeste;
    \item \textbf{Newton} unifica infine la fisica terrestre e quella celeste, mostrando che una sola legge di forza rende conto tanto delle tre leggi di Keplero quanto della caduta dei gravi.
\end{itemize}

\section{Legge di Gravitazione Universale}

Sulla base dei lavori di Keplero, Newton riconobbe la natura della forza responsabile del moto dei pianeti: essa è \textbf{attrattiva}, \textbf{centrale}, proporzionale al prodotto delle masse e inversamente proporzionale al quadrato della distanza.

\begin{teorema}{Legge di Gravitazione Universale}
Due masse puntiformi $m_1$ e $m_2$ poste a distanza $r$ si attraggono con una forza diretta lungo la congiungente, di modulo
\begin{equation}
    F = G \, \frac{m_1 m_2}{r^2}
\end{equation}
Detto $\vec{r}_{12} = \vec{r}_2 - \vec{r}_1$ il vettore che va da $m_1$ a $m_2$, e $\hat{r}_{12}$ il corrispondente versore, la forza esercitata da $m_1$ \emph{su} $m_2$ si scrive
\begin{equation}
    \vec{F}_{12} = - \, \frac{G m_1 m_2}{|\vec{r}_{12}|^3} \, \vec{r}_{12}
                 = - \, \frac{G m_1 m_2}{r^2} \, \hat{r}_{12}
\end{equation}
mentre quella esercitata da $m_2$ su $m_1$ è, per il terzo principio,
\begin{equation}
    \vec{F}_{21} = -\vec{F}_{12} = + \, \frac{G m_1 m_2}{r^2} \, \hat{r}_{12}
\end{equation}
\end{teorema}

\begin{attenzione}{Il Significato del Segno Meno}
Il segno negativo in $\vec{F}_{12}$ non è una convenzione arbitraria: esprime il carattere \textbf{attrattivo} della forza. Il versore $\hat{r}_{12}$ punta da $m_1$ verso $m_2$, e la forza che $m_1$ esercita su $m_2$ è diretta in verso opposto, cioè verso $m_1$. È questa la differenza formale fondamentale rispetto alla forza di Coulomb, che ha la stessa struttura matematica ma può assumere entrambi i segni.
\end{attenzione}

Essendo centrale e proporzionale a $1/r^2$, la forza gravitazionale rientra esattamente nella classe studiata nel capitolo precedente: tutto ciò che vi abbiamo dimostrato --- conservazione del momento angolare, planarità dell'orbita, legge delle aree, riduzione a particella equivalente e potenziale efficace --- si applica qui senza modifiche. In particolare, l'energia potenziale gravitazionale si ottiene dalla formula generale $U(r) = \alpha/r$ ponendo $\alpha = -Gm_1m_2$:
\begin{equation}
    U(r) = - \, \frac{G m_1 m_2}{r}
\end{equation}
con il riferimento fissato all'infinito, $U(+\infty) = 0$. L'energia potenziale è sempre \textbf{negativa}, e tende a zero solo a distanza infinita: è il segno caratteristico di un legame.

\subsection{La Bilancia di Cavendish (1798)}

La legge di Newton contiene una costante universale $G$ il cui valore non può essere dedotto per via teorica, ma solo misurato. L'impresa riuscì a \textbf{Henry Cavendish} nel 1798, con una bilancia a torsione di straordinaria sensibilità.

L'apparato consiste in un'asticella leggera di lunghezza $L$, con due piccole masse $m$ alle estremità, sospesa al centro a un sottile filo di torsione. Avvicinando due grandi masse di piombo $M$, l'attrazione gravitazionale produce una coppia che ruota l'asticella di un angolo $\theta$, finché la torsione elastica del filo non la equilibra:
\begin{equation}
    \tau_{\text{tors}} = c \, \theta = F \cdot L
    \qquad \Longrightarrow \qquad
    c \, \theta = \left( G \, \frac{M m}{r^2} \right) L
\end{equation}
dove $c$ è la costante di torsione del filo, misurabile a parte dal periodo delle oscillazioni libere dell'asticella.

\begin{figure}[ht]
    \centering
    \begin{tikzpicture}[>=stealth, scale=1.05, font=\small]
        % Posizione ruotata dell'asticella
        \draw[dashed, thick, mypen!45] (192:2.4) -- (12:2.4);
        \draw[thick, boxese] (1.55,0) arc (0:12:1.55);
        \node[boxese, anchor=west, xshift=2pt] at (1.6,0.17) {$\theta$};

        % Asticella
        \draw[very thick, mypen] (-2.4,0) -- (2.4,0);
        \draw[<->, thin, mypen!70] (-2.4,1.72) -- (2.4,1.72);
        \node[mypen, anchor=south] at (0,1.74) {$L$};

        % Masse piccole
        \filldraw[boxoss!45, draw=boxoss, very thick] (-2.4,0) circle (0.16);
        \filldraw[boxoss!45, draw=boxoss, very thick] ( 2.4,0) circle (0.16);
        \node[boxoss, anchor=east, xshift=-4pt] at (-2.4,0) {$m$};
        \node[boxoss, anchor=west, xshift=4pt]  at ( 2.4,0) {$m$};

        % Masse grandi
        \filldraw[notered!40, draw=notered, very thick] (-2.4,-1.15) circle (0.32);
        \filldraw[notered!40, draw=notered, very thick] ( 2.4, 1.15) circle (0.32);
        \node[notered] at (-2.4,-1.15) {$M$};
        \node[notered] at ( 2.4, 1.15) {$M$};

        % Forze attrattive: formano una coppia
        \draw[->, very thick, boxatt] (-2.4,-0.24) -- (-2.4,-0.74);
        \draw[->, very thick, boxatt] ( 2.4, 0.24) -- ( 2.4, 0.74);
        \node[boxatt, anchor=east, xshift=-4pt] at (-2.45,-0.5) {$F$};
        \node[boxatt, anchor=west, xshift=4pt]  at ( 2.45, 0.5) {$F$};

        % Distanza r
        \draw[<->, thin, mypen!70] (-1.95,0) -- (-1.95,-1.15);
        \node[mypen, anchor=west, xshift=2pt] at (-1.92,-0.58) {$r$};

        % Filo di torsione, visto dall'alto, e specchio
        \draw[line width=1.8pt, boxese] (-0.33,0) -- (0.33,0);
        \filldraw[white, draw=mypen, thick] (0,0) circle (0.11);
        \draw[thin, mypen] (-0.078,-0.078) -- (0.078,0.078);
        \draw[thin, mypen] (-0.078,0.078) -- (0.078,-0.078);
        \node[mypen, anchor=east] at (-0.6,1.12) {filo di torsione ($c$)};
        \draw[thin, mypen!65] (-0.57,1.06) -- (-0.12,0.15);
        \node[boxese, anchor=east] at (-0.55,-0.72) {specchio};
        \draw[thin, boxese!70] (-0.52,-0.67) -- (-0.22,-0.08);

        % Raggio luminoso
        \draw[->, thick, boxatt!85] (-1.25,-2.25) -- (-0.1,-0.12);
        \draw[->, thick, boxatt!85] (0.1,-0.12) -- (1.25,-2.25);
        \node[boxatt, anchor=north] at (-1.3,-2.3) {sorgente};
        \node[boxatt, anchor=north] at (1.35,-2.3) {scala graduata};
    \end{tikzpicture}
    \caption{Schema della bilancia di torsione di Cavendish, vista dall'alto. Le due grandi masse $M$ sono disposte da lati opposti dell'asticella, così che le attrazioni sulle masse $m$ formino una \emph{coppia}: essa torce il filo di un angolo $\theta$, finché la reazione elastica non la equilibra. Un raggio luminoso riflesso da uno specchio solidale con l'asticella amplifica la rotazione fino a renderla leggibile su una scala graduata distante.}
    \label{fig:cavendish}
\end{figure}

Misurando l'angolo $\theta$ --- amplificato dal raggio di luce riflesso dallo specchio su una scala graduata distante --- si ricava il valore
\begin{equation}
    G = 6{,}674 \times 10^{-11} \; \frac{\mathrm{N} \cdot \mathrm{m}^2}{\mathrm{kg}^2}
\end{equation}

\begin{osservazione}{Perché la Misura è Così Difficile}
$G$ è, di gran lunga, la meno nota fra le costanti fondamentali: la conosciamo con appena quattro o cinque cifre significative, contro le dodici e più di altre costanti. La ragione è l'estrema debolezza dell'interazione gravitazionale fra masse di laboratorio: l'attrazione fra le sfere di Cavendish era dell'ordine di $10^{-7} \, \mathrm{N}$, e non esiste modo di schermare la gravità dei corpi circostanti come si fa con i campi elettrici.
\end{osservazione}

\section{Massa e Campo Gravitazionale}

\subsection{Principio di Equivalenza}

La massa compare nella meccanica in due ruoli concettualmente indipendenti:
\begin{itemize}
    \item la \textbf{massa inerziale} $m_i$ misura la resistenza del corpo all'accelerazione, ed è quella che compare nel secondo principio, $F = m_i a$;
    \item la \textbf{massa gravitazionale} $m_g$ misura l'intensità con cui il corpo interagisce con il campo gravitazionale, ed è quella che compare nella legge di Newton, $F = G M m_g / r^2$.
\end{itemize}

Nulla, a priori, impone che le due coincidano: la prima è una proprietà dinamica, la seconda una sorta di ``carica'' gravitazionale, analoga alla carica elettrica. Eppure l'esperienza mostra che tutti i corpi cadono con la stessa accelerazione $g$. Scriviamo l'equazione del moto di un corpo in caduta libera sulla superficie terrestre:
\begin{equation}
    m_{i} \, a = G \, \frac{M_T \, m_{g}}{R_T^2}
    \qquad \Longrightarrow \qquad
    a = \left( G \, \frac{M_T}{R_T^2} \right) \frac{m_g}{m_i}
\end{equation}

Perché risulti $a = g$ per \emph{qualunque} corpo, indipendentemente dalla sua composizione, il rapporto $m_g/m_i$ deve essere una costante universale; scegliendo opportunamente le unità la si pone uguale a uno.

\begin{teorema}{Principio di Equivalenza (forma debole)}
Massa inerziale e massa gravitazionale sono proporzionali fra loro e, con una scelta opportuna delle unità, numericamente identiche:
\begin{equation}
    m_g \equiv m_i \equiv m
\end{equation}
Di conseguenza l'accelerazione di caduta libera è la stessa per tutti i corpi e vale
\begin{equation}
    g = G \, \frac{M_T}{R_T^2}
\end{equation}
\end{teorema}

\begin{osservazione}{Una Coincidenza che non è una Coincidenza}
Nella meccanica newtoniana l'identità $m_g = m_i$ resta un fatto sperimentale privo di spiegazione --- per quanto verificato con precisione altissima, oggi meglio di una parte su $10^{15}$. Sarà Einstein a elevarla da coincidenza a \emph{principio}, facendone il punto di partenza della relatività generale: se tutti i corpi cadono allo stesso modo, allora la caduta non è l'effetto di una forza ma una proprietà dello spazio-tempo stesso, e un osservatore in caduta libera non può distinguere la propria condizione dall'assenza di gravità.
\end{osservazione}

\subsection{Campo Gravitazionale e Principio di Sovrapposizione}

Conviene separare la sorgente del campo dal corpo che lo subisce, definendo una grandezza che dipenda solo dalla prima.

\begin{definizione}{Campo Gravitazionale}
Si definisce \textbf{campo gravitazionale} generato da una massa $M$ il vettore
\begin{equation}
    \vec{g} = \frac{\vec{F}_g}{m} = - \, \frac{GM}{r^2} \, \hat{r}
\end{equation}
cioè la forza per unità di massa che agirebbe su un corpo di prova posto in quel punto. Esso si misura in $\mathrm{N/kg}$, dimensionalmente equivalenti a $\mathrm{m/s^2}$.
\end{definizione}

Essendo la legge di Newton lineare nelle masse, i campi prodotti da più sorgenti si sommano vettorialmente, senza che l'una interferisca con l'altra. Per $n$ masse $M_i$ vale dunque il \textbf{principio di sovrapposizione}:
\begin{equation}
    \vec{g}_{\text{tot}} = \sum_{i} \vec{g}_i = - \, G \sum_{i} \frac{M_i}{r_i^2} \, \hat{r}_i
\end{equation}

\begin{attenzione}{Campo $\vec{g}$ e Costante $G$}
Si faccia attenzione a non confondere i due simboli: $G$ è la costante di gravitazione universale, uno scalare fissato una volta per tutte, mentre $\vec{g}$ è un campo vettoriale che varia da punto a punto. Sulla superficie terrestre $|\vec{g}| \approx 9{,}81 \, \mathrm{m/s^2}$, ma il suo valore dipende dalla quota e --- come visto nel capitolo sui sistemi rotanti --- anche dalla latitudine.
\end{attenzione}

\section{Moto in Campo Centrale e Leggi di Keplero}

Applichiamo ora al caso gravitazionale l'apparato del potenziale efficace sviluppato nel capitolo precedente. La forza è centrale e vale $F(r) = -k/r^2$ con $k = GmM$; l'energia totale della particella equivalente si scrive allora
\begin{equation}
    E = \frac{1}{2} \mu \dot{r}^2 + U_{\text{eff}}(r)
      = \frac{1}{2} \mu \dot{r}^2 + \frac{L^2}{2\mu r^2} - \frac{GmM}{r}
\end{equation}

\begin{osservazione}{Quando la Massa Ridotta si Può Ignorare}
Per un sistema Sole--pianeta la massa ridotta vale $\mu = \frac{mM}{m+M}$; essendo $M \gg m$, si ha $\mu \approx m$ con ottima approssimazione (per il sistema Sole--Terra l'errore è di circa tre parti su un milione). È per questo che nella pratica astronomica si confonde senza danno la particella equivalente con il pianeta stesso, e il centro di massa con il Sole. La distinzione torna invece essenziale per le stelle binarie, dove le due masse sono confrontabili.
\end{osservazione}

Il potenziale efficace kepleriano ha una forma caratteristica: il termine centrifugo $L^2/2\mu r^2$ domina a piccole distanze creando una barriera repulsiva, mentre il termine attrattivo $-GmM/r$ domina a grandi distanze. Fra i due si apre una buca di potenziale, il cui minimo corrisponde all'orbita circolare.

\subsection{Orbite Circolari e Minimo di Energia}

Per un'orbita circolare di raggio $r^*$ il raggio non varia, $\dot{r} = 0$, e il valore $r^*$ coincide con il minimo del potenziale efficace. Imponiamo dunque l'annullamento della derivata:
\begin{equation}
    \frac{dU_{\text{eff}}}{dr} = 0
    \quad \Longrightarrow \quad
    - \frac{L^2}{\mu r^3} + \frac{GmM}{r^2} = 0
    \quad \Longrightarrow \quad
    r^* = \frac{L^2}{G m M \mu}
\end{equation}

Sostituendo $r^*$ nell'espressione del potenziale efficace si ottiene l'energia corrispondente. Conviene osservare che, in $r = r^*$, il termine centrifugo vale esattamente metà del modulo di quello attrattivo:
\begin{equation}
    \frac{L^2}{2\mu (r^*)^2} = \frac{1}{2}\,\frac{GmM}{r^*}
\end{equation}
da cui
\begin{equation}
    E^* = U_{\text{eff}}(r^*)
        = \frac{1}{2}\,\frac{GmM}{r^*} - \frac{GmM}{r^*}
        = - \, \frac{GmM}{2 r^*}
\end{equation}

\begin{osservazione}{Energia Negativa e Stati Legati}
$E^*$ è negativa, e rappresenta la minima energia compatibile con il momento angolare $L$ assegnato. Il segno negativo dell'energia totale è la firma di uno \textbf{stato legato}: per liberare il corpo e portarlo all'infinito occorrerebbe fornirgli almeno l'energia $|E^*|$. Si noti inoltre che $E^* = \frac{1}{2}U(r^*)$, cioè l'energia totale è pari a metà dell'energia potenziale --- un caso particolare del \emph{teorema del viriale}.
\end{osservazione}

\subsection{Velocità di Fuga}

La condizione $E \geq 0$ separa i moti limitati da quelli illimitati e definisce una velocità caratteristica. Un corpo lanciato dalla superficie di un pianeta di massa $M$ e raggio $R$ riesce ad allontanarsi indefinitamente se la sua energia totale è non negativa:
\begin{equation}
    E = \frac{1}{2} m v^2 - \frac{GmM}{R} \geq 0
    \qquad \Longrightarrow \qquad
    v \geq \sqrt{\frac{2GM}{R}}
\end{equation}

\begin{definizione}{Velocità di Fuga}
Si definisce \textbf{velocità di fuga} la minima velocità iniziale che consente a un corpo di allontanarsi indefinitamente dalla sorgente del campo:
\begin{equation}
    v_{\text{fuga}} = \sqrt{\frac{2GM}{R}} = \sqrt{2gR}
\end{equation}
Per la Terra essa vale circa $11{,}2 \, \mathrm{km/s}$. Si osservi che non dipende dalla massa del corpo lanciato né dalla direzione del lancio, ma solo dal campo della sorgente.
\end{definizione}

\subsection{Classificazione delle Orbite ed Eccentricità}

Risolvendo l'equazione del moto radiale si ottiene la traiettoria in coordinate polari, che risulta essere in ogni caso una \textbf{conica} con un fuoco nel centro di forza:
\begin{equation}
    r(\phi) = \frac{p}{1 + \epsilon \cos\phi}
    \qquad \text{con} \qquad
    p = \frac{L^2}{G m M \mu}
    \, , \qquad
    \epsilon = \sqrt{1 + \frac{2 E L^2}{\mu \,(G m M)^2}}
\end{equation}
dove $p$ è il \textbf{semilato retto} ed $\epsilon$ l'\textbf{eccentricità}. Il legame fra $\epsilon$ ed $E$ è diretto: il segno dell'energia determina il tipo di conica.

\begin{figure}[ht]
    \centering
    \begin{tikzpicture}[>=stealth, scale=1.1, font=\small]
        \def\p{1.2}
        % Assi
        \draw[->, thin, mypen!45] (-3.7,0) -- (1.75,0);
        \node[mypen, anchor=north, yshift=-2pt] at (1.6,-0.06) {$\phi = 0$};
        \draw[thin, mypen!45] (0,-2.75) -- (0,2.75);

        % Iperbole  (eps = 1.6)
        \draw[very thick, boxatt, domain=-110:110, samples=160, smooth]
            plot ({\p/(1+1.6*cos(\x))*cos(\x)}, {\p/(1+1.6*cos(\x))*sin(\x)});
        % Parabola  (eps = 1)
        \draw[very thick, boxnota, domain=-112:112, samples=160, smooth]
            plot ({\p/(1+1.0*cos(\x))*cos(\x)}, {\p/(1+1.0*cos(\x))*sin(\x)});
        % Ellisse   (eps = 0.6)
        \draw[very thick, boxese, domain=0:360, samples=200, smooth]
            plot ({\p/(1+0.6*cos(\x))*cos(\x)}, {\p/(1+0.6*cos(\x))*sin(\x)});
        % Circonferenza (eps = 0)
        \draw[very thick, boxoss] (0,0) circle (\p);

        % Fuoco comune
        \filldraw[notered] (0,0) circle (2.2pt);
        \node[notered, anchor=east, xshift=-4pt, yshift=-3pt] at (0,0) {$M$};

        % Legenda, a destra e ben staccata dalle curve
        \begin{scope}[shift={(2.35,0)}]
            \foreach \i/\col/\txt in {%
                0/boxoss/{$\epsilon = 0$ \; circonferenza, \; $E = E^*$},
                1/boxese/{$0 < \epsilon < 1$ \; ellisse, \; $E^* < E < 0$},
                2/boxnota/{$\epsilon = 1$ \; parabola, \; $E = 0$},
                3/boxatt/{$\epsilon > 1$ \; iperbole, \; $E > 0$}}
            {
                \draw[very thick, \col] (0,{1.5-0.62*\i}) -- (0.34,{1.5-0.62*\i});
                \node[\col, anchor=west] at (0.44,{1.5-0.62*\i}) {\txt};
            }
        \end{scope}
    \end{tikzpicture}
    \caption{Le quattro famiglie di traiettorie in campo kepleriano, tutte con lo stesso semilato retto $p$ e un fuoco comune nel centro di forza $M$. Al crescere dell'eccentricità --- e quindi dell'energia --- l'orbita passa da chiusa ad aperta.}
    \label{fig:classificazione_orbite}
\end{figure}

\begin{formulario}{Classificazione delle Orbite}
\begin{itemize}[leftmargin=1.4em]
    \item $E > 0 \; \Longrightarrow \; \epsilon > 1$: \textbf{iperbole}, moto illimitato (il corpo proviene dall'infinito e vi ritorna);
    \item $E = 0 \; \Longrightarrow \; \epsilon = 1$: \textbf{parabola}, caso limite, corrispondente alla velocità di fuga;
    \item $E^* < E < 0 \; \Longrightarrow \; 0 < \epsilon < 1$: \textbf{ellisse}, moto limitato e periodico;
    \item $E = E^* \; \Longrightarrow \; \epsilon = 0$: \textbf{circonferenza}, orbita di minima energia.
\end{itemize}
\end{formulario}

\subsection{Prima e Seconda Legge di Keplero}

\begin{teorema}{I Legge di Keplero}
Le orbite dei pianeti sono \textbf{ellissi}, delle quali il Sole occupa uno dei due fuochi.
\end{teorema}

\begin{teorema}{II Legge di Keplero (legge delle aree)}
Il raggio vettore che congiunge il Sole al pianeta spazza \textbf{aree uguali in tempi uguali}.
\end{teorema}

La seconda legge è già stata dimostrata nel capitolo precedente, dove si è visto che essa discende dalla sola \emph{centralità} della forza, attraverso la conservazione del momento angolare, e non dalla sua particolare dipendenza da $r$. Vale dunque per ogni campo centrale, gravitazionale o meno. La prima legge, invece, è una proprietà specifica del potenziale $1/r$.

\begin{figure}[ht]
    \centering
    \begin{tikzpicture}[>=stealth, scale=0.95, font=\small]
        \def\semia{3.4}
        \def\semib{2.4}
        \def\focc{2.4}

        % Ellisse e assi
        \draw[very thick, boxese] (0,0) ellipse [x radius=\semia, y radius=\semib];
        \draw[thin, mypen!40] (-\semia-0.5,0) -- (\semia+0.5,0);
        \draw[thin, mypen!40] (0,-\semib-0.3) -- (0,\semib+0.3);

        % Fuochi e centro
        \filldraw[notered] (\focc,0) circle (2.4pt);
        \filldraw[mypen!55] (-\focc,0) circle (2.0pt);
        \node[mypen!75, anchor=south, yshift=3pt] at (-\focc,0) {$F_2$};
        \filldraw[mypen] (0,0) circle (1.6pt);
        \node[mypen, anchor=north east, xshift=-2pt, yshift=-1pt] at (0,0) {$C$};

        % Perielio e afelio
        \filldraw[mypen] (\semia,0) circle (1.8pt);
        \node[mypen, anchor=west, xshift=4pt] at (\semia,0) {perielio};
        \filldraw[mypen] (-\semia,0) circle (1.8pt);
        \node[mypen, anchor=east, xshift=-4pt] at (-\semia,0) {afelio};

        % Semiasse maggiore, dentro l'ellisse
        \draw[<->, very thick, boxoss] (0,-0.62) -- (\semia,-0.62);
        \node[boxoss, anchor=north] at ({\semia/2},-0.66) {$a$};

        % Semiasse minore
        \draw[<->, very thick, boxatt] (-1.3,0) -- (-1.3,{\semib*sqrt(1-1.3*1.3/(\semia*\semia))});
        \node[boxatt, anchor=east, xshift=-3pt] at (-1.3,{\semib*sqrt(1-1.3*1.3/(\semia*\semia))/2}) {$b$};

        % Distanza focale
        \draw[<->, very thick, boxnota] (0,0.95) -- (\focc,0.95);
        \node[boxnota, anchor=south] at ({\focc/2},0.99) {$c = \epsilon a$};

        % Pianeta e raggio vettore
        \def\tp{122}
        \coordinate (P) at ({\semia*cos(\tp)}, {\semib*sin(\tp)});
        \filldraw[boxoss!45, draw=boxoss, very thick] (P) circle (0.15);
        \draw[very thick, notered] (\focc,0) -- (P)
            node[pos=0.72, anchor=south east, notered] {$r$};
        \node[notered, anchor=north, yshift=-3pt] at (\focc,0) {$F_1$ (Sole)};

        % Quote impilate sotto l'ellisse
        \draw[<->, thick, notered!85] (-\semia,-3.05) -- (\focc,-3.05);
        \node[notered, anchor=north] at ({(\focc-\semia)/2},-3.09) {$r_{\max} = a(1+\epsilon)$};
        \draw[<->, thick, notered!85] (\focc,-3.75) -- (\semia,-3.75);
        \node[notered, anchor=west, xshift=4pt] at (\semia,-3.75) {$r_{\min} = a(1-\epsilon)$};
        \draw[dotted, thin, mypen!55] (-\semia,0) -- (-\semia,-3.0);
        \draw[dotted, thin, mypen!55] (\focc,0) -- (\focc,-3.7);
        \draw[dotted, thin, mypen!55] (\semia,0) -- (\semia,-3.7);
    \end{tikzpicture}
    \caption{Geometria dell'orbita ellittica. Il Sole occupa il fuoco $F_1$; l'altro fuoco $F_2$ è vuoto. La distanza minima $r_{\min}$ si raggiunge al perielio, quella massima $r_{\max}$ all'afelio, e il semiasse maggiore $a$ ne è la semisomma.}
    \label{fig:geometria_ellisse}
\end{figure}

Le grandezze caratteristiche dell'ellisse si ricavano direttamente dall'equazione della conica, valutandola agli estremi:
\begin{equation}
    r_{\min} = \frac{p}{1+\epsilon} = a\,(1-\epsilon)
    \qquad \text{(perielio)}
    \qquad\qquad
    r_{\max} = \frac{p}{1-\epsilon} = a\,(1+\epsilon)
    \qquad \text{(afelio)}
\end{equation}
e il semiasse maggiore è la loro semisomma:
\begin{equation}
    a = \frac{r_{\min} + r_{\max}}{2} = \frac{p}{1-\epsilon^2} = \frac{GmM}{2|E|}
\end{equation}

\begin{osservazione}{Il Semiasse Maggiore Dipende solo dall'Energia}
L'ultima uguaglianza merita attenzione: $a$ è determinato dalla sola energia totale, e non dal momento angolare. Due orbite con la stessa energia ma momenti angolari diversi hanno lo stesso semiasse maggiore --- e dunque, come vedremo, lo stesso periodo --- pur avendo forme molto diverse, l'una quasi circolare e l'altra molto allungata. È il momento angolare, attraverso $p$, a fissare l'eccentricità.
\end{osservazione}

\subsection{Terza Legge di Keplero: Derivazione Analitica}

\begin{teorema}{III Legge di Keplero}
Il quadrato del periodo di rivoluzione è proporzionale al cubo del semiasse maggiore dell'orbita:
\begin{equation}
    T^2 = \frac{4\pi^2 a^3}{G\,(M+m)}
\end{equation}
e, nel limite $M \gg m$, il rapporto $a^3/T^2$ è lo stesso per tutti i corpi orbitanti attorno alla medesima sorgente.
\end{teorema}

\emph{Dimostrazione.} L'area dell'ellisse vale $A = \pi a b$. D'altra parte, per la seconda legge la velocità areolare è costante e pari a $\dot{A} = L/2\mu$; integrando su un periodo completo:
\begin{equation}
    A = \int_{0}^{T} \dot{A} \, dt = \frac{L}{2\mu}\,T = \pi a b
    \qquad \Longrightarrow \qquad
    T^2 = \frac{4\pi^2 \mu^2 a^2 b^2}{L^2}
\end{equation}

Introduciamo ora le due relazioni geometriche che legano $b$ e $L$ ai parametri dell'orbita, $b^2 = a^2(1-\epsilon^2)$ e $L^2 = \mu \, G m M \, p = \mu \, G m M \, a(1-\epsilon^2)$:
\begin{equation}
    T^2 = \frac{4\pi^2 \mu^{\cancel{2}} a^2 \left[a^{\cancel{2}}\cancel{(1-\epsilon^2)}\right]}
               {\cancel{\mu} \, G M m \, \cancel{a}\cancel{(1-\epsilon^2)}}
        = \frac{4\pi^2 \mu \, a^3}{G M m}
\end{equation}
Infine, sostituendo la definizione di massa ridotta $\mu = \dfrac{Mm}{M+m}$, il prodotto $Mm$ si semplifica:
\begin{equation}
    T^2 = \frac{4\pi^2 a^3}{G\,(M+m)}
\end{equation}
$\square$

Per il Sistema Solare si ha $M_{\text{Sole}} \gg m_{\text{pianeta}}$, sicché $M + m \approx M_{\text{Sole}}$ e
\begin{equation}
    \frac{a^3}{T^2} \approx \frac{G\,M_{\text{Sole}}}{4\pi^2} = \text{costante per tutti i pianeti}
\end{equation}
che è esattamente la forma empirica in cui Keplero enunciò la legge, senza poterne dare giustificazione.

\subsection{Verifica Newtoniana: da Keplero a \texorpdfstring{$1/r^2$}{1 su r quadro}}

Il ragionamento può essere percorso anche all'indietro, ed è così che Newton procedette: assumendo la terza legge come dato sperimentale, se ne deduce la forma della forza. Consideriamo per semplicità un'orbita circolare, per cui $a = r$, e scriviamo la terza legge come $T^2 = K r^3$. La velocità angolare vale allora
\begin{equation}
    \omega^2 = \left( \frac{2\pi}{T} \right)^{\!2} = \frac{4\pi^2}{K r^3}
\end{equation}
e l'accelerazione centripeta necessaria a mantenere l'orbita è
\begin{equation}
    a_c = \omega^2 r = \frac{4\pi^2}{K r^3} \, r = \frac{4\pi^2}{K} \, \frac{1}{r^2}
\end{equation}

Poiché $F = m \, a_c$, la forza che agisce sul pianeta deve necessariamente andare come $1/r^2$. La terza legge di Keplero \emph{impone} dunque la dipendenza quadratica inversa: è questo il passaggio con cui il moto dei pianeti cessa di essere un fatto astronomico e diventa un problema di dinamica.

\section{Teorema dei Gusci}

Tutto quanto precede tratta i corpi celesti come punti materiali. Resta da giustificare questa approssimazione: perché è lecito sostituire a un pianeta, che è un corpo esteso di raggio migliaia di chilometri, un punto che ne concentra tutta la massa nel centro? La risposta è il \textbf{teorema dei gusci}, che Newton stesso dovette dimostrare prima di poter pubblicare la legge di gravitazione.

\subsection{Dimostrazione Integrale}

Calcoliamo la forza fra un guscio sferico omogeneo di raggio $R$ e massa $M$ e un punto materiale di massa $m$ posto a distanza $r$ dal centro. Per simmetria la forza risultante è diretta lungo la congiungente centro--punto, e le componenti trasversali si elidono a coppie; basta quindi sommare le sole componenti assiali.

Suddividiamo il guscio in anelli infinitesimi individuati dall'angolo polare $\theta$. Ciascun anello ha massa
\begin{equation}
    dM = \sigma \left(2\pi R^2 \sin\theta \, d\theta\right)
    \qquad \text{con} \qquad \sigma = \frac{M}{4\pi R^2}
\end{equation}
e tutti i suoi punti distano ugualmente $s$ dal punto $P$. La componente assiale del contributo dell'anello è
\begin{equation}
    dF = \frac{G m \, dM}{s^2} \cos\psi
       = \frac{G m \, \sigma \, 2\pi R^2 \sin\theta \, d\theta}{s^2}
         \left[ \frac{r - R\cos\theta}{s} \right]
\end{equation}

\begin{figure}[ht]
    \centering
    \begin{tikzpicture}[>=stealth, scale=1.3, font=\small]
        \def\R{1.6}
        \def\rr{3.9}
        \def\th{52}

        % Guscio
        \draw[very thick, mypen] (0,0) circle (\R);
        \draw[dashed, thin, mypen!40] (0,0) ellipse [x radius=\R, y radius=0.42];

        % Anello infinitesimo
        \coordinate (Q) at ({\R*sin(\th)}, {\R*cos(\th)});
        \draw[very thick, boxese] ({\R*sin(\th)},{\R*cos(\th)}) arc (0:-180:{\R*sin(\th)} and 0.3);
        \draw[very thick, boxese, dashed] ({\R*sin(\th)},{\R*cos(\th)}) arc (0:180:{\R*sin(\th)} and 0.3);
        \fill[boxese] (Q) circle (1.6pt);
        \node[boxese, anchor=south, yshift=3pt] at (Q) {$dM$};

        % Asse polare e angolo theta
        \draw[thick, mypen!55] (0,0) -- (0,\R+0.55);
        \draw[thick, mypen!80] (0,0) -- (Q);
        \draw[thick, mypen] (0,0.82) arc (90:{90-\th}:0.82);
        \node[mypen] at ({1.06*sin(\th/2)}, {1.06*cos(\th/2)}) {$\theta$};
        \node[mypen, anchor=west, xshift=3pt] at ({\R*sin(\th)*0.55},{\R*cos(\th)*0.55}) {$R$};

        % Asse centro-punto
        \draw[thick, mypen!45] (0,0) -- (\rr,0);
        \filldraw[mypen] (0,0) circle (1.7pt);
        \node[mypen, anchor=north east, xshift=-1pt] at (0,-0.06) {$O$};
        \filldraw[notered] (\rr,0) circle (2.2pt);
        \node[mypen, anchor=west, xshift=5pt] at (\rr,0) {$P$ \, ($m$)};
        \draw[<->, thin, mypen!70] (0,-0.95) -- (\rr,-0.95);
        \node[mypen, anchor=north] at ({\rr/2},-0.98) {$r$};

        % Distanza s
        \draw[very thick, boxoss] (Q) -- (\rr,0);
        \node[boxoss, anchor=south west] at ({(\R*sin(\th)+\rr)/2 - 0.1},{\R*cos(\th)/2}) {$s$};

        % Angolo psi in P
        \draw[thick, boxatt] ({\rr-0.72},0)
            arc (180:{180-atan(\R*cos(\th)/(\rr-\R*sin(\th)))}:0.72);
        \node[boxatt, anchor=south, yshift=1pt] at ({\rr-1.02},0.28) {$\psi$};

        % Forza elementare
        \draw[->, very thick, notered] (\rr,0) -- ({\rr-0.95},0);
        \node[notered, anchor=north, yshift=-3pt] at ({\rr-0.5},-0.04) {$dF$};
    \end{tikzpicture}
    \caption{Geometria per il teorema dei gusci. Il guscio è suddiviso in anelli infinitesimi individuati dall'angolo $\theta$; tutti i punti dell'anello distano $s$ dal punto $P$, e le componenti trasversali delle forze si elidono per simmetria, lasciando la sola componente assiale $dF \propto \cos\psi$.}
    \label{fig:shell_theorem}
\end{figure}

Conviene ora cambiare variabile d'integrazione, passando da $\theta$ a $s$. Per il teorema di Carnot applicato al triangolo $OQP$:
\begin{equation}
    s^2 = r^2 + R^2 - 2rR\cos\theta
\end{equation}
e differenziando entrambi i membri a $r$ e $R$ fissati:
\begin{equation}
    2s \, ds = 2rR\sin\theta \, d\theta
    \qquad \Longrightarrow \qquad
    \sin\theta \, d\theta = \frac{s \, ds}{rR}
\end{equation}
Dalla stessa relazione si ricava inoltre $\cos\theta = \dfrac{r^2 + R^2 - s^2}{2rR}$, e quindi
\begin{equation}
    r - R\cos\theta = r - \frac{r^2 + R^2 - s^2}{2r} = \frac{r^2 - R^2 + s^2}{2r}
\end{equation}

Sostituendo tutto nell'espressione di $dF$, e usando $\sigma \, 2\pi R^2 = M/2$, si ottiene la forma notevolmente semplice
\begin{equation}
    dF = \frac{G m M}{4 r^2 R} \left[ 1 + \frac{r^2 - R^2}{s^2} \right] ds
\end{equation}
la cui primitiva è immediata:
\begin{equation}
    \int \left[ 1 + \frac{r^2-R^2}{s^2} \right] ds = s - \frac{r^2-R^2}{s}
\end{equation}

\paragraph{Punto esterno al guscio ($r > R$).}
La distanza $s$ varia fra un minimo $r-R$ e un massimo $r+R$. Ricordando che $r^2 - R^2 = (r-R)(r+R)$:
\begin{align}
    F &= \frac{GmM}{4r^2 R} \left[ s - \frac{r^2 - R^2}{s} \right]_{r-R}^{r+R} \notag \\
      &= \frac{GmM}{4r^2 R} \Big[ (r+R) - \underbrace{\tfrac{(r-R)(r+R)}{r+R}}_{=\,r-R}
         - (r-R) + \underbrace{\tfrac{(r-R)(r+R)}{r-R}}_{=\,r+R} \Big] \notag \\
      &= \frac{GmM}{4r^2 R} \Big[ 2(r+R) - 2(r-R) \Big]
       = \frac{GmM}{4r^2 \cancel{R}} \cdot 4\cancel{R}
       = \frac{GmM}{r^2}
\end{align}

\paragraph{Punto interno al guscio ($r < R$).}
Ora la distanza $s$ varia fra $R-r$ e $R+r$, e questa volta $r^2 - R^2 = -(R-r)(R+r)$ è negativo:
\begin{align}
    F &= \frac{GmM}{4r^2 R} \left[ s - \frac{r^2 - R^2}{s} \right]_{R-r}^{R+r} \notag \\
      &= \frac{GmM}{4r^2 R} \Big[ (R+r) + \underbrace{\tfrac{(R-r)(R+r)}{R+r}}_{=\,R-r}
         - (R-r) - \underbrace{\tfrac{(R-r)(R+r)}{R-r}}_{=\,R+r} \Big] \notag \\
      &= \frac{GmM}{4r^2 R} \Big[ (R+r) + (R-r) - (R-r) - (R+r) \Big] = 0
\end{align}

\begin{teorema}{Teorema dei Gusci}
Un guscio sferico omogeneo di massa $M$ esercita su un punto materiale:
\begin{itemize}[leftmargin=1.4em]
    \item \textbf{all'esterno} ($r > R$), la stessa forza che eserciterebbe un punto materiale di massa $M$ collocato nel centro,
    \begin{equation*}
        F = \frac{GmM}{r^2}
    \end{equation*}
    \item \textbf{all'interno} ($r < R$), forza \textbf{nulla} in ogni punto, qualunque sia la posizione.
    \begin{equation*}
        F = 0
    \end{equation*}
\end{itemize}
\end{teorema}

\begin{osservazione}{Perché l'Interno è Esattamente Nullo}
Il risultato interno non è un'approssimazione né un effetto di simmetria banale: il punto non deve stare al centro. Preso un punto qualsiasi dentro il guscio, i coni opposti che lo hanno per vertice intercettano sulla superficie due calotte di area proporzionale al quadrato della rispettiva distanza; l'attrazione maggiore della calotta vicina è compensata esattamente dalla massa maggiore di quella lontana, proprio perché la forza va come $1/r^2$. È dunque un risultato che dipende in modo essenziale dall'esponente $2$: per qualunque altra potenza la cancellazione non avverrebbe.
\end{osservazione}

\subsection{Il Caso della Sfera Piena}

Una sfera piena si tratta come una successione di gusci concentrici, applicando a ciascuno il teorema.

All'\textbf{esterno} ($r > R_T$) ogni guscio agisce come se la sua massa fosse concentrata nel centro; sommando su tutti i gusci si ottiene che l'intera sfera equivale a un punto materiale di massa $M$ posto nel centro. È questo il risultato che legittima a posteriori l'aver trattato i pianeti come punti.

All'\textbf{interno} ($r < R_T$) i gusci più esterni del punto non esercitano alcuna forza, e contribuisce soltanto la massa $M(r)$ contenuta nella sfera di raggio $r$. Per una densità $\rho$ uniforme:
\begin{equation}
    F(r) = \frac{G m \, M(r)}{r^2}
         = \frac{G m \left(\rho \, \tfrac{4}{3}\pi r^3\right)}{r^2}
         = \frac{4}{3}\pi G m \rho \, r \; \propto \; r
\end{equation}

\begin{figure}[ht]
    \centering
    \begin{tikzpicture}[>=stealth, x=1.5cm, y=1.1cm, font=\small]
        % Assi
        \draw[->, thick, mypen] (0,0) -- (5.2,0) node[right] {$r$};
        \draw[->, thick, mypen] (0,0) -- (0,3.0) node[above] {$|\vec{g}(r)|$};

        % Interno: lineare
        \draw[very thick, boxoss] (0,0) -- (1.8,2.3);
        % Esterno: 1/r^2, raccordato in r = R_T
        \draw[very thick, notered, domain=1.8:5.0, samples=120, smooth]
            plot (\x, {2.3*(1.8/\x)^2});

        % Raggio terrestre
        \draw[dashed, thin, mypen!60] (1.8,0) -- (1.8,2.3);
        \draw[dashed, thin, mypen!60] (0,2.3) -- (1.8,2.3);
        \node[mypen, anchor=north] at (1.8,-0.06) {$R_T$};
        \node[mypen, anchor=east] at (-0.06,2.3) {$g$};

        % Etichette
        \node[boxoss, anchor=east] at (1.55,1.05) {$\propto r$};
        \node[notered, anchor=south west] at (3.0,0.95) {$\propto 1/r^2$};
        \node[mypen!75, anchor=north, align=center] at (0.9,-0.25) {interno};
        \node[mypen!75, anchor=north, align=center] at (3.4,-0.25) {esterno};
    \end{tikzpicture}
    \caption{Modulo del campo gravitazionale generato da una sfera piena omogenea in funzione della distanza dal centro. Il campo cresce linearmente all'interno, raggiunge il massimo $g$ sulla superficie e decresce come $1/r^2$ all'esterno.}
    \label{fig:campo_sfera_piena}
\end{figure}

Il campo cresce dunque \textbf{linearmente} dal centro fino alla superficie, dove raggiunge il valore massimo $g = GM/R_T^2$, per poi decrescere come $1/r^2$. In particolare, al centro esatto della Terra il campo gravitazionale è \textbf{nullo}: tutta la massa circostante attrae in modo perfettamente simmetrico.

\chapter{Corpo rigido e oscillazioni.pdf}


% --- Batch Pagine 1-10 ---
% Errore batch 0-10

\chapter{Termodinamica}
\section{Variabili Termodinamiche}

I sistemi costituiti da un numero di elementi dell'ordine di $10^{23}$ sono irrisolvibili dal punto di vista pratico con sistemi di equazioni (es. $10^{23}$ equazioni). In questi sistemi (descritti microscopicamente) una minima variazione può compromettere le condizioni iniziali determinando un comportamento caotico. 

\medskip
Quindi dobbiamo sfruttare una \textbf{descrizione macroscopica} con pochi ($< 10$) \textbf{parametri globali}, il che la rende più semplice e predittiva, seppur meno precisa.

\begin{definizione}{Meccanica Statistica e Termodinamica}
La \textbf{meccanica statistica} stabilisce relazioni fra grandezze microscopiche e macroscopiche di tipo probabilistico.

\medskip
La \textbf{termodinamica} stabilisce relazioni fra grandezze macroscopiche su basi sperimentali tramite i principi termodinamici, aventi validità generale.
\end{definizione}

\subsection{Variabili Termodinamiche}
\begin{definizione}{Variabili termodinamiche}
Le variabili termodinamiche sono grandezze macroscopiche $A, B, C$ che descrivono globalmente un sistema tramite un'\textbf{equazione di stato} $F(A,B,C) = 0$.
\end{definizione}

\begin{itemize}
    \item \textbf{Massa totale} (sistema isolato) -- Meccaniche
    \item \textbf{Volume} (sistema deformabile) -- Meccaniche
    \item \textbf{Pressione} -- Nuova grandezza
    \item \textbf{Temperatura} -- Nuova grandezza
    \item \textbf{Quantità di calore} -- Nuova grandezza
\end{itemize}

Queste variabili sono \textbf{indipendenti dal tempo} e \textbf{definite in ogni punto}, cioè indipendenti dall'istante di partenza ma dipendenti dall'intervallo di tempo considerato.

\subsection{Pressione}
\begin{definizione}{Pressione}
La pressione è il rapporto fra il modulo della forza ortogonale a una superficie e l'area della superficie stessa:
\begin{equation}
p = \frac{||\vec{F}_{\perp}||}{S} = \frac{|\vec{F} \cdot \hat{n}|}{S}
\end{equation}
\end{definizione}

\medskip
\begin{itemize}
    \item Grandezza scalare
    \item Unità di misura nel SI: $1 \text{ Pa} = \frac{1 \text{ N}}{1 \text{ m}^2}$ (Pascal)
    \item $1 \text{ atm} = 101325 \text{ Pa} \simeq 101.3 \text{ kPa}$ (atmosfera)
    \item $1 \text{ bar} = 10^5 \text{ Pa}$
    \item $1 \text{ torr} = 133.3 \text{ Pa}$ (Torricelli)
\end{itemize}

\vspace{0.3cm}
\begin{figure}[ht]
\centering
\begin{tikzpicture}[scale=1]
    % Parallelogram representing the surface
    \fill[cyan!15, opacity=0.8] (0,0) -- (4,0) -- (5,1.5) -- (1,1.5) -- cycle;
    \draw[thick, mypen] (0,0) -- (4,0) -- (5,1.5) -- (1,1.5) -- cycle;
    \node[mypen, font=\bfseries] at (4.3, 0.4) {$\Delta S$};
    
    % Punto di applicazione
    \coordinate (P) at (2.5, 0.75);
    \filldraw[mypen] (P) circle (1.5pt);
    
    % Normal vector
    \draw[->, >=stealth, thick, mypen] (P) -- (2.5, 2.8) node[right] {$\hat{n}$};
    
    % Force vector
    \draw[->, >=stealth, thick, mygreen] (P) -- (3.8, 2.3) node[above right] {$\vec{F}$};
    
    % Proiezione normale di F
    \draw[dashed, mypen] (3.8, 2.3) -- (2.5, 2.3);
    \draw[->, >=stealth, thick, blue] (P) -- (2.5, 2.3) node[left, midway] {$F_\perp$};
\end{tikzpicture}
\caption{Rappresentazione vettoriale della forza $\vec{F}$ esercitata dal fluido su un elemento di superficie $\Delta S$; solo la componente normale $F_\perp$ definisce la pressione idrostatica.}
\end{figure}
\vspace{0.3cm}

\subsection{Stati di Aggregazione}
\begin{itemize}
    \item \textbf{Solido}: possiede forma e volume propri.
    \item \textbf{Fluido liquido}: possiede volume proprio, ma non forma propria.
    \item \textbf{Fluido aeriforme (gas)}: non possiede né volume né forma propri.
\end{itemize}

\begin{definizione}{Fluido Perfetto o Ideale}
Un fluido è detto ideale se presenta assenza di viscosità (attrito interno). In altre parole, considerando il fluido come un insieme di strati sovrapposti, quando due di questi scorrono l'uno sull'altro non si genera alcun attrito.
\end{definizione}
\begin{itemize}
    \item L'\textbf{acqua} è un fluido molto poco viscoso.
    \item L'\textbf{olio} è un fluido molto viscoso.
\end{itemize}

\section{Pressione in un Fluido Ideale}
In un fluido perfetto all'equilibrio, preso un punto P e una superficie $\Delta S$ intorno a P, nel limite per $\Delta S \to 0$ la pressione $F/\Delta S$ è la stessa in ogni direzione (Principio di Pascal).

\begin{figure}[ht]
    \centering
    \begin{tikzpicture}[scale=1.1]
        % Faccia frontale (triangolo rettangolo)
        \fill[cyan!15, opacity=0.8] (0,0) -- (3,0) -- (0,2) -- cycle;
        \draw[thick, mypen] (0,0) -- (3,0) -- (0,2) -- cycle;
        
        % Spigoli posteriori (prospettiva)
        \draw[dashed, mypen] (0,0) -- (1,0.6);
        \draw[dashed, mypen] (3,0) -- (4,0.6);
        \draw[dashed, mypen] (0,2) -- (1,2.6);
        \draw[dashed, mypen] (1,0.6) -- (4,0.6) -- (1,2.6) -- cycle;
        
        % Spigoli visibili superiori
        \draw[thick, mypen] (3,0) -- (4,0.6) -- (1,2.6) -- (0,2);
        
        % Etichetta punto P interno
        \filldraw[mypen] (1.5,0.8) circle (1.5pt) node[below right] {P};
        
        % Vettori pressione sulle facce
        \draw[->, >=stealth, thick, mygreen] (-0.8, 1) -- (0, 1) node[midway, above] {$\vec{F}_x$};
        \draw[->, >=stealth, thick, mygreen] (1.5, -0.8) -- (1.5, 0) node[midway, left] {$\vec{F}_y$};
        \draw[->, >=stealth, thick, notered] (2.8, 1.8) -- (1.8, 1.1) node[midway, above right] {$\vec{F}_n$};
    \end{tikzpicture}
    \caption{Prisma infinitesimo in equilibrio in un fluido ideale: la dimostrazione del Principio di Pascal mostra l'isotropia della pressione nel limite per il volume che tende a zero.}
\end{figure}

\medskip
Se la regione A esercita $\vec{F}$ su $\Delta S$, allora per il terzo principio della dinamica la regione B esercita $-\vec{F}$ su $\Delta S$.
\begin{itemize}
    \item Se il fluido è \textbf{viscoso}, allora $\vec{F} \not\parallel \hat{n}$.
    \item Se il fluido è \textbf{ideale}, allora $\vec{F} \parallel \hat{n}$.
\end{itemize}

All'equilibrio la somma delle forze applicate sul prisma su ogni sua superficie è nulla: $\sum \vec{F}_s = 0$.
Sia $\vec{P} = -mg \hat{z}$ il peso del fluido:
\begin{equation}
\vec{F}_S + \vec{F}_1 + \vec{F}_2 + \vec{F}_3 + \vec{F}_4 + \vec{P} = 0
\end{equation}

Scomponendo lungo gli assi:
\begin{equation*}
\begin{cases}
    \hat{x}: & F_4 - F_3 = 0 \\
    \hat{y}: & p_y b^2 \sin\theta - p b^2 \sin\theta = 0 \\
    \hat{z}: & p_z b^2 \cos\theta - p b^2 \cos\theta - mg = 0
\end{cases}
\implies
\begin{cases}
    p_y = p \\
    p_z = p + \frac{mg}{b^2 \cos\theta}
\end{cases}
\end{equation*}

Poiché $mg = \rho V_{\text{prisma}} g = \rho \left( \frac{b^2 \sin\theta \cos\theta \cdot b}{2} \right) g = \rho \frac{b^3 \sin\theta \cos\theta}{2} g$, sostituendo in $p_z$:
\begin{equation}
p_z = p + \frac{\rho b^3 \sin\theta \cos\theta g}{2 b^2 \cos\theta} = p + \frac{1}{2} \rho b \sin\theta g
\end{equation}

Nel limite per $\Delta S \to 0$, che equivale a far tendere $b \to 0$:
\begin{equation}
\lim_{b \to 0} \left( p + \frac{1}{2} \rho b \sin\theta g \right) = p
\end{equation}

Di conseguenza, ruotando il prisma di $90^\circ$ si dimostra analogamente per l'altro asse che:
\begin{equation}
\boxed{p = p_y = p_z = p_x}
\end{equation}

\vspace{0.3cm}
\section{Legge di Stevino}
\begin{definizione}{Legge di Stevino}
In un fluido all'equilibrio, la pressione esercitata su uno strato del fluido è determinata soltanto dal peso degli strati sovrastanti (che si trovano a una profondità $z$ minore).
\end{definizione}

\medskip
All'equilibrio $\sum \vec{F}_z = 0$. Sia $\vec{R}$ la forza di reazione degli strati a $z' > z$ e $\vec{P}$ la forza peso degli strati a $z' \le z$. Si ottiene:
\begin{equation}
\vec{P} + \vec{R} = 0 \implies \vec{R} = -m\vec{g} \implies R = mg = \rho V g = \rho S z g = p S
\end{equation}

Da cui la pressione:
\begin{equation}
p = \frac{R}{S} = \frac{\rho g S z}{S} = \rho g z
\end{equation}

Se si considera anche la pressione atmosferica o esterna esercitata al livello $z=0$, allora si ottiene la formula generale:
\begin{equation}
\boxed{p(z) = p(0) + \rho g z}
\end{equation}

\vspace{0.3cm}
\section{Legge di Archimede}
\begin{definizione}{Spinta di Archimede}
Un corpo immerso in un fluido riceve una spinta dal basso verso l'alto pari al peso del volume di fluido spostato dal corpo.
\end{definizione}

\medskip
La forza risultante lungo l'asse $z$ si ricava come differenza tra la pressione sulla faccia inferiore e quella sulla faccia superiore:
\begin{align*}
F_z &= F_{\text{up}} - F_{\text{down}} = p(z_2)S - p(z_1)S - mg \\
&= (p(0) + \rho_{\text{fluido}} g z_2)S - (p(0) + \rho_{\text{fluido}} g z_1)S - mg \\
&= \rho_{\text{fluido}} g S (z_2 - z_1) - mg \\
&= \rho_{\text{fluido}} g S (z_2 - z_1) - \rho_{\text{corpo}} g S (z_2 - z_1) \\
&= (\rho_{\text{fluido}} - \rho_{\text{corpo}}) g S (z_2 - z_1) = (\rho_{\text{fluido}} - \rho_{\text{corpo}}) g V_{\text{corpo}}
\end{align*}

Dunque la forza risultante (tenendo conto anche del peso del corpo) è:
\begin{equation}
\boxed{F_z = (\rho_{\text{fluido}} - \rho_{\text{corpo}}) g V_{\text{corpo}}}
\end{equation}

\begin{itemize}
    \item Se $\rho_{\text{fluido}} < \rho_{\text{corpo}} \implies \vec{F}_z \parallel -\hat{z}$, cioè il corpo va a fondo.
    \item Se $\rho_{\text{fluido}} > \rho_{\text{corpo}} \implies \vec{F}_z \parallel +\hat{z}$, cioè il corpo galleggia.
\end{itemize}

Nel caso di galleggiamento, il corpo risale finché non raggiunge una situazione di equilibrio in cui la forza totale è nulla ($F_z = 0$). L'equilibrio si ha quando:
\begin{equation}
\rho_{\text{fluido}} g V_{\text{immerso}} = mg
\end{equation}
dove $V_{\text{immerso}}$ è il volume della parte del corpo effettivamente immersa nel fluido. Pertanto, uguagliando le masse e considerando l'altezza d'immersione:
\begin{equation}
\rho_{\text{fluido}} S (z_2 - \bar{z}) = \rho_{\text{corpo}} S (z_2 - z_1) \implies \boxed{(z_2 - \bar{z}) = (z_2 - z_1) \frac{\rho_{\text{corpo}}}{\rho_{\text{fluido}}}}
\end{equation}
che rappresenta la profondità di immersione.

\section{Misura della Pressione}
La pressione si misura con i \textbf{manometri} (ad es. per gli pneumatici) e i \textbf{barometri} (ad es. per l'atmosfera).
La pressione atmosferica fu misurata per la prima volta da Evangelista Torricelli nel 1643, il quale inventò il barometro.

\begin{figure}[ht]
    \centering
    \begin{tikzpicture}
        % Barometro di Torricelli
        \draw[thick, mypen] (0,0) -- (3,0) -- (3,1) -- (0,1) -- cycle; % vaschetta
        \draw[thick, mypen, fill=blue!30] (0.1,0.1) rectangle (2.9,0.8); % liquido in vaschetta
        \draw[thick, mypen] (1.3,0.5) -- (1.3,4) -- (1.7,4) -- (1.7,0.5); % tubo
        \draw[thick, mypen, fill=blue!30] (1.3,0.5) rectangle (1.7,3); % liquido in tubo
        \node at (2.5,0.5) {A \quad B};
        \node at (1.5,3.5) {vuoto};
        \draw[->, >=stealth, mypen] (3.5, 2) -- (1.8, 3) node[right, xshift=1cm] {aria};
    \end{tikzpicture}
    \caption{Barometro di Torricelli}
\end{figure}

\medskip
L'esperimento utilizza il mercurio e si basa sui seguenti principi:
\begin{itemize}
    \item La densità del mercurio è $\rho_{Hg} = 13.7 \cdot 10^3 \text{ kg/m}^3$ ($\simeq 14 \rho_{H_2O}$), essendo un metallo liquido molto pesante.
    \item L'altezza della colonna di mercurio all'equilibrio è $h_3 = 76 \text{ cm}$.
    \item $p_A = p_B$: se le pressioni fossero diverse ci sarebbero forze non nulle che farebbero entrare o uscire mercurio dal capillare, ma essendo il sistema all'equilibrio le pressioni si equivalgono.
    \item $p_B = \rho_{Hg} g h_1 + \rho_{Hg} g h_2 + p_{\text{atm}}$
    \item $p_A = \rho_{Hg} g h_1 + \rho_{Hg} g h_2 + \rho_{Hg} g h_3$
\end{itemize}

Uguagliando le due pressioni si ottiene:
\begin{align*}
p_{\text{atm}} &= \rho_{Hg} g h_3 = 13.7 \cdot 10^3 \cdot 9.81 \cdot 0.76 \text{ Pa} \simeq 101 \text{ kPa} \\
&\implies \boxed{p_{\text{atm}} = \rho_{Hg} g h_3 \simeq 101 \text{ kPa} \simeq 760 \text{ Torr}}
\end{align*}

\begin{osservazione}{Il peso dell'aria}
L'aria, per quanto impercettibile, ha un peso e si può "maneggiare", cioè si può mettere e togliere: lo spazio vuoto che si viene a creare nel capillare quando scende il mercurio ne è la dimostrazione.
\end{osservazione}


\section{Temperatura}
\begin{definizione}{Temperatura}
La temperatura è la misura della variazione di volume (o altra proprietà termometrica) di un corpo posto a contatto con una sorgente di calore.
\end{definizione}

\medskip
La misura si effettua con i \textbf{termometri}. Per tararli, si scelgono dei punti fissi $P_1$ e $P_2$ a cui si attribuiscono determinati valori di temperatura; questi punti devono essere \textbf{facilmente riproducibili} e \textbf{controllabili}.

A pressione $p_{\text{atm}}$ costante, i punti di riferimento scelti sono:
\begin{itemize}
    \item $P_1$: punto di fusione dell'acqua (equilibrio tra acqua e ghiaccio fondente).
    \item $P_2$: punto di ebollizione dell'acqua (equilibrio tra acqua e vapore).
\end{itemize}

\subsection{Scala Celsius}
Vengono assegnati i valori $P_1 \to 0^\circ\text{C}$ e $P_2 \to 100^\circ\text{C}$.

\subsection{Scala Fahrenheit}
Vengono assegnati i valori $P_1 \to 32^\circ\text{F}$ e $P_2 \to 212^\circ\text{F}$.

\begin{figure}[ht]
    \centering
    \begin{tikzpicture}
        \draw[->, >=stealth, mypen] (0,0) -- (5,0) node[right] {$^\circ\text{C}$};
        \draw[->, >=stealth, mypen] (0,0) -- (0,3) node[above] {$^\circ\text{F}$};
        \draw[mypen] (0,0.5) node[left] {32} -- (0.1,0.5);
        \draw[mypen] (0,2.5) node[left] {212} -- (0.1,2.5);
        \draw[mypen] (3,0) node[below] {100} -- (3,0.1);
        \draw[mypen] (0,0) node[below right] {0};
        \draw[thick, blue] (-1,-0.1) -- (4,3.1);
    \end{tikzpicture}
    \caption{Relazione tra le scale Celsius e Fahrenheit}
\end{figure}

\medskip
La formula per la conversione da $^\circ\text{C}$ a $^\circ\text{F}$ è la seguente:
\begin{equation}
    \boxed{T(^\circ\text{F}) = 32 + \frac{9}{5} T(^\circ\text{C})}
\end{equation}

\section{Sostanze Termometriche e Termometri}
Come sostanza di cui misurare la variazione di volume si usano generalmente gas o liquidi, per la loro elevata capacità di dilatarsi termicamente.

\subsection{Termometro a Mercurio}
Si segna l'altezza della colonna di mercurio quando è inserita nell'acqua bollente ($100^\circ\text{C}$) e nel ghiaccio fondente ($0^\circ\text{C}$). Successivamente si divide in 100 intervalli identici per creare la \textbf{scala centigrada}.

Affinché il termometro sia reattivo, il volume del capillare deve essere molto minore di quello del bulbo ($V_{\text{capillare}} \ll V_{\text{bulbo}}$) e il bulbo non può essere troppo grande, altrimenti il liquido impiega troppo tempo per dilatarsi uniformemente.

Nella scelta del fluido termometrico si deve tenere conto anche delle sue temperature di congelamento ed ebollizione:
\begin{itemize}
    \item \textbf{Mercurio (Hg)}: congelamento a $-39^\circ\text{C}$ (si solidificherebbe con un clima artico) ed ebollizione a $357^\circ\text{C}$.
    \item \textbf{Alcol}: congelamento a $-70^\circ\text{C}$ ed ebollizione a $78.5^\circ\text{C}$ (si vaporizzerebbe prima dell'acqua).
\end{itemize}

\vspace{0.3cm}
\section{Principio Zero della Termodinamica}
\begin{definizione}{Principio Zero}
Due corpi che sono separatamente in equilibrio termico con un terzo corpo, sono in equilibrio termico anche tra di loro.
\end{definizione}

\begin{figure}[ht]
    \centering
    \begin{tikzpicture}[scale=0.95]
        % Parete adiabatica esterna e di separazione tra A e B
        \fill[pattern=north east lines, mypen!30] (0,0) rectangle (6,3);
        \fill[white] (0.2,0.2) rectangle (5.8,2.8);
        \fill[pattern=north east lines, mypen!30] (2.9,1.4) rectangle (3.1,2.8);
        
        % Sistema A e B
        \draw[thick, mypen, fill=cyan!10] (0.5,1.5) rectangle (2.8,2.6) node[midway, font=\bfseries] {Sistema A ($T_A$)};
        \draw[thick, mypen, fill=orange!10] (3.2,1.5) rectangle (5.5,2.6) node[midway, font=\bfseries] {Sistema B ($T_B$)};
        
        % Sistema C (Termometro/riferimento in contatto diatermico con A e B)
        \draw[thick, mypen, fill=green!10] (0.5,0.4) rectangle (5.5,1.3) node[midway, font=\bfseries] {Sistema C ($T_C$) - Parete Diatermica};
        
        % Frecce di equilibrio termico
        \draw[<->, >=stealth, thick, notered] (1.65, 1.4) -- (1.65, 1.6) node[midway, left] {Eq.};
        \draw[<->, >=stealth, thick, notered] (4.35, 1.4) -- (4.35, 1.6) node[midway, right] {Eq.};
    \end{tikzpicture}
    \caption{Illustrazione del Principio Zero: i sistemi A e B, separati da una parete adiabatica, sono singolarmente in equilibrio termico con il terzo sistema C (grazie alla parete diatermica). Ne consegue che A e B sono in equilibrio termico tra loro ($T_A = T_B = T_C$).}
\end{figure}

All'inizio si ha $T_1 \neq T_2$; se li poniamo a contatto termico fra loro (isolandoli dall'esterno), dopo un certo tempo $\Delta t$ essi raggiungono la stessa temperatura finale $T = T_1' = T_2'$, trovandosi in uno stato di \textbf{equilibrio termico}.
Grazie alla proprietà transitiva, se $T_A = T_B$ e $T_B = T_C$, allora $T_A = T_C$.

La temperatura è una grandezza fondamentale, la cui unità di misura nel SI è il Kelvin (K).

\vspace{0.3cm}
\section{Termometro a Gas}
Ci sono termometri a gas di due tipi: a pressione o a volume costante. Nonostante la maggiore precisione, sono anche più ingombranti dei termometri a liquido.

\subsection{Termometro a Pressione Costante}
Un rubinetto collega due capillari fra loro e con un serbatoio (isolato). Se lasciamo il rubinetto aperto e scaldiamo il gas, questo si dilata e spinge il mercurio dall'altra parte.
Quindi facciamo defluire il mercurio nel serbatoio in modo da azzerare il dislivello ($h_1 - h_2 = 0$).

All'equilibrio, la pressione nel punto A può essere calcolata da entrambe le parti:
\begin{align*}
p_A(\text{Dx}) &= p_{\text{gas}} + \rho_{Hg} g h_2 \\
p_A(\text{Sx}) &= p_{\text{atm}} + \rho_{Hg} g h_1
\end{align*}

Dunque si ha la pressione $p_{\text{gas}} = p_{\text{atm}}$.

Misurando il volume del gas si ottiene:
\begin{equation}
    \boxed{V(T_C) = V_0 + \alpha V_0 T_C}
\end{equation}
dove $V_0$ è il volume a $T_C=0^\circ\text{C}$ e $p=1\text{ atm}$.
Il coefficiente $\alpha$ calcolato sperimentalmente vale:
\begin{equation}
    \alpha = \frac{V_{100} - V_0}{100 V_0} (^\circ\text{C})^{-1} = \boxed{ \left( 273.15 \pm 0.01 \, ^\circ\text{C} \right)^{-1} \simeq \frac{1}{273.15} \, (^\circ\text{C})^{-1} }
\end{equation}
Questo valore è comune a gran parte dei gas a pressione atmosferica e temperatura molto più elevata di quella di liquefazione.

\subsection{Termometro a Volume Costante}
Alzando e abbassando il serbatoio di mercurio a seconda dei casi variamo la pressione in funzione della temperatura, mantenendo il volume costante:
\begin{equation}
    \boxed{p(T_C) = p_0 + \alpha p_0 T_C}
\end{equation}
dove $p_0$ è la pressione a $T_C=0^\circ\text{C}$. Il coefficiente $\alpha$ risulta identico, comprensivo dell'incertezza sperimentale: $\alpha = \frac{p_{100} - p_0}{100 p_0} (^\circ\text{C})^{-1} = \boxed{ \left( 273.15 \pm 0.01 \, ^\circ\text{C} \right)^{-1} }$.

\vspace{0.3cm}
\section{Scala Kelvin e Gas Perfetti}
\begin{definizione}{Gas Perfetto}
Le due formule trovate con i termometri a gas a pressione e volume costante sono chiamate \textbf{Leggi di Gay-Lussac}. Un gas che segue queste due leggi per qualsiasi $p, T$ è detto \textbf{gas perfetto} o ideale.
\end{definizione}

Il gas perfetto in natura non esiste, ma molti dei gas nobili hanno un comportamento simile in ampi intervalli di $p$ e $T$.

\medskip
Se calcoliamo i valori per cui $V(T_C) = 0$ e $p(T_C) = 0$, otteniamo $T_C = -1/\alpha \simeq -273.15^\circ\text{C}$.
Esiste quindi una temperatura al di sotto della quale la materia si comporta in modo differente. Si definisce pertanto una \textbf{scala assoluta} detta \textbf{scala Kelvin} che ha come valore $0\text{ K} = -1/\alpha$, detto \textbf{zero assoluto}.

\begin{equation}
\boxed{T(\text{K}) = T(^\circ\text{C}) + 273.15}
\end{equation}
La conversione da $^\circ\text{C}$ a K mantiene lo stesso intervallo: $1^\circ\text{C} = 1\text{ K}$.

Scrivendo le leggi di Gay-Lussac rispetto alla scala Kelvin si ottiene:
\begin{align}
    V(T) &= \alpha V_0 T \quad \text{a } p=\text{costante} \\
    p(T) &= \alpha p_0 T \quad \text{a } V=\text{costante}
\end{align}

Per un gas perfetto a $T=0\text{ K} \implies p=0, V=0$.
Con la meccanica statistica si può dimostrare che l'energia cinetica media $\langle K \rangle \propto T$. Ciò significa che per $T \to 0$ il sistema non può scambiare energia cedendola a un altro sistema, al massimo può solo riceverla, poiché ha raggiunto l'energia minima possibile.
Con "scala assoluta" inoltre si intende che essa è indipendente dalla sostanza termometrica (scala $\Theta$, legata al Secondo Principio della Termodinamica).

\vspace{0.3cm}
\subsection{Derivazione dell'Equazione di Stato, Condizioni STP e Legge di Dalton}
Combinando le due leggi di Gay-Lussac con la legge di Boyle-Mariotte ($p V = \text{costante}$ a $T = \text{costante}$), si ricava l'equazione di stato generale che collega le variabili termodinamiche di una massa gassosa. Consideriamo di passare da uno stato iniziale $A(p_0, V_0, T_0)$ a uno stato finale $B(p, V, T)$ attraverso uno stato intermedio ausiliario $A^*(p_0, V^*, T)$ mediante una trasformazione isobara ($A \to A^*$) e una successiva isoterma ($A^* \to B$):
\begin{align}
    \text{Isobara } A \to A^* &: \quad \frac{V^*}{V_0} = \frac{T}{T_0} \implies V^* = V_0 \frac{T}{T_0} \\
    \text{Isoterma } A^* \to B &: \quad p V = p_0 V^* \implies p V = p_0 V_0 \frac{T}{T_0} \implies \frac{p V}{T} = \frac{p_0 V_0}{T_0} = \text{costante}
\end{align}
Per una mole di gas ($n = 1\text{ mol}$), tale costante universale è indicata con $R$ (\textbf{Costante Universale dei Gas}). Generalizzando a un numero di moli $n = m / M_{\text{molare}}$, si ottiene l'\textbf{Equazione di Stato dei Gas Perfetti}:
\begin{equation}
    \boxed{ p V = n R T = \frac{N}{N_A} R T = N k_B T }
\end{equation}
dove $N$ è il numero totale di molecole, $N_A = 6.022 \cdot 10^{23}\text{ mol}^{-1}$ è il \textbf{Numero di Avogadro} e $k_B = R / N_A = 1.381 \cdot 10^{-23}\text{ J/K}$ è la \textbf{Costante di Boltzmann}.

\begin{definizione}{Condizioni STP e Legge di Avogadro}
Le \textbf{Condizioni Standard di Temperatura e Pressione (STP)} corrispondono a $T_0 = 273.15\text{ K}$ ($0^\circ\text{C}$) e $p_0 = 101.325\text{ kPa} = 1\text{ atm}$. 
Per la \textbf{Legge di Avogadro}, volumi uguali di gas diversi, alle medesime condizioni di temperatura e pressione, contengono lo stesso numero di molecole. Nelle condizioni STP, una mole ($n=1\text{ mol}$) di qualsiasi gas perfetto occupa un volume molare standard pari a:
\begin{equation}
    V_{\text{mol}} = \frac{R T_0}{p_0} = 22.414\text{ L} = 22.414 \cdot 10^{-3}\text{ m}^3
\end{equation}
I valori numerici della costante universale dei gas $R$ nelle diverse unità di misura sono pertanto:
\begin{equation}
    R = 8.314 \, \frac{\text{J}}{\text{mol} \cdot \text{K}} = 0.08205 \, \frac{\text{L} \cdot \text{atm}}{\text{mol} \cdot \text{K}} = 1.987 \, \frac{\text{cal}}{\text{mol} \cdot \text{K}}
\end{equation}
\end{definizione}

\begin{definizione}{Legge di Dalton delle Pressioni Parziali}
In un recipiente di volume $V$ contenente una miscela di k gas chimicamente inerti a temperatura $T$, la pressione totale esercitata dalla miscela sulle pareti è uguale alla somma delle \textbf{pressioni parziali} che ciascun gas eserciterebbe se occupasse da solo l'intero volume:
\begin{equation}
    \boxed{ p_{\text{tot}} = \sum_{i=1}^k p_i = \sum_{i=1}^k \frac{n_i R T}{V} = \frac{R T}{V} \sum_{i=1}^k n_i = \frac{n_{\text{tot}} R T}{V} }
\end{equation}
La pressione parziale dell'i-esimo componente è direttamente proporzionale alla sua frazione molare $x_i = n_i / n_{\text{tot}}$, ovvero $p_i = x_i \, p_{\text{tot}}$.
\end{definizione}

\vspace{0.3cm}
\section{Punto Triplo}
\begin{definizione}{Punto Triplo}
Il punto triplo è la combinazione di pressione e temperatura in cui in una sostanza coesistono i tre stati di aggregazione della materia in equilibrio termodinamico.
\end{definizione}

\medskip
Per l'acqua ($\text{H}_2\text{O}$) questo si verifica a $T \simeq 0.01^\circ\text{C}$ (temperatura alla quale \textbf{deve esistere il ghiaccio fondente}) e $p \simeq 4.58 \text{ Torr}$ (pressione alla quale \textbf{deve esistere l'acqua bollente}): la simultaneità di tali coordinate fisiche garantisce la coesistenza in equilibrio dinamico di ghiaccio, acqua liquida e vapore.
L'unità di misura K (Kelvin) della scala termodinamica assoluta si definisce come la frazione $1/273.16$ della temperatura termodinamica del punto triplo dell'acqua:
\begin{equation}
    \boxed{ 1\text{ K} = \frac{T_{\text{punto triplo}}}{273.16} }
\end{equation}
il fatto che sia 16 deriva dal fatto che c'è una miscela differente di isotopi per i quali il punto di congelamento differisce lievemente.

\begin{osservazione}{Significato quantistico di $0\text{ K}$}
Considerando i livelli quantizzati di energia $E_0, E_1, E_2, \dots$ di un sistema, lo \textbf{stato fondamentale} (di minima energia) corrisponde a $E_0$.
Per un sistema di $N$ particelle si ha quindi $E_{\text{tot}} = N \cdot E_0 = E_{0,\text{minima}}$.
Appena una particella passa a uno stato eccitato ($E_i > E_0$) allora $E_{\text{tot}} > E_0$.
Se il sistema è isolato in un bagno termico a temperatura $T=0\text{ K}$, non vi è alcuno scambio termico verso l'esterno che possa abbassare ulteriormente l'energia, né dall'esterno che la innalzi. Si conclude quindi che $0\text{ K}$ è la temperatura a cui ogni sistema si trova nel suo stato fondamentale.
\end{osservazione}

\vspace{0.3cm}
\section{Dilatazione Termica}
\begin{definizione}{Dilatazione termica}
La dilatazione termica è la variazione di ogni dimensione (lineare, superficiale, volumica) di un corpo al variare della temperatura ($T \uparrow \implies V \uparrow$).
\end{definizione}

\medskip
Se $V(0^\circ\text{C}) = V_0$, la variazione può essere descritta come $V(T_C) = V_0 F(T_C)$, dove la funzione $F(0) = 1$ è adimensionale.
Sviluppando in serie di Taylor nell'intorno dello zero si ottiene:
\begin{equation}
F(T_C) = F(0) + \left(\frac{\partial F}{\partial T_C}\right)_0 T_C + \frac{1}{2} \left(\frac{\partial^2 F}{\partial T_C^2}\right)_0 T_C^2 + o(T_C^3) \simeq 1 + \beta_0 T_C
\end{equation}
in quanto sperimentalmente si osserva che la relazione è lineare. Di conseguenza:
\begin{align}
\beta_0 &= \left(\frac{\partial F}{\partial T_C}\right)_0 = \frac{1}{V_0} \left(\frac{\partial V(T_C)}{\partial T_C}\right)_0 
\end{align}
rappresenta il \textbf{coefficiente di dilatazione termica} a $T_C = 0^\circ\text{C}$ e pressione costante.

\subsection{Dilatazione Termica Lineare ($\beta_0 = \lambda$)}
% Disegno di una sbarra di lunghezza l
\begin{figure}[ht]
\centering
\begin{tikzpicture}[scale=1]
\fill[pattern=north east lines, mypen!20] (0,0) rectangle (4,0.3);
\draw[thick, mypen] (0,0) rectangle (4,0.3);
\draw[<->, >=stealth, thick, notered] (0,0.6) -- (4,0.6) node[midway, above, font=\bfseries] {$\ell(T)$};
\draw[dashed, mypen] (0,0.3) -- (0,0.8);
\draw[dashed, mypen] (4,0.3) -- (4,0.8);
\end{tikzpicture}
\caption{Dilatazione termica lineare di una sbarra metallica all'aumentare della temperatura.}
\end{figure}

Il coefficiente lineare $\lambda$ è definito come $\lambda = \frac{1}{\ell_0} \left. \left( \frac{\partial \ell}{\partial T_C} \right) \right|_{T_C = 0^\circ\text{C}}$.
\begin{equation}
\frac{\partial \ell}{\ell} = \lambda \partial T_C \implies \int_{\ell_0}^{\ell(T_C)} \frac{\partial \ell}{\ell} = \ln \left( \frac{\ell(T_C)}{\ell_0} \right) = \int_0^{T_C} \lambda \partial T_C' = \lambda T_C
\end{equation}

Da cui:
\begin{equation}
\ell(T_C) = \ell_0 e^{\lambda T_C} \stackrel{\lambda T_C \ll 1}{\simeq} \ell_0 (1 + \lambda T_C)
\end{equation}
Nel campo di esistenza dei solidi questa approssimazione è accurata. 
Inoltre $\Delta \ell = \ell_0 \lambda T_C \implies \frac{\Delta \ell}{\ell_0} \ll 1 \implies \Delta \ell \ll \ell_0$ poiché $T_C \ll \lambda^{-1} \ll T_{\text{fusione}}$.

\subsection{Dilatazione Termica Superficiale ($\beta_0 = 2\lambda$)}
% Disegno di un rettangolo a x b
\begin{figure}[ht]
\centering
\begin{tikzpicture}[scale=1.1]
\fill[cyan!15, opacity=0.8] (0,0) rectangle (2.5,1.8);
\draw[thick, mypen] (0,0) rectangle (2.5,1.8);
\draw[<->, >=stealth, mypen] (0,-0.2) -- (2.5,-0.2) node[midway, below] {$a$};
\draw[<->, >=stealth, mypen] (-0.2,0) -- (-0.2,1.8) node[midway, left] {$b$};
\node[mypen, font=\bfseries] at (1.25, 0.9) {$S = a b$};
\end{tikzpicture}
\caption{Dilatazione termica superficiale di una lamina rettangolare: la variazione di area è proporzionale al coefficiente $\beta_0 = 2\lambda$.}
\end{figure}

Sviluppando in serie e trascurando i termini di ordine superiore:
\begin{align*}
S(T_C) &= a(T_C) \cdot b(T_C) = a_0 b_0 (1 + \lambda T_C)^2 \\
&= S_0 (1 + \lambda T_C)^2 = S_0 (1 + \lambda^2 T_C^2 + 2\lambda T_C) \\
&\simeq S_0 (1 + 2\lambda T_C)
\end{align*}
da cui si evince che il coefficiente di dilatazione superficiale è $2\lambda$.

\subsection{Dilatazione Termica Volumica ($\beta_0 = 3\lambda$)}
% Disegno di un parallelepipedo a x b x c
\begin{figure}[ht]
\centering
\begin{tikzpicture}[scale=1]
% Parallepipedo in prospettiva
\fill[orange!15, opacity=0.8] (0,0) -- (2.5,0) -- (2.5,1.5) -- (0,1.5) -- cycle;
\fill[orange!25, opacity=0.8] (0,1.5) -- (0.8,2.1) -- (3.3,2.1) -- (2.5,1.5) -- cycle;
\fill[orange!35, opacity=0.8] (2.5,0) -- (3.3,0.6) -- (3.3,2.1) -- (2.5,1.5) -- cycle;
\draw[thick, mypen] (0,0) -- (2.5,0) -- (2.5,1.5) -- (0,1.5) -- cycle;
\draw[thick, mypen] (0,1.5) -- (0.8,2.1) -- (3.3,2.1) -- (2.5,1.5);
\draw[thick, mypen] (2.5,0) -- (3.3,0.6) -- (3.3,2.1);
\draw[dashed, mypen] (0,0) -- (0.8,0.6) -- (3.3,0.6);
\draw[dashed, mypen] (0.8,0.6) -- (0.8,2.1);
\node[below] at (1.25, 0) {$a$};
\node[right] at (3.3, 1.35) {$c$};
\node[above left] at (0.4, 1.8) {$b$};
\node[mypen, font=\bfseries] at (4.5, 1) {$V = a b c$};
\end{tikzpicture}
\caption{Dilatazione termica volumica di un solido geometrico: la dilatazione tridimensionale è regolata da $\beta_0 = 3\lambda$.}
\end{figure}

Sviluppando analogamente si ottiene:
\begin{align*}
V(T_C) &= a(T_C) \cdot b(T_C) \cdot c(T_C) = a_0 b_0 c_0 (1 + \lambda T_C)^3 \\
&= V_0 (1 + \lambda T_C)^3 = V_0 (1 + 3\lambda^2 T_C + 3\lambda T_C^2 + \lambda^3 T_C^3) \\
&\simeq V_0 (1 + 3\lambda T_C)
\end{align*}
da cui si evince che il coefficiente di dilatazione volumica è $3\lambda$.

\vspace{0.3cm}
\subsection{Interpretazione Quantitativa a Livello Molecolare}
% Grafico curva dell'energia potenziale molecolare
\begin{figure}[ht]
\centering
\begin{tikzpicture}[scale=1]
\draw[->, >=stealth, mypen] (-0.5, 0) -- (4, 0) node[right] {$r$};
\draw[->, >=stealth, mypen] (0, -3) -- (0, 2) node[above] {$U(r)$};
\draw[thick, blue] plot [smooth] coordinates {(0.2, 1.5) (0.3, 0) (0.5, -2.5) (1, -1.5) (2, -0.5) (3.5, -0.1)};
\draw[dashed, mypen] (0.5, -2.5) -- (0, -2.5) node[left] {$E_{min}$};
\draw[dashed, mypen] (0.5, -2.5) -- (0.5, 0) node[above] {$r_0$};
\draw[dashed, mypen] (0.15, -1) -- (1.5, -1);
\draw[dashed, mypen] (0.2, -1.8) -- (1, -1.8);
\node at (-0.3, -1) {$E_2$};
\node at (-0.3, -1.8) {$E_1$};
\node at (0.8, -1.6) {$P_1$};
\node at (1.2, -0.8) {$P_2$};
\node at (0.5, -2.7) {$P_0$};
\end{tikzpicture}
\caption{Curva dell'energia potenziale molecolare}
\end{figure}
Sia $U_{\text{eff}}(r)$ l'energia potenziale nel caso generale e $r_0$ la distanza di equilibrio stabile (punto $P_0$).
Se aumentiamo l'energia del sistema al valore $E_1 > E_{min}$, si ottengono due punti di inversione del moto.
In prima approssimazione $U(r)$ è una parabola per piccole oscillazioni, ma all'aumentare dell'energia, a causa della forma asimmetrica della buca di potenziale, il punto medio dell'oscillazione si sposta verso destra.
Questo significa che l'equilibrio dinamico si raggiunge a una distanza media $r > r_0$; di conseguenza il volume complessivo aumenta poiché le molecole si allontanano mediamente l'una dall'altra. Se l'energia fornita è sufficientemente elevata, le molecole si separano definitivamente (causando ad esempio la fusione della sostanza).

\section{Calore e Quantità di Calore}
\subsection{Stessa Sostanza}

% Due corpi
\begin{figure}[ht]
\centering
\begin{tikzpicture}[scale=1]
\fill[pattern=north east lines, mypen!20] (-0.3,-0.3) rectangle (3.3,1.5);
\draw[thick, mypen] (-0.3,-0.3) rectangle (3.3,1.5);
\draw[thick, fill=cyan!20] (0,0) rectangle (1.4,1.2) node[midway, font=\bfseries] {$m_1, T_1$};
\draw[thick, fill=cyan!40] (1.6,0) rectangle (3,1.2) node[midway, font=\bfseries] {$m_2, T_2$};
\node[above] at (1.5, 1.5) {\footnotesize Contenitore adiabatico (nessun dispersione termica)};
\end{tikzpicture}
\caption{Contatto termico tra due blocchi della stessa sostanza a temperature iniziali differenti all'interno di un involucro adiabatico.}
\end{figure}

Consideriamo due corpi della stessa sostanza con massa e temperatura rispettivamente $m_1, T_1$ e $m_2, T_2$, con $T_1 < T_2$.
Se posti a contatto e isolati dall'esterno, essi raggiungono l'equilibrio termico a una temperatura $T$ intermedia ($T_1 < T < T_2$).

Sperimentalmente si osserva che vale la relazione:
\begin{equation}
m_1 (T - T_1) = m_2 (T_2 - T)
\end{equation}
da cui si ricava che $T$ è la media pesata (con le masse) delle temperature iniziali:
\begin{equation}
T = \frac{m_1 T_1 + m_2 T_2}{m_1 + m_2}
\end{equation}

Generalizzando per $n$ corpi, si ottiene:
\begin{equation}
\sum_{i=1}^n m_i (T - T_i) = 0 \implies T = \frac{\sum m_i T_i}{\sum m_i}
\end{equation}

\vspace{0.3cm}
\subsection{Sostanze Diverse}

% Due corpi
\begin{figure}[ht]
\centering
\begin{tikzpicture}[scale=1]
\fill[pattern=north east lines, mypen!20] (-0.3,-0.3) rectangle (3.3,1.5);
\draw[thick, mypen] (-0.3,-0.3) rectangle (3.3,1.5);
\draw[thick, fill=cyan!20] (0,0) rectangle (1.4,1.2) node[midway, font=\bfseries] {$C_1, T_1$};
\draw[thick, fill=orange!30] (1.6,0) rectangle (3,1.2) node[midway, font=\bfseries] {$C_2, T_2$};
\node[above] at (1.5, 1.5) {\footnotesize Scambio di calore interno: $Q_1 + Q_2 = 0$};
\end{tikzpicture}
\caption{Contatto termico tra due corpi di sostanze diverse (con capacità termiche $C_1$ e $C_2$) posti in isolamento termico rispetto all'ambiente.}
\end{figure}

Consideriamo ora due corpi di sostanze diverse, caratterizzati da massa $m_i$, temperatura $T_i$ e da un coefficiente specifico $c_i$:
\begin{itemize}
    \item $c_i$: \textbf{Calore specifico}, un coefficiente caratteristico della \textit{sostanza} (indipendente dalla massa).
    \item $C_i$: \textbf{Capacità termica}, un coefficiente caratteristico del \textit{corpo} (dipende dalla massa, $C_i = m_i c_i$).
\end{itemize}

Se posti a contatto e isolati dall'esterno, l'equilibrio termico si raggiunge a $T_1 < T < T_2$.
Sperimentalmente si osserva la seguente relazione:
\begin{equation}
m_1 c_1 (T - T_1) = m_2 c_2 (T_2 - T) \iff C_1 (T - T_1) = C_2 (T_2 - T)
\end{equation}

Generalizzando per $n$ corpi, la temperatura finale risulta essere la media pesata con le capacità termiche:
\begin{equation}
\sum_{i=1}^n C_i (T - T_i) = 0 \implies T = \frac{\sum C_i T_i}{\sum C_i}
\end{equation}

\begin{definizione}{Quantità di Calore}
Si definisce \textbf{quantità di calore} $Q_i$ ceduta o assorbita da un corpo:
\begin{equation}
Q_i = C_i (T - T_i) = c_i m_i (T - T_i)
\end{equation}
La convenzione dei segni implica:
\begin{itemize}
    \item Se $T > T_i \implies Q_i > 0$: il corpo \textbf{assorbe} calore.
    \item Se $T < T_i \implies Q_i < 0$: il corpo \textbf{cede} calore.
\end{itemize}
\end{definizione}

\medskip
Considerando l'equazione dell'equilibrio termico $\sum m_i c_i (T - T_i) = 0$ come un sistema omogeneo, abbiamo bisogno di fissare arbitrariamente un calore specifico di riferimento per determinare gli altri. Convenzionalmente si utilizza l'acqua distillata.

\begin{definizione}{Caloria}
La \textbf{chilocaloria} (kcal) è la quantità di calore necessaria per innalzare la temperatura di $1\text{ kg}$ di acqua distillata da $14.5^\circ\text{C}$ a $15.5^\circ\text{C}$ alla pressione di $1\text{ atm}$.
In base a questa definizione, il calore specifico dell'acqua è $c_{\text{acqua}} = 1 \, \frac{\text{kcal}}{\text{kg}^\circ\text{C}}$.
\end{definizione}

\begin{osservazione}{Equivalente meccanico del calore}
La chilocaloria non è una grandezza fisica indipendente, poiché il calore è una forma di energia (la cui unità di misura nel SI è il Joule). Tramite l'esperimento di Joule si è stabilita la conversione:
\begin{equation}
1\text{ kcal} = 4184\text{ J} \implies 1\text{ cal} = 4.184\text{ J}
\end{equation}
(Si nota anche che una frigoria corrisponde a una caloria negativa).
\end{osservazione}

\vspace{0.3cm}
\section{Cambiamenti di Stato}
I cambiamenti di stato (a condizioni standard) sono determinati dalla temperatura:
\begin{itemize}
    \item \textbf{Temperatura di Ebollizione e Liquefazione} ($T_E = T_L$):
    \begin{itemize}
        \item Se $T_{\text{corpo}} > T_E \implies$ gas.
        \item Se $T_{\text{corpo}} < T_E \implies$ liquido.
    \end{itemize}
    \item \textbf{Temperatura di Fusione e Solidificazione} ($T_F = T_S$):
    \begin{itemize}
        \item Se $T_{\text{corpo}} > T_F \implies$ liquido.
        \item Se $T_{\text{corpo}} < T_F \implies$ solido.
    \end{itemize}
    \item \textbf{Temperatura di Brinamento e Sublimazione} ($T_B = T_{\text{sub}}$) [più rari]:
    \begin{itemize}
        \item Se $T_{\text{corpo}} > T_B \implies$ gas.
        \item Se $T_{\text{corpo}} < T_B \implies$ solido.
    \end{itemize}
\end{itemize}

\subsection{Calore Latente}

\begin{definizione}{Calore Latente ($\lambda$)}
Il \textbf{calore latente} è la quantità di calore fornita (o sottratta) a un corpo per passare da uno stato di aggregazione all'altro agendo sui legami intermolecolari, riferita all'unità di massa:
\begin{equation}
    \boxed{ \Delta Q = m \lambda \implies \lambda = \frac{\Delta Q}{m} }
\end{equation}
\end{definizione}

Un tipico esempio è quello del calore latente di fusione $(\Delta Q)_{\text{fusione}}$ o di solidificazione $(\Delta Q)_{\text{solidificazione}}$ dell'acqua ($\text{H}_2\text{O}$).

% Grafico T-Q e T-t
\begin{figure}[ht]
\centering
\begin{tikzpicture}
% Primo grafico T-Q
\draw[->, >=stealth, mypen] (-0.5, 0) -- (4, 0) node[right] {$Q$};
\draw[->, >=stealth, mypen] (0, -1) -- (0, 3) node[above] {$T$};
\draw[mypen] (0, 1) -- (-0.1, 1) node[left] {$0^\circ$};
\draw[thick, mygreen] (-0.5, -0.5) -- (0.5, 1);
\draw[thick, mygreen] (0.5, 1) -- (2.5, 1);
\draw[thick, mygreen] (2.5, 1) -- (3.5, 2.5);
\node at (1.5, 0.7) {$(\Delta Q)_F$};
\draw[dashed, mypen] (1.5, 1) -- (1.5, 0);

% Secondo grafico T-t
\begin{scope}[xshift=6cm]
\draw[->, >=stealth, mypen] (-0.5, 0) -- (4, 0) node[right] {$t$};
\draw[->, >=stealth, mypen] (0, -1) -- (0, 3) node[above] {$T$};
\draw[mypen] (0, 1) -- (-0.1, 1) node[left] {$T_F$};
\draw[thick, blue] (-0.5, -0.5) to[out=45, in=200] (0.5, 1);
\draw[thick, blue] (0.5, 1) -- (2.5, 1);
\draw[thick, blue] (2.5, 1) to[out=20, in=260] (3.5, 2.5);
\node at (1.5, 0.7) {$(\Delta t)_F$};
\draw[dashed, mypen] (1.5, 1) -- (1.5, 0);
\end{scope}
\end{tikzpicture}
\caption{Grafici Temperatura-Calore e Temperatura-tempo per un cambiamento di stato}
\end{figure}

In generale si ha che $(\Delta Q) \propto t$.
Nella fase in cui si fornisce il calore di fusione $(\Delta Q)_{\text{fusione}}$, la temperatura non aumenta e rimane fissa a $T_F = 0^\circ\text{C}$ poiché coesistono ghiaccio e acqua mentre sta avvenendo il passaggio di stato (si tratta di un \textbf{processo isotermo}).
Analizzando i processi inversi (come la solidificazione), il calore scambiato è negativo poiché viene sottratta energia al sistema:
\begin{align}
(\Delta Q)_{\text{fusione}} &= -(\Delta Q)_{\text{solidificazione}} \implies \lambda_F = -\lambda_S = \frac{Q_F}{m} = -\frac{Q_S}{m} \\
(\Delta Q)_{\text{evaporazione}} &= -(\Delta Q)_{\text{liquefazione}} \implies \lambda_E = -\lambda_L = \frac{Q_E}{m} = -\frac{Q_L}{m}
\end{align}

\begin{definizione}{Termostato o Bagno Termico}
Un corpo con massa infinita ($m \to \infty$) e quindi capacità termica infinita ($C \to \infty$) è detto \textbf{termostato}, \textbf{serbatoio}, \textbf{sorgente di calore} o \textbf{bagno termico}. 
La sua caratteristica fondamentale è mantenere la temperatura costante, indipendentemente dal calore scambiato.
\end{definizione}

\medskip
Infatti, se mettiamo in contatto due corpi con capacità termiche $C_1$ e $C_2$, assumendo $C_1 \ll C_2 \to \infty$, la temperatura di equilibrio sarà:
\begin{equation}
T_{\text{equil}} = \frac{C_1 T_1 + C_2 T_2}{C_1 + C_2} \simeq \frac{C_2 T_2}{C_2} = T_2
\end{equation}

\vspace{0.3cm}
\section{Propagazione del Calore}
Il calore si propaga in tre modi differenti:
\begin{itemize}
    \item \textbf{Convezione} (Legge di Newton)
    \item \textbf{Conduzione} (Legge di Fourier)
    \item \textbf{Irraggiamento} (Legge di Stefan-Boltzmann)
\end{itemize}

\subsection{Convezione}
\begin{definizione}{Convezione}
La \textbf{convezione} è un meccanismo di trasmissione del calore caratterizzato da un macroscopico \textbf{trasporto di materia}, tipico dei fluidi (liquidi e gas), in cui il movimento di masse fluide trasferisce energia termica da una regione all'altra.
\end{definizione}

\begin{figure}[ht]
\centering
\begin{tikzpicture}[scale=1]
\fill[cyan!10] (0,0) rectangle (2.5, 2.2);
\draw[thick, mypen] (0,0) rectangle (2.5, 2.2);
\draw[->, >=stealth, thick, notered] (0.6, 0.4) to[out=90, in=180] (1.25, 1.8) to[out=0, in=90] (1.9, 0.4);
\draw[->, >=stealth, thick, blue] (1.9, 0.4) to[out=270, in=0] (1.25, 0.1) to[out=180, in=270] (0.6, 0.4);
\draw[->, >=stealth, thick, notered, line width=1.5pt] (1.25, -0.6) -- (1.25, 0) node[midway, right] {$Q$};
\fill[pattern=north east lines, notered!30] (-0.2, -0.6) rectangle (2.7, 0);
\draw[thick, notered] (-0.2, 0) -- (2.7, 0);
\node[font=\bfseries, mypen] at (3.5, 1.1) {\parbox{3cm}{\centering Moti convettivi nel fluido}};
\end{tikzpicture}
\caption{Moti convettivi in un fluido riscaldato dal basso: il fluido caldo, meno denso, sale verso l'alto ed è sostituito dal fluido freddo che scende, instaurando un ciclo termico continuo con trasporto di materia.}
\end{figure}

All'interno di un fluido riscaldato dal basso si generano delle correnti che ciclicamente si scaldano (salendo verso l'alto per via della minore densità) e si raffreddano (scendendo nuovamente).

\begin{definizione}{Legge di Newton (empirica)}
La quantità di calore scambiata nell'unità di tempo fra l'unità di superficie di un corpo e il fluido circostante è proporzionale alla differenza di temperatura:
\begin{equation}
\boxed{Q = h S t \Delta T}
\end{equation}
dove $h$ è il \textbf{coefficiente di conduttività esterna} (determinato empiricamente).
\end{definizione}

\subsection{Conduzione}
\begin{definizione}{Conduzione}
La \textbf{conduzione} è un meccanismo di propagazione del calore in un mezzo materiale che avviene tramite urti e interazioni a livello molecolare \textbf{senza alcun trasporto di materia} (a differenza della convezione).
\end{definizione}

\begin{figure}[ht]
\centering
\begin{tikzpicture}[scale=1]
\fill[pattern=north east lines, notered!25] (0,0) rectangle (0.6, 2);
\draw[thick, mypen, fill=orange!10] (0.6, 0) rectangle (3.4, 2);
\fill[pattern=north west lines, cyan!25] (3.4, 0) rectangle (4, 2);
\draw[->, >=stealth, thick, mypen] (0.6, -0.5) -- (3.8, -0.5) node[right] {$x$};
\node[font=\bfseries] at (0.3, 2.3) {Sorgente Calda ($T_2$)};
\node[font=\bfseries] at (3.7, 2.3) {Sorgente Fredda ($T_1$)};
\draw[<->, >=stealth, thick, mypen] (1.5, 1.5) -- (2.5, 1.5) node[midway, above] {$\Delta x$};
\draw[->, >=stealth, thick, notered, line width=1.2pt] (1.2, 0.8) -- (2.8, 0.8) node[midway, below, font=\bfseries] {Flusso $Q$};
\node at (0.6, -0.8) {0};
\node at (3.4, -0.8) {$L$};
\end{tikzpicture}
\caption{Conduzione termica lungo un conduttore cilindrico isolato lateralmente di lunghezza $L$, posto tra due serbatoi a temperature fisse $T_2 > T_1$.}
\end{figure}

La conduzione termica si verifica fra due solidi a contatto oppure all'interno di uno stesso solido.
Consideriamo due zone a temperature $T_2$ e $T_1$ con $T_2 > T_1$. Fissiamo un asse $x$ con origine sulla superficie che divide il corpo dalla zona a $T_2$.

\begin{figure}[ht]
\centering
\begin{tikzpicture}[scale=1]
\draw[->, >=stealth, thick, mypen] (0,0) -- (3.5,0) node[right] {$x$};
\draw[->, >=stealth, thick, mypen] (0,0) -- (0,2.8) node[above] {$T(x)$};
\draw[dashed, mypen] (0,2.2) -- (2.8,2.2);
\draw[dashed, mypen] (0,0.6) -- (2.8,0.6);
\draw[dashed, mypen] (2.8,0) -- (2.8,0.6);
\node[left, font=\bfseries, notered] at (0, 2.2) {$T_2$};
\node[left, font=\bfseries, blue] at (0, 0.6) {$T_1$};
\node[below] at (2.8, 0) {$L$};
\draw[thick, mygreen, line width=1.5pt] (0,2.2) -- (2.8, 0.6);
\end{tikzpicture}
\caption{Andamento lineare della temperatura $T(x)$ all'interno di un conduttore omogeneo in regime stazionario, caratterizzato da un gradiente termico costante $\frac{dT}{dx} = \frac{T_1-T_2}{L} < 0$.}
\end{figure}

Spostandoci dalla zona a $T_2$ verso quella a $T_1$, la temperatura del corpo diminuisce gradualmente fino a raggiungere l'estremità opposta; di conseguenza esiste un \textbf{gradiente di temperatura} negativo: $\frac{dT}{dx} < 0$.

\begin{definizione}{Legge di Fourier}
La quantità di calore scambiata nell'unità di tempo fra due solidi (o due strati contigui) è proporzionale alla superficie di contatto e al gradiente termico. La formula sperimentale per un flusso su una lunghezza $L$ in regime stazionario è:
\begin{equation}
\boxed{Q = \frac{k S t \Delta T}{L}}
\end{equation}
dove $k$ è il \textbf{coefficiente di conduttività interna} (dipendente dal materiale).
\begin{itemize}
    \item Se $k$ è \textbf{elevato}, il materiale è un buon conduttore termico.
    \item Se $k$ è \textbf{ridotto}, il materiale è un buon isolante termico.
\end{itemize}
\end{definizione}

\medskip
Il ragionamento alla base della formula infinitesima di Fourier considera le seguenti proporzionalità per il flusso di calore $\delta Q$:
\begin{itemize}
    \item $\delta Q \propto S$ (superficie di contatto)
    \item $\delta Q \propto -\frac{dT}{dx}$ (gradiente termico, il segno meno indica il flusso verso temperature inferiori)
    \item $\delta Q \propto k$ (conduttività)
\end{itemize}

\begin{figure}[ht]
\centering
\begin{tikzpicture}[scale=1.2]
\draw[thick, fill=orange!20] (0,0) rectangle (0.6, 1.5);
\draw[dashed, mypen] (0,0) -- (0,-0.4);
\draw[dashed, mypen] (0.6,0) -- (0.6,-0.4);
\draw[<->, >=stealth, thick, mypen] (0,-0.3) -- (0.6,-0.3) node[midway, below] {$dx$};
\node[left, font=\bfseries] at (0, 0.75) {$T + dT$};
\node[right, font=\bfseries] at (0.6, 0.75) {$T$};
\draw[->, >=stealth, thick, notered, line width=1.5pt] (-0.5, 0.75) -- (1.2, 0.75) node[right, font=\bfseries] {$\delta Q$};
\end{tikzpicture}
\caption{Bilancio termico su uno strato infinitesimo di spessore $dx$: il calore fluisce nel verso delle temperature decrescenti ($\delta Q \propto -dT/dx$).}
\end{figure}

Da queste considerazioni si ricava la forma differenziale della Legge di Fourier:
\begin{equation}
\boxed{\frac{\delta Q}{dt} = -k S \frac{dT}{dx}}
\end{equation}
da cui, integrando e supponendo $k$ uniforme, si ottiene la formula precedente.

Se invece abbiamo una conduzione fra materiali differenti in serie (ovvero $k$ non uniforme):

\begin{figure}[ht]
\centering
\begin{tikzpicture}[scale=1]
\fill[pattern=north east lines, notered!20] (-0.4,0) rectangle (0,2.5);
\draw[thick, fill=orange!10] (0,0) rectangle (1.2,2.5) node[midway] {$k_1$};
\draw[thick, fill=yellow!15] (1.2,0) rectangle (2.6,2.5) node[midway] {$k_2$};
\draw[thick, fill=green!10] (2.6,0) rectangle (4,2.5) node[midway] {$k_3$};
\fill[pattern=north west lines, cyan!20] (4,0) rectangle (4.4,2.5);

\draw[<->, >=stealth] (0,-0.3) -- (1.2,-0.3) node[midway, below] {$d_1$};
\draw[<->, >=stealth] (1.2,-0.3) -- (2.6,-0.3) node[midway, below] {$d_2$};
\draw[<->, >=stealth] (2.6,-0.3) -- (4,-0.3) node[midway, below] {$d_3$};

\node[left, font=\bfseries, notered] at (0, 2.8) {$T_2$};
\node[font=\bfseries] at (1.2, 2.8) {$T_A$};
\node[font=\bfseries] at (2.6, 2.8) {$T_B$};
\node[right, font=\bfseries, blue] at (4, 2.8) {$T_1$};

\draw[->, >=stealth, thick, notered, line width=1.5pt] (0.8, -1) -- (3.2, -1) node[midway, below, font=\bfseries] {Flusso stazionario $Q_{\text{TOT}}$};
\end{tikzpicture}
\caption{Conduzione termica attraverso una parete composta da tre strati di materiale differente in serie ($T_2 > T_A > T_B > T_1$). In regime stazionario il flusso di calore $Q_{\text{TOT}}$ è uniforme attraverso ogni strato.}
\end{figure}

\begin{align}
Q_1 &= \frac{k_1 S (T_2 - T_A) t}{d_1} \\
Q_2 &= \frac{k_2 S (T_A - T_B) t}{d_2} \\
Q_3 &= \frac{k_3 S (T_B - T_1) t}{d_3}
\end{align}

In condizioni stazionarie il calore totale fluito attraverso ciascuno strato deve essere lo stesso: $Q_1 = Q_2 = Q_3 = Q_{\text{TOT}}$.
Esprimendo le differenze di temperatura e sommandole:
\begin{equation}
(T_2 - T_A) + (T_A - T_B) + (T_B - T_1) = T_2 - T_1
\end{equation}
da cui si ottiene la relazione per il calore totale ceduto da $T_2$ a $T_1$:
\begin{equation}
\boxed{Q_{\text{TOT}} = \frac{S (T_2 - T_1) t}{\frac{d_1}{k_1} + \frac{d_2}{k_2} + \frac{d_3}{k_3}}}
\end{equation}
Questa relazione è alla base del principio dei doppi vetri: la quantità rilevante per l'isolamento è il rapporto $\frac{d_i}{k_i}$.

\vspace{0.3cm}
\subsection{Irraggiamento}
\begin{definizione}{Irraggiamento}
L'\textbf{irraggiamento} è la trasmissione di energia termica mediante emissione e propagazione di \textbf{onde elettromagnetiche} da parte di qualsiasi corpo a temperatura maggiore dello zero assoluto ($T > 0\text{ K}$). A differenza della conduzione e della convezione, avviene anche nel vuoto assoluto senza necessità di alcun mezzo materiale in interposizione.
\end{definizione}

\begin{definizione}{Legge di Stefan-Boltzmann}
La quantità di calore irradiata da un corpo nero nell'unità di tempo è proporzionale alla quarta potenza della sua temperatura assoluta:
\begin{equation}
\boxed{Q = S t \sigma T^4}
\end{equation}
dove $\sigma = 5.670 \cdot 10^{-8} \, \frac{\text{W}}{\text{m}^2 \text{K}^4}$ è la \textbf{costante di Stefan-Boltzmann}.
(Per i corpi non neri la costante dipende dall'emissività della sostanza).
\end{definizione}

\begin{figure}[ht]
\centering
\begin{tikzpicture}[scale=1.1]
% Cavità del corpo nero (sezione di sfera con foro)
\fill[gray!20] (0,0) circle (1.5);
\fill[pattern=north east lines, mypen!40] (0,0) circle (1.5);
\fill[white] (0,0) circle (1.2);
\fill[white] (1.1,-0.3) rectangle (1.6,0.3);
\draw[thick, mypen] (1.3,0.3) arc[start angle=12, end angle=348, radius=1.35];

% Raggio incidente e riflessioni interne
\draw[->, >=stealth, thick, notered, line width=1.2pt] (2.5, 0.1) -- (1.3, 0.1) node[midway, above] {Radiazione incidente};
\draw[->, >=stealth, thick, notered] (1.3, 0.1) -- (-0.8, 0.7);
\draw[->, >=stealth, thick, notered] (-0.8, 0.7) -- (-0.5, -0.9);
\draw[->, >=stealth, thick, notered] (-0.5, -0.9) -- (0.7, -0.6);
\draw[->, >=stealth, thick, notered] (0.7, -0.6) -- (0.2, 0.9);
\draw[->, >=stealth, thick, notered] (0.2, 0.9) -- (-1.1, -0.2);

\node[font=\bfseries, mypen] at (0, -1.8) {Cavità isoterma (Corpo Nero ideale)};
\end{tikzpicture}
\caption{Modello fisico di Corpo Nero realizzato tramite una cavità con una piccola apertura: ogni raggio di luce che entra subisce innumerevoli riflessioni interne venendo completamente assorbito dalle pareti, senza possibilità di uscirne.}
\end{figure}

\vspace{0.3cm}
\section{Concetti Fondamentali della Termodinamica}

\begin{definizione}{Sistema Termodinamico e Ambiente}
Un \textbf{sistema termodinamico} è una porzione di materia (con composizione chimica nota) separata dall'ambiente circostante tramite una superficie di controllo. 
Se le proprietà chimico-fisiche sono le stesse in tutto il sistema, esso è detto \textbf{omogeneo}. 
Si distinguono tre tipologie di sistemi:
\begin{itemize}
    \item \textbf{Aperti}: scambiano sia materia che energia con l'ambiente.
    \item \textbf{Chiusi}: scambiano solo energia, non materia.
    \item \textbf{Isolati}: non scambiano né materia né energia.
\end{itemize}
\end{definizione}

\begin{figure}[ht]
\centering
\begin{tikzpicture}
    % Sistema generico
    \draw[thick, fill=cyan!10, mypen] (0,0) circle (2);
    \node at (0,0) {Sistema};
    % Frecce di scambio
    \draw[->, >=stealth, thick, notered] (2.5, 1) -- (1.5, 0.5) node[midway, above] {Calore $Q$};
    \draw[<-, >=stealth, thick, mypen] (2.5, -1) -- (1.5, -0.5) node[midway, below] {Lavoro $L$};
\end{tikzpicture}
\caption{Rappresentazione di un sistema termodinamico aperto/chiuso.}
\end{figure}

Lo \textbf{stato del sistema} è definito da un insieme di grandezze macroscopiche dette \textbf{variabili di stato} (pressione $p$, volume $V$, temperatura $T$). Esse sono legate tra loro dall'equazione di stato $f(p, V, T) = 0$. Ad esempio, per un gas perfetto vale $p V = n R T$.

\begin{definizione}{Equilibrio Termodinamico}
Un sistema si trova in stato di \textbf{equilibrio termodinamico} se soddisfa contemporaneamente tre condizioni:
\begin{enumerate}
    \item \textbf{Equilibrio Meccanico}: La risultante delle forze interne e delle forze esterne è nulla ($\sum \vec{F}_{\text{int}} = 0$, $\sum \vec{F}_{\text{ext}} = 0$).
    \item \textbf{Equilibrio Termico}: La temperatura $T$ è uniforme in tutto il sistema.
    \item \textbf{Equilibrio Chimico}: Assenza di reazioni chimiche macroscopiche e di spostamento netto di materia.
\end{enumerate}
In questo stato la situazione è stazionaria, è descrivibile da poche variabili di stato ed è indipendente dalla storia pregressa del sistema.
\end{definizione}

\vspace{0.3cm}
\subsection{Trasformazioni Termodinamiche}
Una trasformazione termodinamica è un processo durante il quale alcune variabili di stato variano da una situazione iniziale $A$ (di equilibrio) a una situazione finale $B$ (di equilibrio).
Se lo stato finale coincide con lo stato iniziale ($A \equiv B$), la trasformazione è detta \textbf{ciclica}.

Inoltre, se la trasformazione soddisfa le seguenti condizioni:
\begin{itemize}
    \item È una successione continua di stati di equilibrio infinitamente vicini (così che le variabili di stato siano definite in ogni istante).
    \item È infinitamente lenta (il sistema ha il tempo per riequilibrarsi).
    \item Le cause che la generano sono infinitamente piccole.
\end{itemize}
allora la trasformazione è definita \textbf{reversibile} (cioè invertibile per il sistema e per l'ambiente senza lasciare alcuna traccia). Altrimenti, la trasformazione è \textbf{irreversibile} (ogni processo reale caratterizzato da attrito meccanico, turbolenza o gradienti finiti di temperatura è irreversibile).

\begin{osservazione}{Criterio Globale di Reversibilità e Ciclicità}
Una distinzione cruciale tra processi reversibili e irreversibili risiede negli effetti sull'ambiente:
\begin{itemize}
    \item \textbf{Processo Reversibile}: Effettuando la trasformazione diretta $A \to B$ e successivamente la trasformazione inversa $B \to A$, sia il \textbf{sistema} che l'\textbf{ambiente} tornano esattamente nei rispettivi stati iniziali, senza che nell'universo rimanga alcuna traccia residua del ciclo compiuto.
    \item \textbf{Processo Irreversibile}: Anche se è sempre possibile riportare il \textbf{sistema} allo stato iniziale $A$ (facendogli compiere un ciclo chiuso per cui $\Delta U_{\text{sist}} = 0$), l'\textbf{ambiente} subirà inevitabilmente una modificazione permanente (ad esempio un dispendio netto di lavoro trasformato in calore dissipato nel termostato freddo).
\end{itemize}
\end{osservazione}

\begin{osservazione}{Accoppiamento Termico e Catena Infinita di Termostati}
Consideriamo un corpo di capacità termica $C$ inizialmente a temperatura $T_1$. Se lo poniamo bruscamente a contatto con un singolo termostato a temperatura $T_2 > T_1$, il corpo assorbe una quantità di calore $Q = C(T_2 - T_1)$. Se successivamente lo raffreddiamo ponendolo a contatto con un termostato a $T_1$, esso cede il calore $Q' = -C(T_2 - T_1)$. 
Per il corpo il processo è ciclico ($\Delta U = 0$). Tuttavia, per l'ambiente si è verificato un trasferimento diretto di calore dalla sorgente calda $T_2$ alla sorgente fredda $T_1$ senza alcuna produzione di lavoro, un processo fortemente \textbf{irreversibile}.
L'unico modo concettuale per rendere lo scambio termico \textbf{quasi-reversibile} consiste nell'accoppiare il corpo con una \textbf{catena infinita di termostati ausiliari} a temperature $T_1, T_1+\delta T, T_1+2\delta T, \dots, T_2$, con differenze infinitesime $\delta T \to 0$, in modo che in ogni istante il corpo e il termostato in contatto siano in equilibrio termico approssimato.
\end{osservazione}

\begin{osservazione}{Rappresentazione nel Piano di Clapeyron ($p, V$)}
Nel diagramma di Clapeyron, una trasformazione può essere rappresentata da una \textbf{curva continua} e ben definita \textbf{solo se è reversibile}, poiché solo in tal caso ogni punto intermedio corrisponde a uno stato di equilibrio macroscopico con pressione $p$ e temperatura $T$ uniformi in tutto il volume. 
Per una trasformazione \textbf{irreversibile}, durante il processo intermedio il sistema non è in equilibrio e le variabili di stato non hanno valori uniformi o definiti; pertanto, \textbf{la curva $p(V)$ non esiste}. Graficamente, una trasformazione irreversibile tra due stati di equilibrio $A$ e $B$ si indica convenzionalmente con una \textbf{linea spezzata o tratteggiata}.
\end{osservazione}

\begin{figure}[ht]
\centering
\begin{tikzpicture}[scale=1]
% Espansione
\fill[pattern=north east lines, mypen!20] (0,0) rectangle (0.6, 1.6);
\draw[thick, mypen] (0,0) rectangle (0.6, 1.6);
\draw[thick, mypen, fill=cyan!10] (0.6, 0) rectangle (2.2, 1.6);
\draw[thick, fill=gray!30] (2.2, 0) rectangle (2.4, 1.6);
\draw[->, >=stealth, thick, mypen, line width=1.2pt] (2.4, 0.8) -- (3.6, 0.8) node[right, font=\bfseries] {$F$};
\draw[<-, >=stealth, thick, notered] (0.6, 1.3) -- (1.8, 1.3) node[right] {$pS$};
\draw[<-, >=stealth, thick, mygreen] (0.6, 0.3) -- (1.8, 0.3) node[right] {$F_A$};
\node[above, font=\bfseries] at (2.8, 1.2) {Vuoto};
\node[below, font=\bfseries, blue] at (1.4, -0.2) {Espansione ($dV > 0$)};
\node[right] at (4, 0.8) {
    \parbox{7cm}{\small
    \textbf{Equilibrio dinamico:}\\
    $pS = F + F_A \implies F = pS - F_A < pS$ \\
    \textit{La forza motrice esterna richiesta è minore rispetto al caso ideale senza attrito.}
    }
};

% Compressione
\begin{scope}[yshift=-3.8cm]
\fill[pattern=north east lines, mypen!20] (0,0) rectangle (0.6, 1.6);
\draw[thick, mypen] (0,0) rectangle (0.6, 1.6);
\draw[thick, mypen, fill=orange!10] (0.6, 0) rectangle (2.2, 1.6);
\draw[thick, fill=gray!30] (2.2, 0) rectangle (2.4, 1.6);
\draw[<-, >=stealth, thick, mypen, line width=1.2pt] (2.4, 0.8) -- (3.6, 0.8) node[right, font=\bfseries] {$F$};
\draw[->, >=stealth, thick, notered] (0.6, 1.3) -- (1.8, 1.3) node[right] {$pS$};
\draw[<-, >=stealth, thick, mygreen] (0.6, 0.3) -- (1.8, 0.3) node[right] {$F_A$};
\node[above, font=\bfseries] at (2.8, 1.2) {Vuoto};
\node[below, font=\bfseries, notered] at (1.4, -0.2) {Compressione ($dV < 0$)};
\node[right] at (4, 0.8) {
    \parbox{7cm}{\small
    \textbf{Equilibrio dinamico:}\\
    $pS + F_A = F \implies F = pS + F_A > pS$ \\
    \textit{La forza motrice esterna richiesta è maggiore rispetto al caso ideale senza attrito.}
    }
};
\end{scope}
\end{tikzpicture}
\caption{Confronto meccanico tra espansione e compressione di un gas in un cilindro con pistone in presenza di forza d'attrito viscoso $F_A$: l'istèresi meccanica dimostra l'irreversibilità dei processi reali.}
\end{figure}

Di conseguenza, $F_{\text{ESP}} \neq F_{\text{COM}}$, quindi in presenza di attrito il percorso di andata è sempre diverso dal percorso di ritorno, confermando l'irreversibilità del processo reale.
\section{Primo Principio della Termodinamica}

\begin{definizione}{Calore, Lavoro ed Energia Interna}
Il Primo Principio della Termodinamica è un'estensione del principio di conservazione dell'energia. L'energia interna $U$ di un sistema è una \textbf{funzione di stato}; la sua variazione $\Delta U$ tra due stati $A$ e $B$ dipende solo dagli stati estremi e non dal percorso seguito:
\begin{equation}
    \Delta U = U_B - U_A = Q - L
\end{equation}
dove:
\begin{itemize}
    \item $Q$ è il calore scambiato dal sistema (positivo se assorbito, negativo se ceduto).
    \item $L$ è il lavoro compiuto (positivo se compiuto dal sistema sull'ambiente, negativo se subito).
\end{itemize}
\end{definizione}

In forma differenziale si scrive:
\begin{equation}
    dU = \delta Q - \delta L
\end{equation}

\begin{osservazione}{Differenziali Inesatti ($\delta Q$ e $\delta L$)}
Si nota l'uso del simbolo $\delta$ (o đ) per calore e lavoro, in contrapposizione al differenziale esatto $dU$. Questo indica che $Q$ e $L$ \textbf{non sono funzioni di stato}: le quantità infinitesime $\delta Q$ e $\delta L$ non sono i differenziali di una funzione del solo punto (stato), ma dipendono dal cammino seguito lungo la trasformazione. L'integrale $\int_A^B \delta L = L_{AB}$ rappresenta il lavoro lungo una specifica traiettoria, mentre $\int_A^B dU = U_B - U_A$ è lo stesso per qualunque cammino.
\end{osservazione}

\subsection{Lavoro delle Forze di Pressione (Trasformazioni Reversibili e Irreversibili)}
Consideriamo un fluido racchiuso all'interno di un cilindro di sezione di area $S$ munito di un pistone mobile. Se il pistone si sposta di un tratto infinitesimo $dx$, la forza esterna che si oppone all'espansione è $F_e = p_e S$, dove $p_e$ è la pressione esterna. Il lavoro infinitesimo compiuto dal sistema contro l'ambiente esterno è:
\begin{equation}
    \delta L = F_e \, dx = p_e S \, dx = p_e \, dV
\end{equation}
(Il lavoro compiuto dalle forze esterne sul sistema è di segno opposto, $\delta L_e = -p_e dV$).

\begin{itemize}
    \item \textbf{Trasformazione Reversibile}: In ogni istante il sistema è in equilibrio meccanico con l'ambiente, per cui la pressione interna del gas è uniforme e bilancia esattamente quella esterna ($p = p_e$). Il lavoro totale per una variazione finita di volume da $V_A$ a $V_B$ è dato dall'integrale:
    \begin{equation}
        \boxed{ L_{\text{rev}} = \int_{V_A}^{V_B} p(V) \, dV }
    \end{equation}
    Nel piano di Clapeyron ($p, V$), questo integrale rappresenta geometricamente l'\textbf{area sottesa dalla curva} della trasformazione.
    \item \textbf{Trasformazione Irreversibile}: Durante una trasformazione irreversibile non esiste una pressione $p$ interna uniforme definita, per cui l'integrale $\int p \, dV$ perde di significato rispetto alle coordinate del sistema. Il lavoro può essere calcolato \textbf{solo valutando le forze esterne}, ovvero $L_{\text{irr}} = \int p_e \, dV$. Se ad esempio un gas si espande liberamente nel vuoto ($p_e = 0$), il lavoro compiuto è rigorosamente nullo ($L = 0$), nonostante il volume aumenti da $V_A$ a $V_B$.
\end{itemize}

\begin{figure}[ht]
\centering
\begin{tikzpicture}
    % Asse p-V
    \draw[->, >=stealth, thick, mypen] (0,0) -- (0,4) node[above] {$p$};
    \draw[->, >=stealth, thick, mypen] (0,0) -- (5,0) node[right] {$V$};
    
    % Area sotto la curva per il lavoro
    \fill[blue!10, opacity=0.5] (1,0) -- (1,3) to[out=-60, in=150] (4,1) -- (4,0) -- cycle;
    
    % Trasformazione
    \draw[thick, blue, ->, >=stealth] (1,3) to[out=-60, in=150] (4,1);
    
    % Punti A e B
    \filldraw[mypen] (1,3) circle (2pt) node[above right] {$A$};
    \filldraw[mypen] (4,1) circle (2pt) node[above right] {$B$};
    
    \node at (2.5, 1) {$L = \int p \, dV$};
\end{tikzpicture}
\caption{Lavoro termodinamico di una trasformazione reversibile come area sottesa nel piano di Clapeyron ($p-V$).}
\end{figure}

\vspace{0.3cm}
\section{Calorimetria e Capacità Termica}

\begin{definizione}{Calore Specifico Molare}
Oltre al calore specifico per unità di massa ($c = \frac{1}{m}\frac{\delta Q}{dT}$), per i gas è più utile riferirsi al numero di moli $n$. Si definiscono due \textbf{calori specifici molari}:
\begin{itemize}
    \item $c_v$: calore specifico molare a volume costante.
    \item $c_p$: calore specifico molare a pressione costante.
\end{itemize}
Per solidi e liquidi la dilatazione termica è trascurabile ($dV \simeq 0$), sicché il lavoro di espansione è nullo e le due capacità termiche coincidono ($c_p \simeq c_v \simeq c$). Per i gas, invece, la differenza è notevole ed è regolata dalla \textbf{relazione di Mayer}:
\begin{equation}
    c_p - c_v = R
\end{equation}
\end{definizione}

\begin{definizione}{Calori Specifici dei Gas Perfetti, Gradi di Libertà ed Energia Interna}
In base al teorema di equipartizione dell'energia della meccanica statistica, ogni grado di libertà molecolare contribuisce per $\frac{1}{2}R$ al calore specifico molare a volume costante $c_v$. Definendo il rapporto dei calori specifici $\gamma = \frac{c_p}{c_v} > 1$, si ha:
\begin{itemize}
    \item \textbf{Gas Monoatomici} (3 gradi di libertà traslazionali):
    \begin{equation}
        c_v = \frac{3}{2}R, \quad c_p = \frac{5}{2}R, \quad \gamma = \frac{5}{3} \simeq 1.67
    \end{equation}
    \item \textbf{Gas Biatomici} (3 traslazionali + 2 rotazionali a temperatura ambiente):
    \begin{equation}
        c_v = \frac{5}{2}R, \quad c_p = \frac{7}{2}R, \quad \gamma = \frac{7}{5} = 1.40
    \end{equation}
    \item \textbf{Gas Poliatomici} (3 traslazionali + 3 rotazionali):
    \begin{equation}
        c_v = 3R, \quad c_p = 4R, \quad \gamma = \frac{4}{3} \simeq 1.33
    \end{equation}
\end{itemize}
Per un gas perfetto, l'energia interna $U$ è funzione esclusiva della temperatura (Legge di Joule). La variazione di energia interna tra due stati a temperature $T_A$ e $T_B$ vale pertanto:
\begin{equation}
    \boxed{ \Delta U = n c_v \Delta T = n c_v (T_B - T_A) }
\end{equation}
Questa relazione fondamentale è valida per \textbf{qualunque trasformazione} (isocora, isobara, isoterma, adiabatica o persino irreversibile), poiché $U$ è una funzione di stato.
\end{definizione}

\begin{definizione}{Legge di Dulong e Petit per i Solidi}
Per i solidi cristallini a temperatura ambiente (al di sopra della temperatura caratteristica di Debye), ogni atomo oscilla attorno alla propria posizione di equilibrio reticolare comportandosi come un oscillatore armonico tridimensionale (3 gradi di libertà cinetici e 3 potenziali). Il calore specifico molare approssima un valore universale indipendente dalla natura chimica della sostanza:
\begin{equation}
    \boxed{ C_{\text{mol}} \simeq 6 \cdot \frac{1}{2}R = 3R \simeq 24.9 \, \frac{\text{J}}{\text{mol} \cdot \text{K}} \simeq 6 \, \frac{\text{cal}}{\text{mol} \cdot \text{K}} }
\end{equation}
\end{definizione}

\begin{figure}[ht]
\centering
\begin{tikzpicture}
    % Calorimetro
    \draw[thick, mypen] (0,0) -- (0,3) -- (3,3) -- (3,0) -- cycle;
    \draw[thick, mypen] (0.2,0.2) -- (0.2,2.8) -- (2.8,2.8) -- (2.8,0.2) -- cycle;
    
    % Termometro
    \draw[thick, mypen, fill=notered] (1.5, 0.5) circle (0.15);
    \draw[thick, mypen] (1.4, 0.5) -- (1.4, 4) -- (1.6, 4) -- (1.6, 0.5);
    \node at (1.5, 4.2) {$T$};
    
    % Acqua
    \fill[cyan!20] (0.2, 0.2) rectangle (2.8, 2);
    
    \node at (1.5, -0.5) {Calorimetro};
\end{tikzpicture}
\caption{Schema semplificato di un calorimetro.}
\end{figure}

\vspace{0.3cm}
\section{Tipi di Trasformazioni Termodinamiche}
Un sistema può passare da uno stato iniziale a uno finale attraverso diverse trasformazioni notevoli:
\begin{itemize}
    \item \textbf{Isoterma} ($T = \text{costante}$): 
    \begin{equation}
        p V = \text{cost.} \implies L = n R T \ln\left(\frac{V_B}{V_A}\right), \quad \Delta U = 0
    \end{equation}
    
    \item \textbf{Isocora} ($V = \text{costante}$):
    \begin{equation}
        \frac{p}{T} = \text{cost.} \implies L = 0, \quad Q = \Delta U = n c_v \Delta T
    \end{equation}
    
    \item \textbf{Isobara} ($p = \text{costante}$):
    \begin{equation}
        \frac{V}{T} = \text{cost.} \implies L = p \Delta V, \quad Q = n c_p \Delta T
    \end{equation}
    
    \item \textbf{Adiabatica} ($Q = 0$):
    \begin{equation}
        p V^\gamma = \text{cost.} \quad \text{con} \quad \gamma = \frac{c_p}{c_v}
    \end{equation}
    Il lavoro si esprime come $L = -\Delta U = -n c_v \Delta T$.
\end{itemize}

\begin{definizione}{Espansione Libera di Joule (Adiabatica Irreversibile)}
L'\textbf{espansione libera di Joule} è un processo in cui un gas perfetto, contenuto in un recipiente rigidamente isolato termicamente ($Q = 0$), si espande nel vuoto all'apertura di una valvola di comunicazione. Poiché l'espansione avviene contro una pressione esterna nulla ($p_e = 0$), il lavoro compiuto è nullo ($L = 0$). Dal Primo Principio della Termodinamica segue:
\begin{equation}
    \Delta U = Q - L = 0 \implies U_B = U_A
\end{equation}
Essendo l'energia interna di un gas perfetto funzione esclusiva della temperatura, si conclude che $T_B = T_A$ ($\Delta T = 0$). 
\textbf{Attenzione}: Sebbene la temperatura iniziale e finale coincidano, questa \textbf{non è una trasformazione isoterma reversibile}. Durante il moto turbolento nel vuoto, il gas attraversa stati di non-equilibrio in cui pressione e temperatura non sono definite uniformemente. Il processo è \textbf{fortemente irreversibile} e comporta un aumento di entropia dell'universo $\Delta S = n R \ln(V_B / V_A) > 0$. Nel diagramma di Clapeyron i punti estremi $A$ e $B$ giacciono sulla medesima iperbole isoterma, ma non possono essere congiunti da alcuna curva continua.
\end{definizione}

\begin{figure}[ht]
\centering
\begin{tikzpicture}[scale=1.1]
    % Griglia di fondo leggera
    \draw[step=1cm, gray!15, very thin] (0,0) grid (5.2,5.2);

    % Assi
    \draw[->, >=stealth, thick, mypen] (0,0) -- (0,5.2) node[above] {$p$};
    \draw[->, >=stealth, thick, mypen] (0,0) -- (5.2,0) node[right] {$V$};
    
    % Punto di intersezione comune A(2, 3) -> p*V = 6 per isoterma; p*V^1.4 = 7.917 per adiabatica
    % Isocora (V = 2)
    \draw[thick, mygreen] (2, 0.8) -- (2, 4.8) node[above] {\small Isocora};
    
    % Isobara (p = 3)
    \draw[thick, orange] (0.8, 3) -- (4.8, 3) node[right] {\small Isobara};
    
    % Isoterma (p = 6/V)
    \draw[thick, blue, domain=1.2:4.8, samples=50] plot (\x, {6/\x}) node[right] {\small Isoterma};
    
    % Adiabatica (p = 7.917 / V^1.4)
    \draw[thick, notered, domain=1.35:4.8, samples=50] plot (\x, {7.917/(\x^1.4)}) node[right] {\small Adiabatica};
    
    % Punto di intersezione e proiezioni
    \draw[dashed, gray] (2,0) node[below] {$V_0$} -- (2,3) -- (0,3) node[left] {$p_0$};
    \filldraw[mypen] (2,3) circle (2pt) node[above right=1pt] {$A$};
\end{tikzpicture}
\caption{Confronto tra le quattro trasformazioni fondamentali nel piano $p-V$, intersecantesi nel medesimo stato iniziale $A(V_0, p_0)$. Si noti come, durante un'espansione ($V > V_0$), la curva adiabatica scenda più rapidamente rispetto all'isoterma a causa dell'assenza di apporto di calore ($\gamma > 1$).}
\end{figure}
\vspace{0.3cm}
\subsection{Derivazione dell'Adiabatica Reversibile}
Nella trasformazione adiabatica non c'è scambio di calore ($Q = 0$), da cui per il primo principio si ottiene $L = -\Delta U$.
In forma infinitesima per un gas perfetto reversibile si ha:
\begin{equation}
    \delta Q = \delta L + dU = p \, dV + n c_v \, dT = 0
\end{equation}

Dall'equazione di stato $pV = nRT$ differenziando si ricava:
\begin{equation}
    dp \, V + p \, dV = n R \, dT \implies n \, dT = \frac{V}{R} dp + \frac{p}{R} dV
\end{equation}

Sostituendo il termine $n \, dT$ nell'equazione del primo principio:
\begin{equation}
    p \, dV + c_v \left( \frac{V}{R} dp + \frac{p}{R} dV \right) = 0
\end{equation}

Raccogliendo i termini, considerando che $R = c_p - c_v$ e dividendo tutto per $c_v/R$:
\begin{equation}
    p \left( \frac{c_v + R}{c_v} \right) dV + V \, dp = 0 \implies p \frac{c_p}{c_v} dV + V \, dp = 0
\end{equation}

Posto il coefficiente adiabatico $\gamma = \frac{c_p}{c_v}$, si ricava:
\begin{equation}
    \gamma \frac{dV}{V} = - \frac{dp}{p}
\end{equation}

Essendo la trasformazione reversibile, è possibile integrare lungo la curva passante per gli stati $A$ e $B$:
\begin{equation}
    \gamma \int_{V_A}^{V_B} \frac{dV}{V} = -\int_{p_A}^{p_B} \frac{dp}{p} \implies \gamma \ln\left(\frac{V_B}{V_A}\right) + \ln\left(\frac{p_B}{p_A}\right) = 0
\end{equation}
Applicando le proprietà dei logaritmi si giunge infine alle \textbf{equazioni di Poisson} per l'adiabatica reversibile:
\begin{equation}
    \boxed{p V^\gamma = \text{cost.}} \quad \quad \boxed{T V^{\gamma-1} = \text{cost.}} \quad \quad \boxed{p^{1-\gamma} T^\gamma = \text{cost.}}
\end{equation}

\vspace{0.3cm}
\section{Secondo Principio e Macchine Termiche}

\begin{definizione}{Enunciati del Secondo Principio}
Esistono due formulazioni storiche equivalenti del Secondo Principio della Termodinamica:
\begin{itemize}
    \item \textbf{Enunciato di Kelvin-Planck}: È impossibile realizzare una trasformazione il cui unico risultato sia quello di assorbire calore da una sola sorgente e trasformarlo integralmente in lavoro.
    \item \textbf{Enunciato di Clausius}: È impossibile realizzare una trasformazione il cui unico risultato sia il passaggio di calore da un corpo a temperatura minore a uno a temperatura maggiore.
\end{itemize}
\end{definizione}

\medskip
Una \textbf{macchina termica} opera ciclicamente tra due sorgenti a temperature $T_c$ (calda) e $T_f$ (fredda). Il suo \textbf{rendimento} $\eta$ è definito come il rapporto tra il lavoro prodotto $L$ e il calore assorbito dalla sorgente calda $Q_c$:
\begin{equation}
    \eta = \frac{L}{Q_c} = \frac{Q_c - |Q_f|}{Q_c} = 1 - \frac{|Q_f|}{Q_c}
\end{equation}

Per il ciclo reversibile di \textbf{Carnot}, che rappresenta il rendimento massimo teorico operando tra due temperature, vale:
\begin{equation}
    \eta_{\text{Carnot}} = 1 - \frac{T_f}{T_c}
\end{equation}

\begin{figure}[ht]
\centering
\begin{tikzpicture}
    % Sorgente Calda
    \draw[thick, mypen, fill=notered!30] (2,4) rectangle (4,5);
    \node at (3,4.5) {Sorgente $T_c$};
    
    % Macchina Termica
    \draw[thick, mypen, fill=gray!20] (3,2.5) circle (0.8);
    \node at (3,2.5) {M.T.};
    
    % Sorgente Fredda
    \draw[thick, mypen, fill=cyan!30] (2,0) rectangle (4,1);
    \node at (3,0.5) {Sorgente $T_f$};
    
    % Frecce
    \draw[->, >=stealth, thick, notered] (3,4) -- (3,3.3) node[midway, right] {$Q_c$};
    \draw[->, >=stealth, thick, cyan] (3,1.7) -- (3,1) node[midway, right] {$Q_f$};
    \draw[->, >=stealth, thick, mypen] (3.8, 2.5) -- (5, 2.5) node[midway, above] {$L$};
\end{tikzpicture}
\caption{Schema di principio di una macchina termica.}
\end{figure}

\vspace{0.3cm}
\subsection{Espansione Libera (Trasformazione Irreversibile)}

Se il gas espande nel vuoto non compie lavoro macroscopico contro forze esterne e, qualora il recipiente sia adiabatico, non scambia calore.

\begin{figure}[ht]
\centering
\begin{tikzpicture}[scale=1.1]
    % Pareti adiabatiche esterne
    \fill[pattern=north east lines, mypen!25] (-0.3,-0.3) rectangle (5.5,1.8);
    \fill[white] (0,0) rectangle (5.2,1.5);
    \draw[thick, mypen] (0,0) rectangle (5.2,1.5);
    
    % Prima scatola A (Stato A)
    \draw[thick, fill=cyan!20] (0,0) rectangle (2.3,1.5);
    \fill[pattern=north west lines, cyan!40] (0,0) rectangle (2.3,1.5);
    \draw[thick, mypen, line width=1.5pt] (2.3,0) -- (2.3,1.5);
    \node[font=\bfseries] at (1.15, 0.75) {Gas ($V_A, T_A$)};
    \node[font=\bfseries, gray] at (3.75, 0.75) {Vuoto ($p=0$)};
    \node[below, font=\bfseries] at (2.6, -0.4) {Stato A (Iniziale): Setto chiuso};

    % Seconda scatola B (Stato B)
    \begin{scope}[yshift=-3cm]
    \fill[pattern=north east lines, mypen!25] (-0.3,-0.3) rectangle (5.5,1.8);
    \fill[white] (0,0) rectangle (5.2,1.5);
    \draw[thick, mypen] (0,0) rectangle (5.2,1.5);
    \draw[thick, fill=cyan!15] (0,0) rectangle (5.2,1.5);
    \draw[dashed, mypen] (2.3,0) -- (2.3,1.5);
    \node[font=\bfseries] at (2.6, 0.75) {Gas espanso all'equilibrio ($V_B > V_A, T_B = T_A$)};
    \node[below, font=\bfseries] at (2.6, -0.4) {Stato B (Finale): Setto rimosso};
    \end{scope}
    
    \draw[->, >=stealth, thick, mypen, line width=1.5pt] (2.6, -0.9) -- (2.6, -1.3) node[midway, right] {Espansione libera};
\end{tikzpicture}
\caption{Esperimento di espansione libera di Joule nel vuoto all'interno di un contenitore rigido e adiabatico: poiché $L=0$ e $Q=0$, per il Primo Principio si ha $\Delta U = 0$ e, per un gas perfetto, la temperatura rimane invariata ($T_B = T_A$).}
\end{figure}

Poiché il gas si espande nel vuoto, la pressione esterna è nulla ($p_e = 0$) e non viene compiuto lavoro ($L = 0$). Inoltre, essendo il recipiente isolato termicamente, non vi è scambio di calore ($Q = 0$). Per il Primo Principio della Termodinamica si ha dunque:
\begin{equation}
    \Delta U = Q - L = 0
\end{equation}
Per un gas perfetto, l'energia interna dipende unicamente dalla temperatura ($U = U(T)$), quindi $\Delta U = 0 \implies \Delta T = 0$. La trasformazione tra lo stato iniziale $A$ e finale $B$ giunge alla stessa temperatura ($T_A = T_B$), comportandosi complessivamente come una trasformazione \textbf{isoterma} che segue la legge di Boyle:
\begin{equation}
    p_A V_A = p_B V_B
\end{equation}

\vspace{0.3cm}
\subsection{Confronto delle Pendenze tra Adiabatica e Isoterma}
Nel piano di Clapeyron ($p-V$), una curva adiabatica reversibile risulta sempre più ripida di una curva isoterma passante per il medesimo punto.

\begin{figure}[ht]
\centering
\begin{minipage}{0.42\textwidth}
\centering
\begin{tikzpicture}[scale=1]
    % Griglia
    \draw[step=1cm, gray!15, very thin] (0,0) grid (4.2,4.2);
    % Assi
    \draw[->, >=stealth, thick, mypen] (0,0) -- (0,4.2) node[left] {$p$};
    \draw[->, >=stealth, thick, mypen] (0,0) -- (4.2,0) node[below] {$V$};
    
    % Curve passanti per P(1.8, 2.2): p*V = 3.96; p*V^1.4 = 5.019
    \draw[thick, blue, domain=1.0:3.8, samples=40] plot (\x, {3.96/\x}) node[right] {\small Isoterma};
    \draw[thick, notered, domain=1.15:3.6, samples=40] plot (\x, {5.019/(\x^1.4)}) node[right] {\small Adiabatica};
    
    % Tangente Isoterma (pendenza -1.22)
    \draw[thin, blue!70, dashed, domain=1.1:2.7] plot (\x, {4.4 - 1.22*\x});
    
    % Tangente Adiabatica (pendenza -1.71)
    \draw[thin, notered!70, dashed, domain=1.25:2.5] plot (\x, {5.28 - 1.71*\x});
    
    % Punto di contatto
    \filldraw[mypen] (1.8, 2.2) circle (1.8pt) node[above right=1pt] {$P$};
    \draw[dashed, gray] (1.8,0) node[below] {$V_0$} -- (1.8,2.2) -- (0,2.2) node[left] {$p_0$};
\end{tikzpicture}
\end{minipage}\hfill
\begin{minipage}{0.54\textwidth}
Differenziando l'equazione dell'adiabatica reversibile $pV^\gamma = \text{cost.}$:
\begin{align*}
    d(pV^\gamma) = 0 &\implies V^\gamma dp + \gamma V^{\gamma-1} p\,dV = 0 \\
    &\implies \left( \frac{dp}{dV} \right)_{\text{adiab}} = - \gamma \frac{p}{V}
\end{align*}
Confrontando con la pendenza dell'isoterma ($pV = \text{cost.} \implies \gamma = 1$):
\begin{align}
    \left( \frac{dp}{dV} \right)_{\text{isot}} &= - \frac{p}{V} \\
    \left( \frac{dp}{dV} \right)_{\text{adiab}} &= - \gamma \frac{p}{V} = \gamma \left( \frac{dp}{dV} \right)_{\text{isot}}
\end{align}
Poiché $\gamma = \frac{c_p}{c_v} > 1$, la pendenza (in modulo) della curva adiabatica nel punto di intersezione $P(V_0, p_0)$ è maggiore di quella dell'isoterma esattamente di un fattore $\gamma$.
\end{minipage}
\caption{Confronto analitico e geometrico delle pendenze tra una curva isoterma (blu) e un'adiabatica reversibile (rossa) passanti per il medesimo stato $P(V_0, p_0)$. Le rette tratteggiate indicano le rispettive tangenti in $P$.}
\end{figure}

\vspace{0.3cm}
\subsection{Macchine Frigorifere}
A differenza della macchina termica, una \textbf{macchina frigorifera} (o pompa di calore) opera assorbendo un lavoro $L < 0$ dall'ambiente per estrarre calore $Q_f > 0$ da una sorgente fredda a temperatura $T_f$ e cederlo come calore $Q_c < 0$ a una sorgente calda a temperatura $T_c$.

\begin{figure}[ht]
\centering
\begin{tikzpicture}
    % Sorgente Calda
    \draw[thick, mypen, fill=notered!30] (2,4) rectangle (4,5);
    \node at (3,4.5) {Sorgente $T_c$};
    
    % Macchina Frigorifera
    \draw[thick, mypen, fill=gray!20] (3,2.5) circle (0.8);
    \node at (3,2.5) {M.F.};
    
    % Sorgente Fredda
    \draw[thick, mypen, fill=cyan!30] (2,0) rectangle (4,1);
    \node at (3,0.5) {Sorgente $T_f$};
    
    % Frecce
    \draw[->, >=stealth, thick, notered] (3,3.3) -- (3,4) node[midway, right] {$|Q_c|$};
    \draw[->, >=stealth, thick, cyan] (3,1) -- (3,1.7) node[midway, right] {$Q_f$};
    \draw[<-, >=stealth, thick, mypen] (3.8, 2.5) -- (5, 2.5) node[midway, above] {$|L|$};
\end{tikzpicture}
\caption{Schema di principio di una macchina frigorifera.}
\end{figure}

Essendo un ciclo, la variazione di energia interna è nulla, per cui il lavoro scambiato equivale alla somma algebrica dei calori: $|L| = |Q_c| - Q_f$.
L'efficienza di questa macchina è descritta dal \textbf{Coefficiente di Prestazione (COP)}, definito come il rapporto tra l'effetto utile (il calore estratto dalla sorgente fredda) e il costo in termini di lavoro assorbito $|L|$:
\begin{equation}
    \text{COP} = \frac{Q_f}{|L|} = \frac{Q_f}{|Q_c| - Q_f}
\end{equation}

\vspace{0.3cm}
\section{Ciclo di Carnot e Teorema di Carnot}

\begin{figure}[ht]
\centering
\begin{minipage}{0.48\textwidth}
\centering
\begin{tikzpicture}[scale=1.1]
    % Griglia
    \draw[step=1cm, gray!15, very thin] (0,0) grid (4.5,4.5);
    % Assi
    \draw[->, >=stealth, thick, mypen] (0,0) -- (0,4.6) node[left] {$p$};
    \draw[->, >=stealth, thick, mypen] (0,0) -- (4.6,0) node[below] {$V$};
    
    % Curve analitiche del ciclo
    % A->B: Isoterma T2 (p = 4.8/V)
    \draw[thick, blue, domain=1.2:3.0, samples=30] plot (\x, {4.8/\x});
    \node[right, blue] at (2.9, 1.8) {$T_2$};
    
    % B->C: Adiabatica espansione (p = 7.45/V^1.4)
    \draw[thick, notered, domain=3.0:4.0, samples=30] plot (\x, {7.45/(\x^1.4)});
    
    % C->D: Isoterma T1 (p = 2.4/V)
    \draw[thick, blue, domain=1.6:4.0, samples=30] plot (\x, {2.4/\x});
    \node[right, blue] at (3.9, 0.4) {$T_1$};
    
    % D->A: Adiabatica compressione (p = 5.13/V^1.4)
    \draw[thick, notered, domain=1.2:1.6, samples=30] plot (\x, {5.13/(\x^1.4)});
    
    % Riempimento area ciclo (Lavoro L)
    \fill[cyan!15, opacity=0.7, domain=1.2:3.0, variable=\x] (1.2,4.0) -- plot (\x, {4.8/\x}) -- (3.0,1.6) -- (3.0,0.8) -- plot[domain=3.0:1.6] (\x, {2.4/\x}) -- (1.6,1.5) -- plot[domain=1.6:1.2] (\x, {5.13/(\x^1.4)}) -- cycle;
    
    % Vertici
    \filldraw[mypen] (1.2, 4.0) circle (1.8pt) node[above right] {$A$};
    \filldraw[mypen] (3.0, 1.6) circle (1.8pt) node[above right] {$B$};
    \filldraw[mypen] (4.0, 0.6) circle (1.8pt) node[below right] {$C$};
    \filldraw[mypen] (1.6, 1.5) circle (1.8pt) node[below left] {$D$};
    
    \node[mypen, font=\bfseries] at (2.4, 2.0) {$L > 0$};
\end{tikzpicture}
\end{minipage}\hfill
\begin{minipage}{0.50\textwidth}
Il \textbf{Ciclo di Carnot} è un ciclo termodinamico reversibile operante tra due sorgenti a temperatura fissa $T_2 > T_1$. Si compone di quattro fasi alternate:
\begin{itemize}
    \item \textbf{$A \to B$}: Espansione isoterma a temperatura $T_2$ \\
    ($Q_2 = Q_{AB} = L_{AB} > 0$)
    \item \textbf{$B \to C$}: Espansione adiabatica reversibile \\
    ($Q = 0 \implies L_{BC} = -\Delta U > 0$)
    \item \textbf{$C \to D$}: Compressione isoterma a temperatura $T_1$ \\
    ($Q_1 = Q_{CD} = L_{CD} < 0$)
    \item \textbf{$D \to A$}: Compressione adiabatica reversibile \\
    ($Q = 0 \implies L_{DA} = -\Delta U < 0$)
\end{itemize}
\end{minipage}
\caption{Rappresentazione nel piano di Clapeyron del ciclo reversibile di Carnot. L'area racchiusa dal ciclo (azzurra) rappresenta il lavoro utile $L = Q_2 - |Q_1|$ prodotto dalla macchina in ogni ciclo.}
\end{figure}

\vspace{0.3cm}
Essendo il ciclo chiuso, $\Delta U_{\text{tot}} = 0$. Il calore assorbito è $Q_2 = Q_{AB}$ e quello ceduto $Q_1 = Q_{CD}$. Dalle leggi delle adiabatiche si dimostra che $\frac{V_B}{V_A} = \frac{V_C}{V_D}$, per cui il rendimento si semplifica in:
\begin{equation}
    \eta_C = \frac{L}{Q_2} = 1 - \frac{|Q_1|}{Q_2} = 1 - \frac{T_1}{T_2}
\end{equation}

\begin{definizione}{Teorema di Carnot}
Date due sorgenti di calore a temperature $T_1$ e $T_2$ (con $T_2 > T_1$):
\begin{enumerate}
    \item Tutte le macchine termiche \textbf{reversibili} operanti tra le due sorgenti hanno lo stesso rendimento $\eta_C = 1 - T_1/T_2$.
    \item Nessuna macchina termica \textbf{irreversibile} può avere un rendimento maggiore: $\eta_i \le \eta_C$.
\end{enumerate}
\end{definizione}

\vspace{0.3cm}
\begin{figure}[ht]
\centering
\begin{minipage}{0.42\textwidth}
\centering
\begin{tikzpicture}[scale=1.05]
    % Sorgente CALDA T2
    \draw[thick, mypen, fill=notered!20, rounded corners=2pt] (0, 3) rectangle (3, 3.8) node[midway, font=\bfseries] {Sorgente Calda $T_2$};
    
    % Macchina Mc (Carnot in senso inverso - Frigorifero)
    \draw[thick, mypen, fill=white] (0.75, 1.5) circle (0.5) node[font=\bfseries] {$M_C$};
    % Calore ceduto a T2 (freccia verso l'alto da Mc a T2)
    \draw[->, >=stealth, thick, notered] (0.75, 2) -- (0.75, 3) node[left, midway] {$Q_{2C}$};
    % Calore assorbito da T1 (freccia verso l'alto da T1 a Mc)
    \draw[->, >=stealth, thick, cyan] (0.75, 0) -- (0.75, 1) node[left, midway] {$Q_{1C}$};
    % Lavoro assorbito da Mc (freccia entrante in Mc da sinistra)
    \draw[->, >=stealth, thick, mypen] (-0.3, 1.5) node[left] {$L_C$} -- (0.25, 1.5);
    
    % Macchina M (Generica - Motore termico)
    \draw[thick, mypen, fill=white] (2.25, 1.5) circle (0.5) node[font=\bfseries] {$M$};
    % Calore assorbito da T2 (freccia verso il basso da T2 a M)
    \draw[->, >=stealth, thick, notered] (2.25, 3) -- (2.25, 2) node[right, midway] {$Q_2$};
    % Calore ceduto a T1 (freccia verso il basso da M a T1)
    \draw[->, >=stealth, thick, cyan] (2.25, 1) -- (2.25, 0) node[right, midway] {$Q_1$};
    % Lavoro prodotto da M (freccia uscente da M verso destra)
    \draw[->, >=stealth, thick, mypen] (2.75, 1.5) -- (3.3, 1.5) node[right] {$L$};
    
    % Sorgente FREDDA T1
    \draw[thick, mypen, fill=cyan!20, rounded corners=2pt] (0, -0.8) rectangle (3, 0) node[midway, font=\bfseries] {Sorgente Fredda $T_1$};
\end{tikzpicture}
\end{minipage}\hfill
\begin{minipage}{0.54\textwidth}
\textbf{Dimostrazione del Teorema di Carnot}:\\
Si accoppia una macchina termica generica $M$ a una macchina reversibile di Carnot $M_C$ operante in senso inverso (come frigorifero). La macchina $M$ produce un lavoro $L > 0$ estraendo calore dalla sorgente $T_2$, mentre $M_C$ assorbe lavoro $L_C < 0$ per rimettere calore nella medesima sorgente.

Regolando le dimensioni dei cicli in modo che il saldo termico con la sorgente calda sia nullo ($Q_2 + Q_{2C} = 0$), il sistema combinato $M + M_C$ scambia calore unicamente con la sorgente fredda $T_1$.
\end{minipage}
\caption{Schema di accoppiamento per la dimostrazione del Teorema di Carnot via assurdo secondo l'enunciato di Kelvin-Planck.}
\end{figure}

\vspace{0.3cm}
Per il Secondo Principio (enunciato di Kelvin-Planck), il lavoro totale prodotto non può essere positivo (non si può estrarre calore da un'unica sorgente $T_1$ e produrre lavoro netto). Quindi $L_{tot} \le 0 \implies L - |L_C| \le 0$. Da questo limite termodinamico deriva la disuguaglianza sui rendimenti:
\begin{itemize}
    \item Se $M$ è \textbf{irreversibile}: si dimostra che $\eta \le \eta_C$.
    \item Se $M$ è \textbf{reversibile}: il ciclo $M + M_C$ è reversibile e può essere invertito, portando a dedurre che $\eta = \eta_C$.
\end{itemize}

\vspace{0.3cm}
\noindent $\implies$ La macchina reversibile è la più efficiente ($Q \leftrightarrow L$) perché, a parità di calore assorbito, cede la minima quantità di calore alla sorgente fredda.

\begin{osservazione}{Irreversibilità e Moto Perpetuo}
Il primo principio non fa distinzioni tra processi reversibili e irreversibili, tuttavia l'esperienza quotidiana mostra che i processi spontanei (come il calore che fluisce da un corpo caldo a uno freddo, o l'energia dissipata per attrito) sono sempre irreversibili. Questo stabilisce una preferenza temporale (la \textit{freccia del tempo}).

Il \textbf{Secondo Principio} codifica queste osservazioni, vietando:
\begin{itemize}
    \item \textbf{Macchine anti-Clausius}: frigoriferi che estraggono calore da una sorgente fredda verso una calda senza richiedere alcun lavoro esterno in ingresso.
    \item \textbf{Moto perpetuo di seconda specie (anti-Kelvin-Planck)}: macchine termiche in grado di trasformare ciclicamente calore in lavoro scambiando con un'unica sorgente termica.
\end{itemize}
Di conseguenza, l'unica possibilità per un dispositivo che compie un ciclo interagendo con \textbf{una sola sorgente termica (ciclo monotermo)} è quella di subire lavoro per trasformarlo integralmente in calore ($L \le 0$ e $Q \le 0$).
\end{osservazione}

\vspace{0.5cm}
\subsection{Esempio: Ciclo monotermo per un gas perfetto}
Consideriamo un ciclo termodinamico compiuto da $n$ moli di un gas perfetto che scambia calore con una sola sorgente a temperatura $T_1$. Il ciclo è costituito da tre trasformazioni:
\begin{itemize}
    \item \textbf{A $\to$ B (Espansione isoterma reversibile a $T_1$)}: Il gas si espande da $V_A$ a $V_B > V_A$ scambiando calore con l'unico termostato a temperatura $T_1$. Il lavoro compiuto è:
    \begin{equation}
        L_{AB} = n R T_1 \ln\left(\frac{V_B}{V_A}\right)
    \end{equation}
    \item \textbf{B $\to$ C (Compressione adiabatica reversibile)}: Il gas viene compresso fino al volume iniziale $V_C = V_A$, portando la temperatura a $T_2 \neq T_1$. Senza scambio di calore ($Q_{BC}=0$), il lavoro compiuto è pari all'opposto della variazione di energia interna:
    \begin{equation}
        L_{BC} = -\Delta U_{BC} = -n c_v (T_2 - T_1)
    \end{equation}
    \item \textbf{C $\to$ A (Raffreddamento isocoro irreversibile)}: Mantenendo il volume costante a $V_A$, il gas viene rimescolato e riportato alla temperatura iniziale $T_1$ cedendo calore. Essendo isocora, non si compie lavoro: $L_{CA} = 0$. (Nota: l'accoppiamento termico tra una temperatura $T_2$ e l'unico termostato a $T_1$ rende la fase irreversibile).
\end{itemize}

\begin{figure}[ht]
\centering
\begin{minipage}{0.48\textwidth}
\centering
\begin{tikzpicture}[scale=1.1]
    % Griglia
    \draw[step=1cm, gray!15, very thin] (0,0) grid (3.8,3.2);
    % Assi
    \draw[->, >=stealth, thick, mypen] (0,0) -- (0,3.3) node[left] {$p$};
    \draw[->, >=stealth, thick, mypen] (0,0) -- (3.8,0) node[below] {$V$};
    
    % Vertici analitici (VA = VC = 1.0, VB = 3.0)
    % Isoterma T1: p = 1.5/V -> A(1.0, 1.5), B(3.0, 0.5)
    % Adiabatica T1->T2: p = 2.33 / V^1.4 -> B(3.0, 0.5), C(1.0, 2.33)
    
    % A->B Isoterma
    \draw[thick, blue, domain=1.0:3.0, samples=30] plot (\x, {1.5/\x});
    % B->C Adiabatica
    \draw[thick, notered, domain=1.0:3.0, samples=30] plot (\x, {2.33/(\x^1.4)});
    % C->A Isocora
    \draw[thick, mygreen, ->, >=stealth] (1.0, 2.33) -- (1.0, 1.5);
    
    % Frecce di direzione sulle curve
    \draw[->, >=stealth, thick, blue] (1.8, {1.5/1.8}) -- (1.9, {1.5/1.9});
    \draw[<-, >=stealth, thick, notered] (1.7, {2.33/(1.7^1.4)}) -- (1.8, {2.33/(1.8^1.4)});
    
    % Punti
    \filldraw[mypen] (1.0, 1.5) circle (1.6pt) node[left] {$A$};
    \filldraw[mypen] (3.0, 0.5) circle (1.6pt) node[right] {$B$};
    \filldraw[mypen] (1.0, 2.33) circle (1.6pt) node[left] {$C$};
    
    \node[below, blue] at (2.0, 0.7) {\small Isoterma ($T_1$)};
    \node[above right, notered] at (1.6, 1.4) {\small Adiabatica};
    \node[left, mygreen] at (1.0, 1.9) {\small Isocora};
\end{tikzpicture}
\end{minipage}\hfill
\begin{minipage}{0.48\textwidth}
Il lavoro totale del ciclo è la somma dei lavori delle singole trasformazioni:
\begin{equation}
    L_{\text{tot}} = L_{AB} + L_{BC} + L_{CA} = n R T_1 \ln\left(\frac{V_B}{V_A}\right) - n c_v (T_2 - T_1)
\end{equation}
Essendo $L_{CA} = 0$, il segno del lavoro totale dipende dal confronto tra il lavoro positivo di espansione isoterma e quello negativo assorbito nella compressione adiabatica.
\end{minipage}
\caption{Ciclo monotermo costituito da un'isoterma reversibile ($A \to B$), un'adiabatica reversibile ($B \to C$) e un'isocora irreversibile ($C \to A$).}
\end{figure}

Per l'adiabatica reversibile $B \to C$ vale la relazione tra temperature e volumi:
\begin{equation}
    \frac{T_2}{T_1} = \left(\frac{V_B}{V_C}\right)^{\gamma-1} = \left(\frac{V_B}{V_A}\right)^{\gamma-1}
\end{equation}
Passando ai logaritmi e ricordando che $\gamma - 1 = \frac{c_p - c_v}{c_v} = \frac{R}{c_v}$, si ottiene:
\begin{equation}
    \ln\left(\frac{T_2}{T_1}\right) = (\gamma - 1) \ln\left(\frac{V_B}{V_A}\right) = \frac{R}{c_v} \ln\left(\frac{V_B}{V_A}\right) \implies n R \ln\left(\frac{V_B}{V_A}\right) = n c_v \ln\left(\frac{T_2}{T_1}\right)
\end{equation}
Sostituendo questa espressione nel lavoro totale $L_{\text{tot}}$:
\begin{align}
    L_{\text{tot}} &= n c_v T_1 \ln\left(\frac{T_2}{T_1}\right) - n c_v (T_2 - T_1) \nonumber \\
    &= n c_v T_1 \left[ \ln\left(\frac{T_2}{T_1}\right) - \left(\frac{T_2}{T_1} - 1\right) \right]
\end{align}
Per la disuguaglianza fondamentale dei logaritmi, per qualsiasi $x > 0$ con $x \neq 1$ vale sempre $\ln(x) < x - 1$, ossia $\ln(x) - (x-1) < 0$. Posto $x = \frac{T_2}{T_1}$, risulta rigorosamente:
\begin{equation}
    \boxed{L_{\text{tot}} < 0}
\end{equation}
Questo dimostra matematicamente che in un ciclo monotermo irreversibile il lavoro totale prodotto è necessariamente negativo: il sistema subisce lavoro netto dall'ambiente per dissiparlo in calore, in perfetto accordo con il Secondo Principio della Termodinamica (se infatti fosse reversibile o producesse lavoro positivo $L > 0$, violerebbe l'enunciato di Kelvin-Planck).




\subsection{Equivalenza degli Enunciati di Kelvin-Planck e di Clausius}
Per dimostrare l'equivalenza logica dei due enunciati ($K \iff C$), è sufficiente far vedere che se uno dei due è falso, allora necessariamente lo è anche l'altro. Procediamo con due dimostrazioni per assurdo:

\vspace{0.3cm}
\noindent \textbf{1. Dimostrazione $K \implies C$ (se vale Kelvin-Planck, deve valere Clausius)}: \\
Supponiamo per assurdo che l'enunciato di Kelvin-Planck sia vero e quello di Clausius sia falso.
Esistono allora una \textbf{macchina anti-Clausius} $C$ (in grado di trasferire calore $Q_1 > 0$ da una sorgente fredda $T_1$ a una calda $T_2$ senza apporto di lavoro esterno) e una normale macchina termica di Carnot $M_C$ operante tra le stesse due sorgenti.

\begin{figure}[ht]
\centering
\begin{minipage}{0.40\textwidth}
\centering
\begin{tikzpicture}[scale=0.95]
    % T2
    \draw[thick, mypen, fill=notered!20, rounded corners=2pt] (0, 3) rectangle (3, 3.8) node[midway, font=\bfseries] {Sorgente $T_2 > T_1$};
    % T1
    \draw[thick, mypen, fill=cyan!20, rounded corners=2pt] (0, -0.8) rectangle (3, 0) node[midway, font=\bfseries] {Sorgente $T_1$};
    
    % Macchina anti-Clausius C
    \draw[thick, mypen, fill=white] (0.75, 1.5) circle (0.5) node[font=\bfseries] {$C$};
    \draw[->, >=stealth, thick, cyan] (0.75, 0) -- (0.75, 1) node[midway, left] {$Q_1$};
    \draw[->, >=stealth, thick, notered] (0.75, 2) -- (0.75, 3) node[midway, left] {$Q_1$};
    
    % Macchina termica di Carnot Mc (Motore normale)
    \draw[thick, mypen, fill=white] (2.25, 1.5) circle (0.5) node[font=\bfseries] {$M_C$};
    \draw[->, >=stealth, thick, notered] (2.25, 3) -- (2.25, 2) node[midway, right] {$Q_2$};
    \draw[->, >=stealth, thick, cyan] (2.25, 1) -- (2.25, 0) node[midway, right] {$Q_1$};
    \draw[->, >=stealth, thick, mypen] (2.75, 1.5) -- (3.7, 1.5) node[right, font=\bfseries] {$L$};
\end{tikzpicture}
\end{minipage}\hfill
\begin{minipage}{0.56\textwidth}
Accoppiando le due macchine in modo che scambino la stessa quantità di calore $Q_1$ con la sorgente fredda $T_1$, analizziamo i bilanci termici per il sistema combinato:
\begin{itemize}
    \item La sorgente fredda $T_1$ cede calore $Q_1$ alla macchina $C$ e riceve la medesima quantità $Q_1$ dalla macchina $M_C \implies Q_{\text{netto}, T_1} = 0$.
    \item La sorgente calda $T_2$ riceve calore $Q_1$ dalla macchina $C$ e cede calore $Q_2 > Q_1$ alla macchina $M_C \implies Q_{\text{netto}, T_2} = Q_2 - Q_1 > 0$.
    \item Il lavoro totale prodotto dal sistema combinato, essendo $\Delta U_{\text{tot}} = 0$ per un ciclo, vale:
    \begin{equation}
        L = Q_2 - Q_1 > 0
    \end{equation}
\end{itemize}
\end{minipage}
\caption{Violazione dell'enunciato di Kelvin-Planck indotta dall'accoppiamento con una ipotetica macchina anti-Clausius $C$.}
\end{figure}

Il sistema combinato si comporta quindi come una macchina termica che compie un ciclo interagendo con \textbf{un'unica sorgente di calore} ($T_2$, dato che la sorgente $T_1$ non subisce variazioni nette di calore) e produce un lavoro interamente positivo ($L > 0$). Ciò costituisce una diretta violazione dell'enunciato di Kelvin-Planck, contraddicendo l'ipotesi iniziale. Dunque:
\begin{equation}
    K \text{ vero} \implies C \text{ vero}
\end{equation}

\vspace{0.3cm}
\noindent \textbf{2. Dimostrazione $C \implies K$ (se vale Clausius, deve valere Kelvin-Planck)}: \\
Supponiamo per assurdo che l'enunciato di Clausius sia vero e quello di Kelvin-Planck sia falso.
Esiste allora una \textbf{macchina anti-Kelvin-Planck} $K$ capace di prelevare calore $Q > 0$ da una sola sorgente $T_1$ e trasformarlo integralmente in lavoro meccanico $L = Q$.

\begin{figure}[ht]
\centering
\begin{minipage}{0.46\textwidth}
\centering
\begin{tikzpicture}[scale=0.9]
    % T1
    \draw[thick, mypen, fill=cyan!20, rounded corners=2pt] (0, 3) rectangle (1.8, 4) node[midway, font=\bfseries] {Sorgente $T_1$};
    
    % Macchina anti-Kelvin K
    \draw[thick, mypen, fill=white] (0.9, 1.5) circle (0.5) node[font=\bfseries] {$K$};
    \draw[->, >=stealth, thick, cyan] (0.9, 3) -- (0.9, 2) node[midway, right] {$Q$};
    \draw[->, >=stealth, thick, mypen] (1.4, 1.5) -- (2.6, 1.5) node[midway, above] {$L$};
    
    % Attrito / Dissipatore
    \draw[thick, dashed, mypen, fill=orange!15, rounded corners=2pt] (2.6, 0.8) rectangle (4.2, 2.2) node[midway, align=center, font=\small] {Dissipatore\\(attrito)};
    \draw[->, >=stealth, thick, notered] (4.2, 1.5) -- (5.2, 1.5) node[midway, above] {$Q$};
    
    % T2
    \draw[thick, mypen, fill=notered!20, rounded corners=2pt] (5.2, 0.8) rectangle (7, 2.2) node[midway, align=center, font=\bfseries] {Sorgente\\$T_2 > T_1$};
\end{tikzpicture}
\end{minipage}\hfill
\begin{minipage}{0.50\textwidth}
Utilizziamo ora il lavoro $L$ prodotto dalla macchina $K$ per generare calore tramite attrito (o resistenza elettrica) all'interno di un corpo o di un bagno termico a temperatura superiore $T_2 > T_1$.
\begin{itemize}
    \item La sorgente fredda $T_1$ cede calore netto $Q$.
    \item Il lavoro meccanico $L = Q$ viene dissipato integralmente in calore all'interno della sorgente calda $T_2$.
    \item Nessun altro lavoro netto entra o esce dall'ambiente esterno.
\end{itemize}
\end{minipage}
\caption{Violazione dell'enunciato di Clausius indotta dall'accoppiamento con una ipotetica macchina anti-Kelvin-Planck $K$.}
\end{figure}

Il risultato complessivo del sistema accoppiato è il puro trasferimento di una quantità di calore $Q > 0$ da un corpo freddo ($T_1$) a un corpo caldo ($T_2$) senza alcun apporto netto di lavoro dall'esterno. Ciò viola l'enunciato di Clausius, contraddicendo l'ipotesi iniziale. Dunque:
\begin{equation}
    C \text{ vero} \implies K \text{ vero}
\end{equation}
Rimane così provata la totale equivalenza fisica tra le due formulazioni del Secondo Principio ($K \iff C$).

\section{Temperatura Termodinamica Assoluta}
Dal Teorema di Carnot sappiamo che il rendimento di una macchina reversibile operante tra due sorgenti dipende unicamente dalle loro temperature. Da ciò si deduce che il rapporto tra i calori scambiati in un ciclo di Carnot dipende solo dalle temperature empiriche $t_1, t_2$ delle due sorgenti:
\begin{equation*}
    \frac{|Q_1|}{|Q_2|} = f(t_1, t_2)
\end{equation*}

Accoppiando opportunamente due macchine di Carnot operanti rispettivamente tra $(t_2, t_0)$ e $(t_0, t_1)$, Lord Kelvin dimostrò che tale funzione universale deve essere necessariamente fattorizzabile nella forma $f(t_1, t_2) = \frac{\theta(t_1)}{\theta(t_2)}$.

Si può quindi definire la \textbf{temperatura termodinamica assoluta} come $T \equiv \theta(t)$. Questa nuova scala è intrinseca e del tutto \textbf{indipendente dal fluido termodinamico} (al contrario dei termometri reali). Fortunatamente, si dimostra che essa coincide con la temperatura definita tramite il termometro a gas perfetto:
\begin{equation}
    \boxed{\frac{|Q_1|}{|Q_2|} = \frac{T_1}{T_2}} \implies \eta_C = 1 - \frac{T_1}{T_2}
\end{equation}

\vspace{0.3cm}
\begin{osservazione}{Conseguenze sulla macchina termica}
Poiché le temperature termodinamiche $T$ sono per definizione strettamente positive ($T > 0$), si ha sempre $|Q_1| > 0$. Pertanto \textbf{è impossibile} che in un ciclo termico il calore ceduto alla sorgente fredda sia nullo. Questo fornisce un fondamento matematico rigoroso all'enunciato di Kelvin-Planck.
\end{osservazione}
\section{Entropia e Disuguaglianza di Clausius}

La \textbf{disuguaglianza di Clausius} fornisce una formulazione quantitativa del Secondo Principio, permettendo di distinguere analiticamente i cicli reversibili da quelli irreversibili.
Partendo da una macchina termica operante tra due sorgenti, si ha:
\begin{align}
    \text{Reversibile:} &\quad \frac{Q_1}{T_1} + \frac{Q_2}{T_2} = 0 \quad (\eta = \eta_C) \\
    \text{Irreversibile:} &\quad \frac{Q_1}{T_1} + \frac{Q_2}{T_2} < 0 \quad (\eta < \eta_C)
\end{align}
Generalizzando a una macchina che scambia calore con $n$ sorgenti termiche (o integrando per trasformazioni continue), si ottiene il teorema di Clausius:
\begin{equation}
    \boxed{\oint \frac{\delta Q}{T} \le 0}
\end{equation}
dove l'uguaglianza vale se e solo se l'intero ciclo è reversibile.

\vspace{0.3cm}
\subsection{Definizione di Entropia}
Se consideriamo una trasformazione \textbf{ciclica e reversibile}, l'integrale nullo implica che:
\begin{equation*}
    \oint_{\text{rev}} \frac{\delta Q}{T} = \int_{A}^{B} \left(\frac{\delta Q}{T}\right)_{\text{rev}} + \int_{B}^{A} \left(\frac{\delta Q}{T}\right)_{\text{rev}} = 0 \implies \int_{A}^{B} \left(\frac{\delta Q}{T}\right)_{\text{rev}} = \int_{A}^{B} \left(\frac{\delta Q}{T}\right)_{\text{rev, cammino 2}}
\end{equation*}
L'integrale lungo un cammino reversibile \textbf{non dipende dal percorso} ma solo dagli stati iniziale e finale. L'integrando $\frac{\delta Q}{T}$ risulta quindi un differenziale esatto. Questa proprietà definisce una nuova funzione di stato, chiamata \textbf{Entropia ($S$)}:
\begin{equation}
    \boxed{ \Delta S = S(B) - S(A) = \int_A^B \left( \frac{\delta Q}{T} \right)_{\text{rev}} }
\end{equation}
(L'unità di misura è J/K). 
Per un processo \textbf{irreversibile} da $A$ a $B$, si chiude il ciclo tornando ad $A$ con un processo reversibile ausiliario:
\begin{equation*}
    \oint \frac{\delta Q}{T} = \int_{A}^{B} \left( \frac{\delta Q}{T} \right)_{\text{irr}} + \int_{B}^{A} \left( \frac{\delta Q}{T} \right)_{\text{rev}} < 0
\end{equation*}
Da cui si ottiene: $\int_{A}^{B} \left( \frac{\delta Q}{T} \right)_{\text{irr}} < S(B) - S(A)$. In generale:
\begin{equation}
    \boxed{ \Delta S \ge \int_{A}^{B} \frac{\delta Q}{T} } \quad \text{(Disuguaglianza per trasformazioni aperte)}
\end{equation}

\begin{figure}[ht]
\centering
\begin{tikzpicture}
% Assi
\draw[->] (0,0) -- (3,0) node[right] {$V$};
\draw[->] (0,0) -- (0,3) node[above] {$P$};
\fill (0.5, 2) circle (2pt) node[above] {$A$};
\fill (2.5, 0.5) circle (2pt) node[right] {$B$};
\draw[->, dashed] (0.5, 2) to[out=-20, in=120] (1.5, 1.2) node[above] {I} to[out=-60, in=170] (2.5, 0.5);
\draw[<-, dashed] (0.5, 2) to[out=-80, in=160] (1.5, 0.8) node[below] {II} to[out=-20, in=180] (2.5, 0.5);
\end{tikzpicture}
\end{figure}

Per un processo \textbf{irreversibile} da $A$ a $B$, confrontando il percorso irreversibile con uno reversibile:
\begin{align*}
    \int_{A}^{B} \left( \frac{\delta Q}{T} \right)_{\text{irr}} &< S(B) - S(A) \\
    \int_{A}^{B} \frac{\delta Q}{T} &\le S(B) - S(A)
\end{align*}

\begin{definizione}{Formulazione Entropica del Secondo Principio della Termodinamica}
Per qualsiasi trasformazione da uno stato $A$ ad uno stato $B$, l'integrale di Clausius è sempre minore o uguale alla variazione di entropia del sistema. L'uguaglianza vale solo per trasformazioni reversibili.
\end{definizione}

\begin{definizione}{Variazione di Entropia per Sorgenti e Termostati}
Un termostato (o sorgente termica) è per definizione un serbatoio di calore a temperatura costante $T_0$ con capacità termica infinita ($C \to \infty$). Poiché lo scambio termico avviene idealmente in modo isotermo, senza gradienti interni di temperatura, la trasformazione per la sorgente si considera sempre reversibile. La sua variazione di entropia quando scambia un calore $Q_{\text{sorgente}}$ è calcolata elementarmente come:
\begin{equation}
    \boxed{ \Delta S_{\text{sorgente}} = \int \frac{\delta Q_{\text{sorgente}}}{T_0} = \frac{Q_{\text{sorgente}}}{T_0} }
\end{equation}
Se una sorgente cede una quantità di calore $Q > 0$ ad un sistema operativo, rispetto al termostato il calore è uscente ($Q_{\text{sorgente}} = -Q$), con conseguente diminuzione di entropia $\Delta S_{\text{sorgente}} = -\frac{Q}{T_0}$.
\end{definizione}

\vspace{0.3cm}
\subsection{Proprietà dell'Entropia}
\begin{itemize}
    \item \textbf{Additività}: un sistema composto da sottosistemi debolmente interagenti ha l'entropia pari alla somma delle singole entropie dei sottosistemi.
    \item \textbf{Estensività}: l'entropia totale di un sistema composto da più sotto-sistemi non interagenti a lungo raggio è la somma delle entropie delle singole parti. Essendo $\delta Q = \sum_i \delta Q_i$ e $T$ la medesima in condizioni di equilibrio termico, si ha:
    \begin{equation}
        \Delta S = \int_A^B \frac{\delta Q}{T} = \sum_i \int_A^B \frac{\delta Q_i}{T} = \sum_i \Delta S_i
    \end{equation}
    
    \item \textbf{Principio di Non Diminuzione dell'Entropia}: in un \textbf{sistema isolato}, non vi è alcuno scambio termico con l'esterno ($\delta Q = 0$). Pertanto:
    \begin{equation}
        \boxed{ \Delta S_{\text{isolato}} \ge 0 }
    \end{equation}
    
    \item \textbf{Entropia dell'Universo}: L'universo termodinamico, definito come l'insieme di Sistema e Ambiente, è per definizione un sistema isolato. Di conseguenza:
    \begin{equation*}
        \Delta S_{\text{univ}} = \Delta S_{\text{sist}} + \Delta S_{\text{amb}} \ge 0
    \end{equation*}
\end{itemize}

\begin{figure}[ht]
\centering
\begin{tikzpicture}
    \draw[->, >=stealth, thick, mypen] (0,0) -- (6,0) node[right] {Tempo $t$};
    \draw[->, >=stealth, thick, mypen] (0,0) -- (0,3.5) node[above] {Entropia $S$};
    \draw[dashed, thick, notered] (0, 3) node[left] {$S_{max}$} -- (5.8, 3);
    \draw[thick, blue] (0.2, 0.5) to[out=60, in=185] (1.5, 2.0) to[out=5, in=190] (3.0, 2.7) to[out=10, in=180] (5.5, 3.0);
    \node[align=center] at (3, 1) {\small I sistemi isolati evolvono spontaneamente\\ \small verso lo stato di massima entropia};
\end{tikzpicture}
\caption{Andamento nel tempo dell'entropia di un sistema isolato.}
\end{figure}

\subsection{Calcolo della Variazione di Entropia}
Poiché l'entropia è una funzione di stato, la sua variazione $\Delta S$ lungo una trasformazione \textbf{irreversibile} si calcola immaginando un cammino \textbf{reversibile arbitrario} che congiunge gli stessi stati $A$ e $B$.
\begin{itemize}
    \item \textbf{Espansione libera di un gas perfetto}: $\Delta S_{\text{gas}} = n R \ln(V_B/V_A) > 0$.
    \item \textbf{Trasformazione Isoterma}: $\Delta S = Q_{AB}/T = n R \ln(V_B/V_A)$.
    \item \textbf{Trasformazione Adiabatica Reversibile}: $\Delta S = 0$ (trasformazione isoentropica).
    \item \textbf{Trasformazione generica di un gas perfetto}: 
    \begin{align}
        \Delta S &= n c_v \ln\left(\frac{T_B}{T_A}\right) + n R \ln\left(\frac{V_B}{V_A}\right) = n c_p \ln\left(\frac{T_B}{T_A}\right) - n R \ln\left(\frac{p_B}{p_A}\right) \\
        &= n c_v \ln\left( \frac{T_B V_B^{\gamma-1}}{T_A V_A^{\gamma-1}} \right) = n c_v \ln\left( \frac{p_B V_B^\gamma}{p_A V_A^\gamma} \right)
    \end{align}
    \item \textbf{Passaggi di stato (a $T$ e $p$ costanti)}: $\Delta S = \frac{m \lambda}{T}$.
    \item \textbf{Solidi e liquidi incomprimibili}: $\Delta S = m c \ln\left(\frac{T_B}{T_A}\right)$.
\end{itemize}

\vspace{0.3cm}
Inoltre, per un gas perfetto è possibile definire la funzione di stato \textbf{Entropia Assoluta} (a meno di una costante additiva di riferimento $S_0$):
\begin{equation}
    S_{\text{gas}}(p, V) = n c_v \ln\left(p V^\gamma\right) + \text{costante}
\end{equation}
Mentre per un solido o liquido a calore specifico $C$ costante: $S(T) = C \ln T + \text{costante}$.

\begin{osservazione}{Esempio: Calcolo dell'Entropia in una Espansione Libera ed Equivalenza del Percorso Ausiliario}
Consideriamo una mole ($n = 1\text{ mol}$) di gas perfetto che si espande liberamente nel vuoto raddoppiando il proprio volume ($V_B = 2V_A$) all'interno di un contenitore rigido adiabatico ($Q=0, L=0$). Come già dimostrato, la temperatura finale coincide con quella iniziale ($T_B = T_A = T_0$). Poiché la trasformazione reale è irreversibile, per calcolare la variazione di entropia del gas $\Delta S_{\text{gas}}$ dobbiamo connettere gli stati di equilibrio $A$ e $B$ tramite una trasformazione reversibile ausiliaria, ad esempio un'\textbf{isoterma reversibile}:
\begin{equation}
    \Delta S_{\text{gas}} = \int_A^B \left(\frac{\delta Q}{T}\right)_{\text{rev}} = \frac{1}{T_0} \int_{V_A}^{2V_A} p \, dV = n R \ln\left(\frac{2V_A}{V_A}\right) = R \ln 2 \simeq 5.76 \, \frac{\text{J}}{\text{K}} > 0
\end{equation}
Per l'ambiente, nella trasformazione reale non vi è stato alcuno scambio termico ($Q_{\text{amb}} = 0 \implies \Delta S_{\text{amb}} = 0$). Di conseguenza, per l'universo:
\begin{equation}
    \Delta S_{\text{univ}} = \Delta S_{\text{gas}} + \Delta S_{\text{amb}} = R \ln 2 + 0 > 0
\end{equation}
Se volessimo ripristinare il gas allo stato iniziale comprimendolo reversibilmente a contatto con un termostato ausiliario a temperatura $T_0$, il lavoro esterno compiuto sul gas verrebbe integralmente ceduto come calore alla sorgente ($Q_{\text{sorgente}} = +T_0 \Delta S_{\text{gas}}$), incrementando permanentemente l'entropia del termostato di $+R \ln 2$. Ciò conferma l'impossibilità di cancellare dall'universo l'irreversibilità generata dall'espansione libera.
\end{osservazione}

\vspace{0.3cm}
\subsection{Il Diagramma Temperatura-Entropia (T-S)}
Il diagramma \textbf{Temperatura-Entropia (T-S)} rappresenta un fondamentale strumento analitico in termodinamica. Poiché per una trasformazione reversibile vale la definizione infinitesima $\delta Q_{\text{rev}} = T\,dS$, l'area sottesa alla curva di una trasformazione in tale piano rappresenta geometricamente il \textbf{calore scambiato}:
\begin{equation}
    \boxed{ Q_{\text{rev}} = \int_A^B T\,dS }
\end{equation}
Per un ciclo termodinamico reversibile chiuso, essendo la variazione di energia interna nulla ($\oint dU = 0$), l'area interna alla linea chiusa del ciclo rappresenta esattamente il \textbf{lavoro netto} compiuto: $L = \oint \delta Q = \oint T\,dS$.

\vspace{0.3cm}
\subsubsection{Rappresentazione delle Trasformazioni e Confronto delle Pendenze}
Nel piano T-S, le quattro trasformazioni fondamentali assumono forme geometriche caratteristiche:
\begin{itemize}
    \item \textbf{Isoterma ($T = \text{cost}$)}: è rappresentata da una segmento rettilineo \textbf{orizzontale}.
    \item \textbf{Adiabatica Reversibile / Isoentropica ($S = \text{cost}$)}: è rappresentata da un segmento rettilineo \textbf{verticale} ($\Delta S = 0$).
    \item \textbf{Isocora e Isobara}: sono rappresentate da curve esponenziali crescenti della forma $T(S) = T_0 e^{\Delta S / n c}$.
\end{itemize}

\begin{figure}[ht]
\centering
\begin{minipage}{0.46\textwidth}
\centering
\begin{tikzpicture}[scale=1.05]
    % Griglia
    \draw[step=1cm, gray!15, very thin] (0,0) grid (4.2,4.2);
    % Assi
    \draw[->, >=stealth, thick, mypen] (0,0) -- (0,4.4) node[left] {$T$};
    \draw[->, >=stealth, thick, mypen] (0,0) -- (4.4,0) node[below] {$S$};
    
    % Punto iniziale comune A(1.2, 1.2)
    % Isoterma (T = 1.2)
    \draw[thick, blue] (1.2, 1.2) -- (4.0, 1.2) node[right] {\small Isoterma};
    
    % Adiabatica (S = 1.2)
    \draw[thick, notered] (1.2, 1.2) -- (1.2, 4.0) node[above] {\small Adiabatica};
    
    % Isocora: T = 1.2 * exp((S-1.2)/1.2) -> ripida
    \draw[thick, mygreen, domain=1.2:2.65, samples=30] plot (\x, {1.2*exp((\x-1.2)/1.2)}) node[above] {\small Isocora};
    
    % Isobara: T = 1.2 * exp((S-1.2)/2.0) -> meno ripida
    \draw[thick, orange, domain=1.2:3.6, samples=30] plot (\x, {1.2*exp((\x-1.2)/2.0)}) node[right] {\small Isobara};
    
    % Punto A
    \filldraw[mypen] (1.2, 1.2) circle (1.8pt) node[below left] {$A(S_0, T_0)$};
    \draw[dashed, gray] (1.2,0) node[below] {$S_0$} -- (1.2,1.2) -- (0,1.2) node[left] {$T_0$};
\end{tikzpicture}
\end{minipage}\hfill
\begin{minipage}{0.50\textwidth}
\textbf{Confronto analitico delle pendenze}:\\
Dalla relazione differenziale del calore scambiato $\delta Q = n c\,dT = T\,dS$, si ricava la pendenza delle curve nel diagramma T-S:
\begin{equation}
    \left( \frac{dT}{dS} \right) = \frac{T}{n c}
\end{equation}
Confrontando una trasformazione isocora ($c = c_v$) con una isobara ($c = c_p$) passanti per lo stesso stato $(S_0, T_0)$:
\begin{equation}
    \left( \frac{dT}{dS} \right)_v = \frac{T}{n c_v} > \frac{T}{n c_p} = \left( \frac{dT}{dS} \right)_p
\end{equation}
Poiché $c_p = c_v + R > c_v$, nel piano T-S \textbf{la curva isocora cresce sempre con una pendenza maggiore rispetto alla curva isobara}.
\end{minipage}
\caption{Confronto nel piano Temperatura-Entropia tra le quattro trasformazioni termodinamiche elementari a partire dal medesimo stato $A$.}
\end{figure}

\vspace{0.3cm}
\subsubsection{Il Ciclo di Carnot nel Piano T-S}
La potenza concettuale del diagramma T-S si manifesta pienamente nell'analisi del \textbf{Ciclo di Carnot}. Poiché è composto esclusivamente da due isoterme (a $T_2$ e $T_1$) e due adiabatiche reversibili (isoentropiche a $S_B$ e $S_A$), nel piano T-S il ciclo di Carnot assume la geometria elementare di un \textbf{rettangolo}.

\begin{figure}[ht]
\centering
\begin{minipage}{0.44\textwidth}
\centering
\begin{tikzpicture}[scale=1.1]
    % Griglia
    \draw[step=1cm, gray!15, very thin] (0,0) grid (4.2,3.8);
    % Assi
    \draw[->, >=stealth, thick, mypen] (0,0) -- (0,4.0) node[left] {$T$};
    \draw[->, >=stealth, thick, mypen] (0,0) -- (4.2,0) node[below] {$S$};
    
    % Rettangolo di Carnot tra S_A=1.2, S_B=3.4 e T1=1.0, T2=3.0
    \fill[cyan!15, opacity=0.8] (1.2, 1.0) rectangle (3.4, 3.0);
    
    % Ciclo
    \draw[thick, blue, ->, >=stealth] (1.2, 3.0) -- (2.4, 3.0);
    \draw[thick, blue] (2.4, 3.0) -- (3.4, 3.0) node[above, midway] {$Q_2 > 0$};
    
    \draw[thick, notered, ->, >=stealth] (3.4, 3.0) -- (3.4, 2.0);
    \draw[thick, notered] (3.4, 2.0) -- (3.4, 1.0);
    
    \draw[thick, blue, ->, >=stealth] (3.4, 1.0) -- (2.3, 1.0);
    \draw[thick, blue] (2.3, 1.0) -- (1.2, 1.0) node[below, midway] {$Q_1 < 0$};
    
    \draw[thick, notered, ->, >=stealth] (1.2, 1.0) -- (1.2, 2.0);
    \draw[thick, notered] (1.2, 2.0) -- (1.2, 3.0);
    
    % Vertici
    \filldraw[mypen] (1.2, 3.0) circle (1.6pt) node[above left] {$A$};
    \filldraw[mypen] (3.4, 3.0) circle (1.6pt) node[above right] {$B$};
    \filldraw[mypen] (3.4, 1.0) circle (1.6pt) node[below right] {$C$};
    \filldraw[mypen] (1.2, 1.0) circle (1.6pt) node[below left] {$D$};
    
    % Coordinate
    \draw[dashed, gray] (1.2, 0) node[below] {$S_A$} -- (1.2, 1.0);
    \draw[dashed, gray] (3.4, 0) node[below] {$S_B$} -- (3.4, 1.0);
    \draw[dashed, gray] (0, 3.0) node[left] {$T_2$} -- (1.2, 3.0);
    \draw[dashed, gray] (0, 1.0) node[left] {$T_1$} -- (1.2, 1.0);
    
    \node[mypen, font=\bfseries] at (2.3, 2.0) {$L > 0$};
\end{tikzpicture}
\end{minipage}\hfill
\begin{minipage}{0.52\textwidth}
Dal grafico rettangolare si deducono per immediata via geometrica tutte le grandezze del ciclo:
\begin{itemize}
    \item \textbf{Calore assorbito} ($A \to B$ a $T_2$ costante):
    \begin{equation}
        Q_2 = \int_{S_A}^{S_B} T_2\,dS = T_2 (S_B - S_A) > 0
    \end{equation}
    \item \textbf{Calore ceduto} ($C \to D$ a $T_1$ costante):
    \begin{equation}
        Q_1 = \int_{S_B}^{S_A} T_1\,dS = - T_1 (S_B - S_A) < 0
    \end{equation}
    \item \textbf{Lavoro netto} (Area del rettangolo azzurro):
    \begin{equation}
        L = Q_2 - |Q_1| = (T_2 - T_1)(S_B - S_A)
    \end{equation}
\end{itemize}
Il rendimento del ciclo si ottiene quindi con una banale semplificazione algebrica:
\begin{equation}
    \eta_C = \frac{L}{Q_2} = \frac{(T_2 - T_1)(S_B - S_A)}{T_2 (S_B - S_A)} = 1 - \frac{T_1}{T_2}
\end{equation}
\end{minipage}
\caption{Il ciclo di Carnot rappresentato nel piano T-S: le due isoterme orizzontali e le due adiabatiche verticali delimitano un perfetto rettangolo, rendendo immediata la deduzione geometrica del rendimento.}
\end{figure}

\subsection{Degradazione dell'Energia e Lavoro Perduto}
Consideriamo due cicli operanti fra gli stessi termostati $T_A, T_B$, di cui uno di Carnot ($M_C$) e uno irreversibile ($M$).

\begin{figure}[ht]
\centering
\begin{tikzpicture}[scale=1.1]
% Sorgente Calda
\fill[pattern=north east lines, notered!25] (0, 3) rectangle (3, 4);
\draw[thick, mypen, fill=orange!15] (0, 3) rectangle (3, 4) node[midway, font=\bfseries] {Sorgente Calda ($T_B$)};
\node[left, font=\bfseries, notered] at (0, 3.5) {$T_B > T_A$};

% Macchina Termica
\draw[thick, mypen, fill=cyan!15] (1.5, 1.5) circle (0.6);
\node[font=\bfseries] at (1.5, 1.5) {$M / M_C$};

% Sorgente Fredda
\fill[pattern=north west lines, cyan!25] (0, -1) rectangle (3, 0);
\draw[thick, mypen, fill=cyan!10] (0, -1) rectangle (3, 0) node[midway, font=\bfseries] {Sorgente Fredda ($T_A$)};
\node[left, font=\bfseries, blue] at (0, -0.5) {$T_B > T_A$};

% Frecce di scambio termico
\draw[->, >=stealth, thick, notered, line width=1.5pt] (1.5, 3) -- (1.5, 2.1) node[midway, right, font=\bfseries] {$Q_B$};
\draw[->, >=stealth, thick, blue, line width=1.5pt] (1.5, 0.9) -- (1.5, 0) node[midway, right, font=\bfseries] {$Q_A$ (o $Q_{AC}$)};

% Lavoro W
\draw[->, >=stealth, thick, mygreen, line width=1.5pt] (2.1, 1.5) -- (3.5, 1.5) node[right, font=\bfseries] {$L_M$ (o $L_{M_C}$)};
\end{tikzpicture}
\caption{Confronto termodinamico tra un motore reversibile di Carnot ($M_C$) e un motore reale irreversibile ($M$) operanti tra i medesimi serbatoi termici $T_B > T_A$: a parità di calore assorbito $Q_B$, il motore irreversibile cede maggior calore alla sorgente fredda ($|Q_A| > |Q_{AC}|$), comportando una perdita di lavoro utile ($L_M < L_{M_C}$) proporzionale all'entropia generata.}
\end{figure}

Entrambe le macchine $M$ e $M_C$ assorbono lo stesso calore $Q_B$ dalla sorgente calda.
Tuttavia, esse cedono calore alla sorgente fredda in modo diverso: $M$ cede $Q_A$, mentre $M_C$ cede $Q_{AC}$.
Essendo $M_C$ reversibile, il suo rendimento è quello di Carnot:
\begin{equation*}
    \eta_C = 1 - \frac{|Q_{AC}|}{Q_B} = 1 - \frac{T_A}{T_B}
\end{equation*}
Per la macchina irreversibile $M$, il rendimento è strettamente minore di quello di Carnot:
\begin{equation*}
    \eta = 1 - \frac{|Q_A|}{Q_B} < 1 - \frac{T_A}{T_B}
\end{equation*}

Poiché entrambe compiono cicli completi, la variazione di entropia delle macchine è nulla ($\Delta S_M = \Delta S_{M_C} = 0$). 
La variazione di entropia dell'Universo per il ciclo irreversibile è quindi data solo dalle sorgenti:
\begin{equation*}
    \Delta S_{\text{univ}} = \Delta S_{T_A} + \Delta S_{T_B} = - \frac{Q_A}{T_A} - \frac{Q_B}{T_B} > 0
\end{equation*}
Mentre per il ciclo reversibile $M_C$ risulta $\Delta S_{\text{univ}} = 0$.

Possiamo ora calcolare la differenza tra il lavoro prodotto dalla macchina reversibile ($L_{M_C}$) e quello prodotto dalla macchina reale ($L_M$):
\begin{align*}
    L_M &= Q_A + Q_B = (Q_A - Q_{AC}) + (Q_{AC} + Q_B) \\
    &= (Q_A - Q_{AC}) + L_{M_C}
\end{align*}
Sostituendo $Q_{AC} = - \frac{T_A}{T_B} Q_B$ (dal ciclo di Carnot), otteniamo:
\begin{align*}
    L_M &= L_{M_C} + Q_A + \frac{T_A}{T_B} Q_B = L_{M_C} + T_A \left( \frac{Q_A}{T_A} + \frac{Q_B}{T_B} \right) \\
    &= L_{M_C} - T_A \Delta S_{\text{univ}}
\end{align*}
Da cui si ricava infine che:
\begin{equation}
    \boxed{ L_M = L_{M_C} - T_A \Delta S_{\text{univ}} < L_{M_C} }
\end{equation}
La differenza $\boxed{ L_{M_C} - L_M = T_A \Delta S_{\text{univ}} }$ esprime il teorema di \textbf{Gouy-Stodola}: il lavoro meccanico perduto a causa delle irreversibilità è direttamente proporzionale all'entropia generata nell'universo.

\subsection{Casi Tipici di Irreversibilità}
L'irreversibilità di un processo in un sistema isolato si dimostra verificando che la variazione di entropia dell'Universo (Sistema + Ambiente) sia strettamente positiva ($\Delta S_{\text{univ}} > 0$).

\begin{itemize}
    \item \textbf{CASO 1: Riscaldamento per attrito} \\
    Un corpo è trascinato su una superficie scabra. Tutto il lavoro speso per vincere l'attrito si trasforma in variazione di energia interna ($\Delta U = L_{\text{attrito}}$), innalzando la temperatura del sistema (corpo + superficie). Essendo isolato dall'esterno ($Q = 0$), l'ambiente non subisce variazioni. Per calcolare l'entropia, si congiungono gli stati mediante un riscaldamento isocoro reversibile:
    \begin{equation}
        \Delta S_{\text{univ}} = \Delta S_{\text{sist}} = C \ln\left( \frac{T_f}{T_i} \right) > 0 \quad (\text{poiché } T_f > T_i)
    \end{equation}

    \item \textbf{CASO 2: Espansione libera di un gas} \\
    Come analizzato in precedenza, l'espansione nel vuoto di un gas perfetto ($Q=0, L=0, \Delta U=0$) avviene a temperatura costante, ma il volume aumenta. Poiché il sistema è isolato termicamente:
    \begin{equation}
        \Delta S_{\text{univ}} = \Delta S_{\text{sist}} = nR \ln\left( \frac{V_B}{V_A} \right) > 0 \quad (\text{poiché } V_B > V_A)
    \end{equation}

    \item \textbf{CASO 3: Scambio termico per conduzione} \\
    Consideriamo due corpi solidi isolati, con uguale capacità termica $C$ ma temperature iniziali $T_1 \neq T_2$. Essi vengono posti in contatto e raggiungono l'equilibrio a $T_f = \frac{T_1 + T_2}{2}$. L'entropia dell'universo (che coincide col sistema dei due corpi) è la somma delle due variazioni, calcolate su isocore reversibili:
    \begin{equation}
        \Delta S_{\text{univ}} = C \ln\left( \frac{T_f}{T_1} \right) + C \ln\left( \frac{T_f}{T_2} \right) = C \ln\left( \frac{T_f^2}{T_1 T_2} \right)
    \end{equation}
    Ricordando che la media aritmetica è strettamente maggiore della media geometrica per numeri disuguali, si ha $T_f^2 > T_1 T_2$. Dunque il logaritmo è positivo e $\Delta S_{\text{univ}} > 0$.

    \item \textbf{CASO 4: Ciclo monotermo} \\
    In un ciclo termodinamico che scambia calore con una sola sorgente a temperatura $T$, il sistema torna allo stato iniziale e dunque $\Delta S_{\text{sist}} = 0$. Il Primo Principio impone che il lavoro eseguito sia pari al calore assorbito ($L = Q$). Per il principio di Kelvin-Planck, tale lavoro non può essere positivo ($L \le 0 \implies Q \le 0$). Il calore assorbito dalla sorgente è quindi $-Q \ge 0$, da cui:
    \begin{equation}
        \Delta S_{\text{univ}} = \Delta S_{\text{sist}} + \Delta S_{\text{amb}} = 0 - \frac{Q}{T} \ge 0
    \end{equation}
    Se il ciclo produce attriti e richiede lavoro per essere completato ($L < 0$), allora $\Delta S_{\text{univ}} > 0$ e il ciclo è irreversibile.
\end{itemize}

\vspace{0.5cm}
\begin{flushright}
\textit{Fine del Capitolo 6 -- Termodinamica}
\end{flushright}

\chapter{Relatività Ristretta}

\section{Moti Relativi e Sistemi di Riferimento}

Nello studio dei moti relativi si scelgono sempre \textbf{due sistemi di riferimento} e si descrive lo stesso fenomeno fisico dall'uno e dall'altro, confrontando i risultati.

\begin{definizione}{Sistema Fisso e Sistema in Moto}
\begin{itemize}[leftmargin=1.4em]
    \item \textbf{Sistema ``fisso'' (assoluto)} $S$: è per costruzione un sistema \textbf{inerziale}, rispetto al quale si misurano le grandezze dette \emph{assolute}.
    \item \textbf{Sistema ``in moto'' (relativo)} $S'$: non è necessariamente inerziale. Si muove rispetto a $S$ con velocità $\vec{V}$ e accelerazione $\vec{A}$, dette rispettivamente \textbf{velocità di trascinamento} (o \textbf{boost}) e \textbf{accelerazione di trascinamento}.
\end{itemize}
\end{definizione}

\begin{figure}[ht]
    \centering
    \begin{tikzpicture}[>=stealth, scale=1.0]
        % Sistema S
        \draw[->, thick, mypen] (0,0) -- (2.2,0) node[right, font=\small] {$x$};
        \draw[->, thick, mypen] (0,0) -- (0,2.4) node[above, font=\small] {$y$};
        \draw[->, thick, mypen] (0,0) -- (-1.1,-1.1) node[below left, font=\small] {$z$};
        \node[below left, font=\bfseries\small, mypen] at (0,0) {$S$};

        % Sistema S'
        \begin{scope}[xshift=4.2cm]
            \draw[->, thick, boxese] (0,0) -- (2.2,0) node[right, font=\small] {$x'$};
            \draw[->, thick, boxese] (0,0) -- (0,2.4) node[above, font=\small] {$y'$};
            \draw[->, thick, boxese] (0,0) -- (-1.1,-1.1) node[below left, font=\small] {$z'$};
            \node[below left, font=\bfseries\small, boxese] at (0,0) {$S'$};
        \end{scope}

        % Boost
        \draw[->, very thick, notered] (4.2,-1.6) -- (5.9,-1.6)
            node[midway, below, font=\small\bfseries] {$\vec{V}$ (boost)};

        % Punto materiale P
        \coordinate (P) at (6.2,2.0);
        \fill[notered] (P) circle (2.2pt) node[above right, font=\small] {$P$};
        \draw[->, thick, mypen, dashed] (0,0) -- (P) node[pos=0.62, above left, font=\small] {$\vec{r}$};
        \draw[->, thick, boxese, dashed] (4.2,0) -- (P) node[pos=0.5, below right, font=\small] {$\vec{r}\,'$};
        \draw[->, thick, notered] (0,0) -- (4.2,0) node[pos=0.5, below, font=\small] {$\vec{R}$};
    \end{tikzpicture}
    \caption{I due sistemi di riferimento adottati nello studio dei moti relativi: il sistema assoluto $S$ e il sistema relativo $S'$, traslato di $\vec{R}$ e animato dalla velocità di trascinamento $\vec{V}$.}
\end{figure}

\subsection{Relatività Galileiana: $S$ e $S'$ Entrambi Inerziali}

La composizione classica delle velocità lega la velocità $\vec{u}$ di un corpo misurata in $S$ a quella $\vec{u}\,'$ misurata in $S'$:
\begin{equation}
    \vec{u} = \vec{u}\,' + \vec{V}
\end{equation}
dove
\begin{itemize}
    \item $\vec{u}$ è la velocità del corpo vista da $S$;
    \item $\vec{u}\,'$ è la velocità del corpo vista da $S'$;
    \item $\vec{V}$ è la velocità di $S'$ rispetto a $S$.
\end{itemize}

Se $\vec{V} = \text{cost}$ (ossia $S'$ è a sua volta inerziale), l'integrazione nel tempo è immediata:
\begin{equation*}
    \int \vec{u} \, dt = \int \vec{u}\,' \, dt + \vec{V} \int dt
\end{equation*}
e si ottengono le \textbf{trasformazioni di Galileo}.

\begin{definizione}{Trasformazioni di Galileo}
Per un boost $\vec{V} = (V_x, V_y, V_z)$ costante, le coordinate spazio-temporali nei due sistemi sono legate da
\begin{equation}
    \begin{cases}
        x = x' + V_x t \\
        y = y' + V_y t \\
        z = z' + V_z t \\
        t = t'
    \end{cases}
\end{equation}
Il tempo è un \textbf{parametro assoluto}, identico nei due sistemi e indipendente dallo stato di moto dell'osservatore.
\end{definizione}

Derivando ulteriormente la legge di composizione delle velocità si ottiene quella delle accelerazioni:
\begin{equation}
    \vec{a} = \vec{a}\,' + \vec{A}
\end{equation}
dove $\vec{a}$ e $\vec{a}\,'$ sono le accelerazioni del corpo misurate rispettivamente in $S$ e in $S'$, e $\vec{A}$ è l'accelerazione di $S'$ rispetto a $S$. Poiché $\vec{V}$ è costante, si ha $\vec{A} = \vec{0}$ e dunque
\begin{equation}
    \vec{a} = \vec{a}\,' \implies m\vec{a} = m\vec{a}\,' \implies \vec{F} = \vec{F}\,'
\end{equation}

\begin{teorema}{Invarianza Galileiana del Secondo Principio}
Il secondo principio della dinamica è invariante \textbf{in forma e in valore} nel passaggio fra due sistemi di riferimento inerziali collegati da una trasformazione di Galileo.
\end{teorema}

\subsection{Il Caso di $S'$ Non Inerziale: le Forze Apparenti}

Se la velocità di trascinamento non è costante, $S'$ non è inerziale e la derivazione della legge di composizione delle velocità fornisce
\begin{equation}
    \frac{d\vec{u}}{dt} = \frac{d\vec{u}\,'}{dt} + \frac{d\vec{V}}{dt}
    \qquad \text{ovvero} \qquad
    \vec{a} = \vec{a}\,' + \vec{A}, \quad \text{con } \vec{A} \neq \vec{0}
\end{equation}

Moltiplicando per la massa si ottiene $m\vec{a} = m\vec{a}\,' + m\vec{A}$, cioè
\begin{equation}
    \vec{F} = \vec{F}\,' + m\vec{A}
    \qquad \Longrightarrow \qquad
    \vec{F} + \vec{F}_{\text{app}} = \vec{F}\,'
\end{equation}

\begin{definizione}{Forza Apparente}
Si definisce \textbf{forza apparente} (o fittizia) il termine
\begin{equation}
    \vec{F}_{\text{app}} = -m\vec{A}
\end{equation}
Essa non deriva da alcuna interazione fisica fra corpi, ma è il puro effetto cinematico del moto accelerato del sistema di riferimento: per questo motivo \textbf{non verifica il terzo principio della dinamica} (non esiste alcuna reazione uguale e contraria).
\end{definizione}

La relazione $\vec{F} + \vec{F}_{\text{app}} = \vec{F}\,'$ costituisce la \textbf{generalizzazione del secondo principio} ai sistemi di riferimento non inerziali.

\subsection{Il Principio di Relatività di Galilei}

\begin{teorema}{Principio di Relatività di Galilei}
Le leggi della meccanica sono \textbf{invarianti sotto trasformazioni di Galileo} in tutti i sistemi di riferimento inerziali.
\end{teorema}

Da esso discendono tre conseguenze fondamentali:
\begin{enumerate}
    \item \textbf{Invarianza degli intervalli temporali.} Da $t_a = t'_a$ e $t_b = t'_b$ segue immediatamente
    \begin{equation}
        \Delta t = \Delta t'
    \end{equation}
    \item \textbf{Invarianza degli intervalli spaziali.} Da $x_a = x'_a + V t_a$ e $x_b = x'_b + V t_b$ segue
    \begin{equation}
        \Delta x = \Delta x' + V \Delta t
    \end{equation}
    e, misurando le posizioni nel medesimo istante ($\Delta t = 0$), $\Delta x = \Delta x'$: le lunghezze sono le stesse per tutti gli osservatori.
    \item \textbf{Inesistenza di un sistema assoluto privilegiato.} Tutte le velocità sono relative ai vari sistemi; nessuno di essi può essere eletto a riferimento ``vero''.
\end{enumerate}

\section{Il Periodo di Transizione: Maxwell, Michelson e Morley}

\subsection{L'Invarianza della Velocità della Luce}

Attraverso le sue equazioni, \textbf{Maxwell} dimostra che la velocità della luce nel vuoto, $c \approx 3{,}00 \times 10^8 \, \mathrm{m/s}$, è definita da
\begin{equation}
    c = \frac{1}{\sqrt{\varepsilon_0 \mu_0}}
\end{equation}
dove $\varepsilon_0$ è la costante dielettrica del vuoto e $\mu_0$ la permeabilità magnetica del vuoto.

Il punto cruciale è che $c$ risulta \textbf{costante}, poiché dipende unicamente da due costanti universali e non dalle trasformazioni galileiane. Essa è la velocità di propagazione nel vuoto di una \textbf{perturbazione} dei campi $\vec{E}$ (elettrico) e $\vec{B}$ (magnetico), cioè di un'\textbf{onda elettromagnetica}.

\subsection{L'Ipotesi dell'Etere Cosmico}

D'Alembert aveva dimostrato che le onde \emph{meccaniche} necessitano di un mezzo per propagarsi. Per analogia si ipotizzò dunque l'esistenza dell'\textbf{etere cosmico}, un mezzo dalle proprietà singolari:
\begin{itemize}
    \item permea l'intero universo;
    \item è volatile;
    \item ha densità nulla.
\end{itemize}
Il suo unico scopo sarebbe stato quello di far propagare la luce.

Resta però un secondo problema, ben più grave: il fatto che $c$ sia costante è in accordo con Maxwell ma \textbf{non} con Galileo. \emph{Chi ha ragione?} Vengono formulate tre ipotesi alternative:

\begin{enumerate}[label=\textbf{\alph*)}]
    \item Il principio di relatività è \textbf{valido solo per la meccanica} e non per l'elettromagnetismo: nell'elettrodinamica esiste un sistema assoluto e privilegiato (l'etere). Si può applicare Galileo e dimostrarlo sperimentalmente.
    \item Il principio è \textbf{valido per meccanica ed elettromagnetismo}, ma le equazioni di Maxwell sono sbagliate e la luce non ha velocità costante. Si può applicare Galileo e dimostrarlo sperimentalmente.
    \item Il principio è \textbf{valido per meccanica ed elettromagnetismo}, ma sono le equazioni di Galileo e Newton a essere sbagliate: è necessario rivedere le leggi di trasformazione, applicarne di nuove e dimostrarle sperimentalmente.
\end{enumerate}

\begin{osservazione}{Il ``Vento d'Etere''}
Identificare il sistema di riferimento etereo significa dimostrare che la velocità della luce vale $c' = c \pm v$ a seconda dell'orientazione, rispetto al sistema della Terra (dove $v$ è la velocità della Terra rispetto all'etere, cioè la velocità di trascinamento), mentre vale esattamente $c$ rispetto al sistema dell'etere.

Di conseguenza, l'osservatore terrestre avvertirebbe un \textbf{vento d'etere} di velocità $|\vec{v}\,|$, uguale e opposta a quella della Terra rispetto all'etere, in cui essa si muove con i propri moti di rotazione e rivoluzione.
\end{osservazione}

Per ottenere una verifica conclusiva e priva di ambiguità era necessario misurare gli effetti al \textbf{secondo ordine} in $\beta = v/c$:
\begin{equation*}
    \text{I ordine:} \quad \frac{v}{c} = 10^{-4}
    \qquad\qquad
    \text{II ordine:} \quad \left(\frac{v}{c}\right)^{2} = 10^{-8}
\end{equation*}
avendo posto $v \approx 3 \cdot 10^4 \, \mathrm{m/s}$ e $c \approx 3 \cdot 10^8 \, \mathrm{m/s}$. La precisione di $10^{-8}$ è il livello necessario per trovare una risposta definitiva e confermare una delle ipotesi precedenti.

\subsection{L'Esperimento di Michelson e Morley}

\begin{notastorica}{1881--1887}
Nel 1881, e poi in modo definitivo nel 1887 insieme a Morley, \textbf{Michelson} conduce l'esperimento sfruttando l'\textbf{interferometria}, l'unica tecnica dell'epoca capace di raggiungere la sensibilità del secondo ordine richiesta.
\end{notastorica}

\begin{definizione}{Frange d'Interferenza}
Le \textbf{frange d'interferenza} sono l'alternanza fra zone di luce e zone di ombra osservate al telescopio, dovute allo sfasamento temporale $\Delta t$ fra i due raggi ricombinati. Le due cause dello sfasamento sono:
\begin{equation*}
    \ell_1 \neq \ell_2
    \qquad \text{e} \qquad
    c'_1 \neq c'_2
\end{equation*}
cioè la disuguaglianza dei bracci e la diversa velocità apparente della luce lungo i due percorsi.
\end{definizione}

\begin{figure}[ht]
    \centering
    \begin{tikzpicture}[>=stealth, scale=1.15]
        % --- Sorgente S
        \fill[boxatt!75] (-3.1,0) circle (3pt);
        \foreach \a in {-40,-20,0,20,40} {
            \draw[boxatt!75, thick] (\a:0.16) ++(-3.1,0) -- ++(\a:0.22);
        }
        \node[above, font=\small\bfseries, boxatt] at (-3.1,0.3) {$S$};
        \draw[very thick, boxatt!75] (-2.75,0) -- (-0.4,0);

        % --- Specchio semiriflettente M (a 45 gradi)
        \draw[line width=2.2pt, mypen!75] (-0.42,-0.42) -- (0.42,0.42);
        \node[below left, font=\small\bfseries, mypen] at (-0.2,-0.28) {$M$};

        % --- Braccio 1: orizzontale, verso M1
        \draw[very thick, boxoss] (0,0) -- (2.9,0);
        \draw[line width=2.4pt, mypen!75] (3.05,-0.6) -- (3.05,0.6);
        \node[right, font=\small\bfseries, mypen] at (3.15,0) {$M_1$};
        \draw[->, thick, boxoss] (1.15,0.3) -- (1.95,0.3) node[right, font=\scriptsize] {$1$};
        \draw[->, thick, boxoss] (1.95,-0.3) -- (1.15,-0.3) node[left, font=\scriptsize] {$1'$};
        \node[below, font=\small, boxoss] at (1.55,-0.62) {$\ell_1$};

        % --- Braccio 2: verticale, verso M2
        \draw[very thick, boxese] (0,0) -- (0,2.6);
        \draw[line width=2.4pt, mypen!75] (-0.6,2.75) -- (0.6,2.75);
        \node[above, font=\small\bfseries, mypen] at (0,2.85) {$M_2$};
        \draw[->, thick, boxese] (0.3,1.05) -- (0.3,1.85) node[above, font=\scriptsize] {$2$};
        \draw[->, thick, boxese] (-0.3,1.85) -- (-0.3,1.05) node[below, font=\scriptsize] {$2'$};
        \node[left, font=\small, boxese] at (-0.62,1.5) {$\ell_2$};

        % --- Telescopio T
        \draw[very thick, notered] (0,0) -- (0,-1.7);
        \draw[line width=1.2pt, mypen!75, fill=mypen!8] (-0.42,-2.25) rectangle (0.42,-1.7);
        \node[below, font=\small\bfseries, mypen] at (0,-2.32) {$T$};

        % --- Vento d'etere
        \draw[->, line width=1.4pt, boxatt] (-3.1,-1.5) -- (-1.8,-1.5);
        \node[below, font=\scriptsize\bfseries, boxatt, align=center] at (-2.45,-1.58) {vento d'etere \; $\vec{v}$};

        % --- Frange d'interferenza
        \begin{scope}[xshift=5.3cm, yshift=-0.3cm]
            \draw[line width=1.1pt, mypen!75] (0,0) circle (1.0);
            \begin{scope}
                \clip (0,0) circle (0.98);
                \foreach \i in {-0.85,-0.55,-0.25,0.05,0.35,0.65} {
                    \fill[mypen!45] (\i,-1) rectangle (\i+0.15,1);
                }
            \end{scope}
            \node[above, font=\scriptsize\bfseries, align=center] at (0,1.1) {frange d'interferenza};
        \end{scope}
        \draw[->, thick, mypen!60, dashed] (0.55,-2.0) to[out=8, in=205] (4.2,-0.75);
    \end{tikzpicture}
    \caption{Schema dell'interferometro di Michelson. La sorgente $S$ invia il raggio sullo specchio semiriflettente $M$, che lo divide nei due raggi $1$ e $2$ diretti agli specchi $M_1$ e $M_2$ lungo i bracci $\ell_1$ e $\ell_2$; i raggi riflessi $1'$ e $2'$ si ricombinano nel telescopio $T$, dove producono le frange d'interferenza.}
\end{figure}

\paragraph{Calcolo dei tempi di propagazione.}
Si prova a isolare la seconda causa, calcolando esplicitamente i due tempi di percorrenza.

\textbf{Raggio 1 --- parallelo al vento d'etere.} Il percorso $1 + 1' \to 2\ell_1$ viene compiuto in parte concorde e in parte discorde al vento d'etere di velocità $v$:
\begin{equation}
    t_1 = \frac{\ell_1}{c+v} + \frac{\ell_1}{c-v}
    = \frac{\ell_1(c-v) + \ell_1(c+v)}{c^2 - v^2}
    = \boxed{\frac{2c\,\ell_1}{c^2 - v^2}}
\end{equation}
calcolato nel sistema della Terra (dove $t_1 = t'_1$).

\textbf{Raggio 2 --- ortogonale al vento d'etere.} Il percorso $2 + 2' \to 2\ell_2$ richiede alcuni passaggi intermedi: nel tempo $t_2/2$ la luce percorre l'ipotenusa di un triangolo rettangolo i cui cateti sono il braccio $\ell_2$ e lo spostamento $v\,t_2/2$ dell'apparato:
\begin{equation}
    \left(\frac{c\,t_2}{2}\right)^{\!2} = \ell_2^2 + \left(\frac{v\,t_2}{2}\right)^{\!2}
    \implies (c^2 - v^2)\, t_2^2 = 4\ell_2^2
    \implies t_2^2 = \frac{4\ell_2^2}{c^2 - v^2}
\end{equation}
da cui
\begin{equation}
    t_2 = \boxed{\frac{2\ell_2}{c} \cdot \frac{1}{\sqrt{1 - \dfrac{v^2}{c^2}}}}
\end{equation}
calcolato nel sistema dell'etere (dove $t_2 = t'_2$).

Secondo Galileo lo sfasamento vale allora
\begin{equation}
    \Delta t = \Delta t' = t_2 - t_1 = \frac{2}{c}
    \left[ \frac{\ell_2}{\sqrt{1 - \dfrac{v^2}{c^2}}} - \frac{\ell_1}{1 - \dfrac{v^2}{c^2}} \right]
\end{equation}

\paragraph{Rotazione dell'interferometro.}
Si ruota ora l'apparato di $\theta = \pi/2$: il raggio $1$ diventa ortogonale al vento d'etere e il raggio $2$ parallelo. Di conseguenza $\ell_1$ e $\ell_2$ si scambiano, e con essi $t_1$ e $t_2$:
\begin{equation}
    \Delta t^{\,r} = \Delta t'^{\,r} = t_2^{\,r} - t_1^{\,r} = t_1 - t_2 = \frac{2}{c}
    \left[ \frac{\ell_2}{1 - \dfrac{v^2}{c^2}} - \frac{\ell_1}{\sqrt{1 - \dfrac{v^2}{c^2}}} \right]
\end{equation}

\paragraph{Differenza dei tempi di propagazione.}
La grandezza effettivamente misurabile è la variazione dello sfasamento indotta dalla rotazione. Posto $\beta = v/c$ e sviluppando in serie per $x \ll 1$
\begin{equation*}
    \frac{1}{1-x} \approx 1 + x
    \qquad\qquad
    \frac{1}{\sqrt{1-x}} \approx 1 + \frac{x}{2}
\end{equation*}
si ottiene
\begin{align}
    \Delta t^{\,r} - \Delta t
    &= \frac{2}{c} \left[ \frac{\ell_2}{1-\beta^2} - \frac{\ell_1}{\sqrt{1-\beta^2}}
       - \frac{\ell_2}{\sqrt{1-\beta^2}} + \frac{\ell_1}{1-\beta^2} \right] \notag \\[4pt]
    &\approx \frac{2}{c} \left[ \ell_2(1+\beta^2) - \ell_1\left(1+\frac{\beta^2}{2}\right)
       - \ell_2\left(1+\frac{\beta^2}{2}\right) + \ell_1(1+\beta^2) \right] \notag \\[4pt]
    &= \boxed{\frac{(\ell_1 + \ell_2)}{c} \cdot \left(\frac{v}{c}\right)^{\!2}}
\end{align}

\begin{osservazione}{Conclusioni dell'Esperimento}
Il fattore $\Delta t^{\,r} - \Delta t$ di sfasamento \textbf{non fornisce alcuna prova} che l'etere esista: ruotando l'interferometro, i due raggi non si sfasano a sufficienza.

Detto $\Delta N$ il numero di frange che attraversano l'origine del riferimento del telescopio durante la rotazione, il risultato sperimentale è
\begin{equation*}
    \Delta N = \frac{1}{100} \approx 0
    \qquad \text{contro il valore previsto} \qquad
    \Delta N = \frac{4}{10}
\end{equation*}
(per bracci di lunghezza $\approx 11 \, \mathrm{m}$). Data la piccolezza di $\Delta N$, si è concluso che \textbf{non c'è nessuno spostamento}. L'esperimento è stato ripetuto nei successivi cinquant'anni, ma la conclusione è sempre rimasta la stessa.
\end{osservazione}

Di fronte a questo risultato furono avanzate tre interpretazioni:
\begin{enumerate}[label=\textbf{\alph*)}]
    \item \textbf{Trascinamento dell'etere:} $\Delta N = 0$ è corretto, le ipotesi sono corrette e l'etere è ``locale'', trascinato da ogni corpo di massa finita; pertanto $v = 0$ e non esiste alcun vento d'etere.
    \item \textbf{Contrazione di Lorentz--Fitzgerald:} $\Delta N = 0$ è corretto, le ipotesi sono corrette, ma nel corso del moto fra Terra ed etere le lunghezze parallele a $\vec{v}$ si contraggono di un fattore $\sqrt{1-\beta^2}$, mentre quelle perpendicolari restano invariate.
    \item \textbf{Necessità di riformulare le trasformazioni:} $\Delta N = 0$ è corretto, ma le ipotesi sono sbagliate e l'etere \textbf{non esiste} (non c'è vento d'etere perché $v=0$ non ha senso: non c'è nulla rispetto a cui misurarla).
\end{enumerate}

Le interpretazioni \textbf{a)} e \textbf{b)} risultano contraddittorie o non applicabili ai diversi esperimenti condotti: l'unica plausibile è la \textbf{c)}.

\section{I Postulati della Relatività Ristretta}

\begin{teorema}{Postulati di Einstein}
La relatività ristretta poggia su due soli postulati:
\begin{enumerate}
    \item \textbf{Invarianza delle leggi fisiche:} le leggi della fisica sono invarianti in tutti i sistemi di riferimento inerziali.
    \item \textbf{Invarianza di $c$:} la velocità della luce nel vuoto ha lo stesso valore $c$ in tutti i sistemi di riferimento inerziali.
\end{enumerate}
\end{teorema}

Le conseguenze ribaltano punto per punto quelle galileiane:
\begin{enumerate}
    \item \textbf{Dilatazione degli intervalli temporali:} da $t_a = t'_a/\gamma$ e $t_b = t'_b/\gamma$ segue $\Delta t > \Delta t'$.
    \item \textbf{Contrazione degli intervalli spaziali:} da $x_a = x'_a/\gamma$ e $x_b = x'_b/\gamma$ segue $\Delta x < \Delta x'$.
    \item \textbf{Inesistenza di un sistema assoluto privilegiato:} tutte le velocità sono relative ai vari sistemi. Questo è l'unico punto che la relatività ristretta eredita intatto dalla relatività galileiana.
\end{enumerate}

\subsection{La Relatività della Simultaneità}

\begin{osservazione}{La Simultaneità Dipende dal Sistema di Riferimento}
Si sincronizzano gli orologi in modo che $t_a = t_b = 0$, mediante un segnale inviato dall'origine $O$.
\begin{itemize}[leftmargin=1.4em]
    \item \textbf{$S$ e $S'$ in quiete:} $O$ vede arrivare il segnale da $AA'$ nello stesso momento di quello da $BB'$ ($t_a = t_b$); lo stesso vale per $O'$ ($t'_a = t'_b$). Gli eventi $A$, $B$, $A'$, $B'$ sono \textbf{simultanei} sia in $S$ sia in $S'$.
    \item \textbf{$S$ in quiete, $S'$ in moto con velocità $\vec{v}$:} $O$ vede ancora arrivare i due segnali nello stesso momento ($t_a = t_b$), ma $O'$ vede arrivare \textbf{prima} $BB'$ e \textbf{poi} $AA'$ ($t'_a > t'_b$). Gli eventi $A$ e $B$ sono simultanei in $S$, mentre $A'$ e $B'$ \textbf{non} lo sono in $S'$.
\end{itemize}
\end{osservazione}

\begin{figure}[ht]
    \centering
    \begin{tikzpicture}[>=stealth, scale=0.95]
        % --- Pannello (a): entrambi in quiete
        \begin{scope}
            \draw[thick, mypen] (0,0) rectangle (4.2,0.75);
            \draw[thick, boxese] (0,0.75) rectangle (4.2,1.5);
            \node[left, font=\small\bfseries, mypen] at (-0.1,0.375) {$S$};
            \node[left, font=\small\bfseries, boxese] at (-0.1,1.125) {$S'$};
            \foreach \x/\lab in {0.5/A, 2.1/O, 3.7/B} {
                \fill[mypen] (\x,0.375) circle (1.6pt) node[below, font=\scriptsize] {$\lab$};
            }
            \foreach \x/\lab in {0.5/{A'}, 2.1/{O'}, 3.7/{B'}} {
                \fill[boxese] (\x,1.125) circle (1.6pt) node[above, font=\scriptsize] {$\lab$};
            }
            \draw[->, thick, notered] (2.1,1.125) -- (0.75,1.125);
            \draw[->, thick, notered] (2.1,1.125) -- (3.45,1.125);
            \node[font=\small, align=center] at (2.1,-0.95) {(a) $S$ e $S'$ in quiete\\$t_a = t_b$, \; $t'_a = t'_b$};
        \end{scope}
        % --- Pannello (b): S' in moto
        \begin{scope}[xshift=6.4cm]
            \draw[thick, mypen] (0,0) rectangle (4.2,0.75);
            \draw[thick, boxese] (0,0.75) rectangle (4.2,1.5);
            \node[left, font=\small\bfseries, mypen] at (-0.1,0.375) {$S$};
            \node[left, font=\small\bfseries, boxese] at (-0.1,1.125) {$S'$};
            \foreach \x/\lab in {0.5/A, 2.1/O, 3.7/B} {
                \fill[mypen] (\x,0.375) circle (1.6pt) node[below, font=\scriptsize] {$\lab$};
            }
            \foreach \x/\lab in {0.5/{A'}, 2.1/{O'}, 3.7/{B'}} {
                \fill[boxese] (\x,1.125) circle (1.6pt) node[above, font=\scriptsize] {$\lab$};
            }
            \draw[->, thick, notered] (2.1,1.125) -- (0.75,1.125);
            \draw[->, thick, notered] (2.1,1.125) -- (3.45,1.125);
            \draw[->, very thick, boxatt] (4.45,1.125) -- (5.35,1.125)
                node[above, font=\scriptsize\bfseries, midway] {$\vec{v}$};
            \node[font=\small, align=center] at (2.1,-0.95) {(b) $S'$ in moto\\$t_a = t_b$, \; $t'_a > t'_b$};
        \end{scope}
    \end{tikzpicture}
    \caption{La relatività della simultaneità. Due eventi simultanei nel sistema $S$ non lo sono più nel sistema $S'$ in moto: l'osservatore $O'$ va incontro al segnale proveniente da $B$ e si allontana da quello proveniente da $A$.}
\end{figure}

\section{Le Trasformazioni di Lorentz}

\subsection{Derivazione}

Consideriamo l'istante $t = t' = 0$, nel quale le origini $O$ e $O'$ dei due sistemi inerziali coincidono, e supponiamo che da esse venga inviato un segnale luminoso verso un punto $A$. Il segnale si propaga con velocità $c$ in entrambi i sistemi (secondo postulato), per cui
\begin{equation}
    \overline{OA}^2 = \vec{r}^{\;2} = c^2 t^2 = x^2 + y^2 + z^2
    \qquad\qquad
    \overline{O'A}^2 = \vec{r}^{\;\prime 2} = c^2 t'^2 = x'^2 + y'^2 + z'^2
\end{equation}

Poiché il boost è diretto lungo $\hat{x}$, le direzioni trasverse non sono coinvolte: $y' = y$ e $z' = z$. Le due coordinate rimanenti devono invece trasformarsi in modo lineare e generale:
\begin{equation*}
    x' = ax + by + dz + ft
    \qquad\qquad
    t' = gx + hy + mz + nt
\end{equation*}
Dato che $x'$ ``scorre'' solo lungo $x$, mentre $y$ e $z$ restano inalterate, i coefficienti $b$, $d$, $h$, $m$ sono nulli e resta
\begin{equation*}
    x' = ax + ft
    \qquad\qquad
    t' = gx + nt
\end{equation*}

Sappiamo inoltre che un punto solidale con $S'$, che ha $x = vt$ in $S$, ha $x' = 0$ in $S'$; quindi
\begin{equation*}
    x' = 0 = ax + ft \implies a(vt) = -ft \implies f = -av
\end{equation*}
e sostituendo si ottiene la forma cercata
\begin{equation}
    x' = ax - avt = a(x - vt)
\end{equation}

Imponiamo ora l'invarianza del fronte d'onda luminoso, sostituendo $x'$ e $t'$ in $\vec{r}^{\;\prime 2} = c^2t'^2$:
\begin{equation}
    c^2 (gx + nt)^2 = a^2 (x - vt)^2 + y^2 + z^2
\end{equation}
Sviluppando i quadrati:
\begin{equation*}
    (cg)^2 x^2 + (cn)^2 t^2 + 2c^2 gn \, xt = a^2 x^2 + (av)^2 t^2 - 2a^2 v \, xt + y^2 + z^2
\end{equation*}
e raccogliendo i termini omologhi:
\begin{equation}
    \underbrace{\left(c^2 n^2 - a^2 v^2\right)}_{\to \; c^2} t^2 =
    \underbrace{\left(a^2 - c^2 g^2\right)}_{\to \; 1} x^2 + y^2 + z^2
    \underbrace{- \, 2\left(c^2 gn + a^2 v\right)}_{\to \; 0} xt
\end{equation}

Il confronto con $c^2t^2 = x^2+y^2+z^2$ fornisce il sistema dei tre vincoli:
\begin{equation}
    \begin{cases}
        c^2 n^2 - a^2 v^2 = c^2 & \implies a^2 = c^2(n^2-1)/v^2 \\[2pt]
        a^2 - c^2 g^2 = 1 & \implies a^2 = 1 + c^2 g^2 \\[2pt]
        c^2 gn + a^2 v = 0 & \implies g = -a^2 v / (c^2 n)
    \end{cases}
\end{equation}

Sostituendo la terza nella seconda e poi nella prima, si ricava $c^2 = (c^2 - v^2)n^2$, ossia
\begin{equation}
    n = \frac{1}{\sqrt{1 - v^2/c^2}}
\end{equation}
e di conseguenza
\begin{equation}
    a^2 = \frac{c^2}{v^2}\left(\frac{c^2}{c^2-v^2} - 1\right)
        = \frac{c^2}{v^2}\left(\frac{\cancel{c^2} - \cancel{c^2} + v^2}{c^2-v^2}\right)
        = \frac{c^2}{c^2-v^2}
    \implies a = \frac{1}{\sqrt{1-v^2/c^2}} = n
\end{equation}
\begin{equation}
    g = -\frac{a^2 v}{c^2 n} = -\frac{a^2 v}{c^2 a} = -\frac{1}{\sqrt{1-v^2/c^2}} \cdot \frac{v}{c^2}
\end{equation}

\begin{definizione}{Fattore di Lorentz}
Si definisce \textbf{fattore di Lorentz} la quantità adimensionale
\begin{equation}
    \gamma = \frac{1}{\sqrt{1 - \dfrac{v^2}{c^2}}} = \frac{1}{\sqrt{1-\beta^2}}
    \qquad \text{con} \qquad \beta = \frac{v}{c}
\end{equation}
Esso soddisfa sempre $\gamma \geq 1$, con $\gamma = 1$ se e solo se $v = 0$.
\end{definizione}

\begin{formulario}{Trasformazioni di Lorentz}
\begin{minipage}[t]{0.47\textwidth}
\centering \textbf{Da $S$ a $S'$ (dirette)}
\begin{equation*}
    \begin{cases}
        x' = \gamma\,(x - vt) \\
        y' = y \\
        z' = z \\
        t' = \gamma\left(t - \dfrac{v}{c^2}x\right)
    \end{cases}
\end{equation*}
\end{minipage}
\hfill
\begin{minipage}[t]{0.47\textwidth}
\centering \textbf{Da $S'$ a $S$ (inverse)}
\begin{equation*}
    \begin{cases}
        x = \gamma\,(x' + vt') \\
        y = y' \\
        z = z' \\
        t = \gamma\left(t' + \dfrac{v}{c^2}x'\right)
    \end{cases}
\end{equation*}
\end{minipage}
\vspace{4pt}
\end{formulario}

\begin{osservazione}{Simmetria delle Trasformazioni}
Le trasformazioni inverse si ottengono da quelle dirette scambiando gli apici e sostituendo $v \to -v$: nessuno dei due sistemi è privilegiato rispetto all'altro, in pieno accordo con il primo postulato.
\end{osservazione}

\subsection{Trasformazioni Relativistiche per la Velocità}

\paragraph{Caso $\vec{u} \parallel \hat{x}$ (velocità del corpo parallela al boost).}
Siano $\vec{V}$ il boost di $S'$ rispetto a $S$, $\vec{u}\,'$ la velocità del corpo $A$ rispetto a $S'$ e $\vec{u}$ quella rispetto a $S$. Col passare del tempo si ha $x' = u' t'$; sostituendo le trasformazioni di Lorentz $x' = \gamma(x-vt)$ e $t' = \gamma\left(t - \frac{v}{c^2}x\right)$:
\begin{equation*}
    \cancel{\gamma}\,(x - vt) = \cancel{\gamma}\left(t - \frac{v}{c^2}x\right) u'
    \implies x - vt = u't - \frac{u'v}{c^2}x
    \implies x\left(1 + \frac{u'v}{c^2}\right) = (u'+v)\,t
\end{equation*}
da cui
\begin{equation*}
    x = \frac{(u'+v)}{1 + \dfrac{vu'}{c^2}}\, t
\end{equation*}
Sapendo che nel sistema $S$ vale $x = ut$, eguagliamo le due espressioni:
\begin{equation}
    \boxed{\;u_x = \frac{u'_x + v}{1 + \dfrac{v\,u'_x}{c^2}}\;}
    \qquad \text{lungo } \hat{x}
\end{equation}

\paragraph{Caso $\vec{u} \perp \hat{x}$ (velocità del corpo perpendicolare al boost).}
Col passare del tempo $\Delta y' = u'_y \Delta t'$, ma $\Delta y' = \Delta y$ perché il boost non ha componente lungo $\hat{y}$; quindi $u'_y \Delta t' = u_y \Delta t$. Sostituendo la trasformazione del tempo:
\begin{align*}
    u'_y \left[ \gamma\left(t_2 - \frac{v}{c^2}x_2\right) - \gamma\left(t_1 - \frac{v}{c^2}x_1\right) \right] &= u_y (t_2 - t_1) \\
    u'_y \, \gamma \left( \Delta t - \frac{v}{c^2}\Delta x \right) &= u_y \, \Delta t
\end{align*}
e dividendo per $\Delta t$:
\begin{equation}
    \boxed{\;u_y = \gamma \, u'_y \left(1 - \frac{v}{c^2}u_x\right)\;}
    \qquad \text{lungo } \hat{y}
\end{equation}
Per $u_z$ si segue lo stesso ragionamento, ottenendo
\begin{equation}
    \boxed{\;u_z = \gamma \, u'_z \left(1 - \frac{v}{c^2}u_x\right)\;}
    \qquad \text{lungo } \hat{z}
\end{equation}

\begin{attenzione}{Le Componenti Trasverse Cambiano Comunque}
Anche se le \emph{lunghezze} trasverse restano invariate ($\Delta y' = \Delta y$), le \emph{velocità} trasverse si trasformano: ciò accade perché a cambiare è l'intervallo di tempo $\Delta t$ al denominatore, non lo spostamento al numeratore.
\end{attenzione}

\subsection{Invarianti e Velocità Limite}

Abbiamo visto che $c^2t^2 = x^2+y^2+z^2$ e $c^2t'^2 = x'^2+y'^2+z'^2$: dividendo per $t^2$ e $t'^2$ si ricava
\begin{equation}
    \frac{c^2t^2}{t^2} = \frac{x^2}{t^2} + \frac{y^2}{t^2} + \frac{z^2}{t^2} = u_x^2 + u_y^2 + u_z^2
    \qquad\qquad
    \frac{c^2t'^2}{t'^2} = u_x'^2 + u_y'^2 + u_z'^2
\end{equation}
poiché il segnale emesso da $O$ è luminoso e dunque si muove a $c$. Ciò vale sia se $\vec{u}$ è costante sia se non lo è: $c$ è \textbf{invariante}.

Dalle trasformazioni relativistiche per la velocità emerge dunque che se $u = c$ allora anche $u' = c$. Inoltre:
\begin{itemize}
    \item per $v \ll c$ si ha $\gamma \to 1$, e le trasformazioni di Lorentz si riducono a quelle di Galileo (\textbf{principio di corrispondenza});
    \item per $v \to c$ si ha $\gamma \to \infty$;
    \item per $v > c$ (impossibile) si avrebbe $\gamma \to n/i$, cioè ci si sposterebbe nel campo complesso $\mathbb{C}$.
\end{itemize}
Pertanto $c$ è la \textbf{velocità limite} dell'universo.

\begin{figure}[ht]
    \centering
    \begin{tikzpicture}[>=stealth, scale=1.0]
        \begin{scope}
            \draw[->, thick, mypen] (0,0) -- (4.6,0) node[right, font=\small] {$v/c$};
            \draw[->, thick, mypen] (0,0) -- (0,3.4) node[above, font=\small] {$\gamma$};
            \draw[dashed, mypen!55] (0,0.8) -- (4.3,0.8);
            \draw[dashed, notered!70] (4.0,0) -- (4.0,3.3) node[above, font=\scriptsize, notered] {asintoto};
            \node[left, font=\scriptsize] at (0,0.8) {$1$};
            \node[below, font=\scriptsize] at (4.0,-0.05) {$1$};
            \draw[very thick, boxoss, domain=0:0.968, samples=160, smooth, variable=\t]
                plot ({\t*4.0}, {0.8/sqrt(1-\t*\t)});
        \end{scope}
        \begin{scope}[xshift=7.0cm]
            \draw[->, thick, mypen] (0,0) -- (4.6,0) node[right, font=\small] {$v/c$};
            \draw[->, thick, mypen] (0,0) -- (0,3.4) node[above, font=\small] {$1/\gamma$};
            \draw[dashed, mypen!55] (0,2.6) -- (4.0,2.6);
            \node[left, font=\scriptsize] at (0,2.6) {$1$};
            \node[below, font=\scriptsize] at (4.0,-0.05) {$1$};
            \draw[very thick, boxese, domain=0:1, samples=160, smooth, variable=\t]
                plot ({\t*4.0}, {2.6*sqrt(max(1-\t*\t,0))});
        \end{scope}
    \end{tikzpicture}
    \caption{Andamento del fattore di Lorentz $\gamma$ e del suo reciproco $1/\gamma$ in funzione di $\beta = v/c$. Il fattore $\gamma$ diverge per $v \to c$, mentre $1/\gamma$ (il fattore di contrazione) si annulla nello stesso limite.}
\end{figure}

\begin{definizione}{Intervallo Spazio-Temporale}
Si definisce \textbf{intervallo spazio-temporale} la quantità
\begin{equation}
    \Delta \sigma^2 = c^2 \Delta t^2 - \Delta x^2 - \Delta y^2 - \Delta z^2
\end{equation}
Essa è \textbf{invariante} sotto trasformazioni di Lorentz: pur non essendo più assoluti né il tempo né lo spazio separatamente, esiste una loro combinazione che tutti gli osservatori inerziali misurano identica.
\end{definizione}

\section{Conseguenze Cinematiche}

\subsection{Contrazione delle Lunghezze}

Consideriamo un oggetto che ha lunghezza $L_0$ in $S'$ e lunghezza $L$ in $S$, con
\begin{equation*}
    L_0 = x'_B - x'_A
    \qquad\qquad
    L = x_B - x_A
\end{equation*}
Sostituendo le trasformazioni di Lorentz per $x'_B$ e $x'_A$:
\begin{align*}
    x'_B = \gamma(x_B - vt_B) \, , \qquad & x'_A = \gamma(x_A - vt_A) \\
    x'_B - x'_A = \gamma\big(x_B - x_A - v(t_B - t_A)\big) \implies & \; L_0 = \gamma\left(L - v\Delta t\right)
\end{align*}
Ma $\Delta t = 0$, perché le posizioni $x_B$ e $x_A$ vengono misurate nel medesimo istante, per cui
\begin{equation}
    \boxed{\;L = \frac{L_0}{\gamma}\;}
\end{equation}

\begin{definizione}{Lunghezza Propria e Lunghezza Contratta}
\begin{itemize}[leftmargin=1.4em]
    \item La lunghezza del corpo misurata nel sistema \textbf{in quiete rispetto al corpo} ($S'$) si dice \textbf{lunghezza propria} $L_0$: essa è definibile univocamente.
    \item La lunghezza del corpo misurata nel sistema \textbf{in moto rispetto al corpo} ($S$) si dice \textbf{lunghezza contratta} $L$.
\end{itemize}
Vale sempre $L \leq L_0$: la lunghezza propria è la massima misurabile.
\end{definizione}

\begin{attenzione}{Solo le Lunghezze Parallele al Boost si Contraggono}
La contrazione riguarda esclusivamente la direzione parallela al boost; per quelle ortogonali vale
\begin{equation*}
    L^{\perp} = L_0^{\perp}
\end{equation*}
\end{attenzione}

\subsection{Dilatazione dei Tempi}

Consideriamo un evento $E$ che in $S'$ ha luogo nello stesso posto, mentre in $S$ ha luogo in due posizioni diverse: $E(A) = $ inizio ed $E(B) = $ fine. Detto $T_0 = t'_B - t'_A$ l'intervallo misurato in $S'$ e $T = t_B - t_A$ quello misurato in $S$, sostituiamo le trasformazioni di Lorentz:
\begin{align*}
    t_B = \gamma\left(t'_B + \frac{v}{c^2}x'_B\right) \, , \qquad & t_A = \gamma\left(t'_A + \frac{v}{c^2}x'_A\right) \\
    t_B - t_A = \gamma\left(t'_B - t'_A + \frac{v}{c^2}(x'_B - x'_A)\right) \implies & \; T = \gamma\left(T_0 + \frac{v}{c^2}\Delta x'\right)
\end{align*}
Ma $\Delta x' = 0$, perché $x'_A = x'_B$ nel momento in cui l'evento comincia e finisce, per cui
\begin{equation}
    \boxed{\;T = \gamma \, T_0\;}
\end{equation}

\begin{definizione}{Tempo Proprio e Tempo Dilatato}
\begin{itemize}[leftmargin=1.4em]
    \item Il tempo misurato nel sistema \textbf{in quiete rispetto all'evento} ($S'$) si dice \textbf{tempo proprio} $T_0$: esso è definibile univocamente.
    \item Il tempo misurato nel sistema \textbf{in moto rispetto all'evento} ($S$) si dice \textbf{tempo dilatato} $T$.
\end{itemize}
Vale sempre $T \geq T_0$: il tempo proprio è il minimo misurabile.
\end{definizione}

\begin{osservazione}{Dualità fra le Due Conseguenze}
Contrazione delle lunghezze e dilatazione dei tempi sono \textbf{due facce dello stesso fenomeno}: entrambe discendono dalle trasformazioni di Lorentz, ed entrambe si ottengono annullando l'intervallo coniugato ($\Delta t = 0$ per le lunghezze, $\Delta x' = 0$ per i tempi). La grandezza propria --- quella misurata in quiete rispetto all'oggetto della misura --- è sempre l'unica definibile univocamente.
\end{osservazione}

\subsection{Verifica Sperimentale: il Decadimento del Muone}

\begin{definizione}{Muone}
Il \textbf{muone} $\mu$ è un leptone instabile che decade secondo
\begin{equation}
    \mu \longrightarrow e^{+} + \nu_e + \bar{\nu}_{\mu}
\end{equation}
dove $e^{+}$ è il positrone, $\nu_e$ l'elettrone neutrino e $\bar{\nu}_{\mu}$ il muone antineutrino.
\end{definizione}

Il decadimento è un processo statistico governato da due eventi:
\begin{itemize}
    \item \textbf{I evento:} creazione, con $N_0$ muoni prodotti;
    \item \textbf{II evento:} decadimento, con legge $N(t) = N_0 \, e^{-t/\tau}$.
\end{itemize}
La probabilità che il muone \textbf{non} sia ancora decaduto all'istante $t$ è $P = e^{-t/\tau}$, dove $\tau$ è la vita media. Distinguiamo due intervalli:
\begin{enumerate}
    \item vita media \textbf{a riposo}: $\tau_0 = 2{,}2 \, \mu\mathrm{s}$ (sistema di quiete);
    \item vita media \textbf{in moto}: $\tau_{\text{LAB}} = \gamma \, \tau_0$ (sistema in moto con velocità $\vec{v}$).
\end{enumerate}

\begin{figure}[ht]
    \centering
    \begin{tikzpicture}[>=stealth, scale=0.92]
        % --- Sistema LAB
        \begin{scope}
            \node[font=\small\bfseries, mypen] at (-0.7,2.4) {LAB};
            \draw[thick, fill=mypen!8] (0.4,0.2) rectangle (1.3,2.1);
            \node[above, font=\scriptsize\bfseries] at (0.85,2.15) {$B$};
            \draw[thick, fill=mypen!8] (3.7,0.2) rectangle (4.6,2.1);
            \node[above, font=\scriptsize\bfseries] at (4.15,2.15) {$R$};
            \draw[->, thick, notered] (-0.55,1.15) -- (0.3,1.15) node[midway, above, font=\scriptsize] {$P$};
            \foreach \x/\y in {1.7/1.6, 2.3/1.1, 2.9/1.7, 2.1/0.7, 3.2/1.0} {
                \fill[boxese] (\x,\y) circle (1.3pt);
                \node[font=\tiny, boxese] at (\x+0.16,\y+0.16) {$\mu$};
            }
            \draw[->, thick, boxatt] (1.9,0.55) -- (3.1,0.55) node[midway, above, font=\scriptsize] {$\vec{v}$};
            \draw[<->, thick, mypen] (1.3,-0.45) -- (3.7,-0.45);
            \node[font=\small, mypen] at (2.5,-0.8) {$L_0$};
            \node[font=\scriptsize, align=center] at (2.5,-1.55) {estremi fermi:\\$t_{\text{LAB}} = L_0/v$};
        \end{scope}
        % --- Sistema di quiete
        \begin{scope}[xshift=7.4cm]
            \node[font=\small\bfseries, mypen] at (-0.7,2.4) {QUIETE};
            \draw[thick, fill=mypen!8] (0.9,0.2) rectangle (1.8,2.1);
            \node[above, font=\scriptsize\bfseries] at (1.35,2.15) {$B$};
            \draw[thick, fill=mypen!8] (3.2,0.2) rectangle (4.1,2.1);
            \node[above, font=\scriptsize\bfseries] at (3.65,2.15) {$R$};
            \foreach \x/\y in {2.1/1.6, 2.5/1.1, 2.8/1.7, 2.3/0.7} {
                \fill[boxese] (\x,\y) circle (1.3pt);
                \node[font=\tiny, boxese] at (\x+0.16,\y+0.16) {$\mu$};
            }
            \draw[->, thick, boxatt] (1.8,2.45) -- (0.9,2.45) node[midway, above, font=\scriptsize] {$-\vec{v}$};
            \draw[->, thick, boxatt] (4.1,2.45) -- (3.2,2.45);
            \draw[<->, thick, mypen] (1.8,-0.45) -- (3.2,-0.45);
            \node[font=\small, mypen] at (2.5,-0.8) {$L$};
            \node[font=\scriptsize, align=center] at (2.5,-1.55) {lunghezza contratta:\\$\tau_0' = L/v = L_0/(\gamma v)$};
        \end{scope}
    \end{tikzpicture}
    \caption{L'esperimento del muone letto nei due sistemi di riferimento. I protoni $P$ colpiscono il bersaglio $B$ producendo muoni, che percorrono la distanza fino al rivelatore $R$. Nel LAB la distanza è propria ($L_0$) e il tempo dilatato; nel sistema di quiete del muone la distanza è contratta e il tempo è proprio.}
\end{figure}

\begin{esempio}{Le Due Letture del Medesimo Conteggio}
Indichiamo con $P$ i protoni incidenti, $B$ il bersaglio, $R$ il rivelatore, $\vec{v}$ la velocità relativa fra LAB e sistema di quiete, e $\mu$ i muoni prodotti (con $N(t) = N_0$ all'istante iniziale).

\textbf{(1) Nel sistema LAB} la distanza $L_0$ è misurata fra estremi fermi (lunghezza propria) e il tempo di volo è $t_{\text{LAB}} = L_0/v$; la vita media è però dilatata, $\tau_{\text{LAB}} = \gamma \tau_0$:
\begin{equation}
    N(t)\big|_{\text{LAB}} = N_0 \, e^{-\frac{t_{\text{LAB}}}{\tau_{\text{LAB}}}}
    = N_0 \, e^{-\frac{L_0}{v} \cdot \frac{1}{\gamma \tau_0}}
\end{equation}

\textbf{(2) Nel sistema di quiete del muone} la vita media è quella propria $\tau_0$, ma la distanza appare contratta, $L = L_0/\gamma$, e dunque il tempo di volo è $t_0 = L/v = L_0/(\gamma v)$:
\begin{equation}
    N(t)\big|_{\text{quiete}} = N_0 \, e^{-\frac{t_0}{\tau_0}}
    = N_0 \, e^{-\frac{L_0}{\gamma v} \cdot \frac{1}{\tau_0}}
\end{equation}

I due conteggi coincidono:
\begin{equation}
    \boxed{\;N(t)\big|_{\text{LAB}} = N(t)\big|_{\text{quiete}} = N_0 \, e^{-\frac{L_0}{\gamma v}\cdot\frac{1}{\tau_0}}\;}
\end{equation}
I due osservatori attribuiscono la sopravvivenza dei muoni a cause diverse --- l'uno alla dilatazione del tempo, l'altro alla contrazione dello spazio --- ma \textbf{misurano lo stesso numero di particelle}, come deve essere per un fatto oggettivo.
\end{esempio}

\subsection{Causalità e Successione Temporale}

Due eventi $E(1)$ ed $E(2)$ avvengono in $S$ nell'ordine $t_1 < t_2$, cioè $t_2 - t_1 > 0$, e sono connessi da un segnale che viaggia a velocità $u \leq c$. Vale quindi
\begin{equation}
    x_2 - x_1 = u\,(t_2 - t_1) \leq c\,(t_2-t_1) = c\,\Delta t
    \qquad \implies \qquad \Delta x \leq c\,\Delta t
\end{equation}

Vediamo ora quanto vale $\Delta t'$ in un altro sistema inerziale:
\begin{equation}
    \Delta t' = t'_2 - t'_1 = \gamma\left(t_2 - t_1 - \frac{v}{c^2}x_2 + \frac{v}{c^2}x_1\right)
    = \gamma\left(\Delta t - \frac{v}{c^2}\Delta x\right)
\end{equation}

Sappiamo che $\Delta t > 0$: ci chiediamo se $\Delta t'$ possa essere negativo, cioè se l'ordine dei due eventi possa invertirsi. Imponendo $\Delta t' < 0$:
\begin{equation*}
    \gamma\left(\Delta t - \frac{v}{c^2}\Delta x\right) < 0
    \implies \Delta t < \frac{v}{c^2}\Delta x = \frac{v}{c^2}\, u \,\Delta t
    \implies u\,v > c^2
\end{equation*}
il che richiederebbe $u > c$ e $v > c$: \textbf{impossibile}, perché il segnale dovrebbe viaggiare più velocemente della luce.

\begin{teorema}{Conservazione dell'Ordine Causale}
L'ordine temporale di due eventi causalmente connessi \textbf{si mantiene} in ogni sistema di riferimento inerziale: se $t_2 > t_1$, allora $t'_2 > t'_1$ per qualunque osservatore.
\end{teorema}

\subsection{Il Cono di Luce di Minkowski}

\begin{definizione}{Cono di Luce}
Il \textbf{cono di luce di Minkowski} è il cono dello spazio-tempo delimitato dalle linee di universo di un corpo che viaggia a $c$. Esso separa tre regioni:
\begin{itemize}[leftmargin=1.4em]
    \item \textbf{PA --- Passato assoluto} ($t<0$): le posizioni occupate e gli eventi che influenzano il presente.
    \item \textbf{FA --- Futuro assoluto} ($t>0$): le possibili posizioni che il corpo può assumere.
    \item \textbf{PR --- Presente} ($t=0$): la posizione attuale, da cui si va avanti.
\end{itemize}
\end{definizione}

\begin{figure}[ht]
    \centering
    \begin{tikzpicture}[>=stealth, scale=1.0]
        % Regioni
        \fill[boxoss!12] (0,0) -- (2.6,2.6) -- (-2.6,2.6) -- cycle;
        \fill[boxatt!12] (0,0) -- (2.6,-2.6) -- (-2.6,-2.6) -- cycle;
        % Assi
        \draw[->, thick, mypen] (-3.6,0) -- (3.6,0) node[right, font=\small] {$x$};
        \draw[->, thick, mypen] (0,-3.1) -- (0,3.1) node[above, font=\small] {$ct$};
        % Coni di luce
        \draw[very thick, notered] (-2.9,-2.9) -- (2.9,2.9);
        \draw[very thick, notered] (2.9,-2.9) -- (-2.9,2.9);
        \node[right, font=\scriptsize, notered] at (2.55,2.95) {$x = ct$};
        \node[left, font=\scriptsize, notered] at (-2.55,2.95) {$x = -ct$};
        % Ellissi di base
        \draw[dashed, mypen!60] (0,2.6) ellipse (2.6 and 0.42);
        \draw[dashed, mypen!60] (0,-2.6) ellipse (2.6 and 0.42);
        % Etichette
        \node[font=\small\bfseries, boxoss] at (-1.15,1.95) {FA};
        \node[font=\scriptsize, boxoss] at (-1.15,1.55) {futuro assoluto};
        \node[font=\small\bfseries, boxatt] at (0,-1.85) {PA};
        \node[font=\scriptsize, boxatt] at (0,-1.42) {passato assoluto};
        \node[font=\scriptsize, mypen!80, align=center] at (2.55,0.55) {altrove\\(non connesso)};
        \node[font=\scriptsize, mypen!80, align=center] at (-2.55,0.55) {altrove\\(non connesso)};
        \fill[notered] (0,0) circle (2.4pt);
        \node[below right, font=\scriptsize\bfseries, notered] at (0.08,-0.06) {PR};
        % Linea di universo
        \draw[very thick, boxese, dashed] (0,0) .. controls (0.45,1.0) and (0.6,1.8) .. (0.95,2.6);
    \end{tikzpicture}
    \caption{Il cono di luce di Minkowski. Solo gli eventi interni al cono sono causalmente connessi con l'evento nell'origine; la regione esterna (``altrove'') raccoglie gli eventi per raggiungere i quali occorrerebbe una velocità superiore a $c$. La curva tratteggiata in verde è la \emph{linea di universo} di un corpo che accelera restando sempre all'interno del cono.}
\end{figure}

\section{Dinamica Relativistica}

\subsection{Il Quadrivettore Posizione}

\begin{definizione}{Quadrivettore Posizione}
Il \textbf{quadrivettore posizione} individua la posizione di un corpo nello spazio-tempo:
\begin{equation}
    S = (ct,\, x,\, y,\, z) = (ct,\, \vec{x}\,) = (\underbrace{x_0}_{\text{comp. temporale}},\, \underbrace{\vec{x}}_{\text{comp. spaziali}})
\end{equation}
Le coordinate $(x,y,z)$ e il tempo $t$ \textbf{non sono più indipendenti}: $ct$ è la quarta dimensione, e il fattore $c$ le conferisce omogeneità dimensionale con le altre tre (essendo $[c\,t] = \mathrm{m/s} \cdot \mathrm{s} = \mathrm{m}$).
\end{definizione}

\subsection{Trasformazioni di Lorentz in Forma Matriciale}

Comprimendo le quattro trasformazioni di Lorentz per $x$, $y$, $z$ e $t$ possiamo scrivere
\begin{equation}
    (x_0, \vec{x}\,) = \Lambda \, (x'_0, \vec{x}\,')
\end{equation}
dove $\Lambda$ è la \textbf{matrice di Lorentz}, che racchiude le trasformazioni. Partendo da
\begin{equation*}
    \begin{cases}
        ct = \gamma\,(ct' + \beta x') \\
        x = \gamma\,(x' + vt') \\
        y = y' \\
        z = z'
    \end{cases}
    \qquad \text{con } \beta = \frac{v}{c}
\end{equation*}
si ottiene esplicitamente
\begin{equation}
    \begin{pmatrix} x_0 \\ x \\ y \\ z \end{pmatrix}
    =
    \begin{pmatrix} ct \\ x \\ y \\ z \end{pmatrix}
    = \Lambda
    \begin{pmatrix} x'_0 \\ x' \\ y' \\ z' \end{pmatrix}
    =
    \begin{pmatrix}
        \gamma c t' + \beta\gamma x' \\
        \gamma v t' + \gamma x' \\
        y' \\ z'
    \end{pmatrix}
\end{equation}

\begin{formulario}{Matrici di Lorentz}
\begin{minipage}[t]{0.47\textwidth}
\centering \textbf{Da $S'$ a $S$}
\begin{equation*}
    \Lambda =
    \begin{pmatrix}
        \gamma & \beta\gamma & 0 & 0 \\
        \beta\gamma & \gamma & 0 & 0 \\
        0 & 0 & 1 & 0 \\
        0 & 0 & 0 & 1
    \end{pmatrix}
\end{equation*}
\end{minipage}
\hfill
\begin{minipage}[t]{0.47\textwidth}
\centering \textbf{Da $S$ a $S'$}
\begin{equation*}
    \Lambda^{-1} =
    \begin{pmatrix}
        \gamma & -\beta\gamma & 0 & 0 \\
        -\beta\gamma & \gamma & 0 & 0 \\
        0 & 0 & 1 & 0 \\
        0 & 0 & 0 & 1
    \end{pmatrix}
\end{equation*}
\end{minipage}
\vspace{4pt}
\end{formulario}

\subsection{Il Quadriprodotto Scalare}

Dati due quadrivettori $A = (A_0, \vec{A}\,) = (A_0, A_x, A_y, A_z)$ e $B = (B_0, \vec{B}\,) = (B_0, B_x, B_y, B_z)$, il loro prodotto scalare si costruisce con la \textbf{metrica di Minkowski}:
\begin{equation}
    A \cdot B =
    \begin{pmatrix} A_0 & A_x & A_y & A_z \end{pmatrix}
    \begin{pmatrix}
        1 & 0 & 0 & 0 \\
        0 & -1 & 0 & 0 \\
        0 & 0 & -1 & 0 \\
        0 & 0 & 0 & -1
    \end{pmatrix}
    \begin{pmatrix} B_0 \\ B_x \\ B_y \\ B_z \end{pmatrix}
    = A_0 B_0 - \vec{A} \cdot \vec{B}
    \qquad \text{in } S
\end{equation}
e analogamente $A' \cdot B' = A'_0 B'_0 - \vec{A}\,' \cdot \vec{B}\,'$ in $S'$.

\begin{teorema}{Invarianza del Quadriprodotto Scalare}
Il quadriprodotto scalare è \textbf{invariante sotto trasformazioni di Lorentz}:
\begin{equation}
    A \cdot B = A' \cdot B'
\end{equation}
\end{teorema}

\emph{Dimostrazione.} Applicando le trasformazioni al prodotto, e osservando che $A_y B_y = A'_y B'_y$ e $A_z B_z = A'_z B'_z$ (le componenti trasverse non cambiano), basta tenere conto di $A_0 B_0$ e $A_x B_x$:
\begin{align*}
    A \cdot B &= A_0 B_0 - A_x B_x - \cancel{A_y B_y} - \cancel{A_z B_z} \\
    &= \gamma^2 (A'_0 + \beta A'_x)(B'_0 + \beta B'_x) - \gamma^2 (A'_x + \beta A'_0)(B'_x + \beta B'_0) \\
    &= \gamma^2 \Big( A'_0 B'_0 + \cancel{\beta A'_0 B'_x} + \cancel{\beta B'_0 A'_x} + \beta^2 A'_x B'_x
       - A'_x B'_x - \cancel{\beta B'_0 A'_x} - \cancel{\beta A'_0 B'_x} - \beta^2 A'_0 B'_0 \Big) \\
    &= \gamma^2 \left[ A'_0 B'_0 (1-\beta^2) - A'_x B'_x (1-\beta^2) \right]
     = \gamma^2 (1-\beta^2)\,(A'_0 B'_0 - A'_x B'_x)
\end{align*}
Dato che $\gamma^2 = \dfrac{1}{1-\beta^2}$, si ha $\gamma^2(1-\beta^2) = 1$ e dunque
\begin{equation}
    \boxed{\;A_0 B_0 - A_x B_x = A'_0 B'_0 - A'_x B'_x\;}
\end{equation}
$\square$

\subsection{Norma Quadra e Categorie di Quadrivettori}

\begin{definizione}{Norma Quadra}
La \textbf{norma quadra} di un quadrivettore (detta anche \emph{intervallo spazio-temporale di Minkowski}) è il quadriprodotto scalare del quadrivettore con sé stesso:
\begin{equation}
    A \cdot A = A^2 = A_0 A_0 - \vec{A} \cdot \vec{A} = A_0^2 - |\vec{A}|^2
    \qquad \text{con } |\vec{A}|^2 = A_x^2 + A_y^2 + A_z^2
\end{equation}
Con lo stesso calcolo svolto per il quadriprodotto si dimostra che essa è invariante:
\begin{equation}
    A_0^2 - A_x^2 = A_0'^2 - A_x'^2
\end{equation}
\end{definizione}

La norma quadra, essendo invariante, permette di classificare i quadrivettori in tre \textbf{categorie}:
\begin{equation*}
    A^2 > 0 \; \to \; \textbf{tempo}
    \qquad\qquad
    A^2 < 0 \; \to \; \textbf{spazio}
    \qquad\qquad
    A^2 = 0 \; \to \; \textbf{luce}
\end{equation*}

I quattro quadrivettori fondamentali della meccanica relativistica sono quello di \textbf{posizione}, di \textbf{luce}, di \textbf{velocità} e di \textbf{energia--impulso}. Per ciascuno si confronta la norma nel sistema di quiete del corpo ($A'$) con quella nel sistema in moto ($A$).

\paragraph{Quadrivettore posizione: $S = (ct, x, y, z) = (x_0, \vec{x}\,)$.}
Nel sistema in quiete rispetto al corpo si ha $A' = (A'_0, 0, 0, 0)$; nel sistema in moto, ricordando $t = \gamma t'$,
\begin{equation*}
    A = (A_0,\, -\gamma\beta c t',\, 0,\, 0) = (A_0,\, -\gamma v t',\, 0,\, 0) = (A_0,\, -vt,\, 0,\, 0)
\end{equation*}
Quindi $A'^2 = A_0'^2$ e
\begin{equation}
    A^2 = A_0^2 - (vt)^2 = c^2t^2 - v^2t^2 = (c^2-v^2)\,t^2
\end{equation}
Dato che $v < c$ si ha $c^2 - v^2 > 0$, dunque $A^2 > 0$: \textbf{categoria tempo}.

\paragraph{Quadrivettore luce.}
La luce non ha un sistema di quiete: si muove costantemente a $c$ in ogni sistema inerziale. Il quadrivettore luce nel sistema in moto ha quindi
\begin{equation}
    A^2 = A_0^2 - A_x^2 = c^2t^2 - c^2t^2 = 0
\end{equation}
poiché $c = c'$. Si ha $A^2 = 0$: \textbf{categoria luce}.

\paragraph{Quadrivettore velocità.}
Esso si scrive $V = \gamma\,(c, v_x, v_y, v_z) = \gamma\,(c, \vec{v}\,)$.
Nel sistema di quiete $A' = \left(\frac{d(ct')}{dt'}, 0, 0, 0\right) = (c, 0, 0, 0)$; nel sistema in moto $A = \gamma\,(c, v, 0, 0)$. Quindi $A'^2 = c^2$ e
\begin{equation}
    A^2 = \gamma^2 c^2 - \gamma^2 v^2 = \gamma^2 (c^2 - v^2)
        = (c^2-v^2) \cdot \frac{c^2}{c^2-v^2} = c^2
\end{equation}
Dato che $c^2 > 0$, si ha $A^2 > 0$: \textbf{categoria tempo}.

\paragraph{Quadrivettore energia--impulso: $W = \gamma m\,(c, v_x, v_y, v_z) = \left(\dfrac{E}{c}, \vec{p}\,\right)$.}
Nel sistema di quiete $A' = (mc, 0, 0, 0)$; nel sistema in moto $A = (\gamma mc, \gamma mv, 0, 0)$. Quindi $A'^2 = m^2c^2$ e
\begin{equation}
    A^2 = \gamma^2 m^2 c^2 - \gamma^2 m^2 v^2 = \gamma^2 m^2 (c^2-v^2)
        = m^2 (c^2 - v^2) \cdot \frac{c^2}{c^2-v^2} = m^2 c^2
\end{equation}
Dato che $m^2c^2 > 0$, si ha $A^2 > 0$: \textbf{categoria tempo}.

\begin{attenzione}{Invarianza $\neq$ Conservazione}
Due concetti da non confondere mai:
\begin{itemize}[leftmargin=1.4em]
    \item \textbf{Invarianza:} stesso valore in \emph{sistemi di riferimento} diversi.
    \item \textbf{Conservazione:} stesso valore in \emph{istanti} diversi.
\end{itemize}
Una grandezza può essere invariante e non conservata, conservata e non invariante, oppure entrambe le cose (come il quadrimpulso totale in un urto isolato).
\end{attenzione}

\subsection{Energia Relativistica Totale}

Sappiamo che il quadrivettore energia--impulso vale $W = \left(\frac{E}{c}, \vec{p}\,\right)$ e che la sua norma quadra è $W^2 = \left(\frac{E}{c}\right)^2 - \vec{p}^{\;2} = (mc)^2$. Possiamo dunque ricavarci
\begin{equation*}
    \left(\frac{E}{c}\right)^{\!2} = (mc)^2 + \vec{p}^{\;2}
    \implies E^2 = m^2c^4 + \vec{p}^{\;2}c^2
    \implies E^2 = c^2\left(m^2c^2 + \vec{p}^{\;2}\right)
\end{equation*}
\begin{equation}
    \boxed{\;E = \sqrt{m^2c^4 + \vec{p}^{\;2}c^2}\;}
\end{equation}

Esiste anche un secondo modo per esprimere l'energia relativistica, sfruttando la definizione relativistica di quantità di moto $\vec{p} = \gamma m \vec{v}$:
\begin{align*}
    E &= \sqrt{m^2c^4 + \vec{p}^{\;2}c^2} = \sqrt{m^2c^4 + \gamma^2 m^2 v^2 c^2}
    \implies E^2 = m^2c^2\left(c^2 + \gamma^2 v^2\right) \\
    E^2 &= m^2c^2\left(c^2 + \frac{c^2v^2}{c^2-v^2}\right)
       = m^2c^2\left(\frac{c^4 - \cancel{c^2v^2} + \cancel{c^2v^2}}{c^2-v^2}\right)
       = m^2c^4\left(\frac{c^2}{c^2-v^2}\right) = \gamma^2 m^2 c^4
\end{align*}
da cui la celebre relazione

\begin{formulario}{Energia Relativistica Totale}
\begin{equation}
    E = \gamma\, m c^2
    \qquad\qquad\qquad
    E = \sqrt{m^2c^4 + \vec{p}^{\;2}c^2}
\end{equation}
Le due espressioni sono equivalenti: la prima è comoda quando si conosce la velocità, la seconda quando si conosce l'impulso (ed è l'unica utilizzabile per le particelle di massa nulla).
\end{formulario}

\begin{osservazione}{Conservazione del Quadrimpulso negli Urti}
Se il quadrimpulso si conserva in $S$, allora si conserva anche in $S'$. Per una reazione $A + B \to C + D$:
\begin{equation}
    \begin{pmatrix} E_A/c \\ \vec{p}_A \end{pmatrix} + \begin{pmatrix} E_B/c \\ \vec{p}_B \end{pmatrix}
    \longrightarrow
    \begin{pmatrix} E_C/c \\ \vec{p}_C \end{pmatrix} + \begin{pmatrix} E_D/c \\ \vec{p}_D \end{pmatrix}
\end{equation}
Si separano \textbf{sempre} la componente temporale $(E_i/c)$ da quelle spaziali $(\vec{p}_i)$, ottenendo due equazioni indipendenti:
\begin{equation}
    \begin{cases}
        \dfrac{E_A}{c} + \dfrac{E_B}{c} = \dfrac{E_C}{c} + \dfrac{E_D}{c} \\[8pt]
        \vec{p}_A + \vec{p}_B = \vec{p}_C + \vec{p}_D
    \end{cases}
\end{equation}
Applicando a ciascun termine la trasformazione di Lorentz $\frac{E_i}{c} = \gamma\left(\frac{E'_i}{c} + \beta p'_{i,x}\right)$ e sommando, i fattori $\gamma$ si semplificano da entrambi i lati e si ritrova la medesima legge di conservazione in $S'$.
\end{osservazione}

\subsection{Energia Cinetica Relativistica e Difetto di Massa}

\begin{definizione}{Energia Cinetica Relativistica}
L'energia cinetica relativistica è la differenza fra l'energia totale e l'energia a riposo:
\begin{equation}
    T = K_{\text{rel}} = E - E_0 = \gamma m c^2 - mc^2 = (\gamma - 1)\,mc^2
\end{equation}
dove $E$ è l'energia totale ed $E_0 = mc^2$ l'\textbf{energia a riposo}, posseduta da ogni corpo per il solo fatto di avere massa.
\end{definizione}

\begin{esempio}{La Fusione Nucleare nel Sole}
La catena protone--protone che alimenta il Sole si riassume nella reazione
\begin{equation}
    4p \longrightarrow {}^{4}\mathrm{He}^{2+} + 2e^{+} + 2\nu_e
\end{equation}
dove $p$ è il protone, ${}^{4}\mathrm{He}^{2+}$ il nucleo di elio (elio ionizzato), $e^+$ il positrone e $\nu_e$ l'elettrone neutrino.

Misurando le masse si osserva che la quantità di prodotto è \textbf{minore} di quella dei reagenti:
\begin{equation*}
    m(4p) > m({}^{4}\mathrm{He}^{2+}) + 2m(e^+) + 2m(\nu_e)
\end{equation*}
Questo accade perché una parte di $m(4p)$ viene trasformata in energia: è il \textbf{difetto di massa}.
\end{esempio}

\begin{definizione}{Difetto di Massa}
Il difetto di massa è dovuto a una trasformazione di ``energia di massa'' (quella che la massa ha in quanto tale) in ``energia cinetica'' (quella che la massa ha quando si muove). L'energia legata alla massa a riposo è l'\textbf{energia a riposo}:
\begin{equation}
    E_0 = mc^2
    \qquad \Longleftrightarrow \qquad
    \Delta m = \frac{\Delta E}{c^2}
\end{equation}
\end{definizione}

\subsection{Particelle senza Massa}

Per la meccanica classica non ha senso parlare di corpi con $m = 0$, poiché si avrebbe $E = U + \frac{1}{2}mv^2 = 0$. Nella meccanica relativistica, invece, se $m = 0$ allora
\begin{equation}
    E = \gamma m c^2 = \frac{mc^2}{\sqrt{1 - \dfrac{v^2}{c^2}}} =
    \begin{cases}
        0 & \text{se } v < c \\
        \text{indeterminata} & \text{se } v = c
    \end{cases}
\end{equation}
L'indeterminazione si risolve usando l'altra espressione dell'energia:
\begin{equation}
    E = \sqrt{\cancelto{0}{m^2c^4} + \vec{p}^{\;2}c^2} = |\vec{p}\,|\,c
\end{equation}
per cui il quadrimpulso di una particella massless vale
\begin{equation}
    W_{\gamma} = \begin{pmatrix} E/c \\ \vec{p} \end{pmatrix} = \begin{pmatrix} |\vec{p}\,| \\ \vec{p} \end{pmatrix}
\end{equation}
Il \textbf{fotone} --- il quanto di luce --- è l'esempio fondamentale di particella senza massa.

\subsection{Unità di Misura in Meccanica Relativistica}

Si adottano unità di misura particolari per rendere più semplici i calcoli e le formule.

\begin{definizione}{Unità Naturali}
Si pone
\begin{equation}
    c = 1 \qquad (= 2{,}998 \cdot 10^8 \, \mathrm{m/s})
\end{equation}
In queste unità le formule si alleggeriscono notevolmente; per esempio
\begin{equation*}
    \frac{E}{c} \to E \, , \qquad \frac{E}{c^2} \to E \, , \qquad
    E = \sqrt{m^2 + \vec{p}^{\;2}} \, , \qquad E = \gamma m
\end{equation*}
\end{definizione}

\begin{definizione}{Elettronvolt}
L'\textbf{elettronvolt} è l'energia acquistata da una particella di carica $e$ quando attraversa una differenza di potenziale di $1 \, \mathrm{V}$:
\begin{equation}
    1 \, \mathrm{eV} = 1{,}6 \cdot 10^{-19} \, \mathrm{J}
\end{equation}
\end{definizione}

\begin{center}
\renewcommand{\arraystretch}{1.2}
\begin{minipage}[t]{0.42\textwidth}
\centering
\textbf{Multipli e sottomultipli}\\[4pt]
\begin{tabular}{@{}ll@{}}
\toprule
$\mathrm{meV}$ & $= 10^{-3} \, \mathrm{eV}$ \\
$\mathrm{keV}$ & $= 10^{3} \, \mathrm{eV}$ \\
$\mathrm{MeV}$ & $= 10^{6} \, \mathrm{eV}$ \\
$\mathrm{GeV}$ & $= 10^{9} \, \mathrm{eV}$ \\
$\mathrm{TeV}$ & $= 10^{12} \, \mathrm{eV}$ \\
\bottomrule
\end{tabular}
\end{minipage}
\hfill
\begin{minipage}[t]{0.52\textwidth}
\centering
\textbf{Masse di alcune particelle}\\[4pt]
\begin{tabular}{@{}ll@{}}
\toprule
Bosone $Z^0$ & $m_{Z^0} = 91 \, \mathrm{GeV}/c^2$ \\
Protone & $m_p = 938 \, \mathrm{MeV}/c^2$ \\
Neutrone & $m_n = 939{,}57 \, \mathrm{MeV}/c^2$ \\
Elettrone & $m_e = 0{,}511 \, \mathrm{MeV}/c^2$ \\
\bottomrule
\end{tabular}
\end{minipage}
\end{center}

\section{Dinamica degli Urti Relativistici}

\subsection{Confronto fra Urti Classici e Relativistici}

\begin{center}
\renewcommand{\arraystretch}{1.35}
\begin{tabular}{@{}p{0.46\textwidth} p{0.46\textwidth}@{}}
\toprule
\textbf{Meccanica classica} & \textbf{Meccanica relativistica} \\
\midrule
\textbf{Quantità di moto} \newline
$\vec{p}_A + \vec{p}_B = \vec{p}_C + \vec{p}_D$ \quad \checkmark
&
\textbf{Quantità di moto} \newline
$\vec{p}_A + \vec{p}_B = \vec{p}_C + \vec{p}_D$ \quad \checkmark \\
\midrule
\textbf{Quantità di massa} \newline
$m_A + m_B = m_C + m_D$ \quad \checkmark
&
\textbf{Energia (totale)} \newline
$E_A + E_B = E_C + E_D$ \quad \checkmark \\
\midrule
\textbf{Energia (cinetica)} \newline
a) $K_A + K_B = K_C + K_D$ \quad \checkmark \newline
\hspace*{1em} $\Rightarrow$ urto elastico \newline
b) $K_A + K_B \neq K_C + K_D$ \quad \ding{55} \newline
\hspace*{1em} $\Rightarrow$ urto anelastico \newline
\hspace*{1em} (parzialmente o totalmente)
&
\textbf{Quantità di massa ed energia cinetica} \newline
a) $m_A + m_B = m_C + m_D$ \; e \newline
\hspace*{1em} $K_A + K_B = K_C + K_D$ \quad \checkmark \newline
\hspace*{1em} $\Rightarrow$ urto elastico \newline
b) $m_A + m_B \neq m_C + m_D$ \; e \newline
\hspace*{1em} $K_A + K_B \neq K_C + K_D$ \quad \ding{55} \newline
\hspace*{1em} $\Rightarrow$ urto anelastico \\
\bottomrule
\end{tabular}
\end{center}

\begin{attenzione}{Che Cosa Cambia Davvero}
Nella meccanica relativistica la massa \textbf{non} si conserva separatamente: a conservarsi è l'\emph{energia totale}, di cui l'energia a riposo $mc^2$ è una componente. Un urto anelastico relativistico non ``perde'' energia, la converte in massa (o viceversa).
\end{attenzione}

\begin{definizione}{Sistemi di Riferimento per gli Urti}
I moti possono essere studiati in due sistemi privilegiati:
\begin{itemize}[leftmargin=1.4em]
    \item \textbf{Laboratorio (LAB):} il bersaglio è fermo (\emph{target}).
    \item \textbf{Centro di massa (CM):} il CM è fermo e i due corpi si muovono con lo stesso $|\vec{p}\,|$ (\emph{collider}); oppure il CM si muove e i due corpi hanno $\vec{p}$ diversi (\emph{collider asimmetrico}).
\end{itemize}
\end{definizione}

\subsection{Gli Invarianti di Mandelstam}

\paragraph{Invariante $s$ (per quadrivettori di tipo tempo).}
\begin{equation}
    s = (W_A + W_B)^2
\end{equation}

Nel \textbf{LAB} il bersaglio $B$ è fermo, per cui (in unità naturali)
\begin{equation*}
    W_A + W_B \big|_{\text{LAB}} = \begin{pmatrix} E_A + E_B \\ \vec{p}_A + \vec{p}_B \end{pmatrix}
    = \begin{pmatrix} E_A + m_B \\ \vec{p}_A \end{pmatrix}
\end{equation*}
\begin{equation}
    \implies s = E_A^2 + m_B^2 + 2E_A m_B - \vec{p}_A^{\;2}
    = m_A^2 + m_B^2 + 2E_A m_B
\end{equation}
avendo usato $E^2 - p^2 = m^2$.

Nel \textbf{CM} si ha $\vec{p}_A = -\vec{p}_B$, per cui
\begin{equation*}
    W_A + W_B \big|_{\text{CM}} = \begin{pmatrix} E_A + E_B \\ \vec{p}_A + \vec{p}_B \end{pmatrix}
    = \begin{pmatrix} E_A + E_B \\ \vec{0} \end{pmatrix}
    \implies s = E_A^2 + E_B^2 + 2E_A E_B = E_{\text{CM}}^2
\end{equation*}
Infatti, sviluppando esplicitamente:
\begin{align*}
    s &= E_A^2 + E_B^2 + 2E_AE_B - \vec{p}_A^{\;2} - \vec{p}_B^{\;2} - 2\vec{p}_A \cdot \vec{p}_B \\
      &= m_A^2 + m_B^2 + 2E_AE_B - |\vec{p}\,|^2 \cdot 2\,(\cos \pi) \\
      &= m_A^2 + m_B^2 + 2p^2 + 2E_AE_B = E_A^2 + E_B^2 + 2E_AE_B = (E_A + E_B)^2 = E_{\text{CM}}^2
\end{align*}

\begin{osservazione}{Il Significato di $s$}
L'invariante $s$ è il \textbf{quadrato dell'energia totale disponibile nel centro di massa}: è la grandezza che determina quali particelle possono essere prodotte in una reazione. Calcolandolo nel LAB (dove si conoscono i dati sperimentali) e uguagliandolo al valore nel CM (dove si conosce lo stato finale) si risolvono i problemi di soglia senza mai dover cambiare esplicitamente sistema.
\end{osservazione}

\paragraph{Invariante $t$ (per quadrivettori di tipo spazio e tempo).}
\begin{equation}
    t = (W_A - W_C)^2
\end{equation}
Nel CM si ha $W_A = \begin{pmatrix} E_A \\ \vec{p}_A \end{pmatrix}$, $W_B = \begin{pmatrix} E_B \\ -\vec{p}_A \end{pmatrix}$, $W_C = \begin{pmatrix} E_C \\ \vec{p}_C \end{pmatrix}$, $W_D = \begin{pmatrix} E_D \\ -\vec{p}_C \end{pmatrix}$; dalla conservazione $E_A + E_B = E_C + E_D$ segue $2E_A = 2E_C$ per particelle identiche. Quindi:
\begin{align}
    t &= E_A^2 + E_C^2 - 2E_AE_C - \vec{p}_A^{\;2} - \vec{p}_C^{\;2} + 2\vec{p}_A \cdot \vec{p}_C \notag \\
      &= m_A^2 + m_C^2 - 2E_AE_C + 2\,p_A p_C \cos\theta
       = 2m_A^2 - 2E_A^2 + 2|\vec{p}_A|^2 \cos\theta \notag \\
      &= 2m_A^2 - \left(2m_A^2 + 2|\vec{p}_A|^2\right) + 2|\vec{p}_A|^2 \cos\theta
       = -2|\vec{p}_A|^2\,(1 - \cos\theta) \notag \\
      &= \boxed{\,-4\,|\vec{p}_A|^2 \sin^2\!\left(\frac{\theta}{2}\right)}
\end{align}
avendo usato l'identità $1 - \cos\theta = 2\sin^2(\theta/2)$. L'invariante $t$, essendo negativo, è un quadrivettore di \textbf{tipo spazio}: esso misura il \textbf{momento trasferito} nell'urto.

\subsection{Urto Totalmente Anelastico: Classico contro Relativistico}

\begin{esempio}{Confronto Numerico}
Due corpi di masse $m_A = m_B = M = 3{,}0 \cdot 10^4 \, \mathrm{kg}$ urtano in modo totalmente anelastico, con $\vec{v}_A = 50 \, \mathrm{m/s}$ e $\vec{v}_B = 0$.\footnote{Il manoscritto originale riporta $M = 30 \cdot 10^4 \, \mathrm{kg}$; i risultati numerici scritti di seguito nel manoscritto sono però coerenti con $M = 3{,}0 \cdot 10^4 \, \mathrm{kg}$, valore qui adottato.}

\textbf{Nella meccanica classica} la variazione di energia cinetica è
\begin{equation*}
    \Delta K = -\frac{1}{2}\mu \vec{v}_r^{\;2} = -\frac{1}{2}\cdot\frac{M^2}{2M}\vec{v}_r^{\;2}
    = -\frac{M}{4}\vec{v}_r^{\;2} \sim -19 \cdot 10^6 \, \mathrm{J}
\end{equation*}
e \textbf{non} si ha alcuna perdita di massa.

\textbf{Nella meccanica relativistica}, invece, si ha una perdita di massa. Dalla conservazione dell'energia totale:
\begin{equation*}
    \Delta E_{\text{tot}} = \Delta E_c + \Delta E_0 = 0
    \implies \Delta K = -\Delta E_0 = -\gamma \Delta m c^2
\end{equation*}
ma essendo $v \ll c$ si ha $\gamma \to 1$, per cui $\Delta K = -\Delta m c^2$ e
\begin{equation*}
    \Delta m = \frac{|\Delta K|}{c^2} \sim 2 \cdot 10^{-10} \, \mathrm{kg}
\end{equation*}
una massa enormemente più piccola di un granello di sabbia ($\sim 10^{-6} \, \mathrm{kg}$).
\end{esempio}

\begin{osservazione}{Quando il Difetto di Massa Diventa Osservabile}
Se $v \ll c$, allora $\gamma \to 1$ e $\Delta m$ è insignificante e trascurabile: ecco perché la meccanica classica funziona benissimo nella vita quotidiana. Quando invece $v \sim c$ e $\gamma \to +\infty$, $\Delta m$ diventa significativa e misurabile su scala atomica.
\end{osservazione}

\begin{teorema}{Il Quadrimpulso si Conserva Sempre}
Indipendentemente dal tipo di urto, vale
\begin{equation}
    W_A + W_B = W_C + W_D
    \qquad \Longleftrightarrow \qquad
    \begin{cases}
        E_A + E_B = E_C + E_D \\
        \vec{p}_A + \vec{p}_B = \vec{p}_C + \vec{p}_D
    \end{cases}
\end{equation}
con $E = \gamma mc^2 = \sqrt{m^2c^4 + \vec{p}^{\;2}c^2}$ e $\vec{p} = \gamma m \vec{v}$, ossia
\begin{equation}
    \sqrt{\vec{p}_A^{\;2} + m_A^2c^2} + \sqrt{\vec{p}_B^{\;2} + m_B^2c^2}
    = \sqrt{\vec{p}_C^{\;2} + m_C^2c^2} + \sqrt{\vec{p}_D^{\;2} + m_D^2c^2}
\end{equation}
\end{teorema}

Sono inoltre utilissime le due relazioni
\begin{equation}
    \gamma = \frac{E}{mc^2}
    \qquad\qquad
    \beta = \frac{|\vec{v}\,|}{c} \cdot \frac{\gamma m}{\gamma m} = \frac{|\vec{p}\,|}{E/c} = \frac{|\vec{p}\,|\,c}{E}
\end{equation}

\subsection{L'Effetto Compton}

\begin{definizione}{Effetto Compton}
L'\textbf{effetto Compton} è la diffusione di un fotone su un elettrone fermo:
\begin{equation}
    \gamma + e^{-} \longrightarrow e^{-} + \gamma
\end{equation}
con $m_{\gamma} = 0$.
\end{definizione}

\begin{figure}[ht]
    \centering
    \begin{tikzpicture}[>=stealth, scale=1.1]
        % Fotone incidente
        \draw[decorate, decoration={snake, amplitude=1.4pt, segment length=5pt}, very thick, boxatt]
            (-3.0,0) -- (-0.25,0);
        \draw[->, very thick, boxatt] (-0.35,0) -- (-0.15,0);
        \node[above, font=\small, boxatt] at (-1.8,0.18) {$\gamma_i$ \; ($E_{\gamma i}$)};
        % Elettrone bersaglio
        \fill[boxoss] (0,0) circle (2.6pt);
        \node[left, font=\small, boxoss] at (-0.12,-0.62) {$e^{-}_i$ (fermo)};
        % Fotone diffuso
        \draw[decorate, decoration={snake, amplitude=1.4pt, segment length=5pt}, very thick, boxatt]
            (0.2,0.12) -- (2.4,1.55);
        \draw[->, very thick, boxatt] (2.25,1.45) -- (2.5,1.61);
        \node[right, font=\small, boxatt] at (2.5,1.6) {$\gamma_f$ \; ($E_{\gamma f}$)};
        % Elettrone rinculo
        \draw[->, very thick, boxoss] (0.2,-0.12) -- (2.4,-1.35);
        \node[right, font=\small, boxoss] at (2.45,-1.4) {$e^{-}_f$ \; ($E_{ef}$)};
        % Angoli
        \draw[->, thick, mypen] (0.95,0) arc (0:33:0.95);
        \node[font=\small, mypen] at (1.12,0.3) {$\theta_{\gamma}$};
        \draw[->, thick, mypen] (1.25,0) arc (0:-29:1.25);
        \node[font=\small, mypen] at (1.42,-0.36) {$\theta_{e}$};
        % Assi
        \draw[->, thin, mypen!60] (-3.0,-1.9) -- (-2.2,-1.9) node[right, font=\scriptsize] {$\hat{x}$};
        \draw[->, thin, mypen!60] (-3.0,-1.9) -- (-3.0,-1.1) node[above, font=\scriptsize] {$\hat{y}$};
        \draw[dashed, mypen!45] (0,0) -- (3.0,0);
    \end{tikzpicture}
    \caption{Cinematica dell'effetto Compton nel sistema del laboratorio: il fotone incidente $\gamma_i$ viene diffuso di un angolo $\theta_{\gamma}$, mentre l'elettrone inizialmente fermo rincula di un angolo $\theta_e$.}
\end{figure}

Nel LAB i quadrimpulsi iniziali e finali sono
\begin{equation*}
    W_{\text{tot},i} =
    \begin{cases}
        W_{\gamma i} = \begin{pmatrix} E_{\gamma i}/c \\ \vec{p}_{\gamma i} \end{pmatrix} \\[10pt]
        W_{e i} = \begin{pmatrix} E_{e i}/c \\ \vec{p}_{e i} \end{pmatrix} = \begin{pmatrix} m_e c \\ \vec{0} \end{pmatrix}
    \end{cases}
    \qquad\qquad
    W_{\text{tot},f} =
    \begin{cases}
        W_{\gamma f} = \begin{pmatrix} E_{\gamma f}/c \\ \vec{p}_{\gamma f} \end{pmatrix} \\[10pt]
        W_{e f} = \begin{pmatrix} E_{e f}/c \\ \vec{p}_{e f} \end{pmatrix}
    \end{cases}
\end{equation*}
con i versori $\hat{p}_{\gamma i} = \hat{x}$, $\hat{p}_{e i} = \vec{0}$, e
\begin{equation*}
    \hat{p}_{\gamma f} = \cos\theta_{\gamma}\,\hat{x} + \sin\theta_{\gamma}\,\hat{y}
    \qquad\qquad
    \hat{p}_{e f} = \cos\theta_{e}\,\hat{x} - \sin\theta_{e}\,\hat{y}
\end{equation*}

Proiettando la conservazione del quadrimpulso sulle tre coordinate (in unità naturali):
\begin{equation}
    \begin{cases}
        E/c \, : \quad E_{\gamma i} + m_e = E_{\gamma f} + E_{ef} \\
        \hat{x} \, : \quad E_{\gamma i} + 0 = E_{\gamma f}\cos\theta_{\gamma} + p_{ef}\cos\theta_e \\
        \hat{y} \, : \quad 0 + 0 = E_{\gamma f}\sin\theta_{\gamma} - p_{ef}\sin\theta_e
    \end{cases}
\end{equation}
Quadrando e sommando le ultime due equazioni si elimina $\theta_e$:
\begin{equation*}
    p_{ef}^2 = E_{\gamma i}^2 + E_{\gamma f}^2 - 2E_{\gamma i}E_{\gamma f}\cos\theta_{\gamma}
\end{equation*}
e sostituendo in $E_{ef} = \sqrt{m_e^2 + p_{ef}^2}$ nella prima equazione:
\begin{align*}
    m_e^2 + E_{\gamma i}^2 + E_{\gamma f}^2 - 2E_{\gamma i}E_{\gamma f}\cos\theta_{\gamma}
    &= \left(E_{\gamma i} - E_{\gamma f} + m_e\right)^2 \\
    &= E_{\gamma i}^2 + E_{\gamma f}^2 + m_e^2 + 2E_{\gamma i}m_e - 2E_{\gamma f}m_e - 2E_{\gamma i}E_{\gamma f}
\end{align*}
Semplificando i termini uguali:
\begin{equation*}
    E_{\gamma i}E_{\gamma f}\left(\cos\theta_{\gamma} - 1\right) = \left(E_{\gamma f} - E_{\gamma i}\right)m_e
    \implies \frac{1 - \cos\theta_{\gamma}}{m_e} = \frac{E_{\gamma i} - E_{\gamma f}}{E_{\gamma i}E_{\gamma f}}
\end{equation*}

\begin{formulario}{Formula di Compton}
\begin{equation}
    \frac{1}{E_{\gamma f}} - \frac{1}{E_{\gamma i}} = \frac{2\sin^2\!\left(\dfrac{\theta_{\gamma}}{2}\right)}{m_e}
\end{equation}
L'energia ceduta dal fotone è massima per l'angolo di \textbf{retrodiffusione} $\theta_{\gamma} = 180^{\circ}$ ($\pi$), in corrispondenza del quale $\sin^2(\theta_\gamma/2) = 1$.
\end{formulario}

\subsection{Produzione di una Risonanza: $e^+ e^- \to Z^0$}

\begin{esempio}{Energia di Soglia per il Bosone $Z^0$}
Un positrone e un elettrone si annichilano producendo il bosone $Z^0$:
\begin{equation}
    e^+ + e^- \longrightarrow Z^0
\end{equation}
Poiché si lavora in configurazione \emph{collider}, il CM coincide con il LAB e $\hat{p}_+ = +\hat{x}$, $\hat{p}_- = -\hat{x}$, $\hat{p}_Z = 0$. Il $Z^0$ viene prodotto \textbf{a riposo}: l'energia di soglia $E_s$ è quindi l'energia a riposo del $Z^0$, cioè $E_{s,\text{tot}} = E_{0,Z^0}$.

Inoltre $E_+ = E_-$ e $|\vec{p}_+| = |\vec{p}_-|$, perché $e^+$ ed $e^-$ sono particelle identiche a meno della carica. Proiettando la conservazione del quadrimpulso:
\begin{equation*}
    \begin{cases}
        E/c \, : \quad 2E_{\pm} = m_Z \implies E_{\pm} = m_Z/2 \\
        \hat{x} \, : \quad p_+ - p_- = 0 \implies p_+ = p_-
    \end{cases}
\end{equation*}

Da $E_{\pm} = \gamma_{\pm}m_{\pm}$ si ricava il fattore di Lorentz necessario:
\begin{equation*}
    \gamma_{\pm} = \frac{E_{\pm}}{m_{\pm}} = \frac{m_Z}{m_e}\cdot\frac{1}{2} \sim 9 \cdot 10^4
\end{equation*}
e da $\gamma_{\pm}^2 = \dfrac{1}{1-\beta_{\pm}^2}$ segue $1 - \beta_{\pm}^2 = \dfrac{1}{\gamma_{\pm}^2}$, da cui con uno sviluppo di Taylor
\begin{equation*}
    \beta_{\pm} = \sqrt{1 - \frac{1}{\gamma_{\pm}^2}} \approx 1 - \frac{1}{2\gamma_{\pm}^2}
    \implies v_{\pm} = 0{,}999\,c
\end{equation*}
che è la velocità necessaria perché si crei la risonanza.
\end{esempio}

\subsection{Decadimento del Pione Neutro}

\begin{esempio}{$\pi^0 \to \gamma_1 + \gamma_2$}
Il pione neutro $\pi^0$ è una particella instabile che decade in due fotoni. Nel CM il pione è fermo ($\hat{p}_{\pi} = 0$) e i due fotoni sono emessi in direzioni opposte, $\hat{p}_1 = -\hat{x}$ e $\hat{p}_2 = +\hat{x}$. I quadrimpulsi sono
\begin{equation*}
    W_{\text{tot},i} = W_{\pi} = \begin{pmatrix} E_{\pi}/c \\ \vec{p}_{\pi} \end{pmatrix} = \begin{pmatrix} m_{\pi}c \\ \vec{0} \end{pmatrix}
    \qquad\qquad
    W_{\text{tot},f} =
    \begin{cases}
        W_{\gamma 1} = \begin{pmatrix} E_{\gamma 1}/c \\ \vec{p}_{\gamma 1} \end{pmatrix} \\[10pt]
        W_{\gamma 2} = \begin{pmatrix} E_{\gamma 2}/c \\ \vec{p}_{\gamma 2} \end{pmatrix}
    \end{cases}
\end{equation*}
Poiché $\gamma_1$ e $\gamma_2$ sono la stessa particella, $E_{\gamma 1} = E_{\gamma 2}$ e $|\vec{p}_{\gamma 1}| = |\vec{p}_{\gamma 2}|$ (essendo $m_{\gamma} = 0$, vale $E_{\gamma}/c = |\vec{p}_{\gamma}|$). Proiettando:
\begin{equation*}
    \begin{cases}
        E/c \, : \quad m_{\pi} = 2E_{\gamma} \implies E_{\gamma} = m_{\pi}/2 \\
        \hat{x} \, : \quad 0 = -E_{\gamma 1} + E_{\gamma 2} \implies E_{\gamma 1} = E_{\gamma 2}
    \end{cases}
    \qquad \text{nel CM}
\end{equation*}

Ci spostiamo ora dal CM al LAB usando Lorentz per la componente $E/c$:
\begin{align}
    E_1 &= \gamma_1 \left(E_1^{\text{CM}} - \beta_1 \vec{p}_{\gamma 1}^{\;\text{CM}}\right)
         = \gamma_1 |\vec{p}_{\gamma 1}^{\;\text{CM}}| \left(1 - \beta\cos\theta_1\right)
         = \frac{\gamma\, m_{\pi}}{2}\,(1 - \beta) \\
    E_2 &= \gamma_2 \left(E_2^{\text{CM}} + \beta_2 \vec{p}_{\gamma 2}^{\;\text{CM}}\right)
         = \gamma_2 |\vec{p}_{\gamma 2}^{\;\text{CM}}| \left(1 + \beta\cos\theta_2\right)
         = \frac{\gamma\, m_{\pi}}{2}\,(1 + \beta)
\end{align}
nel caso particolare $\theta^{\text{CM}} = 0$. I due fotoni, equivalenti nel CM, appaiono nel LAB uno ``arrossato'' e l'altro ``blueshiftato'': è l'effetto Doppler relativistico.
\end{esempio}

\subsection{La Produzione dell'Antiprotone}

\begin{notastorica}{Segrè e Chamberlain, 1955}
Nel 1955 \textbf{Segrè} e \textbf{Chamberlain} scoprono l'antiprotone grazie all'acceleratore di particelle \emph{Bevatron} del Berkeley National Laboratory (California), vincendo il premio Nobel nel 1959.
\end{notastorica}

La reazione studiata è
\begin{equation}
    \underbrace{p}_{A} + \underbrace{p}_{B} \longrightarrow \underbrace{p}_{1} + \underbrace{p}_{2} + \underbrace{p}_{3} + \underbrace{\bar{p}}_{4}
\end{equation}
Per capire se era possibile produrlo, è stata calcolata l'\textbf{energia di soglia}, la minima perché la reazione si verifichi. I quadrimpulsi totali sono
\begin{equation*}
    W_{\text{tot},i}\big|_{\text{LAB}} = W_A + W_B = \begin{pmatrix} E_A/c + m_B c \\ \vec{p}_A + \vec{0} \end{pmatrix}
    \qquad\qquad
    W_{\text{tot},f}\big|_{\text{CM}} = \begin{pmatrix} (E_1+E_2+E_3+E_4)/c \\ \vec{0} \end{pmatrix}
\end{equation*}
dove tutti gli impulsi finali sono nulli perché si è \textbf{a soglia}: le quattro particelle vengono prodotte ferme nel CM.

Dato che il quadrimpulso si conserva, usiamo l'invariante $s$ di Mandelstam: $s = (W_A+W_B)^2 = (W_1+W_2+W_3+W_4)^2$. Non possiamo però scrivere l'equazione interamente nel LAB (avremmo 16 incognite meno 4 vincoli, cioè 12 variabili): sfruttiamo dunque l'\textbf{invarianza} calcolando $s_i$ nel LAB e $s_f$ nel CM.
\begin{align*}
    s_i\big|_{\text{LAB}} &= E_A^2 - \vec{p}_A^{\;2} + m_B^2 + 2E_A m_B - 2\vec{p}_A \cdot \vec{0}
     = m_A^2 + m_B^2 + 2E_A m_B = 2m_p^2 + 2E_p m_p \\
    s_f\big|_{\text{CM}} &= E_{\text{tot},f}^2 = (4m_p)^2 = 16\,m_p^2 = E_{\text{soglia}}
\end{align*}
Uguagliando:
\begin{equation}
    16\,m_p^2 = 2m_p^2 + 2E_p m_p
    \implies 14\,m_p = 2E_p
    \implies \boxed{E_p = 7\,m_p = 7\,E_{0,p}}
\end{equation}
da cui si ricavano
\begin{equation*}
    K_p = E_p - E_{0,p} = 7m_p - m_p = 6\,m_p
    \qquad
    \gamma_p = \frac{E_p}{m_p} = \frac{7m_p}{m_p} = 7
\end{equation*}
\begin{equation*}
    |\vec{p}_p| = \sqrt{E_p^2 - m_p^2} = m_p\sqrt{7^2 - 1^2} = m_p\sqrt{48}
    \implies \beta = \frac{|\vec{p}_p|}{E_p} = \frac{\sqrt{48}\,m_p}{7\,m_p} = \frac{\sqrt{48}}{7} \approx 0{,}99 \sim 1\,c
\end{equation*}

\subsubsection{La Correzione con l'Impulso di Fermi}

Al tempo, fornire una quantità tale di energia non era possibile; pertanto venne sfruttato l'\textbf{impulso di Fermi} $\vec{p}_F$, l'impulso isotropo posseduto dai nucleoni ($N + P$) del bersaglio in berillio. In questo modo $W_B = \begin{pmatrix} E_B/c \\ \vec{p}_F \end{pmatrix}$ con $p_F \approx 0{,}2\,m_p$, e l'invariante iniziale diventa
\begin{align*}
    s_i &= \underbrace{E_A^2 - \vec{p}_A^{\;2}}_{m_p^2} + \underbrace{E_B^2 - \vec{p}_F^{\;2}}_{m_p^2} + 2E_AE_B - 2\vec{p}_A \cdot \vec{p}_F \\
        &= 2m_p^2 + 2\left(E_AE_B - \vec{p}_A \cdot \vec{p}_F\right)
         = 2m_p^2 + 2\left(E_AE_B + p_A p_F \cdot 1\right) \\
        &= 2m_p^2 + 2E_p\,(m_p + p_F) = 2m_p^2 + 2{,}4\,E_p m_p
\end{align*}
avendo usato $\vec{a}\cdot\vec{b} = a\,b\cos\theta$ con $\theta = \pi$ (l'impulso di Fermi del nucleone è diretto \emph{contro} il proiettile, configurazione più favorevole), e le approssimazioni $E_A \sim |\vec{p}_A| = |\vec{p}_p|$, $E_B \sim m_B$. Uguagliando nuovamente a $s_f = 16 m_p^2$:
\begin{equation}
    16\,m_p^2 = \left(2m_p + 2{,}4\,E_p\right)m_p
    \implies 14\,m_p = 2{,}4\,E_p
    \implies \boxed{E_p = \frac{14}{2{,}4}\,m_p \sim 5{,}83\,m_p}
\end{equation}

\begin{osservazione}{Il Guadagno della Correzione di Fermi}
La correzione di Fermi permette di \textbf{ridurre l'energia} richiesta al fascio da $7\,m_p$ a circa $5{,}83\,m_p$: un risparmio di quasi il $17\%$, decisivo per rendere l'esperimento realizzabile con le tecnologie dell'epoca.
\end{osservazione}

\chapter{Sistemi a Massa Variabile e Sistemi Non Inerziali}

\section{Sistemi a Massa Variabile}

Sappiamo che la formulazione più generale del secondo principio della dinamica è $\vec{F} = \dfrac{d\vec{p}}{dt}$, ma il risultato cambia a seconda che la massa si mantenga costante o no:
\begin{itemize}
    \item \textbf{massa costante:} $\displaystyle \vec{F} = m\,\frac{d\vec{v}}{dt} = m\vec{a}$;
    \item \textbf{massa variabile:} $\displaystyle \vec{F} = \frac{dm}{dt}\,\vec{v} + m\,\frac{d\vec{v}}{dt} = \frac{dm}{dt}\vec{v} + m\vec{a}$.
\end{itemize}

\begin{definizione}{Tasso di Variazione della Massa}
Si definisce \textbf{tasso di variazione della massa} la quantità
\begin{equation}
    \alpha = \frac{dm}{dt}
\end{equation}
e la seconda legge della dinamica assume la forma generale
\begin{equation}
    \vec{F} = \alpha \vec{v} + m\vec{a}
\end{equation}
\end{definizione}

\subsection{Il Nastro Trasportatore}

\begin{esempio}{Nastro Trasportatore con Carico Continuo}
Su un nastro trasportatore viene versata materia a ritmo costante:
\begin{equation}
    \frac{dm}{dt} = \alpha = \text{cost}
    \qquad \implies \qquad
    m(t) = m_0 + \alpha t
\end{equation}
dove $m_0$ è la massa iniziale del nastro.
\end{esempio}

\begin{figure}[ht]
    \centering
    \begin{tikzpicture}[>=stealth, scale=1.0]
        % Tramoggia
        \draw[very thick, mypen] (-0.9,2.8) -- (-0.25,1.6) -- (0.25,1.6) -- (0.9,2.8);
        \foreach \x/\y in {-0.45/2.35, 0.1/2.5, 0.4/2.15, -0.15/2.05} {
            \fill[mygreen!70!black] (\x,\y) circle (1.1pt);
        }
        % Flusso dm
        \foreach \y in {1.35, 1.05, 0.75} {
            \fill[mygreen!70!black] (0,\y) circle (1.2pt);
        }
        \node[right, font=\small, mygreen!50!black] at (0.18,1.05) {$dm$};
        % Nastro
        \draw[very thick, mypen] (-2.6,0) rectangle (2.6,0.45);
        \foreach \x in {-2.0,-1.2,-0.4,0.4,1.2,2.0} {
            \draw[thick, mypen!55] (\x,0) -- (\x,0.45);
        }
        \fill[mygreen!70!black] (-1.6,0.6) circle (2.2pt);
        \fill[mygreen!70!black] (-0.8,0.62) circle (2.6pt);
        \fill[mygreen!70!black] (0.0,0.6) circle (2.0pt);
        % Velocità
        \draw[->, very thick, notered] (0.9,-0.55) -- (2.3,-0.55)
            node[midway, below, font=\small\bfseries] {$\vec{v}(t)$};
        % Motore
        \draw[->, very thick, boxese] (-2.9,0.22) -- (-3.9,0.22);
        \node[left, font=\small, boxese, align=center] at (-3.95,0.22) {motore\\$\vec{F}$};
    \end{tikzpicture}
    \caption{Il nastro trasportatore: la materia versata dalla tramoggia si aggiunge al nastro con tasso costante $\alpha = dm/dt$, facendone crescere la massa e, in assenza di forze esterne, rallentandone il moto.}
\end{figure}

\paragraph{Caso $\vec{F} = 0$ (motore spento).}
La quantità di moto si conserva: $\vec{p}_f = \vec{p}_i = \vec{p}_0$, ovvero $m(t)\,\vec{v}(t) = m_0\vec{v}_0$, da cui
\begin{equation}
    \boxed{\;\vec{v}(t) = \frac{m_0}{m(t)}\,\vec{v}_0 < \vec{v}_0\;}
\end{equation}
poiché $m(t) > m_0$. \textbf{Il nastro rallenta}: la materia che cade deve essere accelerata fino alla velocità del nastro, e l'unica riserva disponibile è la quantità di moto del nastro stesso.

\paragraph{Caso $\vec{F} \neq 0$ (motore acceso).}
Ora $\vec{p}$ varia secondo $d\vec{p} = m\vec{a} + \alpha\vec{v}$. Se vogliamo mantenere $\vec{a} = 0$, cioè $v(t) = v_0$ costante, allora
\begin{equation}
    d\vec{p} = \alpha\vec{v}_0
\end{equation}
prodotta dalla forza $\vec{F}$ del motore.

\begin{formulario}{Nastro Trasportatore a Velocità Costante}
\textbf{Impulso del motore:}
\begin{equation}
    I = \int_0^t \vec{F}(t')\,dt' = \vec{p}_f - \vec{p}_i = m(t)\,v_0 - m_0 v_0
      = \cancel{m_0v_0} + \alpha v_0 t - \cancel{m_0v_0} = \alpha v_0 t
\end{equation}

\textbf{Variazione di energia cinetica:}
\begin{align}
    \frac{dK}{dt} &= \frac{d}{dt}\left(\frac{\vec{p}^{\;2}}{2m}\right)
      = \frac{1}{2m}\frac{d\vec{p}^{\;2}}{dt} + \frac{\vec{p}^{\;2}}{2}\frac{d}{dt}\left(\frac{1}{m}\right)
      = \frac{2\vec{p}}{2m}\frac{d\vec{p}}{dt} + \frac{\vec{p}^{\;2}}{2}\left(-\frac{1}{m^2}\frac{dm}{dt}\right) \notag \\
      &= \frac{\cancel{m}\vec{v}_0}{\cancel{m}}\cdot\vec{F} - \frac{\cancel{m^2}\vec{v}_0^{\;2}}{2}\cdot\frac{\alpha}{\cancel{m^2}}
      = \vec{v}_0 \vec{F} - \frac{1}{2}\alpha \vec{v}_0^{\;2}
\end{align}

\textbf{Potenza del motore:}
\begin{equation}
    P = \frac{dK}{dt} + \frac{1}{2}\alpha \vec{v}_0^{\;2}
\end{equation}

\textbf{Lavoro del motore:}
\begin{align}
    \mathcal{L} &= \int_0^t P(t')\,dt'
      = \int_0^t \frac{dK}{dt'}\,dt' + \int_0^t \frac{1}{2}\alpha \vec{v}_0^{\;2}\,dt'
      = K(t) - K_0 + \frac{1}{2}\alpha t \,\vec{v}_0^{\;2} \notag \\
      &= \frac{(m_0 + \alpha t)\vec{v}_0^{\;2}}{2} - \frac{m_0\vec{v}_0^{\;2}}{2} + \frac{\alpha \vec{v}_0^{\;2} t}{2}
      = \alpha \vec{v}_0^{\;2}\, t
\end{align}
\end{formulario}

\begin{osservazione}{Dove Finisce Metà del Lavoro}
Confrontando il lavoro totale $\mathcal{L} = \alpha v_0^2 t$ con la variazione di energia cinetica $\Delta K = \frac{1}{2}\alpha v_0^2 t$, si vede che \textbf{solo metà} del lavoro compiuto dal motore si ritrova come energia cinetica del sistema. L'altra metà viene dissipata nell'accelerare per attrito la materia versata, che arriva ferma sul nastro e deve essere portata alla velocità $v_0$: è la stessa perdita, irriducibile, che si incontra in un urto totalmente anelastico.
\end{osservazione}

\section{Sistemi Non Inerziali Rotanti}

\subsection{I Moti della Terra}

Rotazione e rivoluzione non hanno molta influenza su fenomeni situati in uno spazio-tempo relativamente piccolo (in cui la Terra può considerarsi \textbf{inerziale}), bensì su fenomeni di durata e distanze maggiori, dove l'effetto risulta più evidente (regime \textbf{non inerziale}).

\begin{esempio}{Le Accelerazioni dei Moti Terrestri}
\textbf{Rotazione} (periodo $T = 24\,\mathrm{h}$, raggio terrestre $r_T = 6{,}3 \cdot 10^6\,\mathrm{m}$):
\begin{equation}
    a_{\omega} = \omega^2 r_T = \left(\frac{2\pi}{24\,\mathrm{h}}\right)^{\!2} \cdot 6{,}3 \cdot 10^6\,\mathrm{m}
    = 0{,}033\,\frac{\mathrm{m}}{\mathrm{s}^2} = 0{,}0034\,g
\end{equation}

\textbf{Rivoluzione} (periodo $T = 365\,\mathrm{d}$, raggio orbitale $r_o = 1{,}5 \cdot 10^{11}\,\mathrm{m}$):
\begin{equation}
    a_{\Omega} = \Omega^2 r_o = \left(\frac{2\pi}{365\,\mathrm{d}}\right)^{\!2} \cdot 1{,}5 \cdot 10^{11}\,\mathrm{m}
    = 0{,}0059\,\frac{\mathrm{m}}{\mathrm{s}^2} = 0{,}0006\,g
\end{equation}

Entrambe sono piccole frazioni di $g$: per questo, nella maggior parte dei problemi di meccanica, il laboratorio terrestre può essere trattato come inerziale.
\end{esempio}

\subsection{Il Fenomeno dei Cicloni}

La rotazione terrestre è in senso \textbf{antiorario} nell'emisfero boreale e \textbf{orario} nell'emisfero australe. In base all'emisfero in cui ci si trova, le correnti d'aria seguono lo stesso moto e si spostano verso il centro del moto, ``convergendo'' ai poli.

I cicloni sono dovuti a differenze di pressione atmosferica elevate: dove essa sarebbe normalmente di $\sim 1000\,\mathrm{mbar}$, nell'occhio del ciclone arriva a $\sim 860\,\mathrm{mbar}$.

\begin{figure}[ht]
    \centering
    \begin{tikzpicture}[>=stealth, scale=1.0]
        % --- Emisfero Nord
        \begin{scope}
            \draw[->, thin, mypen!70] (-1.7,0) -- (1.7,0) node[right, font=\scriptsize] {E};
            \node[left, font=\scriptsize, mypen!70] at (-1.75,0) {W};
            \draw[->, thin, mypen!70] (0,-1.7) -- (0,1.7) node[above, font=\scriptsize] {N};
            \node[below, font=\scriptsize, mypen!70] at (0,-1.75) {S};
            \foreach \a in {0,60,120,180,240,300} {
                \draw[->, very thick, boxoss]
                    (\a:1.45) arc (\a:\a+50:1.45);
            }
            \foreach \a in {30,150,270} {
                \draw[->, thick, boxoss!70] (\a:0.85) arc (\a:\a+55:0.85);
            }
            \node[font=\small\bfseries, align=center] at (0,-2.55) {emisfero NORD\\(antiorario)};
        \end{scope}
        % --- Emisfero Sud
        \begin{scope}[xshift=5.6cm]
            \draw[->, thin, mypen!70] (-1.7,0) -- (1.7,0) node[right, font=\scriptsize] {E};
            \node[left, font=\scriptsize, mypen!70] at (-1.75,0) {W};
            \draw[->, thin, mypen!70] (0,-1.7) -- (0,1.7) node[above, font=\scriptsize] {N};
            \node[below, font=\scriptsize, mypen!70] at (0,-1.75) {S};
            \foreach \a in {0,60,120,180,240,300} {
                \draw[->, very thick, boxatt]
                    (\a:1.45) arc (\a:\a-50:1.45);
            }
            \foreach \a in {30,150,270} {
                \draw[->, thick, boxatt!70] (\a:0.85) arc (\a:\a-55:0.85);
            }
            \node[font=\small\bfseries, align=center] at (0,-2.55) {emisfero SUD\\(orario)};
        \end{scope}
    \end{tikzpicture}
    \caption{Il verso di rotazione dei cicloni nei due emisferi. La deflessione è prodotta dalla forza di Coriolis, che agisce in verso opposto a nord e a sud dell'equatore.}
\end{figure}

\subsection{Trasformazioni per Sistemi Rototraslanti}

Sia $S$ il sistema inerziale e $S'$ il sistema non inerziale, che ruota con velocità angolare $\vec{\omega}$ rispetto a $S$.

\paragraph{Posizione di $P$.}
Dalla somma vettoriale $\overline{OP} = \overline{OO'} + \overline{O'P}$ si ottiene
\begin{equation}
    \vec{r} = \vec{r}\,' + \vec{R}
    \qquad \text{(da $S'$ a $S$)}
    \qquad\qquad
    \vec{r}\,' = \vec{r} - \vec{R}
    \qquad \text{(da $S$ a $S'$)}
\end{equation}
dove $\vec{R}$ è la posizione congiungente relativa fra le due origini $(O, O')$.

\begin{figure}[ht]
    \centering
    \begin{tikzpicture}[>=stealth, scale=1.05]
        % Sistema S
        \draw[->, thick, mypen] (0,0) -- (3.4,0) node[right, font=\small] {$y$};
        \draw[->, thick, mypen] (0,0) -- (0,2.9) node[above, font=\small] {$z$};
        \draw[->, thick, mypen] (0,0) -- (-1.5,-1.5) node[below left, font=\small] {$x$};
        \node[below left, font=\small\bfseries, mypen] at (-0.05,-0.05) {$S$};
        % Origine O'
        \coordinate (Op) at (2.5,1.5);
        \draw[->, very thick, notered] (0,0) -- (Op) node[midway, below right, font=\small] {$\vec{R}$};
        \node[below right, font=\small\bfseries, boxese] at (2.55,1.42) {$O'$};
        % Assi S'
        \draw[->, thick, boxese] (Op) -- ++(1.5,0) node[right, font=\scriptsize] {$y'$};
        \draw[->, thick, boxese] (Op) -- ++(0,1.4) node[above, font=\scriptsize] {$z'$};
        \draw[->, thick, boxese] (Op) -- ++(-0.8,-0.8) node[below left, font=\scriptsize] {$x'$};
        % omega
        \draw[->, very thick, boxatt] (Op) ++(0.35,1.15) arc (200:-20:0.42);
        \node[right, font=\small\bfseries, boxatt] at (3.45,2.75) {$\vec{\omega}$};
        % Punto P
        \coordinate (P) at (3.6,2.5);
        \fill[notered] (P) circle (2.2pt) node[above right, font=\small\bfseries] {$P$};
        \draw[->, thick, mypen, dashed] (0,0) -- (P) node[pos=0.55, above left, font=\small] {$\vec{r}$};
        \draw[->, thick, boxese, dashed] (Op) -- (P) node[pos=0.5, right, font=\small] {$\vec{r}\,'$};
    \end{tikzpicture}
    \caption{Composizione delle posizioni fra il sistema inerziale $S$ e il sistema rototraslante $S'$, quest'ultimo animato dalla velocità angolare $\vec{\omega}$.}
\end{figure}

\paragraph{Velocità di $P$.}
Derivando rispetto al tempo nel sistema $O$:
\begin{align}
    \vec{v} = \left.\frac{d\vec{r}}{dt}\right|_{O}
      &= \left.\frac{d\vec{r}\,'}{dt}\right|_{O} + \left.\frac{d\vec{R}}{dt}\right|_{O} \notag \\
      &= \left.\frac{d\vec{r}\,'}{dt}\right|_{O'} + \vec{\omega}\times\vec{r}\,' + \left.\frac{d\vec{R}}{dt}\right|_{O} \notag \\
      &= \vec{v}\,' + \underbrace{\vec{\omega}\times\vec{r}\,' + \vec{V}}_{\text{velocità di trascinamento}}
\end{align}

\begin{definizione}{Velocità di Trascinamento Locale}
La velocità di trascinamento \textbf{dipende dal punto} $P$ considerato ed è composta da due contributi:
\begin{itemize}[leftmargin=1.4em]
    \item $\vec{V}$: la velocità relativa fra le origini $(O, O')$;
    \item $\vec{\omega}\times\vec{r}\,'$: la velocità dovuta alla rotazione.
\end{itemize}
Invertendo si ottiene la legge di trasformazione da $S$ a $S'$:
\begin{equation}
    \vec{v}\,' = \vec{v} - \vec{\omega}\times\vec{r}\,' - \vec{V}
\end{equation}
\end{definizione}

\paragraph{Accelerazione di $P$.}
Derivando ancora una volta:
\begin{align}
    \vec{a} = \left.\frac{d\vec{v}}{dt}\right|_{O}
      &= \left.\frac{d\vec{v}\,'}{dt}\right|_{O} + \left.\frac{d(\vec{\omega}\times\vec{r}\,')}{dt}\right|_{O} + \left.\frac{d\vec{V}}{dt}\right|_{O} \notag \\
      &= \left.\frac{d\vec{v}\,'}{dt}\right|_{O'} + \vec{\omega}\times\vec{v}\,' + \vec{\omega}\times\left.\frac{d\vec{r}\,'}{dt}\right|_{O} + \left.\frac{d\vec{V}}{dt}\right|_{O} \notag \\
      &= \vec{a}\,' + \vec{\omega}\times\vec{v}\,' + \vec{\omega}\times\left(\vec{v}\,' + \vec{\omega}\times\vec{r}\,'\right) + \vec{A} \notag \\
      &= \vec{a}\,' + \underbrace{2\,\vec{\omega}\times\vec{v}\,'}_{\text{Coriolis}} + \underbrace{\vec{\omega}\times(\vec{\omega}\times\vec{r}\,') + \vec{A}}_{\text{trascinamento}}
\end{align}

\begin{definizione}{Accelerazione di Coriolis e di Trascinamento}
\begin{itemize}[leftmargin=1.4em]
    \item L'\textbf{accelerazione di Coriolis} $2\,\vec{\omega}\times\vec{v}\,'$ compare solo se il corpo \emph{si muove} nel sistema rotante ($\vec{v}\,' \neq 0$).
    \item L'\textbf{accelerazione di trascinamento locale} dipende dal punto $P$ ed è composta da $\vec{A}$ (accelerazione relativa fra le origini) e da $\vec{\omega}\times(\vec{\omega}\times\vec{r}\,')$ (accelerazione dovuta alla rotazione).
\end{itemize}
Invertendo:
\begin{equation}
    \vec{a}\,' = \vec{a} - 2\,\vec{\omega}\times\vec{v}\,' - \vec{\omega}\times(\vec{\omega}\times\vec{r}\,') - \vec{A}
\end{equation}
\end{definizione}

\subsection{Correlazione fra $\vec{g}$ e la Latitudine}

Sia $S(O,x,y,z)$ il sistema inerziale e $S'(O',x',y',z')$ quello rotante solidale con la Terra, con le origini coincidenti ($O \equiv O'$, dunque $\vec{r} = \vec{r}\,'$ e $\vec{V}, \vec{A} = 0$). Per un punto $P$ \textbf{fermo} in $S'$ si ha $\vec{v}\,' = 0$ e $\vec{a}\,' = 0$, per cui i termini di Coriolis si annullano e resta
\begin{equation*}
    \vec{v} = \cancelto{0}{\vec{v}\,'} + \vec{\omega}\times\vec{r}\,'
    \qquad\qquad
    \vec{a} = \cancelto{0}{\vec{a}\,'} + \cancelto{0}{2\vec{\omega}\times\vec{v}\,'} + \vec{\omega}\times(\vec{\omega}\times\vec{r}\,')
\end{equation*}

Supponendo la Terra una sfera perfetta, $|\vec{g}_0| = \text{cost}$, mentre $|\vec{g}\,|$ è l'accelerazione \textbf{locale}, diversa in ogni punto:
\begin{equation}
    \vec{g} = \vec{g}_0 - \vec{\omega}\times(\vec{\omega}\times\vec{r}\,')
\end{equation}
Detta $\theta$ la latitudine, si ha $\vec{\omega} = \omega\hat{z}$ e $\vec{r} = R\,(\cos\theta\,\hat{x} + \sin\theta\,\hat{z})$. Sviluppando il doppio prodotto vettoriale con l'identità $\vec{a}\times(\vec{b}\times\vec{c}) = \vec{b}\,(\vec{a}\cdot\vec{c}) - \vec{c}\,(\vec{a}\cdot\vec{b})$:
\begin{align*}
    \vec{\omega}\times(\vec{\omega}\times\vec{r}\,) &= \vec{\omega}\,(\vec{\omega}\cdot\vec{r}\,) - \vec{r}\,(\vec{\omega}\cdot\vec{\omega})
      = (\omega R \sin\theta)(\omega\hat{z}) - \omega^2\left(R\cos\theta\,\hat{x} + R\sin\theta\,\hat{z}\right) \\
      &= \cancel{\omega^2 R \sin\theta\,\hat{z}} - \omega^2 R\cos\theta\,\hat{x} - \cancel{\omega^2 R\sin\theta\,\hat{z}}
      = -\omega^2 R \cos\theta\,\hat{x}
\end{align*}
da cui
\begin{equation}
    \boxed{\;\vec{g} = \vec{g}_0 + \hat{x}\,\omega^2 R \cos\theta\;}
\end{equation}

\begin{osservazione}{La Deviazione della Verticale}
$\vec{g}$ non ha la stessa direzione di $\vec{g}_0$: è deviata dall'accelerazione ``centrifuga'' dovuta alla rotazione. Più ci si avvicina ai poli (dove il raggio di rotazione si annulla), più $|\vec{g}\,|$ aumenta:
\begin{equation*}
    g = g_0 - \omega^2 R\cos\theta
\end{equation*}
\begin{itemize}[leftmargin=1.4em]
    \item ai \textbf{poli} ($\theta = 90^{\circ}$): $g = g_0$, valore \textbf{massimo};
    \item all'\textbf{equatore} ($\theta = 0^{\circ}$): $g = g_0 - \omega^2 R$, valore \textbf{minimo}.
\end{itemize}
\end{osservazione}

\begin{figure}[ht]
    \centering
    \begin{tikzpicture}[>=stealth, scale=1.0]
        \foreach \i/\lat/\lab in {0/90/{(a) polo}, 1/45/{(b) latitudine media}, 2/0/{(c) equatore}} {
            \begin{scope}[xshift=\i*4.6cm]
                \draw[thick, mypen] (0,0) circle (1.25);
                \draw[dashed, mypen!45] (-1.45,0) -- (1.45,0);
                \draw[->, thick, mypen!60] (0,-1.55) -- (0,1.8) node[above, font=\scriptsize] {$\vec{\omega}$};
                \coordinate (P) at (\lat:1.25);
                \fill[notered] (P) circle (2.2pt);
                % g0: diretta verso il centro
                \draw[->, line width=1.3pt, mypen] (P) -- ($(P)!0.68!(0,0)$);
                \node[font=\scriptsize, mypen] at ($(P)!0.45!(0,0)$) [xshift=-10pt, yshift=-3pt] {$\vec{g}_0$};
                % a_c: orizzontale uscente, modulo proporzionale a cos(latitudine)
                \ifnum\lat=90
                    \node[font=\scriptsize, boxatt, right] at ($(P)+(0.18,0.28)$) {$\vec{a}_c = \vec{0}$};
                \else
                    \draw[->, line width=1.3pt, boxatt] (P) -- ++({1.05*cos(\lat)},0);
                    \node[font=\scriptsize, boxatt] at ($(P)+({1.05*cos(\lat)+0.28},0.22)$) {$\vec{a}_c$};
                \fi
                \node[font=\scriptsize, align=center] at (0,-2.15) {\lab};
            \end{scope}
        }
    \end{tikzpicture}
    \caption{L'accelerazione centrifuga $\vec{a}_c = \omega^2 R\cos\theta$ in funzione della latitudine: nulla ai poli, massima e interamente opposta a $\vec{g}_0$ all'equatore. La somma vettoriale $\vec{g} = \vec{g}_0 + \vec{a}_c$ è la gravità effettivamente misurata.}
\end{figure}

\subsection{La Dinamica in un Riferimento Rotante}

Scrivendo il secondo principio nel sistema inerziale e sostituendo la composizione delle accelerazioni:
\begin{equation}
    \sum\vec{F} = m\vec{a} = m\left(\vec{a}\,' + 2\,\vec{\omega}\times\vec{v}\,' + \vec{\omega}\times(\vec{\omega}\times\vec{r}\,') + \vec{A}\right)
\end{equation}
dove si riconoscono
\begin{align*}
    m\vec{a}\,' &= \textstyle\sum\vec{F}\,' \\
    2m\,\vec{\omega}\times\vec{v}\,' &= \vec{F}_{Co} \quad \to \quad \textbf{forza di Coriolis} \\
    m\,\vec{\omega}\times(\vec{\omega}\times\vec{r}\,') &= \vec{F}_{ce} \quad \to \quad \textbf{forza centrifuga} \\
    m\vec{A} &= \vec{F}_{tr} \quad \to \quad \textbf{forza di trascinamento}
\end{align*}

\begin{formulario}{Secondo Principio in un Riferimento Rotante}
\begin{equation}
    \sum\vec{F}\,' = \sum\vec{F} + \sum\vec{F}_{\text{app}}
    = \sum\vec{F} - \left(2m\,\vec{\omega}\times\vec{v}\,' + m\,\vec{\omega}\times(\vec{\omega}\times\vec{r}\,') + m\vec{A}\right)
\end{equation}
\end{formulario}

\begin{attenzione}{Le Forze Apparenti Esistono Solo nel Sistema Non Inerziale}
Coriolis, centrifuga e trascinamento compaiono \textbf{soltanto} nella descrizione fatta dal sistema non inerziale. Per l'osservatore inerziale esse non esistono: ciò che egli vede è semplicemente un corpo che si muove di moto accelerato.
\end{attenzione}

\subsection{Il Caso della Stazione Spaziale Internazionale}

\begin{esempio}{Assenza di Peso sulla ISS}
La \textbf{ISS} (\emph{International Space Station}) orbita a una quota $d_T = 450\,\mathrm{km}$ con periodo $T = 93\,\mathrm{min}$, cioè
\begin{equation*}
    \frac{1\,\text{giorno}}{93\,\mathrm{min}} = \frac{1440}{93} = 15{,}5 \text{ giri in un giorno}
\end{equation*}
Il suo raggio orbitale è $r_I = r_T + d_T = 6370 + 450 = 6820\,\mathrm{km}$, per cui il campo gravitazionale alla sua quota vale
\begin{equation}
    \frac{g_I}{g} = \frac{r_T^2}{r_I^2} = \frac{r_T^2}{(r_T + d_T)^2}
    = \frac{r_T^2}{r_T^2 + d_T^2 + 2r_Td_T}
    = \frac{1}{1 + \left(\dfrac{d_T}{r_T}\right)^{\!2} + 2\dfrac{d_T}{r_T}}
    \approx 1 - \frac{2d_T}{r_T} \cong 0{,}86
\end{equation}
avendo sviluppato in serie di Taylor e trascurato $(d_T/r_T)^2$ rispetto a $2\,d_T/r_T$.

\textbf{Alla quota della ISS la gravità è ancora l'$86\%$ di quella al suolo}: l'assenza di peso non è dunque dovuta alla lontananza dalla Terra.
\end{esempio}

\begin{figure}[ht]
    \centering
    \begin{tikzpicture}[>=stealth, scale=1.0]
        % Terra
        \shade[ball color=boxoss!25] (0,0) circle (1.5);
        \draw[thick, mypen] (0,0) circle (1.5);
        \node[font=\scriptsize, mypen] at (0,0) {Terra};
        % Orbita
        \draw[dashed, thick, mypen!60] (0,0) circle (2.35);
        % ISS
        \coordinate (I) at (38:2.35);
        \fill[notered] (I) circle (2.6pt);
        \node[right, font=\scriptsize\bfseries, notered] at ($(I)+(0.12,0.1)$) {ISS};
        % raggi
        \draw[->, thick, mypen] (0,0) -- (38:1.5) node[midway, below right, font=\scriptsize] {$r_T$};
        \draw[->, thick, boxatt] (38:1.5) -- (I) node[midway, above left, font=\scriptsize] {$d_T$};
        \draw[->, thick, boxese] (0,0) -- (-160:2.35) node[midway, below, font=\scriptsize] {$r_I$};
        % g_I
        \draw[->, very thick, boxoss] (I) -- ($(I)!0.55!(0,0)$);
        \node[font=\scriptsize, boxoss] at ($(I)!0.42!(0,0)$) [xshift=11pt, yshift=6pt] {$\vec{g}_I$};
    \end{tikzpicture}
    \caption{Geometria orbitale della Stazione Spaziale Internazionale. Alla quota $d_T = 450\,\mathrm{km}$ il campo gravitazionale terrestre vale ancora circa l'$86\%$ del valore al suolo.}
\end{figure}

\paragraph{Un oggetto fermo sulla ISS.}
Consideriamo un oggetto fermo rispetto alla stazione ($\vec{v}\,' = \vec{a}\,' = 0$). Nel \textbf{sistema rotante} solidale con la ISS:
\begin{equation*}
    \sum\vec{F}\,' = \sum\vec{F} + \sum\vec{F}_{\text{app}} = \vec{P}_I + \vec{N} + \vec{F}_{\text{app}}
\end{equation*}
con $\sum\vec{F}\,' = m\vec{a}\,' = 0$, e quindi $\vec{P}_I + \vec{N} = m\vec{g}_I + \vec{N}$. Le forze apparenti valgono
\begin{equation*}
    \vec{F}_{\text{app}} = \cancelto{0}{2m\,\vec{\omega}\times\vec{v}\,'} - m\,\vec{\omega}\times(\vec{\omega}\times\vec{r}\,') - \cancelto{0}{m\vec{A}}
    = m\,\omega^2\vec{r}_I
\end{equation*}
(il termine $\vec{\omega}\times(\vec{\omega}\times\vec{r}\,')$ è trascurabile per il corpo rispetto a $S'$, restando la sola accelerazione centripeta $\omega^2 \vec{r}_I$), da cui
\begin{equation}
    0 = \vec{N} + m\left(\vec{g}_I + \omega^2\vec{r}_I\right)
    \implies \vec{N} = -m\left(\vec{g}_I + \omega^2\vec{r}_I\right)
    \qquad \text{con } \vec{g}_I \parallel -\hat{r}_I \, , \; \vec{r}_I \parallel \hat{r}_I
\end{equation}

\paragraph{L'equazione del moto della stazione.}
Scrivendo il secondo principio per la ISS stessa:
\begin{equation*}
    \sum\vec{F}_I = m_I\,\vec{g}_I = \text{``peso'' della ISS}
    \implies m_I\vec{a}_I = m_I\vec{g}_I \implies \vec{a}_I = \vec{g}_I
\end{equation*}
ovvero, sviluppando l'accelerazione nel sistema rotante,
\begin{equation*}
    \cancelto{0}{\vec{a}\,'_I} + \cancelto{0}{\vec{\omega}\times\vec{v}\,'_I} + \vec{\omega}\times(\vec{\omega}\times\vec{r}\,'_I) + \vec{A} = \vec{g}_I
\end{equation*}
poiché $\vec{a}\,'_I, \vec{v}\,'_I, \vec{r}\,'_I = 0$ (la stazione è ferma rispetto a sé stessa), per cui
\begin{equation}
    \vec{g}_I = \vec{a}_I = \vec{A} = -\omega^2 \vec{r}_I
\end{equation}

\begin{teorema}{L'Assenza di Peso in Orbita}
Siccome $\vec{g}_I = \vec{a}_I = -\omega^2\vec{r}_I$, sostituendo nell'espressione della reazione vincolare si ottiene
\begin{equation}
    \vec{N} = -m\left(-\omega^2\vec{r}_I + \omega^2\vec{r}_I\right) = \vec{0}
\end{equation}
La reazione vincolare è \textbf{nulla} quando la ISS è in orbita: si sperimenta l'\textbf{assenza di gravità}.

Il peso non scompare --- $g_I$ vale ancora l'$86\%$ di $g$ --- ma è interamente ``consumato'' per fornire l'accelerazione centripeta dell'orbita. Stazione e occupanti sono in \textbf{caduta libera} attorno alla Terra: ciò che si annulla non è la gravità, bensì la forza di contatto che, al suolo, ci dà la sensazione di avere un peso.
\end{teorema}

% \include{Chapters/02-StateOfTheArt}
% \include{Chapters/03-Proposal}
% \include{Chapters/04-Evaluation}
% \include{Chapters/05-Conclusions}

%--------------------------------------------------------------
% Bibliography
%--------------------------------------------------------------
% \bibliographystyle{plain} 
% \bibliography{Bibliography} % Entries are in the Bibliography.bib file

%--------------------------------------------------------------
\end{document}
%--------------------------------------------------------------
